% !TeX program = xelatex
% ================================================================
% 《数值线性代数及其应用》（Numerical Linear Algebra and Its Applications, 2e）
% 作者：Xiao-Qing Jin（金小庆）, Yi-Min Wei（魏益民）, Zhi Zhao（赵志）
% 出版社：Science Press（科学出版社），2015
% 中文译本 —— 基于数译(mathtranslations.org)排版模板 mathtranslation.cls v3.3
% 构建：bash tools/build.sh full  （xelatex×2 → biber → xelatex×2）
% 注意：所有 .tex/.cls 须 UTF-8 无 BOM + CRLF 行尾。
% ================================================================
\documentclass[
  BookTitleCN={数值线性代数及其应用},
  BookTitleEN={Numerical Linear Algebra and Its Applications},
  OriginalAuthor={Xiao-Qing Jin, Yi-Min Wei, Zhi Zhao},
  OriginalEdition={Second Edition},
  OriginalPublisher={Science Press},
  OriginalYear={2015},
  Translator={Fingolfin},
  ModelUsed={DeepSeek-V4.1-Flash},
  BibStyle=numeric,
  TranslationDate={2026年11月3日}
]{mathtranslation}

\usepackage{chngcntr}
\usepackage{needspace}% 避免“习题”等小标题落在页末成为孤行

% openright 换章/换篇时会插入空白页。本模板的空白页会沿用上一页的页眉
% 与页码（实测：空白页上印有“iv  《信息与计算科学丛书》序”），故重定义
% \cleardoublepage，令该空白页使用 empty 样式，成为真正的空白页。
\makeatletter
\renewcommand{\cleardoublepage}{%
  \clearpage
  \if@twoside
    \ifodd\c@page\else
      \hbox{}\thispagestyle{empty}\newpage
      \if@twocolumn\hbox{}\newpage\fi
    \fi
  \fi}
\makeatother

% 允许版心稍微拉伸以避免孤立过宽行（不影响数学正确性）
\setlength{\emergencystretch}{3em}

% ---------- 编号体系（待按源书复核后最终确定） ----------
% 公式随章编号：(1.1), (1.2), ...
\numberwithin{equation}{chapter}
% 定理/引理/推论/命题/例/问题 暂保持 cls 默认（按节 X.M），译到各章时据源书调整；
% 若源书按章编号，取消下行注释并按章重映射：
% \counterwithout{theorem}{section}   \counterwithin{theorem}{chapter}
% \counterwithout{lemma}{section}     \counterwithin{lemma}{chapter}
% \counterwithout{corollary}{section} \counterwithin{corollary}{chapter}
% \counterwithout{proposition}{section}\counterwithin{proposition}{chapter}
% \counterwithout{example}{section}   \counterwithin{example}{chapter}
% \counterwithout{problem}{section}   \counterwithin{problem}{chapter}
% \counterwithout{figure}{section}    \counterwithin{figure}{chapter}

% 本项目专用宏（不改动 cls；新宏只加在这里）
\usepackage{mycommand}

% ================================================================
% 标点体例：正文一律使用半角（英文）句读标点。
% 注意：ctex 的 punct 默认为 quanjiao —— 即使源文件写 ASCII 的「, . ; :」，
% 也会被替换成全角字形（实测：PDF 里逗号占一个字宽，文字层还原成「，」）。
% 故必须显式改为 plain（西文标点），ASCII 标点才按拉丁字体半角排印。
\ctexset{punct=plain}
% 正文源文件已由 tools/fix_punct.py 批量转换；此处补上三处由模板/宏
% 生成的标点，使其与正文一致（不改动 mathtranslation.cls 与 mycommand.sty）：
%   1) \newterm 的中英对照括注：原来印作「矩阵（matrix）」，改为「矩阵 (matrix)」；
%   2) 算法标题：原来印作「算法 2.1（高斯消去法）」，改为半角括号（重定义见下）；
%   3) 「注：」的标题冒号：由全角冒号改为半角冒号（重定义 mremark 环境）。
% ================================================================
\makeatletter
\renewcommand{\newterm}[3]{%
  \@ifundefined{MT@seen@#1}{%
    \expandafter\gdef\csname MT@seen@#1\endcsname{}%
    \phantomsection
    \label{term:#1}%
    \protected@write\@auxout{}{%
      \string\MTRegisterTerm{\unexpanded{#2}}{\unexpanded{#3}}{term:#1}%
    }%
  }{}%
  {\termfont #2} (#3)%
}
% 算法环境：与 mycommand.sty 的定义逐字相同，仅把标题两侧的全角括号换成半角
% （含其后不必再补的间距，见 \newterm 口径）。若将来 mycommand.sty 改动此处，
% 需同步本定义。
\renewenvironment{algorithm}[1]{%
  \par\medskip
  \refstepcounter{algorithm}%
  \noindent{\headfont 算法\ \textbf{\thealgorithm} (#1)}\par
  \smallskip
}{%
  \par\medskip
}
% 「注：」标题冒号改半角：mycommand.sty 用 \newtheoremstyle{mtcolon} 把冒号
% 写死在样式里（不改动该文件），故在本地重定义 mremark 环境，间距与字体
% 保持与 mtcolon 相同（上下各 6pt、正文 \theoremfont、标签 \headfont）。
\renewenvironment{mremark}{%
  \par\addvspace{6pt}%
  \noindent{\headfont 注:}\hspace{0.5em}\theoremfont\ignorespaces
}{%
  \par\addvspace{6pt}%
}
\makeatother

% ================================================================
% proof 标题分隔符：v3.3 把分隔符抽成 \proofSeparator，默认 \enspace，
% 即标题排作「证明 」而不再带冒号（模仿 ntheorem 的 nonumberplain）。
% 本书此前一直按「证明：」排版且已定稿，故显式恢复旧体例，使本次升级
% 不改变标题既有观感；\proofColonBack 是 cls 内定义的具名回缩长度。
% 若将来决定统一为半角冒号，改成 \renewcommand{\proofSeparator}{: } 即可。
% ================================================================
\renewcommand{\proofSeparator}{：\hspace*{\proofColonBack}}

% 参考文献（biblatex + biber，numeric 样式，与源书 [n] 引用一致）
\addbibresource{references.bib}
% 忠实还原源书引用编号 [1]..[100]：覆盖 cls 默认的 sorting=nyt，
% 改用 sorting=none，并以 \nocite{*} 让全部 100 条按 .bib 文件顺序（= 源书顺序）编号。
\ExecuteBibliographyOptions{sorting=none}
\nocite{*}

% ================================================================
% 封面（覆盖 cls 的 \makecover，只调版式，不改动 cls 本体）
% cls 原实现有两个版式缺陷：
%   1) 用 \hspace{2cm} 做首行缩进 —— 标题自动换行后，第二行落回左边距，
%      出现“应用”“Applications”“ZHAO”悬出到左侧的现象；
%   2) 作者名套 \textsc —— XeLaTeX + Fandol 无真小型大写，被伪造成
%      “XiAO-QiNG JiN” 这种大小混杂的形态。
% 这里改用 \leftskip 让整段（含换行后的各行）统一缩进，并去掉 \textsc。
% ================================================================
\makeatletter
\renewcommand{\makecover}{%
  \begin{titlepage}
    \thispagestyle{empty}%
    \flushleft
    \rule{\textwidth}{1.2pt} \\[0.1cm]
    \rule{\textwidth}{0.8pt} \\[1.5cm]
    \begingroup
      \leftskip=2cm\relax
      {\fontsize{38pt}{48pt}\selectfont\headfont \BookTitleCN\par}%
    \endgroup
    \vspace{0.7cm}
    \begingroup
      \leftskip=2cm\relax
      \ifblank{\BookTitleEN}{}{%
        {\fontsize{22pt}{28pt}\selectfont\itshape \BookTitleEN\par}%
      }%
    \endgroup
    \vspace{1.2cm}
    \rule{\textwidth}{0.6pt} \\[2cm]
    \begingroup
      \leftskip=2cm\relax
      {\LARGE \OriginalAuthor\par}%
    \endgroup
    \vfill
    \noindent
    \begin{minipage}[b]{0.46\textwidth}
      \raggedright
      \MT@coverlogo
    \end{minipage}%
    \hfill
    \begin{minipage}[b]{0.50\textwidth}
      \raggedleft
      \gdef\MT@joinline{}%
      \MT@join{\OriginalEdition}\MT@join{\OriginalPublisher}\MT@join{\OriginalYear}%
      \ifx\MT@joinline\@empty\else{\small \MT@joinline} \\[0.25cm]\fi
      {\normalsize \Translator\ 翻译及重排} \\[0.25cm]
      \ifdefempty{\ModelUsed}{}{{\small Model: \ModelUsed} \\[0.25cm]}%
      {\small Updated: \ifdefempty{\TranslationDate}{\today}{\TranslationDate}}%
    \end{minipage}
    \vspace*{1cm}
  \end{titlepage}%
}
\makeatother

% ================================================================
% 后物料：参考文献与术语索引为「无编号」章
% v3.2 起 cls 的 \makebibliography / \printterminology 已自动 \backmatter
% 并以 \chapter* + 目录条目 + 页眉排版，故原先为去掉章号而做的两处
% \patchcmd（匹配 \chapter{参考文献} / \chapter{术语索引}）在 v3.3 下必然失配，
% 且会打印 MT-WARN 级警告，已删除；此处不再做任何干预。
% ================================================================

\begin{document}
\makecover

\frontmatter

% 前置材料各篇长短悬殊，末页若按 flushbottom 拉伸会出现 Underfull \vbox，
% 故前置部分改为顶齐，进入正文后恢复底齐。
\raggedbottom

% ---------- 前置材料 ----------
% front/translator-note.tex —— 译者说明
% 说明底本、翻译原则、术语与记号约定、排版与校勘处理。
% 注：每段必须写成一行（不得硬换行），否则 xeCJK 会把换行当成空格，句中留出空隙。
% 全文须控制在一页之内（约 40 行），超出一页会使第二页只剩寥寥数行。
\cleardoublepage
\phantomsection
\chapter*{译者说明}
\addcontentsline{toc}{chapter}{译者说明}
\markboth{译者说明}{译者说明}

\section*{一、关于底本}

本书据科学出版社 2015 年出版的 \textit{Numerical Linear Algebra and Its Applications} (Second Edition) 译出 (「信息与计算科学丛书」第 73 卷, ISBN 978-7-03-046425-5), 作者为金小庆, 魏益民, 赵志. 原书初版于 2004 年; 第二版订正了初版中的若干疏漏, 补入定理~\ref{thm:6-7} 的证明, 新增第 10 章 (空间分数阶扩散方程), 参考文献亦有扩充. 全书共 10 章, 依次讨论线性方程组的直接法与迭代法, Krylov 子空间方法, 非对称与对称特征值问题, 以及上述方法在常微分方程与分数阶扩散方程中的应用.

\section*{二、翻译原则}

\noindent\textbf{1.~忠实原文.} 译文以原书正文为准, 不增删命题, 公式, 例题与习题, 不改动定理的条件与结论; 章节划分与叙述次序均依原书. 原书未设脚注, 译文亦不加脚注.

\smallskip
\noindent\textbf{2.~术语.} 术语以国内通行译名为准, 如 Gaussian elimination 作「高斯消去法」; 首见时以「中文 (English)」附出英文原词, 并收入书末术语索引, 同一概念全书只用一种译名.

\smallskip
\noindent\textbf{3.~人名.} 数学家的人名 (Jacobi, Gauss, Householder, Krylov 等) 保留原文拼写, 不另行音译; 致谢中的人名照录其缩写形式.

\smallskip
\noindent\textbf{4.~文献.} 原书以 [1]--[100] 引用文献, 译文沿用原编号, 与书末「参考文献」一一对应.

\smallskip
\noindent\textbf{5.~编号.} 公式, 定理, 定义, 引理, 推论, 例一律按章编号 (如式 (3.1), 定理 3.1); 习题按章自 1 起连续编号, 均依原书体例.

\smallskip
\noindent\textbf{6.~记号.} 向量用粗正体 ($\mathbf{x}$), 矩阵用大写斜体 ($A$); 转置记作 $^{T}$, 共轭转置记作 $^{*}$; 不等式用 $\leqslant$, $\geqslant$; $\R$, $\C$ 为实数域与复数域; $\norm{\cdot}$ 为范数, 所指依上下文判定; 虚数单位记作 $\ii$.

\section*{三、排版与校勘}

译文以 \LaTeX{} 重排, 用 XeLaTeX 编译, 中文正文采用 Fandol 字体, 数学排版沿用 \texttt{amsmath}; 算法框, 示意图与全部公式均重新排印, 未使用扫描图像. 原书为影印本, 无文字层, 翻译时先以光学字符识别取得初稿, 再逐页比对页面图版核校, 凡定理陈述, 矩阵元素, 上下标与各类编号一律以图版为准; 个别疑似排印疏漏因无从确证作者原意, 译文一仍其旧.

\section*{四、关于译文}

译者力求译文准确, 术语统一, 行文可读, 惟学力有限, 疏漏恐所难免; 译文中的错误概由译者负责, 敬请读者指正. 原书的学术观点, 方法与结论均属原作者.

\vspace{0.8em}
\rightline{译者}
\rightline{2026 年 11 月}
% 译者说明（底本、翻译原则、记号与校勘说明）
% front/series-preface.tex —— 原书丛书序（Preface to the Series）
% 作者：Shi Zhongci，2005 年 7 月。
% 注：每段必须写成一行（不得硬换行），否则 xeCJK 会把换行当成空格。
\cleardoublepage
\phantomsection
\chapter*{《信息与计算科学丛书》序}
\addcontentsline{toc}{chapter}{《信息与计算科学丛书》序}
\markboth{《信息与计算科学丛书》序}{《信息与计算科学丛书》序}

20 世纪 70 年代以来, 科学出版社已陆续出版三十余种「计算方法丛书」. 这套丛书由已故中国科学院院士, 中国科学院计算中心主任冯康先生倡导并主持. 丛书及时反映了计算数学的前沿方向与最新研究成果, 对青年科学工作者和整个研究群体产生了很大影响, 对我国计算数学的发展起了重要作用.

为适应科学发展的新形势, 1998 年, 中华人民共和国教育部将计算数学, 数值算法, 信息科学, 运筹学与控制论等学科合并为「信息与计算科学」这一新学科. 科学出版社亦相应调整了丛书编委会, 并将丛书更名为「信息与计算科学丛书」. 第一届编委会会议于 2004 年 9 月在北京召开, 讨论了新丛书的目标, 方向与内容.

新丛书旨在向高年级本科生, 研究生和从事这些领域研究的科学工作者介绍信息与计算科学的现状与前沿, 因而将以具体而系统的方式阐述信息与计算科学的进展, 并兼及相关交叉学科的发展.

谨向历任编委会成员与编辑部同仁, 以及所有为「计算方法丛书」作出重要贡献的数学家致谢. 正是由于他们的努力, 这套丛书在学术界赢得了卓越声誉. 我真诚希望各位继续支持「信息与计算科学丛书」, 使新丛书同样能推动本世纪信息与计算科学的发展.

\vspace{1.5em}
\rightline{石钟慈}
\rightline{2005 年 7 月}
% 原书丛书序（Shi Zhongci, 2005.7）
% front/preface.tex —— 作者原序（Preface）
% 含：正文、引文、致谢、修订重印说明、第二版说明。原书末尾无署名。
% 注：每段必须写成一行（不得硬换行），否则 xeCJK 会把换行当成空格。
\cleardoublepage
\phantomsection
\chapter*{前言}
\addcontentsline{toc}{chapter}{前言}
\markboth{前言}{前言}

自 1946 年第一台现代计算机诞生以来, 数值线性代数 (又称矩阵计算) 一直是科学与工程计算的核心. 科学与工程中的大多数问题最终都归结为矩阵计算问题, 因此, 研究数值线性代数十分重要. 本书对矩阵计算作初步介绍, 同时也收录了近年来取得的一些新成果. 全书开篇, 我们首先在第 1 章给出数值线性代数的概貌.

第 2 章介绍求解一般线性方程组的基本直接法——高斯消去法. 高斯消去法通常用于求解规模中等, 没有特殊结构的稠密线性方程组, 其运算量为 $\mathcal{O}(n^{3})$, 其中 $n$ 为方程组的规模. 本章还讨论了主元技术.

第 3 章为讨论扰动与误差对数值解的影响, 引入向量范数与矩阵范数并研究其性质, 同时给出浮点运算与部分主元技术的误差分析.

第 4 章研究线性最小二乘问题, 重点讨论超定线性方程组 $A\mathbf{x}=\mathbf{b}$ ($A$ 的行数多于列数) 的最小二乘解的求法. 我们利用一些正交变换与 QR 分解来设计求解最小二乘问题的高效算法.

第 5 章研究求解 $A\mathbf{x}=\mathbf{b}$ 的经典迭代法. 迭代法与高斯消去法等直接法颇不相同. 当 $A$ 相当大时, 基于 $A$ 的 LU 分解的直接法在计算时间与计算机存储两方面都难以承受. 通常, 在大多数大型问题中矩阵是稀疏的; 但稀疏性可能在 LU 分解过程中丢失, 于是在 LU 分解结束时, 存储成为关键问题. 对这类问题, 可以使用一类称为迭代法的方法. 本章只讨论其中若干经典迭代法.

第 6 章介绍另一类迭代法, 即 Krylov 子空间方法. 我们只研究这类方法中的两个版本: 共轭梯度 (CG) 法与广义极小残量 (GMRES) 法. 1952 年提出的 CG 法是求解对称正定线性方程组最著名的迭代法之一; GMRES 方法则提出于 1986 年, 用于求解非对称线性方程组. 本章还简要讨论了预处理技术.

特征值问题在科学计算中格外引人关注. 第 7 章研究非对称特征值问题, 介绍幂法, 反幂法与 QR 方法等若干著名方法.

对称特征值问题具有良好的性质与丰富的数学理论, 是数值线性代数中最有趣的课题之一. 第 8 章研究这一课题, 讨论对称 QR 迭代法, Jacobi 法, 对分法以及分而治之技术.

第 9 章简要研究用边值方法求解带初值的常微分方程组 (ODEs) 的若干近期进展. 这类方法需要求解一个或多个非对称, 大型且稀疏的线性方程组, 因此我们用第 6 章的 GMRES 方法并配以适当的预条件子来求解它们. 主要结果之一是: 若对 $m$ 阶常微分方程组采用 $A_{\nu_1,\nu_2}$-稳定的边值方法, 则预处理后的矩阵可分解为 $I+L$, 其中 $I$ 为单位矩阵, $L$ 的秩至多为 $2m(\nu_1+\nu_2)$. 由此可知, 对预处理后的方程组施用 GMRES 方法, 至多经过 $2m(\nu_1+\nu_2)+1$ 次迭代即收敛. 文中还给出了在若干时滞微分方程上的应用.

第 10 章研究用于求解二维空间分数阶扩散方程的 Crank--Nicolson 加权与移位 Grünwald 差分 (CN-WSGD) 格式. 该格式需要求解一个或多个非对称, 大型且满的线性方程组. 我们使用第 6 章介绍的 GMRES 方法, 并结合重启策略与预处理技术来求解这些方程组. 主要结果之一是: 对预处理后的方程组施用重启 GMRES 方法时, 每个时间层的计算量为 $\mathcal{O}(N^{2}\log N)$, 其中 $N=N_{1}=N_{2}$ 为 $x$ 方向与 $y$ 方向的空间剖分数.

\vspace{0.8em}
\begin{quote}
「如果还有哪个数学分支之于数学科学的基础性, 像微积分与微分方程一样, 那便是数值线性代数. 」

\raggedleft ——\,L. Trefethen, D. Bau III
\end{quote}
\vspace{0.8em}

\noindent\textbf{致谢:} 我们要感谢香港中文大学数学系 Raymond H. F. Chan 教授始终如一的鼓励, 多年的友谊与经费支持; 感谢复旦大学数学科学学院 Z. H. Cao 教授的许多有益讨论与宝贵建议. 我们还要感谢朋友 Z. C. Shi 教授的鼓励与宝贵意见. 当然, 特别的谢意要献给作者生命中两个重要的机构——澳门大学与复旦大学, 它们为本书的写作提供了极好的学术氛围. 本书的大部分写作是在夜晚, 周末与假日完成的. 最后, 还要感谢我们的家人, 他们的关爱, 理解, 鼓励与支持是本书得以完成的必要条件. 谨向他们致以最诚挚的谢意!

\begingroup\raggedright
本书的出版部分得到澳门大学研究基金 (项目号 \mbox{RG024/01-02S/JXQ/FST}, \mbox{RG031/02-03S/JXQ/FST}, \mbox{RG064/03-04S/JXQ/FST}) 与国家自然科学基金项目 10471027 的资助; 上海市教育委员会与复旦大学亦提供了部分经费支持.
\par\endgroup

\medskip
\noindent\textbf{作者关于修订重印的说明:} 在本书的第二次印刷中, 我们订正了第一次印刷中的一些数学上与排印上的小错误. 谨向所有指出这些错误的读者致谢. 第 7 章作了若干补充与修订, 参考文献亦有更新, 书中还增加了一些习题. 本书的第二次印刷得到澳门大学研究基金 \mbox{RG081/04-05S/JXQ/FST} 的资助.

\medskip
\noindent\textbf{作者关于第二版的说明:} 自 2004 年本书问世至今已十余年, 我们收到了许多人的反馈. 同事, 学生与其他读者指出了书中一些数学上与排印上的小错误, 并对如何改进本书提出了建议. 我们尤其感谢上海海事大学数学系 Z. L. Xu 教授, 他提出了许多富有见地的意见与建设性的评论. 本版补入了定理~\ref{thm:6-7} 的证明, 并新增一章 (第 10 章). 第 10 章讨论数值分数阶扩散方程的近期进展. 参考文献有所扩充与更新. 本书的修订得到澳门大学研究基金 \mbox{MYRG098(Y3-L3)-FST13-JXQ}, 国家自然科学基金项目 11271084 以及复旦大学的资助.
% 作者前言（含致谢与两版说明）

\makecontents

\mainmatter
\flushbottom

% chapters/ch01.tex —— 第 1 章 绪论
% 内容权威：源 PDF 物理页 15--22（书页 1--8）。
% 源书本章公式仅 (1.1) 一个编号；定理 1.1；章末习题 15 道。
\chapter{绪论}
\label{ch:introduction}

数值线性代数 (NLA) 又称\newterm{matrixcomputation}{矩阵计算}{matrix computation}. 自 1946 年前后第一台现代计算机问世以来, 它一直是科学与工程计算的核心. 科学与工程中的绝大多数问题, 最终都会转化为数值线性代数问题. 因此, 研究数值线性代数对我们而言极为重要. 本书对数值线性代数作初等介绍, 同时也纳入了近年来取得的一些新成果.

\section{基本记号}
\label{sec:ch1-basicsymbols}

本书通篇采用下列记号.

\begin{itemize}
  \item 令 $\R$ 表示实数集, $\C$ 表示复数集, 且 $\ii \equiv \sqrt{-1}$.

  \item 令 $\R^n$ 表示实 $n$ 维向量的集合, $\C^n$ 表示复 $n$ 维向量的集合. 向量几乎总取列向量.

  \item 令 $\R^{m\times n}$ 表示 $m\times n$ 实\newterm{matrix}{矩阵}{matrix}的线性向量空间, $\C^{m\times n}$ 表示 $m\times n$ 复矩阵的线性向量空间.

  \item 用大写字母如 $A,B,C,\Delta,\Lambda$ 等表示矩阵, 用粗体小写字母如 $\mathbf{x},\mathbf{y},\mathbf{z}$ 等表示向量; 小写字母如 $\alpha,\beta,\lambda$ 等表示标量.

  \item 符号 $a_{ij}$ 表示矩阵 $A$ 的第 $ij$ 个\newterm{entry}{元素}{entry}.

  \item 符号 $A^T$ 表示矩阵 $A$ 的\newterm{transpose}{转置}{transpose}, $A^*$ 表示其\newterm{conjtranspose}{共轭转置}{conjugate transpose}.

  \item 设 $\mathbf{a}_1,\ldots,\mathbf{a}_m\in\R^n$ (或 $\C^n$). 用 $\spn\{\mathbf{a}_1,\ldots,\mathbf{a}_m\}$ 表示由 $\mathbf{a}_1,\ldots,\mathbf{a}_m$ 的所有线性组合构成的线性向量空间.

  \item 令 $\rank(A)$ 表示矩阵 $A$ 的\newterm{rank}{秩}{rank}.

  \item 令 $\dim(S)$ 表示向量空间 $S$ 的维数.

  \item 用 $\det(A)$ 表示矩阵 $A$ 的\newterm{determinant}{行列式}{determinant}; 用 $\diag(a_{11},\cdots,a_{nn})$ 表示如下的 $n\times n$\newterm{diagmatrix}{对角矩阵}{diagonal matrix}:
  \[
    \diag(a_{11},\dots,a_{nn}) = \begin{bmatrix}
      a_{11} & 0      & \cdots & 0      \\
      0      & a_{22} & \ddots & \vdots \\
      \vdots & \ddots & \ddots & 0      \\
      0      & \cdots & 0      & a_{nn}
    \end{bmatrix}.
  \]

  \item 对矩阵 $A=[a_{ij}]$, 符号 $|A|$ 表示元素为 $(|A|)_{ij}=|a_{ij}|$ 的矩阵.

  \item 符号 $I$ 表示\newterm{identity}{单位矩阵}{identity matrix}, 即
  \[
    I = \begin{bmatrix}
      1      & 0      & \cdots & 0      \\
      0      & 1      & \ddots & \vdots \\
      \vdots & \ddots & \ddots & 0      \\
      0      & \cdots & 0      & 1
    \end{bmatrix},
  \]
  而 $\mathbf{e}_i$ 表示第 $i$ 个单位向量, 即 $I$ 的第 $i$ 列.

  \item 符号 $\mathbf{0}$ 表示尺寸适当的\newterm{zeromatrix}{零矩阵}{zero matrix}或零向量.

  \item 用 $\norm{\cdot}$ 表示矩阵或向量的\newterm{norm}{范数}{norm}; 符号 $\norm{\cdot}_1,\norm{\cdot}_2,\norm{\cdot}_\infty$ 分别表示 $p=1,2,\infty$ 时的 $p$ 范数.

  \item 仿照 MATLAB, 在算法中 $A(i,j)$ 表示矩阵 $A$ 的第 $(i,j)$ 个元素; $A(i,:)$ 与 $A(:,j)$ 分别表示 $A$ 的第 $i$ 行与第 $j$ 列; $A(i_1:i_2,k)$ 表示由 $A$ 的第 $k$ 列中从第 $i_1$ 个到第 $i_2$ 个元素构成的列向量; $A(k,j_1:j_2)$ 表示由 $A$ 的第 $k$ 行中从第 $j_1$ 个到第 $j_2$ 个元素构成的行向量; $A(k:l,p:q)$ 表示由 $A$ 的第 $k$ 行至第 $l$ 行, 第 $p$ 列至第 $q$ 列所构成的 $(l-k+1)\times(q-p+1)$ 子矩阵.
\end{itemize}

\section{数值线性代数中的基本问题}
\label{sec:ch1-basicproblems}

数值线性代数包含下列三个主要问题, 它们正是本书的研究对象.

\begin{parenlist}
  \item 求线性方程组
    \[
      A\mathbf{x}=\mathbf{b},
    \]
    的解, 其中 $A$ 为 $n\times n$ 非奇异矩阵, $\mathbf{b}$ 为 $n$ 维向量.

  \item 线性最小二乘问题: 对任意 $m\times n$ 矩阵 $A$ 与 $m$ 维向量 $\mathbf{b}$, 求一个 $n$ 维向量 $\mathbf{x}$, 使得
    \[
      \norm{A\mathbf{x}-\mathbf{b}}_2 = \min_{\mathbf{y}\in\R^n}\norm{A\mathbf{y}-\mathbf{b}}_2.
    \]

  \item 特征值问题: 对任意 $n\times n$ 矩阵 $A$, 求其部分 (或全部) 特征值及相应的特征向量. 我们在此指出, 若存在非零向量 $\mathbf{x}\in\C^n$ 使得
    \[
      A\mathbf{x}=\lambda\mathbf{x},
    \]
    则称复数 $\lambda$ 为 $A$ 的\newterm{eigenvalue}{特征值}{eigenvalue}, 其中 $\mathbf{x}$ 称为 $A$ 的属于 $\lambda$ 的\newterm{eigenvector}{特征向量}{eigenvector}.
\end{parenlist}

除这些主要问题外, 数值线性代数中还有许多其他基本问题, 例如总体最小二乘问题, 矩阵方程, 广义逆, 逆特征值问题, 非线性特征值问题, 奇异值问题以及逆奇异值问题等. 关于这些问题的文献, 我们参见 \cite{src:3,src:6,src:13,src:26,src:27,src:53,src:72,src:73,src:76,src:84,src:86,src:88,src:100}.

\section{为何要研究数值方法?}
\label{sec:ch1-whynumerical}

为回答这一问题, 让我们考虑如下线性方程组
\[
  A\mathbf{x}=\mathbf{b},
\]
其中 $A$ 为 $n\times n$ 非奇异矩阵, $\mathbf{x}=[x_1,x_2,\ldots,x_n]^T$. 若使用著名的 Cramer 法则, 则可得如下解:
\[
  x_1=\frac{\det(A_1)}{\det(A)},\quad
  x_2=\frac{\det(A_2)}{\det(A)},\quad
  \ldots,\quad
  x_n=\frac{\det(A_n)}{\det(A)},
\]
其中 $A_i$ ($i=1,2,\ldots,n$) 是把第 $i$ 列换成向量 $\mathbf{b}$ 后所得的矩阵. 于是需要计算 $n+1$ 个行列式, 即 $\det(A_i)$ ($i=1,2,\ldots,n$) 与 $\det(A)$. 这共需
\[
  [n!(n-1)](n+1)=(n-1)(n+1)!
\]
次乘法. 当 $n=25$ 时, 若用一台每秒运算 $10^{10}$ 次的计算机, 则需要
\[
  \frac{24\times 26!}{10^{10}\times 3600\times 24\times 365}\approx 30.6\times 10^{9}\ \text{年}.
\]
而若采用高斯消去法, 则只需
\[
  \sum_{i=1}^{n}(i-1)(i+1)=\sum_{i=1}^{n}i^{2}-n=\frac{1}{6}n(n+1)(2n+1)-n=\mathcal{O}(n^{3})
\]
次乘法. 于是用同一台计算机, 不到 1 秒即可解出一个 $25\times 25$ 的线性方程组. 由上述讨论可见, 用不同的数值方法求解同一个问题, 效果差别竟如此之大. 因此, 研究数值方法的性质对我们而言是必不可少的.

\section{矩阵分解}
\label{sec:ch1-factorization}

对任意线性方程组 $A\mathbf{x}=\mathbf{b}$, 若能把 $A$ 分解为 $A=LU$, 其中 $L$ 为下三角矩阵, $U$ 为上三角矩阵, 则有
\begin{equation}
  \begin{cases}
    L\mathbf{y}=\mathbf{b},\\
    U\mathbf{x}=\mathbf{y}.
  \end{cases}
  \label{eq:ch1-lu}
\end{equation}
通过回代, 容易解出 \eqref{eq:ch1-lu}, 进而解出 $A\mathbf{x}=\mathbf{b}$. 因此, 矩阵分解是数值线性代数中非常重要的工具. 下面这个定理是线性代数中最基本, 也最有用的结果之一 \cite{src:39}.

\begin{theorem}[Jordan 分解定理]
\label{thm:ch1-jordan}
设 $A\in\C^{n\times n}$, 则存在非奇异矩阵 $X\in\C^{n\times n}$ 使得
\[
  X^{-1}AX=J\equiv\diag(J_1,J_2,\ldots,J_p),
\]
亦即 $A=XJX^{-1}$, 其中 $J$ 称为 $A$ 的\newterm{jordanform}{Jordan 标准形}{Jordan canonical form}, 而
\[
  J_i=\begin{bmatrix}
    \lambda_i & 1         & 0      & \cdots & 0      \\
    0         & \lambda_i & 1      & \ddots & \vdots \\
    \vdots    & 0         & \ddots & \ddots & 0      \\
    \vdots    &           & \ddots & \ddots & 1      \\
    0         & \cdots    & \cdots & 0      & \lambda_i
  \end{bmatrix}\in\C^{n_i\times n_i},\qquad i=1,2,\ldots,p,
\]
称为\newterm{jordanblock}{Jordan 块}{Jordan block}, 且 $n_1+\cdots+n_p=n$. 除对角 Jordan 块的排列次序外, $A$ 的 Jordan 标准形是唯一的. 若 $A\in\R^{n\times n}$ 且仅有实特征值, 则矩阵 $X$ 可取为实矩阵.
\end{theorem}

\section{扰动与误差分析}
\label{sec:ch1-perturbation}

数值算法给出的解很少是绝对正确的. 误差通常有两类: 其一, 输入数据本身带有误差, 它由先前的计算或测量产生; 其二, 算法内部的近似会引入误差. 因此, 我们需要进行扰动与误差分析 \cite{src:12,src:69,src:96}.

\medskip
\noindent (1) 扰动.

给定 $x$, 欲计算函数 $f(x)$ 的值. 设 $x$ 有一个扰动 $\delta x$, 且 $|\delta x|/|x|$ 非常小. 我们希望找到尽可能小的正数 $c(x)$, 使得
\[
  \frac{|f(x+\delta x)-f(x)|}{|f(x)|}\leqslant c(x)\frac{|\delta x|}{|x|}.
\]
此时称 $c(x)$ 为 $f(x)$ 在 $x$ 处的\newterm{condnum}{条件数}{condition number}. 若 $c(x)$ 很大, 则称函数 $f$ 在 $x$ 处是病态的; 若 $c(x)$ 很小, 则称 $f$ 在 $x$ 处是良态的.

\begin{mremark}
一个计算问题是否病态, 与我们所采用的数值方法无关.
\end{mremark}

\medskip
\noindent (2) 误差.

用某种数值方法计算函数 $f$ 在点 $x$ 处的值, 得到 $\hat{y}$. 由于舍入误差 (或截断误差), 通常有
\[
  \hat{y}\neq f(x).
\]
若存在 $\delta x$ 使得
\[
  \hat{y}=f(x+\delta x),\qquad |\delta x|\leqslant\epsilon|x|,
\]
其中 $\epsilon$ 是一个与所用数值方法及计算机密切相关的正常数, 则当 $\epsilon$ 很小时, 称该方法是\newterm{stable}{稳定的}{stable}; 当 $\epsilon$ 很大时, 称该方法是不稳定的.

\begin{mremark}
一个数值方法是否稳定, 与我们所面临的计算问题无关.
\end{mremark}

综合扰动与误差分析, 可得
\[
  \frac{|\hat{y}-f(x)|}{|f(x)|}
  =\frac{|f(x+\delta x)-f(x)|}{|f(x)|}
  \leqslant c(x)\frac{|\delta x|}{|x|}
  \leqslant \epsilon c(x).
\]
因此, 一个数值结果是否准确, 既取决于数值方法的稳定性, 也取决于计算问题的条件数.

\section{运算量与收敛速度}
\label{sec:ch1-cost}

数值算法通常分为两类:
\[
  \begin{cases}
    \text{(i) 直接法;}\\
    \text{(ii) 迭代法.}
  \end{cases}
\]
用直接法, 可在精确算术下经有限步得到计算问题的精确解; 用迭代法, 经有限步只能得到计算问题的近似解.

运算量是衡量算法的重要指标. 一个算法的运算量, 是指算法中所用的「加, 减, 乘, 除」运算的总次数. 需要指出, 算法的速度只部分地取决于运算量. 在现代计算机中, 运算的速度远快于数据传输的速度. 因此, 有时算法的速度主要取决于数据传输的总量.

对直接法, 我们通常以运算量作为衡量算法速度的主要指标; 对迭代法, 则需要考虑
\[
  \begin{cases}
    \text{(i) 每次迭代的运算量;}\\
    \text{(ii) 方法的收敛速度.}
  \end{cases}
\]

对迭代算法所产生的序列 $\{\mathbf{x}_k\}$, 若 $\{\mathbf{x}_k\}\to\mathbf{x}$ (精确解), 且满足
\[
  \norm{\mathbf{x}_k-\mathbf{x}}\leqslant c\norm{\mathbf{x}_{k-1}-\mathbf{x}},\quad k=1,2,\ldots,
\]
其中 $0<c<1$, $\norm{\cdot}$ 为任一向量范数 (详见\autoref{ch:perturbation}), 则称其收敛速度是线性的. 若满足
\[
  \norm{\mathbf{x}_k-\mathbf{x}}\leqslant c\norm{\mathbf{x}_{k-1}-\mathbf{x}}^{p},\quad k=1,2,\ldots,
\]
其中 $0<c<1$, $p>1$, 则称其收敛速度是超线性的.

\begin{exercises}

\item 若上三角矩阵的对角元全为零, 则称之为严格上三角矩阵. 证明: 若 $A$ 为 $n$ 阶严格上三角矩阵, 则 $A^{n}=\mathbf{0}$.

\item 设 $A\in\C^{n\times m}$, $B\in\C^{m\times l}$. 证明
  \[
    \rank(AB)\geqslant\rank(A)+\rank(B)-m.
  \]

\item 设
  \[
    A=\begin{bmatrix} A_{11} & A_{12}\\ A_{21} & A_{22}\end{bmatrix},
  \]
  其中 $A_{ij}$ ($i,j=1,2$) 均为方阵, $\det(A_{11})\neq 0$, 且满足
  \[
    A_{11}A_{21}=A_{21}A_{11}.
  \]
  则
  \[
    \det(A)=\det(A_{11}A_{22}-A_{21}A_{12}).
  \]

\item 证明 $\det(I-\mathbf{u}\mathbf{v}^{*})=1-\mathbf{v}^{*}\mathbf{u}$, 其中 $\mathbf{u},\mathbf{v}\in\C^{m}$ 为列向量.

\item 证明 Hadamard 不等式: 对 $A\in\C^{n\times n}$ 有
  \[
    |\det(A)|\leqslant\prod_{j=1}^{n}\norm{\mathbf{a}_j}_2,
  \]
  其中 $\mathbf{a}_j=A(:,j)$, $\norm{\mathbf{a}_j}_2=\left(\sum_{i=1}^{n}|A(i,j)|^{2}\right)^{\frac{1}{2}}$. 等号何时成立?

\item 设 $B$ 为幂零矩阵, 即存在整数 $k>0$ 使得 $B^{k}=\mathbf{0}$. 证明: 若 $AB=BA$, 则
  \[
    \det(A+B)=\det(A).
  \]

\item 设 $A$ 为 $m\times n$ 矩阵, $B$ 为 $n\times m$ 矩阵. 证明矩阵
  \[
    \begin{bmatrix} AB & \mathbf{0}\\ B & \mathbf{0}\end{bmatrix}
    \quad\text{与}\quad
    \begin{bmatrix} \mathbf{0} & \mathbf{0}\\ B & BA\end{bmatrix}
  \]
  相似. 由此推出: $AB$ 的非零特征值与 $BA$ 的非零特征值相同, 且
  \[
    \det(I_m+AB)=\det(I_n+BA).
  \]

\item 若矩阵 $M\in\C^{n\times n}$ 满足
  \[
    M=M^{*},\qquad \mathbf{x}^{*}M\mathbf{x}>0,
  \]
  对一切 $\mathbf{x}\neq\mathbf{0}\in\C^{n}$ 成立, 则称 $M$ 为 Hermite 正定矩阵. 设 $A$ 与 $B$ 均为 Hermite 正定矩阵.
  \subquest 证明矩阵乘积 $AB$ 的特征值全为正.
  \subquest 证明 $AB$ 为 Hermite 矩阵当且仅当 $A$ 与 $B$ 可交换.

\item 证明任意矩阵 $A\in\C^{n\times n}$ 都可唯一地写成如下形式
  \[
    A=B+\ii C,
  \]
  其中 $B$ 与 $C$ 均为 Hermite 矩阵.

\item 证明: 若 $A$ 为反 Hermite 矩阵, 即 $A^{*}=-A$, 则 $A$ 的全部特征值都落在虚轴上.

\item 设
  \[
    A=\begin{bmatrix} A_{11} & A_{12}\\ A_{21} & A_{22}\end{bmatrix}.
  \]
  假设 $A_{11}$, $A_{22}$ 为方阵, 且 $A_{11}$, $A_{22}-A_{21}A_{11}^{-1}A_{12}$ 均非奇异. 设
  \[
    B=\begin{bmatrix} B_{11} & B_{12}\\ B_{21} & B_{22}\end{bmatrix}
  \]
  为 $A$ 的逆. 证明
  \begin{align*}
    B_{22}&=(A_{22}-A_{21}A_{11}^{-1}A_{12})^{-1}, &
    B_{12}&=-A_{11}^{-1}A_{12}B_{22},\\
    B_{21}&=-B_{22}A_{21}A_{11}^{-1}, &
    B_{11}&=A_{11}^{-1}-B_{12}A_{21}A_{11}^{-1}.
  \end{align*}

\item 设 $A$ 与 $B$ 均为 Hermite 矩阵, 且 $A$ 正定. 证明: $A+B$ 正定当且仅当 $A^{-1}B$ 的全部特征值都大于 $-1$.

\item 设 $A$ 为幂等矩阵, 即 $A^{2}=A$. 证明 $A$ 的每个特征值或为 $0$, 或为 $1$.

\item 设 $A$ 为所有元素都等于 $1$ 的矩阵. 证明 $A$ 可写成 $A=\mathbf{e}\mathbf{e}^{T}$, 其中 $\mathbf{e}^{T}=[1,1,\ldots,1]$, 且 $A$ 为半正定矩阵. 求 $A$ 的特征值与特征向量.

\item 证明任意矩阵 $A\in\C^{n\times n}$ 都有极分解 $A=HQ$, 其中 $H$ 为 Hermite 半正定矩阵, $Q$ 为酉矩阵. 回忆: 若 $M^{-1}=M^{*}$, 则称 $M\in\C^{n\times n}$ 为酉矩阵. 此外, 若 $A$ 非奇异, 则 $H$ 为 Hermite 正定矩阵, 且 $A$ 的极分解是唯一的.

\end{exercises}

% chapters/ch02.tex —— 第 2 章：高斯消去法
% 源书：Chapter 2 Gaussian Elimination，书页 9–22（物理页 23–36）
% 公式编号 (2.1)–(2.6)；定理 2.1–2.4；定义 2.1；引理 2.1；算法 2.1–2.4；习题 1–11
\chapter{高斯消去法}
\label{ch:gaussian}

线性方程组的求解问题在 NLA 中处于核心地位. 求解线性方程组一般有两类方法: 一类称为\newterm{directmethod}{直接法}{direct method}, 另一类称为\newterm{iterativemethod}{迭代法}{iterative method}. 使用直接法, 可在有限步内得到计算问题在精确算术下的精确解; 使用迭代法, 在有限步内只能得到计算问题的近似解.

本章介绍求解一般线性方程组的一种基本直接法——\newterm{gaussianelimination}{高斯消去法}{Gaussian elimination}. 通常, 高斯消去法用于求解规模中等, 无特殊结构的稠密线性方程组.

\section{三角线性方程组与 $LU$ 分解}

我们先研究三角线性方程组.

\subsection{三角线性方程组}
\label{subsec:2-1-1}

考虑如下非奇异下三角线性方程组
\begin{equation}
  L\mathbf{y}=\mathbf{b},
  \label{eq:2-1}
\end{equation}
其中 $\mathbf{b}=[b_1,b_2,\cdots,b_n]^T\in\R^n$ 为已知向量, $\mathbf{y}=[y_1,y_2,\cdots,y_n]^T$ 为未知向量, 而 $L=[l_{ij}]\in\R^{n\times n}$ 由下式给出
\[
  L=
  \begin{bmatrix}
    l_{11} & 0      & \cdots & \cdots & 0 \\
    l_{21} & l_{22} & 0      & \cdots & \vdots \\
    l_{31} & l_{32} & l_{33} & \ddots & \vdots \\
    \vdots & \vdots & \vdots & \ddots & 0 \\
    l_{n1} & l_{n2} & l_{n3} & \cdots & l_{nn}
  \end{bmatrix}
\]
且 $l_{ii}\neq0\ (i=1,2,\cdots,n)$. 由 \eqref{eq:2-1} 的第一个方程, 有
\[
  l_{11}y_1=b_1,
\]
从而
\[
  y_1=\frac{b_1}{l_{11}}.
\]
类似地, 由 \eqref{eq:2-1} 的第二个方程, 有
\[
  y_2=\frac{1}{l_{22}}\bigl(b_2-l_{21}y_1\bigr).
\]
一般地, 若已求得 $y_1,y_2,\cdots,y_{i-1}$, 则由 \eqref{eq:2-1} 的第 $i$ 个方程, 有
\[
  y_i=\frac{1}{l_{ii}}\left(b_i-\sum_{j=1}^{i-1}l_{ij}y_j\right).
\]
该算法称为\newterm{forwardsubstitution}{前代法}{forward substitution}, 需要 $\mathcal{O}(n^2)$ 次运算.

现在考虑如下非奇异上三角线性方程组
\begin{equation}
  U\mathbf{x}=\mathbf{y},
  \label{eq:2-2}
\end{equation}
其中 $\mathbf{x}=[x_1,x_2,\ldots,x_n]^T$ 为未知向量, 而 $U\in\R^{n\times n}$ 由下式给出
\[
  U=
  \begin{bmatrix}
    u_{11} & u_{12} & u_{13} & \cdots & u_{1n} \\
    0      & u_{22} & u_{23} & \cdots & \vdots \\
    0      & 0      & u_{33} & \ddots & \vdots \\
    \vdots & \vdots & \ddots & \ddots & \vdots \\
    0      & \cdots & \cdots & 0      & u_{nn}
  \end{bmatrix}
\]
且 $u_{ii}\neq0\ (i=1,2,\ldots,n)$. 从 \eqref{eq:2-2} 的最后一个方程出发, 可逐步求得 $x_n,x_{n-1},\ldots,x_1$. 其中 $x_n=y_n/u_{nn}$, 而 $x_i$ 由下式给出
\[
  x_i=\frac{1}{u_{ii}}\left(y_i-\sum_{j=i+1}^{n}u_{ij}x_j\right)
\]
对 $i=n-1,\ldots,1$ 计算. 该算法称为\newterm{backwardsubstitution}{回代法}{backward substitution}, 同样需要 $\mathcal{O}(n^2)$ 次运算.

对于一般线性方程组
\begin{equation}
  A\mathbf{x}=\mathbf{b},
  \label{eq:2-3}
\end{equation}
其中 $A\in\R^{n\times n}$ 与 $\mathbf{b}\in\R^n$ 为已知. 若能把矩阵 $A$ 分解为
\[
  A=LU,
\]
其中 $L$ 为下三角矩阵, $U$ 为上三角矩阵, 则可通过以下两个步骤求出 \eqref{eq:2-3} 的解:

\begin{parenlist}
  \item 用前代法求出 $L\mathbf{y}=\mathbf{b}$ 的解 $\mathbf{y}$.
  \item 用回代法求出 $U\mathbf{x}=\mathbf{y}$ 的解 $\mathbf{x}$.
\end{parenlist}

现在面临的问题是如何把矩阵 $A$ 分解为 $A=LU$. 为此引入高斯变换矩阵.

\subsection{高斯变换矩阵}

令
\[
  L_k=I-\mathbf{l}_k\mathbf{e}_k^T,
\]
其中 $I\in\R^{n\times n}$ 为单位矩阵, $\mathbf{l}_k=[0,\ldots,0,l_{k+1,k},\ldots,l_{nk}]^T\in\R^n$, 而 $\mathbf{e}_k\in\R^n$ 为第 $k$ 个单位向量. 则对任意 $k$,
\[
  L_k=
  \begin{bmatrix}
    1      & 0      & \cdots & \cdots & \cdots & 0 \\
    0      & \ddots &        &        &        & \vdots \\
    \vdots &        & 1      &        &        & \vdots \\
    \vdots &        & -l_{k+1,k} & 1  &        & \vdots \\
    \vdots &        & \vdots &        & \ddots & 0 \\
    0      & \cdots & -l_{nk}& \cdots & 0      & 1
  \end{bmatrix}
\]
称为\newterm{gaussiantransform}{高斯变换矩阵}{Gaussian transform matrix}. 这样的矩阵是单位下三角矩阵. 这里指出, \newterm{unittriangular}{单位三角矩阵}{unit triangular matrix}指对角元全为 $1$ 的三角矩阵. 对任意给定向量
\[
  \mathbf{x}=[x_1,x_2,\ldots,x_n]^T\in\R^n,
\]
有
\[
  \begin{aligned}
    L_k\mathbf{x}
    &=[x_1,\ldots,x_k,x_{k+1}-x_kl_{k+1,k},\ldots,x_n-x_kl_{nk}]^T\\
    &=[x_1,\ldots,x_k,0,\ldots,0]^T
  \end{aligned}
\]
只要取
\[
  l_{ik}=\frac{x_i}{x_k},\qquad i=k+1,\ldots,n
\]
且 $x_k\neq0$. 注意 $\mathbf{e}_k^T\mathbf{l}_k=0$, 容易验证
\[
  L_k^{-1}=I+\mathbf{l}_k\mathbf{e}_k^T.
\]

对给定矩阵 $A\in\R^{n\times n}$, 有
\[
  L_kA=(I-\mathbf{l}_k\mathbf{e}_k^T)A=A-\mathbf{l}_k(\mathbf{e}_k^TA)
\]
以及
\[
  \rank\bigl(\mathbf{l}_k(\mathbf{e}_k^TA)\bigr)=1.
\]
因此, $L_kA$ 是矩阵 $A$ 的\newterm{rankonemodification}{秩一修正}{rank-one modification}.

\subsection{$LU$ 分解的计算}

考虑下面这个简单例子. 令
\[
  A=
  \begin{bmatrix}
    1 & 5 & 9 \\
    2 & 4 & 7 \\
    3 & 3 & 10
  \end{bmatrix}.
\]
用高斯变换矩阵
\[
  L_1=
  \begin{bmatrix}
    1  & 0 & 0 \\
    -2 & 1 & 0 \\
    -3 & 0 & 1
  \end{bmatrix},
\]
得到
\[
  L_1A=
  \begin{bmatrix}
    1 & 5  & 9  \\
    0 & -6 & -11\\
    0 & -12& -17
  \end{bmatrix}.
\]
接着用高斯变换矩阵
\[
  L_2=
  \begin{bmatrix}
    1 & 0  & 0 \\
    0 & 1  & 0 \\
    0 & -2 & 1
  \end{bmatrix},
\]
得到
\[
  L_2(L_1A)\equiv U=
  \begin{bmatrix}
    1 & 5  & 9  \\
    0 & -6 & -11\\
    0 & 0  & 5
  \end{bmatrix}.
\]
因而最终有
\[
  A=LU,
\]
其中
\[
  L\equiv(L_2L_1)^{-1}=L_1^{-1}L_2^{-1}=
  \begin{bmatrix}
    1 & 0 & 0 \\
    2 & 1 & 0 \\
    3 & 2 & 1
  \end{bmatrix}.
\]

对一般的 $n\times n$ 矩阵 $A$, 可用 $n-1$ 个高斯变换矩阵 $L_1,L_2,\ldots,L_{n-1}$, 使 $L_{n-1}\cdots L_1A$ 为上三角矩阵. 事实上, 令 $A^{(0)}\equiv A$, 并设已求得 $k-1$ 个高斯变换矩阵 $L_1,\ldots,L_{k-1}\in\R^{n\times n}$, 使得
\[
  A^{(k-1)}=L_{k-1}\cdots L_1A=
  \begin{bmatrix}
    A_{11}^{(k-1)} & A_{12}^{(k-1)}\\
    \mathbf{0}     & A_{22}^{(k-1)}
  \end{bmatrix},
\]
其中 $A_{11}^{(k-1)}$ 为 $(k-1)\times(k-1)$ 上三角矩阵, 且
\[
  A_{22}^{(k-1)}=
  \begin{bmatrix}
    a_{kk}^{(k-1)} & \cdots & a_{kn}^{(k-1)}\\
    \vdots         & \ddots & \vdots\\
    a_{nk}^{(k-1)} & \cdots & a_{nn}^{(k-1)}
  \end{bmatrix}.
\]
若 $a_{kk}^{(k-1)}\neq0$, 则可取高斯变换矩阵
\[
  L_k=I-\mathbf{l}_k\mathbf{e}_k^T,
\]
其中
\[
  \mathbf{l}_k=[0,\ldots,0,l_{k+1,k},\ldots,l_{nk}]^T,
\]
而
\[
  l_{ik}=\frac{a_{ik}^{(k-1)}}{a_{kk}^{(k-1)}},\qquad i=k+1,\ldots,n,
\]
使 $L_kA^{(k-1)}$ 第 $k$ 列的最后 $n-k$ 个元素变为零. 于是有
\[
  A^{(k)}\equiv L_kA^{(k-1)}=
  \begin{bmatrix}
    A_{11}^{(k)} & A_{12}^{(k)}\\
    \mathbf{0}   & A_{22}^{(k)}
  \end{bmatrix},
\]
其中 $A_{11}^{(k)}$ 为 $k\times k$ 上三角矩阵. 经过 $n-1$ 步后, 得到所需的 $A^{(n-1)}$, 它是上三角矩阵. 令
\[
  L=(L_{n-1}\cdots L_1)^{-1},\qquad U=A^{(n-1)},
\]
则有 $A=LU$. 现在说明 $L$ 是单位下三角矩阵. 注意当 $j<i$ 时 $\mathbf{e}_j^T\mathbf{l}_i=0$, 故
\[
  \begin{aligned}
    L
    &=L_1^{-1}\cdots L_{n-1}^{-1}\\
    &=(I+\mathbf{l}_1\mathbf{e}_1^T)(I+\mathbf{l}_2\mathbf{e}_2^T)\cdots(I+\mathbf{l}_{n-1}\mathbf{e}_{n-1}^T)\\
    &=I+\mathbf{l}_1\mathbf{e}_1^T+\cdots+\mathbf{l}_{n-1}\mathbf{e}_{n-1}^T\\
    &=I+[\mathbf{l}_1,\mathbf{l}_2,\ldots,\mathbf{l}_{n-1},\mathbf{0}]\\
    &=
    \begin{bmatrix}
      1      & 0      & \cdots & 0      & 0 \\
      l_{21} & 1      & 0      & \cdots & 0 \\
      l_{31} & l_{32} & 1      & \ddots & \vdots \\
      \vdots & \vdots & \vdots & \ddots & 0 \\
      l_{n1} & l_{n2} & l_{n3} & \cdots & 1
    \end{bmatrix}.
  \end{aligned}
\]
上述 $LU$ 分解的计算过程称为高斯消去法. 于是有如下算法.

\begin{algorithm}{高斯消去法}
  \label{alg:2-1}
\[
  \left\{
  \begin{array}{ll}
    \textbf{for} & k=1:n-1\\
    & A(k+1:n,k)=A(k+1:n,k)/A(k,k)\\
    & A(k+1:n,k+1:n)=A(k+1:n,k+1:n)\\
    & \hphantom{A(k+1:n,k+1:n)={}}-A(k+1:n,k)A(k,k+1:n)\\
    \textbf{end}
  \end{array}
  \right.
\]
\end{algorithm}

高斯消去法的运算量为
\[
  \begin{aligned}
    \sum_{k=1}^{n-1}\Bigl((n-k)+2(n-k)^2\Bigr)
    &=\frac{n(n-1)}{2}+\frac{n(n-1)(2n-1)}{3}\\
    &=\frac{2}{3}n^3+\mathcal{O}(n^2)=\mathcal{O}(n^3).
  \end{aligned}
\]
注意, 高斯消去法要求 $a_{kk}^{(k-1)}\ (k=1,\ldots,n-1)$ 非零. 我们有如下定理.

\begin{theorem}
\label{thm:2-1}
元素 $a_{ii}^{(i-1)}\neq0\ (i=1,\ldots,k)$ 当且仅当 $A$ 的所有\newterm{leadingprincipal}{顺序主子矩阵}{leading principal submatrix} $A_i\ (i=1,\ldots,k)$ 非奇异.
\end{theorem}

\begin{proof}
对 $k$ 作归纳. $k=1$ 时显然成立. 设结论对直到 $k-1$ 的情形成立. 我们要证明: 若 $A_1,\ldots,A_{k-1}$ 非奇异, 则
\[
  A_k\ \text{非奇异}\iff a_{kk}^{(k-1)}\neq0.
\]
由假设知
\[
  a_{ii}^{(i-1)}\neq0,\qquad i=1,\ldots,k-1.
\]
用 $k-1$ 个高斯变换矩阵 $L_1,\ldots,L_{k-1}$, 得到
\begin{equation}
  A^{(k-1)}=L_{k-1}\cdots L_1A=
  \begin{bmatrix}
    A_{11}^{(k-1)} & A_{12}^{(k-1)}\\
    \mathbf{0}     & A_{22}^{(k-1)}
  \end{bmatrix},
  \label{eq:2-4}
\end{equation}
其中 $A_{11}^{(k-1)}$ 为对角元 $a_{ii}^{(i-1)}\ (i=1,\ldots,k-1)$ 非零的上三角矩阵. 因此 $A^{(k-1)}$ 的第 $k$ 个顺序主子矩阵具有如下形式
\[
  \begin{bmatrix}
    A_{11}^{(k-1)} & *\\
    \mathbf{0}     & a_{kk}^{(k-1)}
  \end{bmatrix}.
\]
记 $(L_1)_k,\ldots,(L_{k-1})_k$ 分别为 $L_1,\ldots,L_{k-1}$ 的第 $k$ 个顺序主子矩阵. 利用 \eqref{eq:2-4}, 得到
\[
  (L_{k-1})_k\cdots(L_1)_kA_k=
  \begin{bmatrix}
    A_{11}^{(k-1)} & *\\
    \mathbf{0}     & a_{kk}^{(k-1)}
  \end{bmatrix}.
\]
注意 $L_i\ (i=1,\ldots,k-1)$ 为单位下三角矩阵, 立即可知
\[
  \det(A_k)=a_{kk}^{(k-1)}\det\bigl(A_{11}^{(k-1)}\bigr)\neq0
\]
当且仅当 $a_{kk}^{(k-1)}\neq0$.
\end{proof}

于是有

\begin{theorem}
\label{thm:2-2}
若矩阵 $A\in\R^{n\times n}$ 的所有顺序主子矩阵 $A_i\ (i=1,\cdots,n)$ 均非奇异, 则 $A$ 的 $LU$ 分解存在且唯一.
\end{theorem}

\section{带主元选取的 $LU$ 分解}

在研究主元选取技巧之前, 先看下面这个简单例子:
\[
  \begin{bmatrix}
    0.3\times10^{-11} & 1\\
    1                 & 1
  \end{bmatrix}
  \begin{bmatrix}
    x_1\\ x_2
  \end{bmatrix}
  =
  \begin{bmatrix}
    0.7\\ 0.9
  \end{bmatrix}.
\]
若用 10 位十进制浮点算术施行高斯消去法, 则得到
\[
  \widehat{L}=
  \begin{bmatrix}
    1 & 0\\
    0.3333333333\times10^{12} & 1
  \end{bmatrix}
\]
以及
\[
  \widehat{U}=
  \begin{bmatrix}
    0.3\times10^{-11} & 1\\
    0 & -0.3333333333\times10^{12}
  \end{bmatrix}.
\]
于是计算解为
\[
  \hat{\mathbf{x}}=[0.0000000000,\ 0.7000000000]^T,
\]
与精确解
\[
  \mathbf{x}=[0.20000000000006\ldots,\ 0.69999999999994\ldots]^T
\]
相比, 这个结果并不好.

若仅交换第一个方程与第二个方程, 则得到
\[
  \begin{bmatrix}
    1                 & 1\\
    0.3\times10^{-11} & 1
  \end{bmatrix}
  \begin{bmatrix}
    x_1\\ x_2
  \end{bmatrix}
  =
  \begin{bmatrix}
    0.9\\ 0.7
  \end{bmatrix}.
\]
再次用 10 位十进制浮点算术施行高斯消去法, 得到
\[
  \widehat{L}=
  \begin{bmatrix}
    1 & 0\\
    0.3\times10^{-11} & 1
  \end{bmatrix},
  \qquad
  \widehat{U}=
  \begin{bmatrix}
    1 & 1\\
    0 & 1
  \end{bmatrix}.
\]
于是计算解为
\[
  \hat{\mathbf{x}}=[0.2000000000,\ 0.7000000000]^T,
\]
这个结果很好. 可见需要把置换引入高斯消去法. 先定义置换矩阵.

\begin{definition}
\label{def:2-1}
\newterm{permutationmatrix}{置换矩阵}{permutation matrix} $P$ 是行经过置换的单位矩阵.
\end{definition}

置换矩阵的重要性质包含在下面的引理中. 其证明是直接的.

\begin{lemma}
\label{lem:2-1}
设 $P,P_1,P_2\in\R^{n\times n}$ 为置换矩阵, $X\in\R^{n\times n}$. 则
\begin{romanlist}
  \item $PX$ 等于把 $X$ 的行作相应置换所得矩阵; $XP$ 等于把 $X$ 的列作相应置换所得矩阵.
  \item $P^{-1}=P^T$.
  \item $\det(P)=\pm1$.
  \item $P_1\cdot P_2$ 也是置换矩阵.
\end{romanlist}
\end{lemma}

现在介绍本节的主要定理.

\begin{theorem}
\label{thm:2-3}
若 $A$ 非奇异, 则存在置换矩阵 $P_1$ 与 $P_2$, 单位下三角矩阵 $L$ 以及非奇异上三角矩阵 $U$, 使得
\[
  P_1AP_2=LU.
\]
其中 $P_1$ 与 $P_2$ 只需其一.
\end{theorem}

\begin{proof}
对维数 $n$ 作归纳. $n=1$ 时显然成立. 设结论对 $n-1$ 成立. 若 $A$ 非奇异, 则它必有非零元素. 选取置换矩阵 $P_1'$ 与 $P_2'$, 使 $P_1'AP_2'$ 的 $(1,1)$ 位置非零. 现在写出所求的分解并对未知分量求解:
\begin{align}
  P_1'AP_2'
  &=
  \begin{bmatrix}
    a_{11} & A_{12}\\
    A_{21} & A_{22}
  \end{bmatrix}
  =
  \begin{bmatrix}
    1      & \mathbf{0}\\
    L_{21} & I
  \end{bmatrix}\cdot
  \begin{bmatrix}
    u_{11} & U_{12}\\
    \mathbf{0} & \widetilde{A}_{22}
  \end{bmatrix}\notag\\
  &=
  \begin{bmatrix}
    u_{11} & U_{12}\\
    L_{21}u_{11} & L_{21}U_{12}+\widetilde{A}_{22}
  \end{bmatrix},
  \label{eq:2-5}
\end{align}
其中 $A_{22},\widetilde{A}_{22}$ 为 $(n-1)\times(n-1)$ 矩阵, 而 $L_{21},U_{12}^T$ 为 $(n-1)\times1$ 矩阵.

求解该 $2\times2$ 分块分解的各分量, 得到
\[
  u_{11}=a_{11}\neq0,\qquad U_{12}=A_{12},
\]
以及
\[
  L_{21}u_{11}=A_{21},\qquad A_{22}=L_{21}U_{12}+\widetilde{A}_{22}.
\]
因此得到
\[
  L_{21}=\frac{A_{21}}{a_{11}},\qquad
  \widetilde{A}_{22}=A_{22}-L_{21}U_{12}.
\]
我们希望把归纳假设用于 $\widetilde{A}_{22}$, 但为此需验证
\[
  \det(\widetilde{A}_{22})\neq0.
\]
由于
\[
  \det(P_1'AP_2')=\pm\det(A)\neq0,
\]
并且
\[
  \begin{aligned}
    \det(P_1'AP_2')
    &=\det
    \begin{bmatrix}
      1      & \mathbf{0}\\
      L_{21} & I
    \end{bmatrix}\cdot
    \det
    \begin{bmatrix}
      u_{11} & U_{12}\\
      \mathbf{0} & \widetilde{A}_{22}
    \end{bmatrix}\\
    &=u_{11}\cdot\det(\widetilde{A}_{22}),
  \end{aligned}
\]
可知
\[
  \det(\widetilde{A}_{22})\neq0.
\]
因此, 由归纳假设, 存在置换矩阵 $\widetilde{P}_1$ 与 $\widetilde{P}_2$, 使得
\begin{equation}
  \widetilde{P}_1\widetilde{A}_{22}\widetilde{P}_2=\widetilde{L}\widetilde{U},
  \label{eq:2-6}
\end{equation}
其中 $\widetilde{L}$ 为单位下三角矩阵, $\widetilde{U}$ 为非奇异上三角矩阵. 把 \eqref{eq:2-6} 代入 \eqref{eq:2-5}, 得到
\[
  \begin{aligned}
    P_1'AP_2'
    &=
    \begin{bmatrix}
      1      & \mathbf{0}\\
      L_{21} & I
    \end{bmatrix}
    \begin{bmatrix}
      u_{11} & U_{12}\\
      \mathbf{0} & \widetilde{P}_1^T\widetilde{L}\widetilde{U}\widetilde{P}_2^T
    \end{bmatrix}\\
    &=
    \begin{bmatrix}
      1      & \mathbf{0}\\
      L_{21} & I
    \end{bmatrix}
    \begin{bmatrix}
      1      & \mathbf{0}\\
      \mathbf{0} & \widetilde{P}_1^T\widetilde{L}
    \end{bmatrix}
    \begin{bmatrix}
      u_{11} & U_{12}\\
      \mathbf{0} & \widetilde{U}\widetilde{P}_2^T
    \end{bmatrix}\\
    &=
    \begin{bmatrix}
      1      & \mathbf{0}\\
      L_{21} & \widetilde{P}_1^T\widetilde{L}
    \end{bmatrix}
    \begin{bmatrix}
      u_{11} & U_{12}\widetilde{P}_2\\
      \mathbf{0} & \widetilde{U}
    \end{bmatrix}
    \begin{bmatrix}
      1      & \mathbf{0}\\
      \mathbf{0} & \widetilde{P}_2^T
    \end{bmatrix}\\
    &=
    \begin{bmatrix}
      1      & \mathbf{0}\\
      \mathbf{0} & \widetilde{P}_1^T
    \end{bmatrix}
    \begin{bmatrix}
      1      & \mathbf{0}\\
      \widetilde{P}_1L_{21} & \widetilde{L}
    \end{bmatrix}
    \begin{bmatrix}
      u_{11} & U_{12}\widetilde{P}_2\\
      \mathbf{0} & \widetilde{U}
    \end{bmatrix}
    \begin{bmatrix}
      1      & \mathbf{0}\\
      \mathbf{0} & \widetilde{P}_2^T
    \end{bmatrix},
  \end{aligned}
\]
于是得到 $A$ 的所求分解:
\[
  \begin{aligned}
    P_1AP_2
    &=\left(
      \begin{bmatrix}
        1 & \mathbf{0}\\
        \mathbf{0} & \widetilde{P}_1
      \end{bmatrix}P_1'
    \right)
    A
    \left(
      P_2'
      \begin{bmatrix}
        1 & \mathbf{0}\\
        \mathbf{0} & \widetilde{P}_2
      \end{bmatrix}
    \right)\\
    &=
    \begin{bmatrix}
      1      & \mathbf{0}\\
      \widetilde{P}_1L_{21} & \widetilde{L}
    \end{bmatrix}
    \begin{bmatrix}
      u_{11} & U_{12}\widetilde{P}_2\\
      \mathbf{0} & \widetilde{U}
    \end{bmatrix}.
  \end{aligned}
\]
\end{proof}

这种行列交换策略称为\newterm{completepivoting}{全主元选取}{complete pivoting}. 于是有如下算法.

\begin{algorithm}{全主元高斯消去法}
  \label{alg:2-2}
\[
  \left\{
  \begin{array}{ll}
    \textbf{for} & k=1:n-1\\
    & \text{choose } p,q,\ (k\leqslant p,q\leqslant n)\ \text{such that}\\
    & \hphantom{\text{choose }}\bigl|A(p,q)\bigr|=\max\bigl\{\bigl|A(i,j)\bigr|:i=k:n,\ j=k:n\bigr\}\\
    & A(k,1:n)\leftrightarrow A(p,1:n)\\
    & A(1:n,k)\leftrightarrow A(1:n,q)\\
    & \textbf{if}\ \ A(k,k)\neq0\\
    & \hphantom{\textbf{if}\ \ }A(k+1:n,k)=A(k+1:n,k)/A(k,k)\\
    & \hphantom{\textbf{if}\ \ }A(k+1:n,k+1:n)=A(k+1:n,k+1:n)\\
    & \hphantom{\textbf{if}\ \ A(k+1:n,k+1:n)={}}-A(k+1:n,k)A(k,k+1:n)\\
    & \textbf{else}\\
    & \hphantom{\textbf{if}\ \ }\textbf{stop}\\
    & \textbf{end}\\
    \textbf{end}
  \end{array}
  \right.
\]
\end{algorithm}

注意, 虽然带全主元选取的 $LU$ 分解能克服不带主元选取的 $LU$ 分解的某些缺点, 但全主元选取的代价很高. 通常, 为选取主元而比较矩阵元素就需要 $\mathcal{O}(n^3)$ 次运算.

为了降低选主元的运算代价, 人们提出了带部分主元选取的 $LU$ 分解. 在部分主元选取中, 第 $k$ 步从子矩阵 $A_{22}^{(k-1)}$ 中选取满足
\[
  \left|a_{pk}^{(k-1)}\right|=\max\left\{\left|a_{ik}^{(k-1)}\right|:k\leqslant i\leqslant n\right\}
\]
的 $a_{pk}^{(k-1)}$. 当 $A$ 非奇异时, 带部分主元选取的 $LU$ 分解总可进行到底, 最终得到
\[
  PA=LU.
\]
在该算法中, 为选取主元而比较矩阵元素的运算量为 $\mathcal{O}(n^2)$. 于是有

\begin{algorithm}{部分主元高斯消去法}
  \label{alg:2-3}
\[
  \left\{
  \begin{array}{ll}
    \textbf{for} & k=1:n-1\\
    & \text{choose } p,\ (k\leqslant p\leqslant n)\ \text{such that}\\
    & \hphantom{\text{choose }}\bigl|A(p,k)\bigr|=\max\bigl\{\bigl|A(i,k)\bigr|:i=k:n\bigr\}\\
    & A(k,1:n)\leftrightarrow A(p,1:n)\\
    & \textbf{if}\ \ A(k,k)\neq0\\
    & \hphantom{\textbf{if}\ \ }A(k+1:n,k)=A(k+1:n,k)/A(k,k)\\
    & \hphantom{\textbf{if}\ \ }A(k+1:n,k+1:n)=A(k+1:n,k+1:n)\\
    & \hphantom{\textbf{if}\ \ A(k+1:n,k+1:n)={}}-A(k+1:n,k)A(k,k+1:n)\\
    & \textbf{else}\\
    & \hphantom{\textbf{if}\ \ }\textbf{stop}\\
    & \textbf{end}\\
    \textbf{end}
  \end{array}
  \right.
\]
\end{algorithm}

关于秩揭示 $LU$ 分解, 参见 \cite{src:40}. 关于带或不带主元选取的高斯消去法的概率分析, 可参见 \cite{src:34,src:81,src:99}.

\section{Cholesky 分解}

设 $A\in\R^{n\times n}$ 为\newterm{symmetricpositivedefinite}{对称正定}{symmetric positive definite}矩阵, 即它满足
\[
  A=A^T,\qquad \mathbf{x}^TA\mathbf{x}>0
\]
对所有 $\mathbf{x}\neq\mathbf{0}\in\R^n$ 成立. 我们有

\begin{theorem}
\label{thm:2-4}
设 $A\in\R^{n\times n}$ 为对称正定矩阵. 则存在对角元为正的下三角矩阵 $L\in\R^{n\times n}$, 使得
\[
  A=LL^T.
\]
该分解称为\newterm{choleskyfactorization}{Cholesky 分解}{Cholesky factorization}.
\end{theorem}

\begin{proof}
由于 $A$ 正定, $A$ 的所有主子矩阵都应正定. 由定理 \ref{thm:2-2}, 存在单位下三角矩阵 $\widetilde{L}$ 与上三角矩阵 $U$, 使得
\[
  A=\widetilde{L}U.
\]
令
\[
  D=\diag(u_{11},\ldots,u_{nn}),\qquad \widetilde{U}=D^{-1}U,
\]
其中 $u_{ii}>0\ (i=1,\ldots,n)$. 则有
\[
  \widetilde{U}^TD\widetilde{L}^T=A^T=A=\widetilde{L}D\widetilde{U}.
\]
因此
\[
  \widetilde{L}^T\widetilde{U}^{-1}=D^{-1}\widetilde{U}^{-T}\widetilde{L}D.
\]
注意 $\widetilde{L}^T\widetilde{U}^{-1}$ 是单位上三角矩阵, 而 $D^{-1}\widetilde{U}^{-T}\widetilde{L}D$ 是下三角矩阵. 故
\[
  \widetilde{L}^T\widetilde{U}^{-1}=I=D^{-1}\widetilde{U}^{-T}\widetilde{L}D,
\]
这说明 $\widetilde{U}=\widetilde{L}^T$. 从而
\[
  A=\widetilde{L}D\widetilde{L}^T.
\]
令
\[
  L=\widetilde{L}\cdot\diag\bigl(\sqrt{u_{11}},\ldots,\sqrt{u_{nn}}\bigr).
\]
最终得到
\[
  A=LL^T. \qedhere
\]
\end{proof}

于是, 当矩阵 $A$ 对称正定时, 可通过以下三个步骤求方程组 $A\mathbf{x}=\mathbf{b}$ 的解:

\begin{parenlist}
  \item 求 $A$ 的 Cholesky 分解: $A=LL^T$.
  \item 求 $L\mathbf{y}=\mathbf{b}$ 的解 $\mathbf{y}$.
  \item 求 $L^T\mathbf{x}=\mathbf{y}$ 的解 $\mathbf{x}$.
\end{parenlist}

由定理 \ref{thm:2-4} 可知, Cholesky 分解不需要选取主元. 此外, 可通过比较 $A=LL^T$ 两端对应元素来直接计算 $L$. 于是有如下算法.

\begin{algorithm}{Cholesky 分解}
  \label{alg:2-4}
\[
  \left\{
  \begin{array}{ll}
    \textbf{for} & k=1:n\\
    & A(k,k)=\sqrt{A(k,k)}\\
    & A(k+1:n,k)=A(k+1:n,k)/A(k,k)\\
    & \textbf{for}\ \ j=k+1:n\\
    & \hphantom{\textbf{for}\ \ }A(j:n,j)=A(j:n,j)-A(j:n,k)A(j,k)\\
    & \textbf{end}\\
    \textbf{end}
  \end{array}
  \right.
\]
\end{algorithm}

Cholesky 分解的运算量为 $n^3/3$.

\begin{exercises}
  \item 设 $S,T\in\R^{n\times n}$ 为上三角矩阵, 且
  \[
    (ST-\lambda I)\mathbf{x}=\mathbf{b}
  \]
  为非奇异方程组. 给出一个 $\mathcal{O}(n^2)$ 次运算的算法来计算 $\mathbf{x}$.

  \item 证明对称正定矩阵 $A$ 的 $LDL^T$ 分解是唯一的, 其中 $L$ 为单位下三角矩阵.

  \item 设 $A\in\R^{n\times n}$ 为对称正定矩阵. 给出一个算法, 计算上三角矩阵 $U\in\R^{n\times n}$, 使得 $A=UU^T$.

  \item 设 $A=[a_{ij}]\in\R^{n\times n}$ 为\newterm{strictlydiagdominant}{严格对角占优}{strictly diagonally dominant}矩阵, 即
  \[
    |a_{kk}|>\sum_{\substack{j=1\\ j\neq k}}^{n}|a_{kj}|,\qquad k=1,2,\ldots,n.
  \]
  证明: 严格对角占优矩阵非奇异, 且对角元为正的严格对角占优对称矩阵正定.

  \item 设
  \[
    A=
    \begin{bmatrix}
      A_{11} & A_{12}\\
      A_{21} & A_{22}
    \end{bmatrix}
  \]
  其中 $A_{11}$ 为 $k\times k$ 非奇异矩阵. 则
  \[
    S=A_{22}-A_{21}A_{11}^{-1}A_{12}
  \]
  称为 $A_{11}$ 在 $A$ 中的 \newterm{schurcomplement}{Schur 补}{Schur complement}. 证明: 经过 $k$ 步不带主元选取的高斯消去后, $A_{22}^{(k-1)}=S$.

  \item 设 $A$ 为对称正定矩阵. 在高斯消去第一步结束时, 有
  \[
    \begin{bmatrix}
      a_{11} & \mathbf{a}_1^T\\
      \mathbf{0} & A_{22}
    \end{bmatrix}.
  \]
  证明 $A_{22}$ 也是对称正定的.

  \item 设 $A=[a_{ij}]\in\R^{n\times n}$ 为严格对角占优矩阵. 经过一步高斯消去后, 有
  \[
    \begin{bmatrix}
      a_{11} & \mathbf{a}_1^T\\
      \mathbf{0} & A_{22}
    \end{bmatrix}.
  \]
  证明 $A_{22}$ 也是严格对角占优的.

  \item 证明: 若 $PAQ=LU$ 由带主元选取的高斯消去法得到, 则 $|u_{ii}|\geqslant|u_{ij}|$ 对 $j=i+1,\ldots,n$ 成立.

  \item 设 $H=A+\ii B$ 为 Hermite 正定矩阵, 其中 $A,B\in\R^{n\times n}$.
  \subquest 证明矩阵
  \[
    C=
    \begin{bmatrix}
      A & -B\\
      B & A
    \end{bmatrix}
  \]
  为对称正定矩阵.
  \subquest 如何仅用实数运算求解
  \[
    (A+\ii B)(\mathbf{x}+\ii\mathbf{y})=\mathbf{b}+\ii\mathbf{c},
    \qquad \mathbf{x},\mathbf{y},\mathbf{b},\mathbf{c}\in\R^n\ ?
  \]

  \item 给出一个用带部分主元选取的高斯消去法求解三对角方程组的算法.

  \item 证明: 若奇异矩阵 $A\in\R^{n\times n}$ 具有唯一的 $LU$ 分解, 则 $A_k$ 对 $k=1,2,\ldots,n-1$ 非奇异.
\end{exercises}

% chapters/ch03.tex —— 第 3 章：扰动与误差分析
% 源书：Chapter 3 Perturbation and Error Analysis，书页 23–45（物理页 37–59）
% 公式编号 (3.1)–(3.22)；定理 3.1–3.12；定义 3.1–3.5；引理 3.1–3.2；推论 3.1–3.3；例 3.1；习题 1–15
\chapter{扰动与误差分析}
\label{ch:perturbation}

本章讨论扰动与误差对数值解的影响, 并给出浮点运算与部分主元选取技术的误差分析. 众所周知, 线性向量空间中距离与大小的本质刻画由范数承担. 因此在展开扰动与误差分析之前, 我们需要先引入向量范数与矩阵范数并研究其性质.

\section{向量范数与矩阵范数}
\label{sec:3-1}

先介绍向量范数.

\subsection{向量范数}

设
\[
  \mathbf{x}=[x_1,x_2,\ldots,x_n]^T\in\R^n.
\]

\begin{definition}
\label{def:3-1}
$\R^n$ 上的\newterm{vectornorm}{向量范数}{vector norm}是一个函数, 它对每个 $\mathbf{x}\in\R^n$ 赋予一个实数 $\norm{\mathbf{x}}$, 称为 $\mathbf{x}$ 的范数, 且对一切 $\mathbf{x},\mathbf{y}\in\R^n$ 与一切 $\alpha\in\R$ 满足以下三条性质:
\begin{romanlist}
  \item 若 $\mathbf{x}\neq\mathbf{0}$ 则 $\norm{\mathbf{x}}>0$; 且 $\norm{\mathbf{x}}=0$ 当且仅当 $\mathbf{x}=\mathbf{0}$;
  \item $\norm{\alpha\mathbf{x}}=|\alpha|\cdot\norm{\mathbf{x}}$;
  \item $\norm{\mathbf{x}+\mathbf{y}}\leqslant\norm{\mathbf{x}}+\norm{\mathbf{y}}$.
\end{romanlist}
\end{definition}

一类有用的向量范数是由下式定义的\newterm{pnorm}{$p$ 范数}{$p$-norm}:
\[
  \norm{\mathbf{x}}_p\equiv\left(\sum_{i=1}^{n}|x_i|^p\right)^{\frac{1}{p}},
\]
其中 $1\leqslant p$. 实践中最为常用的 $p$ 范数是
\[
  \norm{\mathbf{x}}_1=\sum_{i=1}^{n}|x_i|,\qquad
  \norm{\mathbf{x}}_2=\left(\sum_{i=1}^{n}|x_i|^2\right)^{\frac{1}{2}},\qquad
  \norm{\mathbf{x}}_\infty=\max_{1\leqslant i\leqslant n}|x_i|.
\]

关于 $\norm{\cdot}_2$ 的 Cauchy-Schwarz 不等式如下:
\[
  \left|\mathbf{x}^T\mathbf{y}\right|\leqslant\norm{\mathbf{x}}_2\norm{\mathbf{y}}_2
\]
对 $\mathbf{x},\mathbf{y}\in\R^n$ 成立; 它是如下 Hölder 不等式 \cite{src:39} 的特例:
\[
  \left|\mathbf{x}^T\mathbf{y}\right|\leqslant\norm{\mathbf{x}}_p\norm{\mathbf{y}}_q,
\]
其中 $1/p+1/q=1$.

$\R^n$ 上向量范数的一个非常重要的性质是: $\R^n$ 上所有向量范数都是等价的, 如下述定理所示 \cite{src:75}.

\begin{theorem}
\label{thm:3-1}
若 $\norm{\cdot}_\alpha$ 与 $\norm{\cdot}_\beta$ 是 $\R^n$ 上的两个范数, 则存在两个正常数 $c_1$ 与 $c_2$, 使得
\[
  c_1\norm{\mathbf{x}}_\alpha\leqslant\norm{\mathbf{x}}_\beta\leqslant c_2\norm{\mathbf{x}}_\alpha
\]
对一切 $\mathbf{x}\in\R^n$ 成立.
\end{theorem}

例如, 若 $\mathbf{x}\in\R^n$, 则有
\[
  \norm{\mathbf{x}}_2\leqslant\norm{\mathbf{x}}_1\leqslant\sqrt{n}\,\norm{\mathbf{x}}_2,
\]
\[
  \norm{\mathbf{x}}_\infty\leqslant\norm{\mathbf{x}}_2\leqslant\sqrt{n}\,\norm{\mathbf{x}}_\infty
\]
以及
\[
  \norm{\mathbf{x}}_\infty\leqslant\norm{\mathbf{x}}_1\leqslant n\norm{\mathbf{x}}_\infty.
\]

我们指出, 对任一向量序列 $\{\mathbf{x}^k\}$, 其中 $\mathbf{x}^k=[x_1^{(k)},\ldots,x_n^{(k)}]^T\in\R^n$, 以及 $\mathbf{x}=[x_1,\ldots,x_n]^T\in\R^n$, 由定理 \ref{thm:3-1} 可证
\[
  \lim_{k\to\infty}\norm{\mathbf{x}^k-\mathbf{x}}=0
  \iff
  \lim_{k\to\infty}\left|x_i^{(k)}-x_i\right|=0
\]
对 $i=1,\ldots,n$ 成立.

\subsection{矩阵范数}

设
\[
  A=[a_{ij}]_{i,j=1}^{n}\in\R^{n\times n}.
\]

下面转向矩阵范数.

\begin{definition}
\label{def:3-2}
\newterm{matrixnorm}{矩阵范数}{matrix norm}是一个函数, 它对每个 $A\in\R^{n\times n}$ 赋予一个实数 $\norm{A}$, 称为 $A$ 的范数, 且对一切 $A,B\in\R^{n\times n}$ 与一切 $\alpha\in\R$ 满足以下四条性质:
\begin{romanlist}
  \item 若 $A\neq0$ 则 $\norm{A}>0$; 且 $\norm{A}=0$ 当且仅当 $A=0$;
  \item $\norm{\alpha A}=|\alpha|\cdot\norm{A}$;
  \item $\norm{A+B}\leqslant\norm{A}+\norm{B}$;
  \item $\norm{AB}\leqslant\norm{A}\cdot\norm{B}$.
\end{romanlist}
\end{definition}

$\R^{n\times n}$ 上矩阵范数的一个重要性质是: $\R^{n\times n}$ 上所有矩阵范数都是等价的. 关于向量范数与矩阵范数之间的关系, 有

\begin{definition}
\label{def:3-3}
若矩阵范数 $\norm{\cdot}_M$ 与向量范数 $\norm{\cdot}_v$ 满足
\[
  \norm{A\mathbf{x}}_v\leqslant\norm{A}_M\norm{\mathbf{x}}_v
\]
对 $A\in\R^{n\times n}$ 与 $\mathbf{x}\in\R^n$ 成立, 则称这两个范数是\newterm{consistentnorm}{相容的}{mutually consistent}.
\end{definition}

对任意向量范数 $\norm{\cdot}_v$, 可按如下自然方式定义一个矩阵范数:
\[
  \norm{A}_M\equiv\max_{\mathbf{x}\neq\mathbf{0}}\frac{\norm{A\mathbf{x}}_v}{\norm{\mathbf{x}}_v}
  =\max_{\norm{\mathbf{x}}_v=1}\norm{A\mathbf{x}}_v.
\]
这样得到的矩阵范数称为由该向量范数\newterm{inducednorm}{诱导}{induced}的范数.

最重要的矩阵范数是由向量 $p$ 范数诱导的矩阵 $p$ 范数, $p=1,2,\infty$. 我们有下述定理.

\begin{theorem}
\label{thm:3-2}
设
\[
  A=[a_{ij}]_{i,j=1}^{n}\in\R^{n\times n}.
\]
则
\begin{romanlist}
  \item $\displaystyle\norm{A}_1=\max_{\norm{\mathbf{x}}_1=1}\norm{A\mathbf{x}}_1
    =\max_{1\leqslant j\leqslant n}\sum_{i=1}^{n}\abs{a_{ij}}$;
  \item $\displaystyle\norm{A}_\infty=\max_{\norm{\mathbf{x}}_\infty=1}\norm{A\mathbf{x}}_\infty
    =\max_{1\leqslant i\leqslant n}\sum_{j=1}^{n}\abs{a_{ij}}$;
  \item $\displaystyle\norm{A}_2=\max_{\norm{\mathbf{x}}_2=1}\norm{A\mathbf{x}}_2
    =\sqrt{\lambda_{\max}(A^TA)}$, 其中 $\lambda_{\max}(A^TA)$ 是 $A^TA$ 的最大特征值.
\end{romanlist}
\end{theorem}

\begin{proof}
我们只给出 (i) 与 (iii) 的证明. 以下总假定 $A\neq0$.

对 (i), 把矩阵 $A$ 按列分块:
\[
  A=[\mathbf{a}_1\mid\mathbf{a}_2\mid\cdots\mid\mathbf{a}_n].
\]
令
\[
  \delta=\norm{\mathbf{a}_{j_0}}_1=\max_{1\leqslant j\leqslant n}\norm{\mathbf{a}_j}_1.
\]
则对任一满足 $\norm{\mathbf{x}}_1=\sum_{i=1}^{n}|x_i|=1$ 的向量 $\mathbf{x}=[x_1,x_2,\ldots,x_n]^T\in\R^n$, 有
\[
  \begin{aligned}
    \norm{A\mathbf{x}}_1
    &=\norm{\sum_{j=1}^{n}x_j\mathbf{a}_j}_1
      \leqslant\sum_{j=1}^{n}|x_j|\cdot\norm{\mathbf{a}_j}_1\\
    &\leqslant\Bigl(\sum_{j=1}^{n}|x_j|\Bigr)\cdot\max_{1\leqslant j\leqslant n}\norm{\mathbf{a}_j}_1
      =\norm{\mathbf{a}_{j_0}}_1=\delta.
  \end{aligned}
\]
令 $\mathbf{e}_{j_0}$ 表示第 $j_0$ 个单位向量, 则
\[
  \norm{A\mathbf{e}_{j_0}}_1=\norm{\mathbf{a}_{j_0}}_1=\delta.
\]
因此
\[
  \norm{A}_1=\max_{\norm{\mathbf{x}}_1=1}\norm{A\mathbf{x}}_1
  =\delta=\max_{1\leqslant j\leqslant n}\norm{\mathbf{a}_j}_1
  =\max_{1\leqslant j\leqslant n}\sum_{i=1}^{n}\abs{a_{ij}}.
\]

对 (iii), 有
\[
  \begin{aligned}
    \norm{A}_2
    &=\max_{\norm{\mathbf{x}}_2=1}\norm{A\mathbf{x}}_2
      =\max_{\norm{\mathbf{x}}_2=1}\bigl[(A\mathbf{x})^T(A\mathbf{x})\bigr]^{1/2}\\
    &=\max_{\norm{\mathbf{x}}_2=1}\bigl[\mathbf{x}^T(A^TA)\mathbf{x}\bigr]^{1/2}.
  \end{aligned}
\]
由于 $A^TA$ 是半正定的, 其特征值可按如下次序排列:
\[
  \lambda_1\geqslant\lambda_2\geqslant\cdots\geqslant\lambda_n\geqslant0.
\]
令
\[
  \mathbf{v}_1,\mathbf{v}_2,\ldots,\mathbf{v}_n\in\R^n
\]
分别表示对应于 $\lambda_1,\lambda_2,\ldots,\lambda_n$ 的标准正交特征向量. 则对任一满足 $\norm{\mathbf{x}}_2=1$ 的向量 $\mathbf{x}\in\R^n$, 有
\[
  \mathbf{x}=\sum_{i=1}^{n}\alpha_i\mathbf{v}_i,\qquad\sum_{i=1}^{n}\alpha_i^2=1.
\]
因此
\[
  \mathbf{x}^TA^TA\mathbf{x}=\sum_{i=1}^{n}\lambda_i\alpha_i^2\leqslant\lambda_1.
\]
另一方面, 令 $\mathbf{x}=\mathbf{v}_1$, 得
\[
  \mathbf{x}^TA^TA\mathbf{x}=\mathbf{v}_1^TA^TA\mathbf{v}_1
  =\mathbf{v}_1^T\lambda_1\mathbf{v}_1=\lambda_1.
\]
于是
\[
  \norm{A}_2=\max_{\norm{\mathbf{x}}_2=1}\norm{A\mathbf{x}}_2
  =\sqrt{\lambda_1}=\sqrt{\lambda_{\max}(A^TA)}.
\]
\end{proof}

关于范数 $\norm{\cdot}_2$, 我们有下述定理.

\begin{theorem}
\label{thm:3-3}
设 $A\in\R^{n\times n}$. 则
\begin{romanlist}
  \item $\displaystyle\norm{A}_2=\max_{\norm{\mathbf{x}}_2=1}\max_{\norm{\mathbf{y}}_2=1}\abs{\mathbf{y}^*A\mathbf{x}}$, 其中 $\mathbf{x},\mathbf{y}\in\C^n$;
  \item $\norm{A^T}_2=\norm{A}_2=\sqrt{\norm{A^TA}_2}$;
  \item $\norm{A}_2=\norm{QAZ}_2$, 对任意正交矩阵 $Q$ 与 $Z$ 成立.
\end{romanlist}
我们回忆, 矩阵 $M\in\R^{n\times n}$ 称为\newterm{orthogonalmatrix}{正交的}{orthogonal}, 若 $M^{-1}=M^T$.
\end{theorem}

\begin{proof}
我们只证 (i). 先引入向量范数 $\norm{\cdot}$ 的\newterm{dualnorm}{对偶范数}{dual norm} $\norm{\cdot}^D$, 其定义为
\[
  \norm{\mathbf{y}}^D=\max_{\norm{\mathbf{x}}=1}\abs{\mathbf{y}^*\mathbf{x}}.
\]
对 $\norm{\cdot}_2$, 由 Cauchy-Schwarz 不等式有
\[
  \abs{\mathbf{y}^*\mathbf{x}}\leqslant\norm{\mathbf{y}}_2\norm{\mathbf{x}}_2,
\]
当 $\mathbf{x}=\mathbf{y}/\norm{\mathbf{y}}_2$ 时取等号. 因此 $\norm{\cdot}_2$ 的对偶范数为
\[
  \norm{\mathbf{y}}_2^D=\max_{\norm{\mathbf{x}}_2=1}\abs{\mathbf{y}^*\mathbf{x}}
  =\max_{\norm{\mathbf{x}}_2=1}\norm{\mathbf{y}}_2\norm{\mathbf{x}}_2=\norm{\mathbf{y}}_2.
\]
所以 $\norm{\cdot}_2$ 是自对偶的. 现在考虑
\[
  \begin{aligned}
    \norm{A}_2
    &=\max_{\norm{\mathbf{x}}_2=1}\norm{A\mathbf{x}}_2
      =\max_{\norm{\mathbf{x}}_2=1}\norm{A\mathbf{x}}_2^D\\
    &=\max_{\norm{\mathbf{x}}_2=1}\max_{\norm{\mathbf{y}}_2=1}\abs{(A\mathbf{x})^*\mathbf{y}}
      =\max_{\norm{\mathbf{x}}_2=1}\max_{\norm{\mathbf{y}}_2=1}\abs{\mathbf{y}^*A\mathbf{x}}.
  \end{aligned}
\]
\end{proof}

另一有用的范数是\newterm{frobeniusnorm}{Frobenius 范数}{Frobenius norm}, 其定义为
\[
  \norm{A}_F\equiv\left(\sum_{j=1}^{n}\sum_{i=1}^{n}\abs{a_{ij}}^2\right)^{\frac{1}{2}}.
\]
$\norm{\cdot}_F$ 最重要的性质之一是: 对任意正交矩阵 $Q$ 与 $Z$,
\[
  \norm{A}_F=\norm{QAZ}_F.
\]

下面把关于范数的讨论推广到 $\C$ 域. 我们指出, 从范数的观点看, $\R$ 域中的矩阵或向量与 $\C$ 域中的矩阵或向量并无本质区别.

\begin{definition}
\label{def:3-4}
设 $A\in\C^{n\times n}$. 则 $A$ 的所有特征值构成的集合称为 $A$ 的\newterm{spectrum}{谱}{spectrum}, 而
\[
  \rho(A)=\max\{|\lambda|:\lambda\text{ 属于 }A\text{ 的谱}\}
\]
称为 $A$ 的\newterm{spectralradius}{谱半径}{spectral radius}.
\end{definition}

关于谱半径与矩阵范数之间的关系, 有

\begin{theorem}
\label{thm:3-4}
设 $A\in\C^{n\times n}$. 则
\begin{romanlist}
  \item 对任意矩阵范数, 有
  \[
    \rho(A)\leqslant\norm{A};
  \]
  \item 对任意 $\epsilon>0$, 存在定义在 $\C^{n\times n}$ 上的范数 $\norm{\cdot}_\epsilon$, 使得
  \[
    \norm{A}_\epsilon\leqslant\rho(A)+\epsilon.
  \]
\end{romanlist}
\end{theorem}

\begin{proof}
对 (i), 设 $\mathbf{x}\in\C^n$ 满足
\[
  \mathbf{x}\neq\mathbf{0},\qquad A\mathbf{x}=\lambda\mathbf{x},\qquad|\lambda|=\rho(A).
\]
则
\[
  \rho(A)\norm{\mathbf{x}\mathbf{e}_1^T}=\norm{\lambda\mathbf{x}\mathbf{e}_1^T}
  =\norm{A\mathbf{x}\mathbf{e}_1^T}\leqslant\norm{A}\cdot\norm{\mathbf{x}\mathbf{e}_1^T}.
\]
故
\[
  \rho(A)\leqslant\norm{A}.
\]

对 (ii), 利用定理 \ref{thm:ch1-jordan}(Jordan 分解定理), 知存在非奇异矩阵 $X\in\C^{n\times n}$, 使得
\[
  X^{-1}AX=
  \begin{bmatrix}
    \lambda_1 & \delta_1 &           &        &                \\
              & \lambda_2 & \delta_2 &        &                \\
              &           & \ddots    & \ddots &                \\
              &           &           & \lambda_{n-1} & \delta_{n-1}\\
              &           &           &        & \lambda_n
  \end{bmatrix},
\]
其中 $\delta_i=1$ 或 $0$. 对任意给定的 $\epsilon>0$, 令
\[
  D_\epsilon=\diag(1,\epsilon,\epsilon^2,\ldots,\epsilon^{n-1}),
\]
则
\[
  D_\epsilon^{-1}X^{-1}AXD_\epsilon=
  \begin{bmatrix}
    \lambda_1 & \epsilon\delta_1 &           &        &                \\
              & \lambda_2 & \epsilon\delta_2 &        &                \\
              &           & \ddots    & \ddots &                \\
              &           &           & \lambda_{n-1} & \epsilon\delta_{n-1}\\
              &           &           &        & \lambda_n
  \end{bmatrix}.
\]
现在定义
\[
  \norm{G}_\epsilon=\norm{D_\epsilon^{-1}X^{-1}GXD_\epsilon}_\infty,
  \qquad G\in\C^{n\times n}.
\]
容易看出, 这个矩阵范数 $\norm{\cdot}_\epsilon$ 实际上是由如下定义的向量范数诱导的:
\[
  \norm{\mathbf{x}}_{XD_\epsilon}=\norm{(XD_\epsilon)^{-1}\mathbf{x}}_\infty,
  \qquad\mathbf{x}\in\C^n.
\]
因此
\[
  \norm{A}_\epsilon=\norm{D_\epsilon^{-1}X^{-1}AXD_\epsilon}_\infty
  =\max_{1\leqslant i\leqslant n}\bigl(|\lambda_i|+|\epsilon\delta_i|\bigr)
  \leqslant\rho(A)+\epsilon,
\]
其中 $\delta_n=0$.
\end{proof}

我们指出, 对任一矩阵序列 $\{A_{(k)}\}$, 其中 $A_{(k)}=[a_{ij}^{(k)}]\in\R^{n\times n}$, 以及 $A=[a_{ij}]\in\R^{n\times n}$,
\[
  \lim_{k\to\infty}\norm{A_{(k)}-A}=0
  \iff
  \lim_{k\to\infty}a_{ij}^{(k)}=a_{ij}
\]
对 $i,j=1,\ldots,n$ 成立.

\begin{theorem}
\label{thm:3-5}
设 $A\in\C^{n\times n}$. 则
\[
  \lim_{k\to\infty}A^k=0\iff\rho(A)<1,
\]
其中 $0$ 是 $n\times n$ 零矩阵.
\end{theorem}

\begin{proof}
先假定
\[
  \lim_{k\to\infty}A^k=0.
\]
设 $\lambda$ 是 $A$ 的特征值且 $\rho(A)=|\lambda|$. 则对任意 $k$, $\lambda^k$ 是 $A^k$ 的特征值. 由定理 \ref{thm:3-4} (i) 知, 对任意 $k$,
\[
  \rho(A)^k=|\lambda|^k=\abs{\lambda^k}\leqslant\rho(A^k)\leqslant\norm{A^k}.
\]
因此
\[
  \lim_{k\to\infty}\rho(A)^k=0,
\]
这说明 $\rho(A)<1$.

反之, 假定 $\rho(A)<1$. 由定理 \ref{thm:3-4} (ii), 存在矩阵范数 $\norm{\cdot}_\epsilon$ 使得 $\norm{A}_\epsilon<1$. 因此
\[
  0\leqslant\norm{A^k}_\epsilon\leqslant\norm{A}_\epsilon^k\to0,
  \qquad k\to\infty,
\]
即
\[
  \lim_{k\to\infty}A^k=0.
\]
\end{proof}

利用定理 \ref{thm:3-5}, 容易证明下述重要定理.

\begin{theorem}
\label{thm:3-6}
设 $A\in\C^{n\times n}$. 则
\begin{romanlist}
  \item $\sum_{k=0}^{\infty}A^k$ 收敛当且仅当 $\rho(A)<1$;
  \item 当 $\sum_{k=0}^{\infty}A^k$ 收敛时, 有
  \[
    \sum_{k=0}^{\infty}A^k=(I-A)^{-1}.
  \]
\end{romanlist}
此外, 存在定义在 $\C^{n\times n}$ 上的范数, 使得对任意 $m$,
\[
  \norm{(I-A)^{-1}-\sum_{k=0}^{m}A^k}\leqslant\frac{\norm{A}^{m+1}}{1-\norm{A}}.
\]
\end{theorem}

\begin{corollary}
\label{cor:3-1}
设 $\norm{\cdot}$ 是定义在 $\C^{n\times n}$ 上的范数, $\norm{I}=1$, 且 $A\in\C^{n\times n}$ 满足 $\norm{A}<1$. 则 $I-A$ 非奇异, 且
\[
  \norm{(I-A)^{-1}}\leqslant\frac{1}{1-\norm{A}}.
\]
\end{corollary}

\begin{proof}
只需注意
\[
  \norm{(I-A)^{-1}}=\norm{\sum_{k=0}^{\infty}A^k}
  \leqslant1+\sum_{k=1}^{\infty}\norm{A}^k=\frac{1}{1-\norm{A}}.
\]
\end{proof}

\section{线性方程组的扰动分析}

先考虑如下简单的线性方程组 $A\mathbf{x}=\mathbf{b}$:
\[
  \begin{bmatrix}
    2.0001 & 1.9999\\
    1.9999 & 2.0001
  \end{bmatrix}
  \begin{bmatrix}
    x_1\\ x_2
  \end{bmatrix}
  =
  \begin{bmatrix}
    4\\ 4
  \end{bmatrix}.
\]
该方程组的解为 $\mathbf{x}=[1,1]^T$. 若 $\mathbf{b}$ 有一个小扰动, 比如
\[
  \delta\mathbf{b}=[1\times10^{-4},\,-1\times10^{-4}]^T,
\]
则方程组变为
\[
  \begin{bmatrix}
    2.0001 & 1.9999\\
    1.9999 & 2.0001
  \end{bmatrix}
  \begin{bmatrix}
    \hat{x}_1\\ \hat{x}_2
  \end{bmatrix}
  =
  \begin{bmatrix}
    4.0001\\ 3.9999
  \end{bmatrix}.
\]
该扰动方程组的解为 $\hat{\mathbf{x}}=[1.5,0.5]^T$. 因此
\[
  \frac{\norm{\mathbf{x}-\hat{\mathbf{x}}}_\infty}{\norm{\mathbf{x}}_\infty}=\frac{1}{2},
  \qquad
  \frac{\norm{\delta\mathbf{b}}_\infty}{\norm{\mathbf{b}}_\infty}=\frac{1}{40000},
\]
即解的相对误差是 $\mathbf{b}$ 上扰动相对误差的 $20000$ 倍.

由此可见, 在求解线性方程组 $A\mathbf{x}=\mathbf{b}$ 时, 需要一个良好的度量来刻画计算解对输入小扰动的敏感程度. 于是引入矩阵的条件数, 它把 $\mathbf{x}$ 的扰动与 $A$, $\mathbf{b}$ 的扰动联系起来.

\begin{definition}
\label{def:3-5}
设 $\norm{\cdot}$ 为任一矩阵范数, $A$ 为非奇异矩阵. $A$ 的条件数定义为
\begin{equation}
  \kappa(A)\equiv\norm{A}\cdot\norm{A^{-1}}.
  \label{eq:3-1}
\end{equation}
\end{definition}

显然, 条件数依赖于所用的矩阵范数. 当 $\kappa(A)$ 小时, 称 $A$ 是\newterm{wellconditioned}{良态的}{well-conditioned}; 当 $\kappa(A)$ 大时, 称 $A$ 是\newterm{illconditioned}{病态的}{ill-conditioned}. 注意对任意 $p$ 范数, 有
\[
  1=\norm{I}=\norm{A\cdot A^{-1}}\leqslant\norm{A}\cdot\norm{A^{-1}}=\kappa(A).
\]

设 $\hat{\mathbf{x}}$ 是 $A\mathbf{x}=\mathbf{b}$ 的精确解 $\mathbf{x}$ 的一个近似. \newterm{errorvector}{误差向量}{error vector}定义为
\[
  \mathbf{e}=\mathbf{x}-\hat{\mathbf{x}},
\]
即
\begin{equation}
  \mathbf{x}=\hat{\mathbf{x}}+\mathbf{e}.
  \label{eq:3-2}
\end{equation}
\newterm{absoluteerror}{绝对误差}{absolute error}由 $\norm{\mathbf{e}}=\norm{\mathbf{x}-\hat{\mathbf{x}}}$ 给出 (对任意向量范数). 若 $\mathbf{x}\neq\mathbf{0}$, 则\newterm{relativeerror}{相对误差}{relative error}定义为
$\norm{\mathbf{e}}/\norm{\mathbf{x}}=\norm{\mathbf{x}-\hat{\mathbf{x}}}/\norm{\mathbf{x}}$.

把 \eqref{eq:3-2} 代入 $A\mathbf{x}=\mathbf{b}$, 得
\[
  A(\hat{\mathbf{x}}+\mathbf{e})=A\hat{\mathbf{x}}+A\mathbf{e}=\mathbf{b}.
\]
因此
\[
  A\hat{\mathbf{x}}=\mathbf{b}-A\mathbf{e}=\hat{\mathbf{b}}.
\]
$\hat{\mathbf{x}}$ 是 $A\hat{\mathbf{x}}=\hat{\mathbf{b}}$ 的精确解, 其中 $\hat{\mathbf{b}}$ 是 $\mathbf{b}$ 的扰动向量. 由 $\mathbf{x}=A^{-1}\mathbf{b}$ 与 $\hat{\mathbf{x}}=A^{-1}\hat{\mathbf{b}}$, 有
\begin{equation}
  \norm{\mathbf{x}-\hat{\mathbf{x}}}=\norm{A^{-1}(\mathbf{b}-\hat{\mathbf{b}})}
  \leqslant\norm{A^{-1}}\cdot\norm{\mathbf{b}-\hat{\mathbf{b}}}.
  \label{eq:3-3}
\end{equation}
类似地,
\[
  \norm{\mathbf{b}}=\norm{A\mathbf{x}}\leqslant\norm{A}\cdot\norm{\mathbf{x}},
\]
即
\begin{equation}
  \frac{1}{\norm{\mathbf{x}}}\leqslant\frac{\norm{A}}{\norm{\mathbf{b}}}.
  \label{eq:3-4}
\end{equation}
结合 \eqref{eq:3-3}, \eqref{eq:3-4} 与 \eqref{eq:3-1}, 得到下述定理, 它以 $A$ 的条件数刻画了向量 $\mathbf{b}$ 的扰动对 $A\mathbf{x}=\mathbf{b}$ 的解的影响.

\begin{theorem}
\label{thm:3-7}
设 $\hat{\mathbf{x}}$ 是 $A\mathbf{x}=\mathbf{b}$ 的精确解 $\mathbf{x}$ 的近似解. 则
\[
  \frac{\norm{\mathbf{x}-\hat{\mathbf{x}}}}{\norm{\mathbf{x}}}
  \leqslant\kappa(A)\frac{\norm{\mathbf{b}-\hat{\mathbf{b}}}}{\norm{\mathbf{b}}}.
\]
\end{theorem}

下一个定理则以 $A$ 的条件数刻画了系数矩阵 $A$ 的扰动对 $A\mathbf{x}=\mathbf{b}$ 的解的影响.

\begin{theorem}
\label{thm:3-8}
设 $A$ 为非奇异矩阵, $\widehat{A}$ 是 $A$ 的扰动矩阵, 且
\[
  \norm{A-\widehat{A}}\cdot\norm{A^{-1}}<1.
\]
若 $A\mathbf{x}=\mathbf{b}$ 且 $\widehat{A}\hat{\mathbf{x}}=\hat{\mathbf{b}}$, 其中 $\hat{\mathbf{b}}$ 是 $\mathbf{b}$ 的扰动向量, 则
\[
  \frac{\norm{\mathbf{x}-\hat{\mathbf{x}}}}{\norm{\mathbf{x}}}
  \leqslant\frac{\kappa(A)}{1-\kappa(A)\norm{A-\widehat{A}}/\norm{A}}
  \left(\frac{\norm{A-\widehat{A}}}{\norm{A}}+\frac{\norm{\mathbf{b}-\hat{\mathbf{b}}}}{\norm{\mathbf{b}}}\right).
\]
\end{theorem}

\begin{proof}
令
\[
  E=A-\widehat{A}\qquad\text{且}\qquad\delta\mathbf{b}=\mathbf{b}-\widehat{\mathbf{b}}.
\]
用 $\widehat{A}\hat{\mathbf{x}}=\hat{\mathbf{b}}$ 减去 $A\mathbf{x}=\mathbf{b}$, 得
\[
  A(\mathbf{x}-\hat{\mathbf{x}})=-E\hat{\mathbf{x}}+\delta\mathbf{b}.
\]
进一步有
\[
  \frac{\norm{\mathbf{x}-\hat{\mathbf{x}}}}{\norm{\mathbf{x}}}
  \leqslant\norm{A^{-1}E}\frac{\norm{\hat{\mathbf{x}}}}{\norm{\mathbf{x}}}
  +\norm{A^{-1}}\frac{\norm{A\mathbf{x}}}{\norm{\mathbf{x}}}\frac{\norm{\delta\mathbf{b}}}{\norm{\mathbf{b}}}.
\]
利用
\[
  \norm{\hat{\mathbf{x}}}\leqslant\norm{\hat{\mathbf{x}}-\mathbf{x}}+\norm{\mathbf{x}}
  \qquad\text{且}\qquad
  \norm{A\mathbf{x}}\leqslant\norm{A}\cdot\norm{\mathbf{x}},
\]
便得
\[
  \frac{\norm{\mathbf{x}-\hat{\mathbf{x}}}}{\norm{\mathbf{x}}}
  \leqslant\norm{A^{-1}E}\frac{\norm{\mathbf{x}-\hat{\mathbf{x}}}}{\norm{\mathbf{x}}}
  +\norm{A^{-1}E}+\norm{A^{-1}}\norm{A}\frac{\norm{\delta\mathbf{b}}}{\norm{\mathbf{b}}},
\]
即
\[
  \bigl(1-\norm{A^{-1}E}\bigr)\frac{\norm{\mathbf{x}-\hat{\mathbf{x}}}}{\norm{\mathbf{x}}}
  \leqslant\norm{A^{-1}E}+\kappa(A)\frac{\norm{\delta\mathbf{b}}}{\norm{\mathbf{b}}}.
\]
由于
\[
  \norm{A^{-1}E}\leqslant\norm{A^{-1}}\cdot\norm{E}
  =\norm{A^{-1}}\cdot\norm{A-\widehat{A}}<1,
\]
得到
\[
  \frac{\norm{\mathbf{x}-\hat{\mathbf{x}}}}{\norm{\mathbf{x}}}
  \leqslant\bigl(1-\norm{A^{-1}E}\bigr)^{-1}
  \left(\norm{A^{-1}E}+\kappa(A)\frac{\norm{\delta\mathbf{b}}}{\norm{\mathbf{b}}}\right).
\]
再利用
\[
  \norm{A^{-1}E}\leqslant\norm{A^{-1}}\cdot\norm{E}
  =\kappa(A)\frac{\norm{E}}{\norm{A}},
\]
最终得到
\[
  \frac{\norm{\mathbf{x}-\hat{\mathbf{x}}}}{\norm{\mathbf{x}}}
  \leqslant\frac{\kappa(A)}{1-\kappa(A)\norm{E}/\norm{A}}
  \left(\frac{\norm{E}}{\norm{A}}+\frac{\norm{\delta\mathbf{b}}}{\norm{\mathbf{b}}}\right).
\]
\end{proof}

定理 \ref{thm:3-7} 与定理 \ref{thm:3-8} 以 $A$ 的条件数给出了 $\mathbf{x}$ 的相对误差的上界. 由这两个定理可知, 若 $A$ 是良态的, 即 $\kappa(A)$ 小, 则当 $A$ 与 $\mathbf{b}$ 两者的相对误差都小时, $\mathbf{x}$ 的相对误差也小.

\begin{corollary}
\label{cor:3-2}
设 $\norm{\cdot}$ 为任一满足 $\norm{I}=1$ 的矩阵范数, $A$ 为非奇异矩阵, 且 $A+\widetilde{A}$ 是 $A$ 的扰动矩阵, 满足
\[
  \norm{A^{-1}\widetilde{A}}<1.
\]
则 $A+\widetilde{A}$ 非奇异, 且
\[
  \frac{\norm{(A+\widetilde{A})^{-1}-A^{-1}}}{\norm{A^{-1}}}
  \leqslant\frac{\kappa(A)}{1-\kappa(A)\norm{\widetilde{A}}/\norm{A}}
  \cdot\frac{\norm{\widetilde{A}}}{\norm{A}}.
\]
\end{corollary}

\begin{proof}
先证明
\[
  \norm{(A+\widetilde{A})^{-1}-A^{-1}}
  \leqslant\frac{\norm{\widetilde{A}}\cdot\norm{A^{-1}}^2}{1-\norm{A^{-1}\widetilde{A}}}.
\]
注意
\[
  A+\widetilde{A}=A(I+A^{-1}\widetilde{A})=A\bigl[I-(-A^{-1}\widetilde{A})\bigr].
\]
令 $F=-A^{-1}\widetilde{A}$ 且 $r=\norm{A^{-1}\widetilde{A}}$. 现在
\[
  (A+\widetilde{A})^{-1}=(I-F)^{-1}A^{-1}.
\]
因此, 注意到 $\norm{F}=r<1$ 并利用推论 \ref{cor:3-1}, 有
\[
  \norm{(A+\widetilde{A})^{-1}}=\norm{(I-F)^{-1}A^{-1}}
  \leqslant\norm{(I-F)^{-1}}\cdot\norm{A^{-1}}
  <\frac{\norm{A^{-1}}}{1-\norm{F}}=\frac{\norm{A^{-1}}}{1-r}.
\]
利用恒等式
\[
  B^{-1}=A^{-1}-B^{-1}(B-A)A^{-1},
\]
有
\[
  (A+\widetilde{A})^{-1}-A^{-1}=-(A+\widetilde{A})^{-1}\widetilde{A}A^{-1}.
\]
于是
\[
  \norm{(A+\widetilde{A})^{-1}-A^{-1}}
  \leqslant\norm{A^{-1}}\cdot\norm{\widetilde{A}}\cdot\norm{(A+\widetilde{A})^{-1}}
  \leqslant\frac{\norm{A^{-1}}^2\norm{\widetilde{A}}}{1-r}.
\]
最后得到
\[
  \begin{aligned}
    \frac{\norm{(A+\widetilde{A})^{-1}-A^{-1}}}{\norm{A^{-1}}}
    &\leqslant\frac{\norm{A^{-1}}\cdot\norm{\widetilde{A}}}{1-r}
     \leqslant\frac{\norm{A^{-1}}\cdot\norm{\widetilde{A}}}{1-\norm{A^{-1}}\cdot\norm{\widetilde{A}}}\\
    &\leqslant\frac{\kappa(A)}{1-\kappa(A)\norm{\widetilde{A}}/\norm{A}}
     \cdot\frac{\norm{\widetilde{A}}}{\norm{A}}.
  \end{aligned}
\]
\end{proof}

\begin{mremark}
除 \eqref{eq:3-1} 所定义的条件数之外, 文献 \cite{src:22} 为数值求解偏微分方程提出了\newterm{effectivecondnum}{有效条件数}{effective condition number}. 详细研究见 \cite{src:57}.
\end{mremark}

\section{浮点算术的误差分析}

在计算机中, \newterm{floatingpoint}{浮点数}{floating point number} $f$ 表示为
\[
  f=\pm\omega\times\beta^J,\qquad L\leqslant J\leqslant U,
\]
其中 $\beta$ 是基数, $J$ 是阶, $\omega$ 是尾数. 通常 $\omega$ 具有如下形式:
\[
  \omega=0.d_1d_2\ldots d_t,
\]
其中 $t$ 是 $\omega$ 的长度 (精度), $d_1\neq0$, 且 $0\leqslant d_i<\beta$ 对 $i=2,\ldots,t$ 成立.

令
\[
  \mathcal{F}=\{0\}\cup\{f:f=\pm\omega\times\beta^J,\ 0\leqslant d_i<\beta,\ d_1\neq0,\ L\leqslant J\leqslant U\}.
\]
则 $\mathcal{F}$ 含
\[
  2(\beta-1)\beta^{t-1}(U-L+1)+1
\]
个浮点数. 这些数对称地分布在区间 $[m,M]$ 与 $[-M,-m]$ 中, 其中
\begin{equation}
  m=\beta^{L-1},\qquad M=\beta^U(1-\beta^{-t}).
  \label{eq:3-5}
\end{equation}
我们指出, $\mathcal{F}$ 只是一个有限集, 它不可能包含这两个区间中的所有实数.

令 $fl(x)$ 表示任一实数 $x$ 的浮点数. 则
\[
  fl(x)=0,\qquad\text{当 }x=0.
\]
若 $m\leqslant|x|\leqslant M$, 则按\newterm{rounding}{舍入}{rounding}, $fl(x)$ 是使
\[
  \abs{fl(x)-x}=\min_{f\in\mathcal{F}}|f-x|
\]
成立的 $f$; 而按\newterm{chopping}{截断}{chopping}, $fl(x)$ 是使
\[
  \abs{fl(x)-x}=\min_{|f|\leqslant|x|}|f-x|
\]
成立的 $f$.

例如, 令 $\beta=10$, $t=3$, $L=0$, $U=2$. 考虑 $x=5.45627$ 的浮点表示. 按舍入, 得 $fl(x)=0.546\times10$; 按截断, 得 $fl(x)=0.545\times10$. 下述定理给出浮点表示相对误差的估计.

\begin{theorem}
\label{thm:3-9}
设 $m\leqslant|x|\leqslant M$, 其中 $m$ 与 $M$ 由 \eqref{eq:3-5} 定义. 则
\[
  fl(x)=x(1+\delta),\qquad|\delta|\leqslant\mu_\varepsilon,
\]
其中 $\mu_\varepsilon$ 是\newterm{machineprecision}{机器精度}{machine precision}, 即
\[
  \mu_\varepsilon=
  \begin{cases}
    \dfrac{1}{2}\beta^{1-t}, & \text{按舍入},\\[8pt]
    \beta^{1-t}, & \text{按截断}.
  \end{cases}
\]
\end{theorem}

\begin{proof}
以下假定 $x>0$. 设 $\alpha$ 为整数且满足
\begin{equation}
  \beta^{\alpha-1}\leqslant x<\beta^{\alpha}.
  \label{eq:3-6}
\end{equation}
由于 $[\beta^{\alpha-1},\beta^{\alpha})$ 中浮点数的阶为 $\alpha$, 所有数
\[
  0.d_1d_2\ldots d_t\times\beta^{\alpha}
\]
都以间距 $\beta^{\alpha-t}$ 分布在该区间内. 对舍入误差, 由 \eqref{eq:3-6} 有
\[
  \abs{fl(x)-x}\leqslant\frac{1}{2}\beta^{\alpha-t}
  =\frac{1}{2}\beta^{\alpha-1}\beta^{1-t}\leqslant\frac{1}{2}x\beta^{1-t},
\]
即
\[
  \frac{\abs{fl(x)-x}}{x}\leqslant\frac{1}{2}\beta^{1-t}.
\]
对截断误差, 有
\[
  \abs{fl(x)-x}\leqslant\beta^{\alpha-t}=\beta^{\alpha-1}\beta^{1-t}\leqslant x\beta^{1-t},
\]
即
\[
  \frac{\abs{fl(x)-x}}{x}\leqslant\beta^{1-t}.
\]
证毕.
\end{proof}

下面考虑基本运算的舍入误差. 设 $a,b\in\mathcal{F}$, 而``$\circ$''表示任一基本运算: ``$+,-,\times,\div$''. 由定理 \ref{thm:3-9} 立即得到

\begin{theorem}
\label{thm:3-10}
有
\[
  fl(a\circ b)=(a\circ b)(1+\delta),\qquad|\delta|\leqslant\mu_\varepsilon.
\]
\end{theorem}

\begin{theorem}
\label{thm:3-11}
若 $|\delta_i|\leqslant\mu_\varepsilon$ 且 $n\mu_\varepsilon\leqslant0.01$, 则
\[
  1-1.01n\mu_\varepsilon\leqslant\prod_{i=1}^{n}(1+\delta_i)
  \leqslant1+1.01n\mu_\varepsilon.
\]
\end{theorem}

\begin{proof}
由于 $|\delta_i|\leqslant\mu_\varepsilon$, 有
\begin{equation}
  (1-\mu_\varepsilon)^n\leqslant\prod_{i=1}^{n}(1+\delta_i)\leqslant(1+\mu_\varepsilon)^n.
  \label{eq:3-7}
\end{equation}
对 $(1-\mu_\varepsilon)^n$ 的下界, 利用 $(1-x)^n$ 的 Taylor 展开, 即
\[
  (1-x)^n=1-nx+\frac{n(n-1)}{2}(1-\theta x)^{n-2}x^2,
\]
有
\[
  1-nx\leqslant(1-x)^n.
\]
因此
\begin{equation}
  1-1.01n\mu_\varepsilon\leqslant1-n\mu_\varepsilon\leqslant(1-\mu_\varepsilon)^n.
  \label{eq:3-8}
\end{equation}
现在估计 $(1+\mu_\varepsilon)^n$ 的上界. 利用 $e^x$ 的 Taylor 展开, 有
\[
  e^x=1+x+\frac{x^2}{2!}+\frac{x^3}{3!}+\cdots
  =1+x+\frac{x}{2}\cdot x\cdot\left(1+\frac{x}{3}+\cdots\right).
\]
因此, 当 $0\leqslant x\leqslant0.01$ 时, 利用 $e^{0.01}<2$ 可知
\begin{equation}
  1+x\leqslant e^x\leqslant1+x+\frac{0.01}{2}xe^x\leqslant1+1.01x.
  \label{eq:3-9}
\end{equation}
令 $x=\mu_\varepsilon$, 由 \eqref{eq:3-9} 的左不等式得
\begin{equation}
  (1+\mu_\varepsilon)^n\leqslant e^{n\mu_\varepsilon}.
  \label{eq:3-10}
\end{equation}
令 $x=n\mu_\varepsilon$, 由 \eqref{eq:3-9} 的右不等式得
\begin{equation}
  e^{n\mu_\varepsilon}\leqslant1+1.01n\mu_\varepsilon.
  \label{eq:3-11}
\end{equation}
结合 \eqref{eq:3-10} 与 \eqref{eq:3-11}, 有
\begin{equation}
  (1+\mu_\varepsilon)^n\leqslant1+1.01n\mu_\varepsilon.
  \label{eq:3-12}
\end{equation}
由 \eqref{eq:3-7}, \eqref{eq:3-8} 与 \eqref{eq:3-12}, 证毕.
\end{proof}

考虑下述例子.

\begin{example}
\label{ex:3-1}
给定 $\mathbf{x},\mathbf{y}\in\R^n$ 与 $n\mu_\varepsilon\leqslant0.01$, 估计
\[
  \abs{fl(\mathbf{x}^T\mathbf{y})-\mathbf{x}^T\mathbf{y}}
\]
的上界.

令
\[
  S_k=fl\left(\sum_{i=1}^{k}x_iy_i\right).
\]
由定理 \ref{thm:3-10}, 有
\[
  S_1=x_1y_1(1+\gamma_1),\qquad|\gamma_1|\leqslant\mu_\varepsilon,
\]
以及
\[
  \begin{aligned}
    S_k&=fl\bigl(S_{k-1}+fl(x_ky_k)\bigr)\\
       &=\bigl[S_{k-1}+x_ky_k(1+\gamma_k)\bigr](1+\delta_k),
       \qquad|\delta_k|,|\gamma_k|\leqslant\mu_\varepsilon.
  \end{aligned}
\]
因此
\[
  \begin{aligned}
    fl(\mathbf{x}^T\mathbf{y})
    &=S_n=\sum_{i=1}^{n}x_iy_i(1+\gamma_i)\prod_{j=i}^{n}(1+\delta_j)\\
    &=\sum_{i=1}^{n}(1+\epsilon_i)x_iy_i,
  \end{aligned}
\]
其中
\[
  1+\epsilon_i=(1+\gamma_i)\prod_{j=i}^{n}(1+\delta_j),
\]
且 $\delta_1=0$. 于是, 若 $n\mu_\varepsilon\leqslant0.01$, 由定理 \ref{thm:3-11} 得
\[
  \abs{fl(\mathbf{x}^T\mathbf{y})-\mathbf{x}^T\mathbf{y}}
  \leqslant\sum_{i=1}^{n}|\epsilon_i|\cdot|x_iy_i|
  \leqslant1.01n\mu_\varepsilon\sum_{i=1}^{n}|x_iy_i|.
\]
\end{example}

在结束本节之前, 简要讨论基本矩阵运算的浮点分析. 回忆以下记号:
\[
  |E|=[|e_{ij}|],
\]
其中 $E=[e_{ij}]\in\R^{n\times n}$, 以及
\[
  |E|\leqslant|F|\iff|e_{ij}|\leqslant|f_{ij}|
\]
对 $i,j=1,2,\ldots,n$ 成立. 设 $A,B\in\R^{n\times n}$ 的元素均属于 $\mathcal{F}$, $\alpha\in\mathcal{F}$. 由定理 \ref{thm:3-10}, 有
\[
  fl(\alpha A)=\alpha A+E,\qquad|E|\leqslant\mu_\varepsilon|\alpha A|,
\]
以及
\[
  fl(A+B)=(A+B)+E,\qquad|E|\leqslant\mu_\varepsilon|A+B|.
\]
由例 \ref{ex:3-1} 还有
\[
  fl(AB)=AB+E,\qquad|E|\leqslant1.01n\mu_\varepsilon|A|\cdot|B|.
\]
注意 $|A|\cdot|B|$ 可能远大于 $|AB|$. 因此 $AB$ 的相对误差未必小. 关于浮点算术数值计算的详细研究, 参见 \cite{src:65}.

\section{部分主元法的误差分析}

我们证明: 若用带部分主元选取的高斯消去法求解 $A\mathbf{x}=\mathbf{b}$, 则计算解 $\tilde{\mathbf{x}}$ 满足
\[
  (A+E)\tilde{\mathbf{x}}=\mathbf{b},
\]
其中 $E$ 是一个误差矩阵. 同时也给出 $E$ 的上界. 先研究 $A$ 的 $LU$ 分解的舍入误差.

\begin{lemma}
\label{lem:3-1}
设 $A\in\R^{n\times n}$ 的元素为浮点数. 假定 $A$ 具有 $LU$ 分解, 且 $6n\mu_\varepsilon\leqslant1$, 其中 $\mu_\varepsilon$ 是机器精度. 则用高斯消去法得到
\[
  \widetilde{L}\widetilde{U}=A+E,
\]
其中
\[
  |E|\leqslant3n\mu_\varepsilon\bigl(|A|+|\widetilde{L}|\cdot|\widetilde{U}|\bigr).
\]
\end{lemma}

\begin{proof}
对 $n$ 作归纳. 显然, $n=1$ 时引理成立. 设引理对 $n-1$ 成立. 现在考虑矩阵 $A\in\R^{n\times n}$:
\[
  A=
  \begin{bmatrix}
    \alpha & \mathbf{w}^T\\
    \mathbf{u} & A_1
  \end{bmatrix},
\]
其中 $A_1\in\R^{(n-1)\times(n-1)}$. 在高斯消去的第一步, 应计算向量 $\tilde{\mathbf{v}}_1=fl(\mathbf{u}/\alpha)$, 并把矩阵 $A_1$ 修改为
\[
  \widetilde{A}_1=fl\bigl(A_1-fl(\tilde{\mathbf{v}}_1\mathbf{w}^T)\bigr).
\]
由定理 \ref{thm:3-10}, 有
\begin{equation}
  \tilde{\mathbf{v}}_1=\mathbf{u}/\alpha+\mathbf{f},
  \qquad|\mathbf{f}|\leqslant\frac{\mu_\varepsilon}{|\alpha|}|\mathbf{u}|
  \label{eq:3-13}
\end{equation}
以及
\begin{equation}
  \widetilde{A}_1=A_1-\tilde{\mathbf{v}}_1\mathbf{w}^T+F,
  \qquad|F|\leqslant(2+\mu_\varepsilon)\mu_\varepsilon
  \bigl(|A_1|+|\tilde{\mathbf{v}}_1|\cdot|\mathbf{w}|^T\bigr).
  \label{eq:3-14}
\end{equation}
对 $\widetilde{A}_1$, 利用归纳假设, 得到 $LU$ 分解: 单位下三角矩阵 $\widetilde{L}_1$ 与上三角矩阵 $\widetilde{U}_1$ 使得
\[
  \widetilde{L}_1\widetilde{U}_1=\widetilde{A}_1+E_1,
\]
其中
\[
  |E_1|\leqslant3(n-1)\mu_\varepsilon
  \bigl(|\widetilde{A}_1|+|\widetilde{L}_1|\cdot|\widetilde{U}_1|\bigr).
\]
于是有
\[
  \widetilde{L}\widetilde{U}
  =
  \begin{bmatrix}
    1 & \mathbf{0}\\
    \tilde{\mathbf{v}}_1 & \widetilde{L}_1
  \end{bmatrix}
  \begin{bmatrix}
    \alpha & \mathbf{w}^T\\
    \mathbf{0} & \widetilde{U}_1
  \end{bmatrix}
  =A+E,
\]
其中
\[
  E=
  \begin{bmatrix}
    \mathbf{0} & \mathbf{0}\\
    \alpha\mathbf{f} & E_1+F
  \end{bmatrix}.
\]
利用 \eqref{eq:3-14}, 得到
\[
  |\widetilde{A}_1|\leqslant(1+2\mu_\varepsilon+\mu_\varepsilon^2)
  \bigl(|A_1|+|\tilde{\mathbf{v}}_1|\cdot|\mathbf{w}|^T\bigr).
\]
因此, 利用条件 $6n\mu_\varepsilon\leqslant1$, 有
\[
  \begin{aligned}
    |E_1+F|
    &\leqslant|E_1|+|F|\\
    &\leqslant3(n-1)\mu_\varepsilon\bigl(|\widetilde{A}_1|+|\widetilde{L}_1|\cdot|\widetilde{U}_1|\bigr)
      +(2+\mu_\varepsilon)\mu_\varepsilon\bigl(|A_1|+|\tilde{\mathbf{v}}_1|\cdot|\mathbf{w}|^T\bigr)\\
    &\leqslant3(n-1)\mu_\varepsilon\Bigl(\bigl(1+2\mu_\varepsilon+\mu_\varepsilon^2\bigr)
      \bigl(|A_1|+|\tilde{\mathbf{v}}_1|\cdot|\mathbf{w}|^T\bigr)
      +|\widetilde{L}_1|\cdot|\widetilde{U}_1|\Bigr)\\
    &\qquad+(2+\mu_\varepsilon)\mu_\varepsilon\bigl(|A_1|+|\tilde{\mathbf{v}}_1|\cdot|\mathbf{w}|^T\bigr)\\
    &\leqslant\mu_\varepsilon\Bigl(3n-1+\bigl[6n+3(n-1)\mu_\varepsilon-5\bigr]\mu_\varepsilon\Bigr)
      \bigl(|A_1|+|\tilde{\mathbf{v}}_1|\cdot|\mathbf{w}|^T\bigr)\\
    &\qquad+3(n-1)\mu_\varepsilon\bigl(|\widetilde{L}_1|\cdot|\widetilde{U}_1|\bigr)\\
    &\leqslant3n\mu_\varepsilon\bigl(|A_1|+|\tilde{\mathbf{v}}_1|\cdot|\mathbf{w}|^T
      +|\widetilde{L}_1|\cdot|\widetilde{U}_1|\bigr).
  \end{aligned}
\]
结合 \eqref{eq:3-13}, 得到
\[
  \begin{aligned}
    |E|
    &=
    \begin{bmatrix}
      \mathbf{0} & \mathbf{0}\\
      |\alpha|\,|\mathbf{f}| & |E_1+F|
    \end{bmatrix}\\
    &\leqslant3n\mu_\varepsilon
    \begin{bmatrix}
      \mathbf{0} & \mathbf{0}\\
      |\mathbf{u}| & |A_1|+|\tilde{\mathbf{v}}_1|\cdot|\mathbf{w}|^T
        +|\widetilde{L}_1|\cdot|\widetilde{U}_1|
    \end{bmatrix}\\
    &\leqslant3n\mu_\varepsilon\left(
    \begin{bmatrix}
      |\alpha| & |\mathbf{w}|^T\\
      |\mathbf{u}| & |A_1|
    \end{bmatrix}
    +
    \begin{bmatrix}
      1 & \mathbf{0}\\
      |\tilde{\mathbf{v}}_1| & |\widetilde{L}_1|
    \end{bmatrix}
    \begin{bmatrix}
      |\alpha| & |\mathbf{w}|^T\\
      \mathbf{0} & |\widetilde{U}_1|
    \end{bmatrix}\right)\\
    &=3n\mu_\varepsilon\bigl(|A|+|\widetilde{L}|\cdot|\widetilde{U}|\bigr).
  \end{aligned}
\]
证毕.
\end{proof}

\begin{corollary}
\label{cor:3-3}
设 $A\in\R^{n\times n}$ 非奇异且元素为浮点数, $6n\mu_\varepsilon\leqslant1$. 假定用带部分主元选取的高斯消去法得到
\[
  \widetilde{L}\widetilde{U}=PA+E,
\]
其中 $\widetilde{L}=[l_{ij}]$ 是满足 $|l_{ij}|\leqslant1$ 的单位下三角矩阵, $\widetilde{U}$ 是上三角矩阵, $P$ 是置换矩阵. 则 $E$ 满足下述不等式:
\[
  |E|\leqslant3n\mu_\varepsilon\bigl(|PA|+|\widetilde{L}|\cdot|\widetilde{U}|\bigr).
\]
\end{corollary}

在得到 $A$ 的 $LU$ 分解之后, 求解 $A\mathbf{x}=\mathbf{b}$ 的问题就化为求解如下两个三角方程组:
\[
  \widetilde{L}\mathbf{y}=P\mathbf{b},\qquad\widetilde{U}\mathbf{x}=\mathbf{y}.
\]
因此, 还需估计求解三角方程组的舍入误差.

\begin{lemma}
\label{lem:3-2}
设 $S\in\R^{n\times n}$ 为非奇异三角矩阵, 元素为浮点数, 且 $1.01n\mu_\varepsilon\leqslant0.01$. 用 \ref{subsec:2-1-1} 小节提出的方法求解 $S\mathbf{x}=\mathbf{b}$, 则所得计算解 $\tilde{\mathbf{x}}$ 满足
\[
  (S+H)\tilde{\mathbf{x}}=\mathbf{b},
\]
其中
\[
  |H|\leqslant1.01n\mu_\varepsilon|S|.
\]
\end{lemma}

\begin{proof}
对 $n$ 作归纳. 不失一般性, 设 $S=L$ 为下三角矩阵. 显然, $n=1$ 时引理成立. 设引理对 $n-1$ 成立. 现在考虑下三角矩阵 $L\in\R^{n\times n}$. 设 $\tilde{\mathbf{x}}$ 是 $L\mathbf{x}=\mathbf{b}$ 的计算解, 并把 $L$, $\mathbf{b}$ 与 $\tilde{\mathbf{x}}$ 分块如下:
\[
  L=
  \begin{bmatrix}
    l_{11} & \mathbf{0}\\
    \mathbf{v}_1 & L_1
  \end{bmatrix},
  \qquad
  \mathbf{b}=
  \begin{bmatrix}
    b_1\\ \mathbf{c}
  \end{bmatrix},
  \qquad
  \tilde{\mathbf{x}}=
  \begin{bmatrix}
    \tilde{x}_1\\ \tilde{\mathbf{y}}
  \end{bmatrix},
\]
其中 $\mathbf{c},\tilde{\mathbf{y}}\in\R^{n-1}$, $L_1\in\R^{(n-1)\times(n-1)}$. 由定理 \ref{thm:3-10}, 有
\begin{equation}
  \tilde{x}_1=fl(b_1/l_{11})=\frac{b_1}{l_{11}(1+\delta_1)},
  \qquad|\delta_1|\leqslant\mu_\varepsilon.
  \label{eq:3-15}
\end{equation}
注意 $\tilde{\mathbf{y}}$ 是 $(n-1)$ 阶方程组
\[
  L_1\mathbf{y}=fl(\mathbf{c}-\tilde{x}_1\mathbf{v}_1)
\]
的计算解. 由归纳假设, 有
\[
  (L_1+H_1)\tilde{\mathbf{y}}=fl(\mathbf{c}-\tilde{x}_1\mathbf{v}_1),
\]
其中
\begin{equation}
  |H_1|\leqslant1.01(n-1)\mu_\varepsilon|L_1|.
  \label{eq:3-16}
\end{equation}
再次由定理 \ref{thm:3-10}, 得到
\[
  fl(\mathbf{c}-\tilde{x}_1\mathbf{v}_1)
  =fl\bigl(\mathbf{c}-fl(\tilde{x}_1\mathbf{v}_1)\bigr)
  =(I+D_\gamma)^{-1}(\mathbf{c}-\tilde{x}_1\mathbf{v}_1-\tilde{x}_1D_\delta\mathbf{v}_1),
\]
其中
\[
  D_\gamma=\diag(\gamma_2,\ldots,\gamma_n),\qquad
  D_\delta=\diag(\delta_2,\ldots,\delta_n),
\]
且
\[
  |\gamma_i|,\,|\delta_i|\leqslant\mu_\varepsilon,\qquad i=2,\ldots,n.
\]
因此
\[
  \tilde{x}_1\mathbf{v}_1+\tilde{x}_1D_\delta\mathbf{v}_1
  +(I+D_\gamma)(L_1+H_1)\tilde{\mathbf{y}}=\mathbf{c},
\]
从而
\[
  (L+H)\tilde{\mathbf{x}}=\mathbf{b},
\]
其中
\[
  H=
  \begin{bmatrix}
    \delta_1l_{11} & \mathbf{0}\\
    D_\delta\mathbf{v}_1 & H_1+D_\gamma(L_1+H_1)
  \end{bmatrix}.
\]
利用 \eqref{eq:3-15}, \eqref{eq:3-16} 与条件 $1.01n\mu_\varepsilon\leqslant0.01$, 有
\[
  \begin{aligned}
    |H|
    &\leqslant
    \begin{bmatrix}
      |\delta_1|\cdot|l_{11}| & \mathbf{0}\\
      |D_\delta|\cdot|\mathbf{v}_1| & |H_1|+|D_\gamma|\bigl(|L_1|+|H_1|\bigr)
    \end{bmatrix}\\
    &\leqslant
    \begin{bmatrix}
      \mu_\varepsilon|l_{11}| & \mathbf{0}\\
      \mu_\varepsilon|\mathbf{v}_1| & |H_1|+\mu_\varepsilon\bigl(|L_1|+|H_1|\bigr)
    \end{bmatrix}\\
    &\leqslant\mu_\varepsilon
    \begin{bmatrix}
      |l_{11}| & \mathbf{0}\\
      |\mathbf{v}_1| & \bigl[1.01(n-1)+1+1.01(n-1)\mu_\varepsilon\bigr]|L_1|
    \end{bmatrix}\\
    &\leqslant1.01n\mu_\varepsilon|L|.
  \end{aligned}
\]
证毕.
\end{proof}

于是得到本节的主要定理.

\begin{theorem}
\label{thm:3-12}
设 $A\in\R^{n\times n}$ 为非奇异矩阵, 元素为浮点数, 且 $1.01n\mu_\varepsilon\leqslant0.01$. 若用带部分主元选取的高斯消去法求解 $A\mathbf{x}=\mathbf{b}$, 则所得计算解 $\tilde{\mathbf{x}}$ 满足
\[
  (A+\delta A)\tilde{\mathbf{x}}=\mathbf{b},
\]
其中
\begin{equation}
  \norm{\delta A}_\infty\leqslant\mu_\varepsilon\bigl(3n+5.04n^3\rho\bigr)\norm{A}_\infty,
  \label{eq:3-17}
\end{equation}
而\newterm{growthfactor}{增长因子}{growth factor}为
\[
  \rho\equiv\frac{1}{\norm{A}_\infty}\max_{i,j,k}\abs{a_{ij}^{(k)}}.
\]
\end{theorem}

\begin{proof}
利用带部分主元选取的高斯消去法, 得到如下两个三角方程组:
\[
  \widetilde{L}\mathbf{y}=P\mathbf{b},\qquad\widetilde{U}\mathbf{x}=\mathbf{y}.
\]
由引理 \ref{lem:3-2}, 计算解 $\tilde{\mathbf{x}}$ 应满足
\[
  (\widetilde{L}+F)(\widetilde{U}+G)\tilde{\mathbf{x}}=P\mathbf{b},
\]
即
\begin{equation}
  \bigl(\widetilde{L}\widetilde{U}+F\widetilde{U}+\widetilde{L}G+FG\bigr)\tilde{\mathbf{x}}=P\mathbf{b},
  \label{eq:3-18}
\end{equation}
其中
\begin{equation}
  |F|\leqslant1.01n\mu_\varepsilon|\widetilde{L}|,\qquad
  |G|\leqslant1.01n\mu_\varepsilon|\widetilde{U}|.
  \label{eq:3-19}
\end{equation}
把 $\widetilde{L}\widetilde{U}=PA+E$ 代入 \eqref{eq:3-18}, 得
\[
  (A+\delta A)\tilde{\mathbf{x}}=\mathbf{b},
\]
其中
\[
  \delta A=P^T\bigl(E+F\widetilde{U}+\widetilde{L}G+FG\bigr).
\]
利用 \eqref{eq:3-19}, 推论 \ref{cor:3-3} 与条件 $1.01n\mu_\varepsilon\leqslant0.01$, 有
\begin{equation}
  \begin{aligned}
    |\delta A|
    &\leqslant P^T\bigl(3n\mu_\varepsilon|PA|+(3n+2.04n)\mu_\varepsilon
      |\widetilde{L}|\cdot|\widetilde{U}|\bigr)\\
    &=n\mu_\varepsilon P^T\bigl(3|PA|+5.04|\widetilde{L}|\cdot|\widetilde{U}|\bigr).
  \end{aligned}
  \label{eq:3-20}
\end{equation}
再次由推论 \ref{cor:3-3}, $\widetilde{L}$ 中各元素的绝对值均不超过 $1$. 因此
\begin{equation}
  \norm{\widetilde{L}}_\infty\leqslant n.
  \label{eq:3-21}
\end{equation}
现在定义
\[
  \rho\equiv\frac{1}{\norm{A}_\infty}\max_{i,j,k}\abs{a_{ij}^{(k)}},
\]
于是有
\begin{equation}
  \norm{\widetilde{U}}_\infty\leqslant n\rho\norm{A}_\infty.
  \label{eq:3-22}
\end{equation}
把 \eqref{eq:3-21} 与 \eqref{eq:3-22} 代入 \eqref{eq:3-20}, 即得 \eqref{eq:3-17}. 证毕.
\end{proof}

\begin{mremark}
与由给定数据引起的初始误差相比, $\norm{\delta A}_\infty$ 通常非常小. 因此, 带部分主元选取的高斯消去法是数值稳定的. 带完全主元选取的高斯消去法中的增长因子可在 \cite{src:33,src:38} 中找到.
\end{mremark}

\begin{exercises}
  \item 设
  \[
    A=
    \begin{bmatrix}
      1 & 0.999999\\
      0.999999 & 1
    \end{bmatrix}.
  \]
  计算 $A^{-1}$, $\det(A)$ 以及 $A$ 的条件数.

  \item 证明 $\norm{AB}_F\leqslant\norm{A}_2\norm{B}_F$ 与 $\norm{AB}_F\leqslant\norm{A}_F\norm{B}_2$.

  \item 证明: 对任意方阵 $A$, 有 $\norm{A}_2^2\leqslant\norm{A}_1\norm{A}_\infty$.

  \item 证明
  \[
    \norm{
    \begin{bmatrix}
      A_{11} & 0\\
      0 & A_{22}
    \end{bmatrix}}_2
    \leqslant
    \norm{
    \begin{bmatrix}
      A_{11} & A_{12}\\
      A_{21} & A_{22}
    \end{bmatrix}}_2.
  \]

  \item 设 $A$ 非奇异. 证明
  \[
    \norm{A^{-1}}_2^{-1}=\min_{\norm{\mathbf{x}}_2=1}\norm{A\mathbf{x}}_2.
  \]

  \item 证明: 若 $S$ 为实矩阵且 $S=-S^T$, 则 $I-S$ 非奇异, 且矩阵
  \[
    (I-S)^{-1}(I+S)
  \]
  是正交矩阵. 这个矩阵称为 $S$ 的 Cayley 变换.

  \item 证明: 若 $A$ 与 $A+E$ 均非奇异, 则
  \[
    \frac{\norm{(A+E)^{-1}-A^{-1}}}{\norm{A^{-1}}}
    \leqslant\norm{(A+E)^{-1}}\cdot\norm{E}.
  \]

  \item 设 $A\in\R^{n\times n}$ 非奇异, 且 $\mathbf{x},\mathbf{y},\mathbf{z}\in\R^n$ 满足 $A\mathbf{x}=\mathbf{b}$ 与 $A\mathbf{y}=\mathbf{b}+\mathbf{z}$. 证明
  \[
    \frac{\norm{\mathbf{z}}_2}{\norm{A}_2}
    \leqslant\norm{\mathbf{x}-\mathbf{y}}_2
    \leqslant\norm{A^{-1}}_2\norm{\mathbf{z}}_2.
  \]

  \item 设 $A=[a_{ij}]$ 为 $m\times n$ 矩阵. 定义
  \[
    \triplenorm{A}_l=\max_{i,j}\abs{a_{ij}}.
  \]
  $\triplenorm{\cdot}_l$ 是矩阵范数吗? 说明理由.

  \item 证明: 若 $X\in\C^{n\times n}$ 非奇异, 则 $\norm{A}_X=\norm{X^{-1}AX}_2$ 定义了一个矩阵范数.

  \item 设 $A=LDL^T\in\R^{n\times n}$ 为对称正定矩阵, 且
  \[
    D=\diag(d_{11},\ldots,d_{nn}).
  \]
  证明
  \[
    \kappa_2(A)\geqslant\frac{\max_i\{d_{ii}\}}{\min_i\{d_{ii}\}}.
  \]

  \item 验证
  \[
    \norm{\mathbf{x}\mathbf{y}^*}_F=\norm{\mathbf{x}\mathbf{y}^*}_2
    =\norm{\mathbf{x}}_2\norm{\mathbf{y}}_2
  \]
  对任意 $\mathbf{x},\mathbf{y}\in\C^n$ 成立.

  \item 证明: 若 $\mathbf{0}\neq\mathbf{v}\in\R^n$ 且 $E\in\R^{n\times n}$, 则
  \[
    \norm{E\left(I-\frac{\mathbf{v}\mathbf{v}^T}{\mathbf{v}^T\mathbf{v}}\right)}_F^2
    =\norm{E}_F^2-\frac{\norm{E\mathbf{v}}_2^2}{\mathbf{v}^T\mathbf{v}}.
  \]

  \item 设 $A=[a_{ij}]\in\R^{n\times n}$, $\mathbf{x}=[x_1,\ldots,x_n]^T\in\R^n$. 证明
  \[
    |A|\cdot|\mathbf{x}|\leqslant\tau\norm{A}_\infty|\mathbf{x}|,
  \]
  其中 $\tau=\max_i|x_i|/\min_i|x_i|$.

  \item 证明 Sherman-Morrison-Woodbury 公式. 设 $U,V$ 为 $n\times k$ 矩阵, $k\leqslant n$, $A$ 为 $n\times n$ 矩阵. 则
  \[
    T=I+V^TA^{-1}U
  \]
  非奇异当且仅当 $A+UV^T$ 非奇异. 此时有
  \[
    (A+UV^T)^{-1}=A^{-1}-A^{-1}UT^{-1}V^TA^{-1}.
  \]
\end{exercises}

% chapters/ch04.tex —— 第 4 章：最小二乘问题
% 源书：Chapter 4 Least Squares Problems，书页 46–59（物理页 60–73）
% 公式编号 (4.1)–(4.8)；定理 4.1–4.9；定义 4.1–4.3；推论 4.1–4.2；算法 4.1–4.2；习题 1–15
\chapter{最小二乘问题}
\label{ch:leastsquares}

本章研究线性最小二乘问题:
\[
  \min_{\mathbf{y}\in\R^n}\norm{A\mathbf{y}-\mathbf{b}}_2,
\]
其中数据矩阵 $A\in\R^{m\times n}$ ($m\geqslant n$) 与观测向量 $\mathbf{b}\in\R^m$ 是给定的. 为构造求解此类问题的高效算法, 我们将介绍一些著名的正交变换与 QR 分解. 关于最小二乘问题的文献, 可参见 \cite{src:9,src:11,src:28,src:36,src:88,src:92,src:93,src:94}.

\section{最小二乘问题}
\label{sec:4-1}

实践中, 若给定 $m$ 个点 $t_1,t_2,\ldots,t_m$ 以及这些点上的数据 $y_1,y_2,\ldots,y_m$, 并在这些点上定义函数 $\phi_1(t),\phi_2(t),\ldots,\phi_n(t)$, 我们便试图找出由
\[
  f(\mathbf{x},t)\equiv\sum_{j=1}^{n}x_j\phi_j(t)
\]
定义的函数 $f(\mathbf{x},t)$, 其中 $\mathbf{x}=[x_1,\ldots,x_n]^T$, 使得由
\[
  r_i(\mathbf{x})\equiv y_i-f(\mathbf{x},t_i)
  =y_i-\sum_{j=1}^{n}x_j\phi_j(t_i),
  \qquad i=1,2,\ldots,m
\]
定义的\newterm{residual}{残差}{residual}尽可能小. 写成矩阵形式, 有
\[
  \mathbf{r}(\mathbf{x})=\mathbf{b}-A\mathbf{x},
\]
其中
\[
  A=
  \begin{bmatrix}
    \phi_1(t_1) & \cdots & \phi_n(t_1)\\
    \vdots      &        & \vdots\\
    \phi_1(t_m) & \cdots & \phi_n(t_m)
  \end{bmatrix},
\]
以及
\[
  \mathbf{b}=[y_1,\ldots,y_m]^T,\qquad
  \mathbf{r}(\mathbf{x})=[r_1(\mathbf{x}),\ldots,r_m(\mathbf{x})]^T.
\]

当 $m=n$ 时, 我们可以要求 $\mathbf{r}(\mathbf{x})=\mathbf{0}$, 此时 $\mathbf{x}$ 可通过求解方程组 $A\mathbf{x}=\mathbf{b}$ 得到. 当 $m>n$ 时, 我们要求 $\mathbf{r}(\mathbf{x})$ 在范数 $\norm{\cdot}_2$ 下达到最小. 为此引入如下最小二乘问题的定义.

\begin{definition}
\label{def:4-1}
设 $A\in\R^{m\times n}$, $\mathbf{b}\in\R^m$. 求 $\mathbf{x}\in\R^n$ 使得
\begin{equation}
  \norm{\mathbf{b}-A\mathbf{x}}_2=\norm{\mathbf{r}(\mathbf{x})}_2
  =\min_{\mathbf{y}\in\R^n}\norm{\mathbf{r}(\mathbf{y})}_2
  =\min_{\mathbf{y}\in\R^n}\norm{\mathbf{b}-A\mathbf{y}}_2.
  \label{eq:4-1}
\end{equation}
该问题称为\newterm{leastsquares}{最小二乘问题}{least squares problem}, $\mathbf{r}(\mathbf{x})$ 称为残差.
\end{definition}

以下只考虑
\[
  \rank(A)=n<m
\]
的情形. 先研究如下方程的解 $\mathbf{x}$:
\begin{equation}
  A\mathbf{x}=\mathbf{b},\qquad A\in\R^{m\times n}.
  \label{eq:4-2}
\end{equation}

矩阵 $A$ 的\newterm{range}{值域}{range}定义为
\[
  \mathcal{R}(A)\equiv\{\mathbf{y}\in\R^m:\mathbf{y}=A\mathbf{x},\ \mathbf{x}\in\R^n\}.
\]
容易看出
\[
  \mathcal{R}(A)=\spn\{\mathbf{a}_1,\ldots,\mathbf{a}_n\},
\]
其中 $\mathbf{a}_i\ (i=1,\ldots,n)$ 是 $A$ 的列向量. $A$ 的\newterm{nullspace}{零空间}{nullspace}定义为
\[
  \mathcal{N}(A)\equiv\{\mathbf{x}\in\R^n:A\mathbf{x}=\mathbf{0}\}.
\]
$\mathcal{N}(A)$ 的维数称为 $A$ 的\newterm{nullity}{零度}{nullity}, 记为 $\operatorname{nullity}(A)$. 子空间 $\mathcal{S}\subset\R^n$ 的\newterm{orthcomplement}{正交补}{orthogonal complement}定义为
\[
  \mathcal{S}^{\perp}\equiv\{\mathbf{y}\in\R^n:
  \mathbf{y}^T\mathbf{x}=0,\ \text{对一切}\ \mathbf{x}\in\mathcal{S}\}.
\]

关于 \eqref{eq:4-2}, 我们有如下定理.

\begin{theorem}
\label{thm:4-1}
方程 \eqref{eq:4-2} 有解 $\iff\rank(A)=\rank([A,\mathbf{b}])$.
\end{theorem}

\begin{theorem}
\label{thm:4-2}
设 $\mathbf{x}$ 是 \eqref{eq:4-2} 的一个特解, 则 \eqref{eq:4-2} 的解集为 $\mathbf{x}+\mathcal{N}(A)$.
\end{theorem}

\begin{corollary}
\label{cor:4-1}
假设方程 \eqref{eq:4-2} 有解, 则解唯一 $\iff$ $\operatorname{nullity}(A)=0$.
\end{corollary}

关于 \eqref{eq:4-1} 的解, 我们有如下基本定理.

\begin{theorem}
\label{thm:4-3}
最小二乘问题 \eqref{eq:4-1} 总有解, 且解唯一当且仅当 $\operatorname{nullity}(A)=0$.
\end{theorem}

\begin{proof}
由于
\[
  \R^m=\mathcal{R}(A)\oplus\mathcal{R}(A)^{\perp},
\]
向量 $\mathbf{b}$ 可唯一地表示为
\[
  \mathbf{b}=\mathbf{b}_1+\mathbf{b}_2,
\]
其中 $\mathbf{b}_1\in\mathcal{R}(A)$, $\mathbf{b}_2\in\mathcal{R}(A)^{\perp}$. 对任意 $\mathbf{x}\in\R^n$, 由于 $\mathbf{b}_1-A\mathbf{x}\in\mathcal{R}(A)$ 且与 $\mathbf{b}_2$ 正交, 故有
\[
  \begin{aligned}
    \norm{\mathbf{r}(\mathbf{x})}_2^2
    &=\norm{\mathbf{b}-A\mathbf{x}}_2^2
     =\norm{(\mathbf{b}_1-A\mathbf{x})+\mathbf{b}_2}_2^2\\
    &=\norm{\mathbf{b}_1-A\mathbf{x}}_2^2+\norm{\mathbf{b}_2}_2^2.
  \end{aligned}
\]
注意 $\norm{\mathbf{r}(\mathbf{x})}_2^2$ 达到最小当且仅当 $\norm{\mathbf{b}_1-A\mathbf{x}}_2^2$ 达到最小. 由于 $\mathbf{b}_1\in\mathcal{R}(A)$, $\norm{\mathbf{r}(\mathbf{x})}_2^2$ 达到最小当且仅当
\[
  A\mathbf{x}=\mathbf{b}_1,
\]
即
\[
  \norm{\mathbf{b}_1-A\mathbf{x}}_2^2=0.
\]
于是由推论 \ref{cor:4-1} 知, $A\mathbf{x}=\mathbf{b}_1$ 的解唯一 (即 \eqref{eq:4-1} 的解唯一) 当且仅当
\[
  \operatorname{nullity}(A)=0. \qedhere
\]
\end{proof}

令
\[
  \mathcal{X}=\{\mathbf{x}\in\R^n:\mathbf{x}\ \text{是 \eqref{eq:4-1} 的解}\}.
\]
我们有

\begin{theorem}
\label{thm:4-4}
向量 $\mathbf{x}\in\mathcal{X}$ 当且仅当
\begin{equation}
  A^TA\mathbf{x}=A^T\mathbf{b}.
  \label{eq:4-3}
\end{equation}
\end{theorem}

\begin{proof}
设 $\mathbf{x}\in\mathcal{X}$. 由定理 \ref{thm:4-3} 知 $A\mathbf{x}=\mathbf{b}_1$, 其中 $\mathbf{b}_1\in\mathcal{R}(A)$, 且
\[
  \mathbf{r}(\mathbf{x})=\mathbf{b}-A\mathbf{x}=\mathbf{b}-\mathbf{b}_1
  =\mathbf{b}_2\in\mathcal{R}(A)^{\perp}.
\]
因此
\[
  A^T\mathbf{r}(\mathbf{x})=A^T\mathbf{b}_2=\mathbf{0}.
\]
把 $\mathbf{r}(\mathbf{x})=\mathbf{b}-A\mathbf{x}$ 代入 $A^T\mathbf{r}(\mathbf{x})=\mathbf{0}$, 即得 \eqref{eq:4-3}.

反之, 设 $\mathbf{x}\in\R^n$ 满足
\[
  A^TA\mathbf{x}=A^T\mathbf{b},
\]
则对任意 $\mathbf{y}\in\R^n$, 有
\[
  \begin{aligned}
    \norm{\mathbf{b}-A(\mathbf{x}+\mathbf{y})}_2^2
    &=\norm{\mathbf{b}-A\mathbf{x}}_2^2
     -2\mathbf{y}^TA^T(\mathbf{b}-A\mathbf{x})+\norm{A\mathbf{y}}_2^2\\
    &=\norm{\mathbf{b}-A\mathbf{x}}_2^2+\norm{A\mathbf{y}}_2^2\\
    &\geqslant\norm{\mathbf{b}-A\mathbf{x}}_2^2.
  \end{aligned}
\]
于是 $\mathbf{x}\in\mathcal{X}$. \qedhere
\end{proof}

因此, 求解最小二乘问题有如下算法:
\begin{parenlist}
  \item 计算 $C=A^TA$ 与 $\mathbf{d}=A^T\mathbf{b}$.
  \item 求 $C=LL^T$ 的 Cholesky 分解.
  \item 求解三角线性方程组: $L\mathbf{y}=\mathbf{d}$ 与 $L^T\mathbf{x}=\mathbf{y}$.
\end{parenlist}

我们指出, 计算 $A^TA$ 的运算量通常为 $\mathcal{O}(n^2m)$, 并且 $A$ 的某些信息可能丢失. 例如, 考虑
\[
  A=
  \begin{bmatrix}
    1 & 1 & 1\\
    \epsilon & 0 & 0\\
    0 & \epsilon & 0\\
    0 & 0 & \epsilon
  \end{bmatrix}.
\]
有
\[
  A^TA=
  \begin{bmatrix}
    1+\epsilon^2 & 1 & 1\\
    1 & 1+\epsilon^2 & 1\\
    1 & 1 & 1+\epsilon^2
  \end{bmatrix}.
\]
假设 $\epsilon=10^{-3}$ 且采用 6 位十进制浮点系统, 则 $1+\epsilon^2=1+10^{-6}$ 被舍入为 $1$, 这意味着 $A^TA$ 是奇异的!

我们指出, \eqref{eq:4-3} 的解 $\mathbf{x}$ 可表示为
\[
  \mathbf{x}=(A^TA)^{-1}A^T\mathbf{b}.
\]
若令
\[
  A^{\dagger}=(A^TA)^{-1}A^T,
\]
则最小二乘解 $\mathbf{x}$ 可写为
\[
  \mathbf{x}=A^{\dagger}\mathbf{b}.
\]
实际上, $n\times m$ 矩阵 $A^{\dagger}$ 就是 $A$ 的 Moore–Penrose 广义逆, 它是唯一的 \cite{src:30,src:39,src:88}. 一般而言, 有

\begin{definition}
\label{def:4-2}
设 $X\in\R^{n\times m}$. 若它满足以下条件:
\[
  AXA=A,\qquad XAX=X,\qquad (AX)^T=AX,\qquad (XA)^T=XA,
\]
则称 $X$ 为 $A$ 的\newterm{moorepenrose}{Moore–Penrose 广义逆}{Moore–Penrose generalized inverse}, 记为 $A^{\dagger}$.
\end{definition}

下面建立最小二乘问题的扰动分析. 假设 $\mathbf{b}$ 上有一个扰动 $\delta\mathbf{b}$, 并记 $\mathbf{x}$ 与 $\mathbf{x}+\delta\mathbf{x}$ 分别为如下最小二乘问题的解:
\[
  \min_{\mathbf{x}}\norm{\mathbf{b}-A\mathbf{x}}_2,
  \qquad
  \min_{\mathbf{x}}\norm{(\mathbf{b}+\delta\mathbf{b})-A\mathbf{x}}_2.
\]
则
\[
  \mathbf{x}=A^{\dagger}\mathbf{b},
\]
且
\[
  \mathbf{x}+\delta\mathbf{x}
  =A^{\dagger}(\mathbf{b}+\delta\mathbf{b})=A^{\dagger}\tilde{\mathbf{b}},
\]
其中 $\tilde{\mathbf{b}}=\mathbf{b}+\delta\mathbf{b}$. 矩阵 $A\in\R^{m\times n}$ 的矩阵 2 范数定义为
\[
  \norm{A}_2\equiv\max_{\mathbf{0}\neq\mathbf{x}\in\R^n}
  \frac{\norm{A\mathbf{x}}_2}{\norm{\mathbf{x}}_2}
  =\max_{\norm{\mathbf{x}}_2=1}\norm{A\mathbf{x}}_2,
\]
这与方阵的矩阵 2 范数相同. 此处有

\begin{theorem}
\label{thm:4-5}
设 $\mathbf{b}_1$ 与 $\tilde{\mathbf{b}}_1$ 分别表示 $\mathbf{b}$ 与 $\tilde{\mathbf{b}}$ 在 $\mathcal{R}(A)$ 上的\newterm{orthprojection}{正交投影}{orthogonal projection}. 若 $\mathbf{b}_1\neq\mathbf{0}$, 则
\[
  \frac{\norm{\delta\mathbf{x}}_2}{\norm{\mathbf{x}}_2}
  \leqslant\kappa_2(A)\frac{\norm{\delta\mathbf{b}_1}_2}{\norm{\mathbf{b}_1}_2},
\]
其中 $\kappa_2(A)=\norm{A}_2\norm{A^{\dagger}}_2$, 且 $\tilde{\mathbf{b}}_1=\mathbf{b}_1+\delta\mathbf{b}_1$.
\end{theorem}

\begin{proof}
设 $\mathbf{b}_2$ 表示 $\mathbf{b}$ 在 $\mathcal{R}(A)^{\perp}$ 上的正交投影, 则 $\mathbf{b}=\mathbf{b}_1+\mathbf{b}_2$ 且 $A^T\mathbf{b}_2=\mathbf{0}$. 注意
\[
  A^{\dagger}\mathbf{b}=A^{\dagger}\mathbf{b}_1+A^{\dagger}\mathbf{b}_2
  =A^{\dagger}\mathbf{b}_1+(A^TA)^{-1}A^T\mathbf{b}_2
  =A^{\dagger}\mathbf{b}_1.
\]
类似地, $A^{\dagger}\tilde{\mathbf{b}}=A^{\dagger}\tilde{\mathbf{b}}_1$. 因此
\begin{align}
  \norm{\delta\mathbf{x}}_2
  &=\norm{A^{\dagger}\mathbf{b}-A^{\dagger}\tilde{\mathbf{b}}}_2
   =\norm{A^{\dagger}(\mathbf{b}_1-\tilde{\mathbf{b}}_1)}_2\notag\\
  &\leqslant\norm{A^{\dagger}}_2\norm{\mathbf{b}_1-\tilde{\mathbf{b}}_1}_2
   =\norm{A^{\dagger}}_2\norm{\delta\mathbf{b}_1}_2.
   \label{eq:4-4}
\end{align}
由于 $A\mathbf{x}=\mathbf{b}_1$, 有
\begin{equation}
  \norm{\mathbf{b}_1}_2\leqslant\norm{A}_2\norm{\mathbf{x}}_2.
  \label{eq:4-5}
\end{equation}
结合 \eqref{eq:4-4} 与 \eqref{eq:4-5}, 证毕.
\end{proof}

我们指出, 条件数 $\kappa_2(A)$ 对最小二乘问题十分重要. 当 $\kappa_2(A)$ 大时, 称该最小二乘问题是病态的; 当 $\kappa_2(A)$ 小时, 称该最小二乘问题是良态的.

\begin{theorem}
\label{thm:4-6}
设 $A$ 的列向量线性无关, 则
\[
  \kappa_2(A)^2=\kappa_2(A^TA).
\]
\end{theorem}

\begin{proof}
由定理 \ref{thm:3-3} (ii) 与所给条件, 有
\[
  \norm{A}_2^2=\norm{A^TA}_2,\qquad
  \norm{A^{\dagger}}_2^2=\norm{A^{\dagger}(A^{\dagger})^T}_2
  =\norm{(A^TA)^{-1}}_2.
\]
因此
\[
  \kappa_2(A)^2=\norm{A}_2^2\norm{A^{\dagger}}_2^2
  =\norm{A^TA}_2\norm{(A^TA)^{-1}}_2=\kappa_2(A^TA). \qedhere
\]
\end{proof}

\section{正交变换}
\label{sec:4-2}

为构造求解最小二乘问题的高效算法, 我们介绍一些著名的正交变换.

\subsection{Householder 变换}
\label{subsec:4-2-1}

首先介绍 Householder 变换的如下定义.

\begin{definition}
\label{def:4-3}
设 $\mathbf{w}\in\R^n$ 且 $\norm{\mathbf{w}}_2=1$. 定义 $H\in\R^{n\times n}$ 如下:
\begin{equation}
  H=I-2\mathbf{w}\mathbf{w}^T.
  \label{eq:4-6}
\end{equation}
矩阵 $H$ 称为 \newterm{householder}{Householder 变换}{Householder transformation}.
\end{definition}

\begin{theorem}
\label{thm:4-7}
设 $H$ 由 \eqref{eq:4-6} 定义, 则 $H$ 具有以下性质:
\begin{romanlist}
  \item $H$ 是对称正交矩阵.
  \item $H^2=I$.
  \item $H$ 称为\newterm{reflection}{反射}{reflection}, 因为 $H\mathbf{x}$ 是 $\mathbf{x}\in\R^n$ 关于过 $\mathbf{0}$ 且垂直于 $\mathbf{w}$ 的平面的反射.
\end{romanlist}
\end{theorem}

\begin{proof}
我们只证 (iii). 注意对任意向量 $\mathbf{x}\in\R^n$, 它可表示为
\[
  \mathbf{x}=\mathbf{u}+\alpha\mathbf{w},
\]
其中 $\mathbf{u}\in\spn\{\mathbf{w}\}^{\perp}$, $\alpha\in\R$. 利用 $\mathbf{u}^T\mathbf{w}=0$ 与 $\mathbf{w}^T\mathbf{w}=1$, 有
\[
  H\mathbf{x}=(I-2\mathbf{w}\mathbf{w}^T)(\mathbf{u}+\alpha\mathbf{w})
  =\mathbf{u}+\alpha\mathbf{w}-2\mathbf{w}\mathbf{w}^T\mathbf{u}
  -2\alpha\mathbf{w}\mathbf{w}^T\mathbf{w}
  =\mathbf{u}-\alpha\mathbf{w}. \qedhere
\]
\end{proof}

\begin{theorem}
\label{thm:4-8}
对任意 $\mathbf{0}\neq\mathbf{x}\in\R^n$, 可以构造单位向量 $\mathbf{w}$, 使得由 \eqref{eq:4-6} 定义的 Householder 变换满足
\[
  H\mathbf{x}=\alpha\mathbf{e}_1,
\]
其中 $\alpha=\pm\norm{\mathbf{x}}_2$.
\end{theorem}

\begin{proof}
注意 $H\mathbf{x}=(I-2\mathbf{w}\mathbf{w}^T)\mathbf{x}
=\mathbf{x}-2(\mathbf{w}^T\mathbf{x})\mathbf{w}$. 令
\[
  \mathbf{w}=\frac{\mathbf{x}-\alpha\mathbf{e}_1}
  {\norm{\mathbf{x}-\alpha\mathbf{e}_1}_2},
\]
则有
\begin{align}
  H\mathbf{x}&=\mathbf{x}-2(\mathbf{w}^T\mathbf{x})\mathbf{w}\notag\\
  &=\mathbf{x}-\frac{2}{\norm{\mathbf{x}-\alpha\mathbf{e}_1}_2^2}
   \bigl[(\mathbf{x}^T-\alpha\mathbf{e}_1^T)\mathbf{x}\bigr]
   (\mathbf{x}-\alpha\mathbf{e}_1)\notag\\
  &=\mathbf{x}
   -\frac{2\bigl(\norm{\mathbf{x}}_2^2-\alpha x_1\bigr)}
   {\norm{\mathbf{x}-\alpha\mathbf{e}_1}_2^2}\mathbf{x}
   +\frac{2\bigl(\norm{\mathbf{x}}_2^2-\alpha x_1\bigr)\alpha}
   {\norm{\mathbf{x}-\alpha\mathbf{e}_1}_2^2}\mathbf{e}_1\notag\\
  &=\left[1-\frac{2\bigl(\norm{\mathbf{x}}_2^2-\alpha x_1\bigr)}
   {\norm{\mathbf{x}-\alpha\mathbf{e}_1}_2^2}\right]\mathbf{x}
   +\frac{2\bigl(\norm{\mathbf{x}}_2^2-\alpha x_1\bigr)\alpha}
   {\norm{\mathbf{x}-\alpha\mathbf{e}_1}_2^2}\mathbf{e}_1,
   \label{eq:4-7}
\end{align}
其中 $x_1$ 是向量 $\mathbf{x}$ 的第一个分量. 令 $\mathbf{x}$ 的系数为零, 便得如下方程:
\[
  1-\frac{2\bigl(\norm{\mathbf{x}}_2^2-\alpha x_1\bigr)}
  {\norm{\mathbf{x}-\alpha\mathbf{e}_1}_2^2}=0.
\]
解此方程得 $\alpha=\pm\norm{\mathbf{x}}_2$. 把它代入 \eqref{eq:4-7}, 便有
\[
  H\mathbf{x}=\mp\norm{\mathbf{x}}_2\mathbf{e}_1. \qedhere
\]
\end{proof}

我们指出, 对任意向量 $\mathbf{0}\neq\mathbf{x}\in\R^n$, 由定理 \ref{thm:4-8}, 可以构造 Householder 矩阵 $H$, 使 $H\mathbf{x}$ 的后 $n-1$ 个分量均为零. 可用以下两步构造 $H$ 的单位向量 $\mathbf{w}$:
\begin{parenlist}
  \item 计算 $\mathbf{v}=\mathbf{x}\pm\norm{\mathbf{x}}_2\mathbf{e}_1$;
  \item 计算 $\mathbf{w}=\mathbf{v}/\norm{\mathbf{v}}_2$.
\end{parenlist}

接下来的自然问题是: $\norm{\mathbf{x}}_2$ 前面的符号如何选取? 通常取
\[
  \mathbf{v}=\mathbf{x}+\sign(x_1)\norm{\mathbf{x}}_2\mathbf{e}_1,
\]
其中 $x_1\neq0$ 是向量 $\mathbf{x}$ 的第一个分量 \cite{src:80}. 由于
\[
  H=I-2\mathbf{w}\mathbf{w}^T
  =I-\frac{2}{\mathbf{v}^T\mathbf{v}}\mathbf{v}\mathbf{v}^T
  =I-\beta\mathbf{v}\mathbf{v}^T,
\]
其中 $\beta=2/\mathbf{v}^T\mathbf{v}$, 我们只需计算 $\beta$ 与 $\mathbf{v}$, 而不必形成 $\mathbf{w}$. 于是有如下算法.

\begin{algorithm}{Householder 变换}
\label{alg:4-1}
\[
  \left\{
  \begin{array}{l}
    \text{\textbf{function}:}\ [\mathbf{v},\beta]=\text{house}(\mathbf{x})\\
    \quad n=\text{length}(\mathbf{x})\\
    \quad\sigma=\mathbf{x}(2:n)^T\mathbf{x}(2:n)\\
    \quad\mathbf{v}(1)=\mathbf{x}(1)+\sign(\mathbf{x}(1))
      \sqrt{\mathbf{x}(1)^2+\sigma}\\
    \quad\mathbf{v}(2:n)=\mathbf{x}(2:n)\\
    \quad\textbf{if}\ \sigma=0\\
    \quad\quad\beta=0\\
    \quad\textbf{else}\\
    \quad\quad\beta=2/(\mathbf{v}(1)^2+\sigma)\\
    \quad\textbf{end}
  \end{array}
  \right.
\]
\end{algorithm}

\subsection{Givens 旋转}
\label{subsec:4-2-2}

\newterm{givens}{Givens 旋转}{Givens rotation}定义如下:
\[
  \begin{aligned}
    G(i,k,\theta)
    &=I+s(\mathbf{e}_i\mathbf{e}_k^T-\mathbf{e}_k\mathbf{e}_i^T)
      +(c-1)(\mathbf{e}_i\mathbf{e}_i^T+\mathbf{e}_k\mathbf{e}_k^T)\\
    &=
    \begin{bmatrix}
      1      &        & \vdots &        & \vdots &        &        \\
             & \ddots & \vdots &        & \vdots &        &        \\
      \cdots & \cdots & c      & \cdots & s      & \cdots & \cdots \\
             &        & \vdots &        & \vdots &        &        \\
      \cdots & \cdots & -s     & \cdots & c      & \cdots & \cdots \\
             &        & \vdots &        & \vdots & \ddots &        \\
             &        & \vdots &        & \vdots &        & 1
    \end{bmatrix},
  \end{aligned}
\]
其中 $c=\cos\theta$, $s=\sin\theta$. 容易证明 $G(i,k,\theta)$ 是正交矩阵.

设 $\mathbf{x}=[x_1,\ldots,x_n]^T\in\R^n$, $\mathbf{y}=G(i,k,\theta)\mathbf{x}$, 则有
\[
  y_i=cx_i+sx_k,\qquad y_k=-sx_i+cx_k,\qquad
  y_j=x_j,\quad j\neq i,k.
\]
若要使 $y_k=0$, 只需取
\[
  c=\frac{x_i}{\sqrt{x_i^2+x_k^2}},\qquad
  s=\frac{x_k}{\sqrt{x_i^2+x_k^2}}.
\]
因此
\[
  y_i=\sqrt{x_i^2+x_k^2},\qquad y_k=0.
\]
我们指出, 对任意向量 $\mathbf{0}\neq\mathbf{x}\in\R^n$, 可以构造作用在 $\mathbf{x}$ 上的 Givens 旋转 $G(i,k,\theta)$, 使 $\mathbf{x}$ 的某个非零分量变为零.

\section{QR 分解}
\label{sec:4-3}

设 $A\in\R^{m\times n}$, $\mathbf{b}\in\R^m$. 由定理 \ref{thm:3-3} (iii), 对任意正交矩阵 $Q$, 有
\[
  \norm{A\mathbf{x}-\mathbf{b}}_2=\norm{Q^T(A\mathbf{x}-\mathbf{b})}_2.
\]
因此, 最小二乘问题
\[
  \min_{\mathbf{x}}\norm{Q^TA\mathbf{x}-Q^T\mathbf{b}}_2
\]
与 \eqref{eq:4-1} 等价. 我们希望找到一个合适的正交矩阵 $Q$, 使原来的最小二乘问题化为更易求解的最小二乘问题. 我们有

\begin{theorem}
\label{thm:4-9}
(QR 分解) 设 $A\in\R^{m\times n}$ 且 $m\geqslant n$, 则 $A$ 具有 QR 分解:
\begin{equation}
  A=Q
  \begin{bmatrix}
    R\\
    \mathbf{0}
  \end{bmatrix},
  \label{eq:4-8}
\end{equation}
其中 $Q\in\R^{m\times m}$ 是正交矩阵, $R\in\R^{n\times n}$ 是对角元非负的上三角矩阵. 当 $m=n$ 且 $A$ 非奇异时, 该分解唯一.
\end{theorem}

\begin{proof}
对 $n$ 作归纳. 当 $n=1$ 时, 由定理 \ref{thm:4-8} 知结论成立. 今设定理对 $\R^{p\times(n-1)}$ ($p\geqslant n-1$) 中的所有矩阵成立. 设 $A\in\R^{m\times n}$ 的第一列为 $\mathbf{a}_1$. 再次由定理 \ref{thm:4-8}, 存在正交矩阵 $Q_1\in\R^{m\times m}$ 使得
\[
  Q_1^T\mathbf{a}_1=\norm{\mathbf{a}_1}_2\mathbf{e}_1.
\]
因此, 有
\[
  Q_1^TA=
  \begin{bmatrix}
    \norm{\mathbf{a}_1}_2 & \mathbf{v}^T\\
    \mathbf{0}            & A_1
  \end{bmatrix}.
\]
对矩阵 $A_1\in\R^{(m-1)\times(n-1)}$, 由归纳假设得
\[
  A_1=Q_2
  \begin{bmatrix}
    R_2\\
    \mathbf{0}
  \end{bmatrix},
\]
其中 $Q_2\in\R^{(m-1)\times(m-1)}$ 是正交矩阵, $R_2$ 是对角元非负的上三角矩阵. 于是令
\[
  Q=Q_1
  \begin{bmatrix}
    1 & \mathbf{0}\\
    \mathbf{0} & Q_2
  \end{bmatrix},
  \qquad
  R=
  \begin{bmatrix}
    \norm{\mathbf{a}_1}_2 & \mathbf{v}^T\\
    \mathbf{0}            & R_2\\
    \mathbf{0}            & \mathbf{0}
  \end{bmatrix}.
\]
则 $Q$ 与 $R$ 即为满足定理条件的矩阵.

当 $A\in\R^{m\times m}$ 非奇异时, 我们要证明 QR 分解唯一. 设
\[
  A=QR=\widetilde{Q}\widetilde{R},
\]
其中 $Q,\widetilde{Q}\in\R^{m\times m}$ 是正交矩阵, $R,\widetilde{R}\in\R^{m\times m}$ 是对角元非负的上三角矩阵. 由于 $A$ 非奇异, $R$ 与 $\widetilde{R}$ 的对角元均为正. 因此, 矩阵
\[
  \widetilde{Q}^TQ=\widetilde{R}R^{-1}
\]
既正交又是对角元为正的上三角矩阵. 于是
\[
  \widetilde{Q}^TQ=\widetilde{R}R^{-1}=I,
\]
即
\[
  \widetilde{Q}=Q,\qquad\widetilde{R}=R. \qedhere
\]
\end{proof}

后面需要 QR 分解的复矩阵版本.

\begin{corollary}
\label{cor:4-2}
设 $A\in\C^{m\times n}$ 且 $m\geqslant n$, 则 $A$ 具有 QR 分解:
\[
  A=Q
  \begin{bmatrix}
    R\\
    \mathbf{0}
  \end{bmatrix},
\]
其中 $Q\in\C^{m\times m}$ 是\newterm{unitarymatrix}{酉矩阵}{unitary matrix}, $R\in\C^{n\times n}$ 是对角元非负的上三角矩阵. 当 $m=n$ 且 $A$ 非奇异时, 该分解唯一.
\end{corollary}

下面用 QR 分解求解最小二乘问题 \eqref{eq:4-1}. 假设 $A\in\R^{m\times n}$ ($m\geqslant n$) 的列线性无关, $\mathbf{b}\in\R^m$, 且 $A$ 具有 QR 分解 \eqref{eq:4-8}. 把 $Q$ 分块为
\[
  Q=[Q_1\mid Q_2],
\]
以及
\[
  Q^T\mathbf{b}=
  \begin{bmatrix}
    Q_1^T\\
    Q_2^T
  \end{bmatrix}\mathbf{b}
  =
  \begin{bmatrix}
    \mathbf{c}_1\\
    \mathbf{c}_2
  \end{bmatrix}.
\]
则
\[
  \norm{A\mathbf{x}-\mathbf{b}}_2^2
  =\norm{Q^TA\mathbf{x}-Q^T\mathbf{b}}_2^2
  =\norm{R\mathbf{x}-\mathbf{c}_1}_2^2+\norm{\mathbf{c}_2}_2^2.
\]
$\mathbf{x}$ 是最小二乘问题 \eqref{eq:4-1} 的解当且仅当它是 $R\mathbf{x}=\mathbf{c}_1$ 的解. 注意, 由于 $R$ 是上三角矩阵, 通过求解 $R\mathbf{x}=\mathbf{c}_1$ 来得到 \eqref{eq:4-1} 的解要容易得多. 于是有如下求解最小二乘问题的算法:
\begin{parenlist}
  \item 计算 $A$ 的 QR 分解.
  \item 计算 $\mathbf{c}_1=Q_1^T\mathbf{b}$.
  \item 求解上三角方程组 $R\mathbf{x}=\mathbf{c}_1$.
\end{parenlist}

最后, 我们讨论如何用 Householder 变换计算 $A$ 的 QR 分解. 设 $m=7$, $n=5$. 假设已求出 Householder 变换 $H_1$ 与 $H_2$, 使得
\[
  H_2H_1A=
  \begin{bmatrix}
    \times & \times & \times & \times & \times\\
    0      & \times & \times & \times & \times\\
    0      & 0      & +      & \times & \times\\
    0      & 0      & +      & \times & \times\\
    0      & 0      & +      & \times & \times\\
    0      & 0      & +      & \times & \times\\
    0      & 0      & +      & \times & \times
  \end{bmatrix}.
\]
现在构造 Householder 变换 $\widetilde{H}_3\in\R^{5\times5}$, 使得
\[
  \widetilde{H}_3
  \begin{bmatrix}
    +\\ +\\ +\\ +\\ +
  \end{bmatrix}
  =
  \begin{bmatrix}
    \times\\ 0\\ 0\\ 0\\ 0
  \end{bmatrix}.
\]
令 $H_3=\diag(I_2,\widetilde{H}_3)$, 得到
\[
  H_3H_2H_1A=
  \begin{bmatrix}
    \times & \times & \times & \times & \times\\
    0      & \times & \times & \times & \times\\
    0      & 0      & \times & \times & \times\\
    0      & 0      & 0      & \times & \times\\
    0      & 0      & 0      & \times & \times\\
    0      & 0      & 0      & \times & \times\\
    0      & 0      & 0      & \times & \times
  \end{bmatrix}.
\]
一般地, 经过 $n$ 个这样的步骤之后, 可把矩阵 $A$ 化为如下形式:
\[
  H_nH_{n-1}\cdots H_1A=
  \begin{bmatrix}
    R\\
    \mathbf{0}
  \end{bmatrix},
\]
其中 $R$ 是对角元非负的上三角矩阵. 令 $Q=H_1\cdots H_n$, 即得
\[
  A=Q
  \begin{bmatrix}
    R\\
    \mathbf{0}
  \end{bmatrix}.
\]
于是有如下算法.

\begin{algorithm}{QR 分解: Householder 变换}
\label{alg:4-2}
\[
  \left\{
  \begin{array}{l}
    \textbf{for}\ j=1:n\\
    \quad[\mathbf{v},\beta]=\text{house}(A(j:m,j))\\
    \quad A(j:m,j:n)=(I_{m-j+1}-\beta\mathbf{v}\mathbf{v}^T)A(j:m,j:n)\\
    \quad\textbf{if}\ j<m\\
    \quad\quad A(j+1:m,j)=\mathbf{v}(2:m-j+1)\\
    \quad\textbf{end}\\
    \textbf{end}
  \end{array}
  \right.
\]
\end{algorithm}

我们指出, QR 分解不仅是求解最小二乘问题的基本工具, 也是求解数值线性代数中其他一些基本问题的重要工具 \cite{src:20,src:23,src:24}.

\begin{exercises}
  \item 设 $A\in\R^{m\times n}$ 列满秩. 证明: 若 $E$ 满足 $\norm{E}_2\leqslant1/\norm{A^{\dagger}}_2$, 则 $A+E$ 也列满秩, 其中 $A^{\dagger}=(A^TA)^{-1}A^T$.

  \item 设 $U=[u_{ij}]$ 为非奇异上三角矩阵. 证明
  \[
    \kappa_\infty(U)\geqslant
    \frac{\underset{i}{\max}|u_{ii}|}{\underset{i}{\min}|u_{ii}|},
  \]
  其中 $\kappa_\infty(U)=\norm{U}_\infty\norm{U^{-1}}_\infty$.

  \item 设 $A\in\R^{m\times n}$, $m\geqslant n$, 且列满秩. 证明
  \[
    \begin{bmatrix}
      I & A\\
      A^T & \mathbf{0}
    \end{bmatrix}
    \begin{bmatrix}
      \mathbf{r}\\
      \mathbf{x}
    \end{bmatrix}
    =
    \begin{bmatrix}
      \mathbf{b}\\
      \mathbf{0}
    \end{bmatrix}
  \]
  有解, 其中 $\mathbf{x}$ 使 $\norm{A\mathbf{x}-\mathbf{b}}_2$ 达到最小.

  \item 设 $\mathbf{x}\in\R^n$, $P$ 是 Householder 变换且满足
  \[
    P\mathbf{x}=\pm\norm{\mathbf{x}}_2\mathbf{e}_1.
  \]
  设 $G_{12},G_{23},\ldots,G_{n-1,n}$ 为 Givens 旋转, 且令
  \[
    Q=G_{12}G_{23}\cdots G_{n-1,n}.
  \]
  若 $Q\mathbf{x}=\pm\norm{\mathbf{x}}_2\mathbf{e}_1$, 问 $P$ 是否等于 $Q$? 给出证明或反例.

  \item 设 $A\in\R^{m\times n}$. 证明 $X=A^{\dagger}$ 在所有 $X\in\R^{n\times m}$ 上使 $\norm{AX-I}_F$ 达到最小. 最小值是多少?

  \item 设
  $\mathbf{x}=
  \begin{bmatrix}
    x_1\\ x_2
  \end{bmatrix}\in\C^2$.
  给出计算如下酉矩阵的算法
  \[
    Q=
    \begin{bmatrix}
      c & \bar{s}\\
      -s & c
    \end{bmatrix},
    \qquad c\in\R,\quad c^2+|s|^2=1,
  \]
  使得
  $Q\mathbf{x}=
  \begin{bmatrix}
    *\\ 0
  \end{bmatrix}$.

  \item 设 $m\times n$ 矩阵 $A$ 具有形式
  \[
    A=
    \begin{bmatrix}
      A_1\\
      A_2
    \end{bmatrix},
  \]
  其中 $A_1$ 是 $n\times n$ 非奇异矩阵, $A_2$ 是 $(m-n)\times n$ 任意矩阵. 证明 $\norm{A^{\dagger}}_2\leqslant\norm{A_1^{-1}}_2$.

  \item 考虑如下著名的病态矩阵
  \[
    A=
    \begin{bmatrix}
      1 & 1 & 1\\
      \epsilon & 0 & 0\\
      0 & \epsilon & 0\\
      0 & 0 & \epsilon
    \end{bmatrix},
    \qquad|\epsilon|\ll1.
  \]
  \par\smallskip
  \hangindent=2em\hangafter=1\noindent\hspace*{2em}(a)\enspace
  选取小的 $\epsilon$ 使 $\rank(A)=3$, 然后计算 $\kappa_2(A)$, 以说明 $A$ 是病态的.
  \par\smallskip
  \hangindent=2em\hangafter=1\noindent\hspace*{2em}(b)\enspace
  对如上给定的 $A$ 与 $\mathbf{b}=[3,\epsilon,\epsilon,\epsilon]^T$, 分别用以下方法求最小二乘解:
  \par\smallskip
  \hangindent=4em\hangafter=1\noindent\hspace*{4em}(i)\enspace 正规方程法;
  \par\smallskip
  \hangindent=4em\hangafter=1\noindent\hspace*{4em}(ii)\enspace QR 方法.

  \item 设 $A=BC$, 其中 $B\in\C^{m\times r}$, $C\in\C^{r\times n}$, 且
  \[
    r=\rank(A)=\rank(B)=\rank(C).
  \]
  证明
  \[
    A^{\dagger}=C^*(CC^*)^{-1}(B^*B)^{-1}B^*.
  \]

  \item 设 $A=U\Sigma V^*\in\C^{m\times n}$, 其中 $U\in\C^{m\times n}$ 满足 $U^*U=I$, $V\in\C^{n\times n}$ 满足 $V^*V=I$, $\Sigma$ 是 $n\times n$ 对角矩阵. 证明
  \[
    A^{\dagger}=V\Sigma^{\dagger}U^*.
  \]

  \item 证明
  \[
    A^{\dagger}=\lim_{\epsilon\to0}(A^*A+\epsilon I)^{-1}A^*
    =\lim_{\epsilon\to0}A^*(AA^*+\epsilon I)^{-1}.
  \]

  \item 证明
  \[
    \mathcal{R}(A)\cap\mathcal{N}(A^*)=\{\mathbf{0}\}.
  \]

  \item 设 $A=[a_{ij}]\in\C^{n\times n}$ 为\newterm{idempotent}{幂等的}{idempotent}矩阵, 则
  \[
    \mathcal{R}(A)\oplus\mathcal{N}(A)=\C^n,\qquad
    \rank(A)=\sum_{i=1}^{n}a_{ii}.
  \]

  \item 设 $A\in\C^{m\times n}$. 证明
  \[
    \mathcal{R}(AA^{\dagger})=\mathcal{R}(AA^*)=\mathcal{R}(A),
  \]
  \[
    \mathcal{R}(A^{\dagger}A)=\mathcal{R}(A^*A)
    =\mathcal{R}(A^{\dagger})=\mathcal{R}(A^*),
  \]
  \[
    \mathcal{N}(AA^{\dagger})=\mathcal{N}(AA^*)
    =\mathcal{N}(A^{\dagger})=\mathcal{N}(A^*),
  \]
  \[
    \mathcal{N}(A^{\dagger}A)=\mathcal{N}(A^*A)=\mathcal{N}(A).
  \]
  因此 $AA^{\dagger}$ 与 $A^{\dagger}A$ 是正交投影算子.

  \item 证明推论 \ref{cor:4-2}.
\end{exercises}

% chapters/ch05.tex —— 第 5 章：经典迭代法
% 源书：Chapter 5 Classical Iterative Methods，书页 60–75（物理页 74–89）
% 公式编号 (5.1)–(5.11)；定理 5.1–5.13；定义 5.1–5.5；引理 5.1；推论 5.1；例 5.1；习题 1–14
\chapter{经典迭代法}
\label{ch:iterative}

我们研究求解 $A\mathbf{x}=\mathbf{b}$ 的经典迭代法. 迭代法最初由 Gauss 于 1823 年, Liouville 于 1837 年和 Jacobi 于 1845 年提出, 它与高斯消去法等直接法有很大不同. 参见 \cite{src:2}.

若矩阵 $A$ 相当大, 则基于 $A$ 的 $LU$ 分解的直接法在计算时间与计算机存储两方面都变得难以承受. 在偏微分方程离散化等实际问题中, 矩阵的阶数可高达数十万. 对这类问题, 直接法变得不切实际. 此外, 大多数大型问题都是稀疏的, 而在 $LU$ 分解过程中稀疏性通常会丢失. 因此, 在 $LU$ 分解结束时我们不得不面对一个含有大量非零元素的大矩阵, 此时存储成为关键问题. 对这类问题, 我们可以使用一类称为迭代法的方法. 本章只讨论若干经典迭代法.

我们指出, 经典迭代法的缺点是收敛速度可能很慢, 甚至可能发散, 而且需要寻找合适的停止准则.

\section{Jacobi 方法与 Gauss-Seidel 方法}
\label{sec:5-1}

\subsection{Jacobi 方法}
\label{subsec:5-1-1}

考虑如下线性方程组
\[
  A\mathbf{x}=\mathbf{b},
\]
其中 $A=[a_{ij}]\in\R^{n\times n}$. 可以把矩阵 $A$ 写成如下形式
\[
  A=D-L-U,
\]
其中
\[
  D=\diag(a_{11},a_{22},\ldots,a_{nn}),
\]
\[
  L=
  \begin{bmatrix}
    0 & & & & \\
    -a_{21} & 0 & & & \\
    -a_{31} & -a_{32} & 0 & & \\
    \vdots & \vdots & \ddots & \ddots & \\
    -a_{n1} & -a_{n2} & \cdots & -a_{n,n-1} & 0
  \end{bmatrix},
\]
以及
\[
  U=
  \begin{bmatrix}
    0 & -a_{12} & -a_{13} & \cdots & -a_{1n}\\
      & 0 & -a_{23} & \cdots & -a_{2n}\\
      & & \ddots & \ddots & \vdots\\
      & & & 0 & -a_{n-1,n}\\
      & & & & 0
  \end{bmatrix}.
\]
于是容易看出
\[
  \mathbf{x}=B_J\mathbf{x}+\mathbf{g},
\]
其中
\[
  B_J=D^{-1}(L+U),\qquad \mathbf{g}=D^{-1}\mathbf{b}.
\]
矩阵 $B_J$ 称为 Jacobi 迭代矩阵. 给定初始向量 $\mathbf{x}_0=\bigl[x_1^{(0)},x_2^{(0)},\ldots,x_n^{(0)}\bigr]^T$ 后, 相应的迭代
\begin{equation}
  \mathbf{x}_k=B_J\mathbf{x}_{k-1}+\mathbf{g},\qquad k=1,2,\ldots,
  \label{eq:5-1}
\end{equation}
称为\newterm{jacobimethod}{Jacobi 方法}{Jacobi method}.

\subsection{Gauss-Seidel 方法}
\label{subsec:5-1-2}

在 Jacobi 方法中, 为计算向量
\[
  \mathbf{x}_{k+1}=\Bigl[x_1^{(k+1)},x_2^{(k+1)},\ldots,x_n^{(k+1)}\Bigr]^T
\]
的各分量, 只用到向量 $\mathbf{x}_k$ 的各分量. 然而注意, 为计算 $x_i^{(k+1)}$, 我们本可以使用已经得到的 $x_1^{(k+1)},x_2^{(k+1)},\ldots,x_{i-1}^{(k+1)}$. 因此, Jacobi 方法的一个自然改进是把 Jacobi 迭代 \eqref{eq:5-1} 改写为
\begin{equation}
  \mathbf{x}_k=(D-L)^{-1}U\mathbf{x}_{k-1}+(D-L)^{-1}\mathbf{b},\qquad k=1,2,\ldots
  \label{eq:5-2}
\end{equation}
其思想是: 每算出一个新分量, 就立即在下一个分量的计算中使用它. 迭代 \eqref{eq:5-2} 称为\newterm{gaussseidelmethod}{Gauss-Seidel 方法}{Gauss-Seidel method}.

注意矩阵 $D-L$ 是以 $a_{11},\ldots,a_{nn}$ 为对角元的下三角矩阵. 由于假定这些元素非零, 故矩阵 $D-L$ 非奇异. 矩阵
\[
  B_{GS}=(D-L)^{-1}U
\]
称为 Gauss-Seidel 迭代矩阵.

\section{收敛性分析}
\label{sec:5-2}

\subsection{收敛定理}
\label{subsec:5-2-1}

通常很难给出一个好的初始近似 $\mathbf{x}_0$. 因此, 最好能有条件保证 Jacobi 方法与 Gauss-Seidel 方法对任意选定的初始近似都收敛.

Jacobi 迭代与 Gauss-Seidel 迭代都可用下式表示
\begin{equation}
  \mathbf{x}_{k+1}=B\mathbf{x}_k+\mathbf{g},\qquad k=0,1,\ldots
  \label{eq:5-3}
\end{equation}
对 Jacobi 迭代, 有
\[
  B_J=D^{-1}(L+U),\qquad \mathbf{g}=D^{-1}\mathbf{b};
\]
而对 Gauss-Seidel 迭代, 有
\[
  B_{GS}=(D-L)^{-1}U,\qquad \mathbf{g}=(D-L)^{-1}\mathbf{b}.
\]
迭代 \eqref{eq:5-3} 称为\newterm{linearstationary}{线性定常迭代}{linear stationary iteration}, 其中 $B\in\R^{n\times n}$ 称为\newterm{iterationmatrix}{迭代矩阵}{iteration matrix}, $\mathbf{g}\in\R^n$ 称为常向量, $\mathbf{x}_0\in\R^n$ 称为初始向量. 下面给出一个收敛性定理.

\begin{theorem}
\label{thm:5-1}
迭代 \eqref{eq:5-3} 对任意初始猜测 $\mathbf{x}_0$ 收敛, 当且仅当 $k\to\infty$ 时 $B^k\to0$.
\end{theorem}

\begin{proof}
由 $\mathbf{x}=B\mathbf{x}+\mathbf{g}$ 与 $\mathbf{x}_{k+1}=B\mathbf{x}_k+\mathbf{g}$, 有
\begin{equation}
  \mathbf{x}-\mathbf{x}_{k+1}=B(\mathbf{x}-\mathbf{x}_k).
  \label{eq:5-4}
\end{equation}
由于它对 $k$ 的任意值都成立, 可以写出
\begin{equation}
  \mathbf{x}-\mathbf{x}_k=B(\mathbf{x}-\mathbf{x}_{k-1}).
  \label{eq:5-5}
\end{equation}
把 \eqref{eq:5-5} 代入 \eqref{eq:5-4}, 得
\[
  \mathbf{x}-\mathbf{x}_{k+1}=B^2(\mathbf{x}-\mathbf{x}_{k-1}).
\]
如此继续 $k$ 次, 可写出
\[
  \mathbf{x}-\mathbf{x}_{k+1}=B^{k+1}(\mathbf{x}-\mathbf{x}_0).
\]
这表明, 对任意选定的 $\mathbf{x}_0$, $\{\mathbf{x}_k\}$ 收敛到解 $\mathbf{x}$ 当且仅当 $k\to\infty$ 时 $B^k\to0$.
\end{proof}

回忆定理 \ref{thm:3-5}: $k\to\infty$ 时 $B^k\to0$ 当且仅当谱半径 $\rho(B)<1$. 由于 $|\lambda_i|\leqslant\norm{B}$, 判断 $\rho(B)<1$ 的一个好办法是用行和范数或列和范数计算 $\norm{B}$, 看是否有 $\norm{B}<1$. 注意其逆不真. 把定理 \ref{thm:5-1} 的结论与上述观察结合起来, 便得下述定理.

\begin{theorem}
\label{thm:5-2}
迭代 \eqref{eq:5-3} 对任意选定的 $\mathbf{x}_0$ 收敛, 当且仅当 $\rho(B)<1$. 此外, 若对某个矩阵范数有 $\norm{B}<1$, 则迭代 \eqref{eq:5-3} 收敛.
\end{theorem}

考虑下面的例子:
\[
  A_1=
  \begin{bmatrix}
    1 & 2 & -2\\
    1 & 1 & 1\\
    2 & 2 & 1
  \end{bmatrix},
  \qquad
  A_2=
  \begin{bmatrix}
    2 & -1 & 1\\
    1 & 1 & 1\\
    1 & 1 & -2
  \end{bmatrix}.
\]
容易验证, 对 $A_1$, 即使 Gauss-Seidel 方法不收敛, Jacobi 方法仍收敛; 对 $A_2$, Jacobi 方法发散而 Gauss-Seidel 方法收敛.

\subsection{收敛的充分条件}
\label{subsec:5-2-2}

现在应用定理 \ref{thm:5-2} 给出一系列准则, 它们保证 Jacobi 方法和 (或)Gauss-Seidel 方法对任意选定的初始近似 $\mathbf{x}_0$ 收敛.

\begin{definition}
\label{def:5-1}
设 $A=[a_{ij}]\in\R^{n\times n}$. 若
\[
  |a_{ii}|>\sum_{\substack{j=1\\ j\neq i}}^{n}|a_{ij}|,
  \qquad i=1,2,\ldots,n,
\]
则称矩阵 $A$ 为严格对角占优的. 若
\[
  |a_{ii}|\geqslant\sum_{\substack{j=1\\ j\neq i}}^{n}|a_{ij}|,
  \qquad i=1,2,\ldots,n,
\]
且其中至少有一个严格不等式成立, 则称矩阵 $A$ 为\newterm{weaklydiagdominant}{弱对角占优的}{weakly diagonally dominant}.
\end{definition}

\begin{theorem}
\label{thm:5-3}
若 $A$ 严格对角占优, 则 Jacobi 方法对任意初始近似 $\mathbf{x}_0$ 收敛.
\end{theorem}

\begin{proof}
由于 $A=[a_{ij}]$ 严格对角占优, 有
\[
  |a_{ii}|>\sum_{\substack{j=1\\ j\neq i}}^{n}|a_{ij}|,
  \qquad i=1,2,\ldots,n.
\]
回忆 Jacobi 迭代矩阵
\[
  B_J=D^{-1}(L+U)
\]
由下式给出:
\[
  B_J=
  \begin{bmatrix}
    0 & -\dfrac{a_{12}}{a_{11}} & \cdots & \cdots & -\dfrac{a_{1n}}{a_{11}}\\[8pt]
    -\dfrac{a_{21}}{a_{22}} & 0 & -\dfrac{a_{23}}{a_{22}} & \cdots & -\dfrac{a_{2n}}{a_{22}}\\[8pt]
    \vdots & \ddots & \ddots & \ddots & \vdots\\[2pt]
    \vdots & & \ddots & \ddots & -\dfrac{a_{n-1,n}}{a_{n-1,n-1}}\\[8pt]
    -\dfrac{a_{n1}}{a_{nn}} & \cdots & \cdots & \dfrac{a_{n,n-1}}{a_{nn}} & 0
  \end{bmatrix}.
\]
我们知道每一行的绝对行和都小于 $1$, 即
\[
  \norm{B_J}_\infty<1.
\]
于是由定理 \ref{thm:5-2} 知, Jacobi 方法收敛.
\end{proof}

\begin{theorem}
\label{thm:5-4}
若 $A=[a_{ij}]\in\R^{n\times n}$ 严格对角占优, 则 Gauss-Seidel 方法对任意初始近似 $\mathbf{x}_0$ 收敛.
\end{theorem}

\begin{proof}
Gauss-Seidel 迭代矩阵由
\[
  B_{GS}=(D-L)^{-1}U
\]
给出. 设 $\lambda$ 是该矩阵的特征值, $\mathbf{x}=[x_1,x_2,\ldots,x_n]^T$ 是相应的特征向量, 其模最大的分量为 $1$. 则由
\[
  B_{GS}\mathbf{x}=\lambda\mathbf{x}
\]
得
\[
  U\mathbf{x}=\lambda(D-L)\mathbf{x},
\]
即
\[
  -\sum_{j=i+1}^{n}a_{ij}x_j=\lambda\sum_{j=1}^{i}a_{ij}x_j,
  \qquad 1\leqslant i\leqslant n,
\]
它可改写为
\begin{equation}
  \lambda a_{ii}x_i=-\lambda\sum_{j=1}^{i-1}a_{ij}x_j-\sum_{j=i+1}^{n}a_{ij}x_j,
  \qquad 1\leqslant i\leqslant n.
  \label{eq:5-6}
\end{equation}
设 $x_k$ 是向量 $\mathbf{x}$ 中模最大的分量 (其模为 $1$). 则由 \eqref{eq:5-6} 得
\[
  |\lambda|\cdot|a_{kk}|\leqslant|\lambda|\sum_{j=1}^{k-1}|a_{kj}|+\sum_{j=k+1}^{n}|a_{kj}|,
\]
即
\[
  |\lambda|\Bigl(|a_{kk}|-\sum_{j=1}^{k-1}|a_{kj}|\Bigr)\leqslant\sum_{j=k+1}^{n}|a_{kj}|,
\]
或
\begin{equation}
  |\lambda|\leqslant
  \frac{\displaystyle\sum_{j=k+1}^{n}|a_{kj}|}
       {\displaystyle|a_{kk}|-\sum_{j=1}^{k-1}|a_{kj}|}.
  \label{eq:5-7}
\end{equation}
由于 $A$ 严格对角占优, 有
\[
  |a_{kk}|-\sum_{j=1}^{k-1}|a_{kj}|>\sum_{j=k+1}^{n}|a_{kj}|.
\]
于是由 \eqref{eq:5-7} 断定 $|\lambda|<1$, 即 $\rho(B_{GS})<1$. 由定理 \ref{thm:5-2} 知, Gauss-Seidel 方法收敛.
\end{proof}

下面讨论对称正定矩阵的 Jacobi 方法与 Gauss-Seidel 方法的收敛性.

\begin{theorem}
\label{thm:5-5}
设 $A=[a_{ij}]\in\R^{n\times n}$ 对称, 且对角元 $a_{ii}>0$ ($i=1,2,\ldots,n$). 则 Jacobi 方法收敛当且仅当 $A$ 与 $2D-A$ 都正定.
\end{theorem}

\begin{proof}
由于
\[
  B_J=D^{-1}(L+U)=D^{-1}(D-A)=I-D^{-1}A,
\]
且
\[
  D=\diag(a_{11},a_{22},\ldots,a_{nn})
\]
中 $a_{ii}>0$ ($i=1,2,\ldots,n$), 故
\[
  B_J=I-D^{-1}A=D^{-1/2}\bigl(I-D^{-1/2}AD^{-1/2}\bigr)D^{1/2}.
\]
容易看出 $I-D^{-1/2}AD^{-1/2}$ 是对称的, 且与 $B_J$ 相似. 于是 $B_J$ 的特征值均为实数.

现在先假设 Jacobi 方法收敛, 则由定理 \ref{thm:5-2} 知 $\rho(B_J)<1$. 矩阵
\[
  I-D^{-1/2}AD^{-1/2}
\]
的特征值的绝对值小于 $1$, 即 $D^{-1/2}AD^{-1/2}$ 的特征值落在 $(0,2)$ 内. 因此 $A$ 正定. 另一方面, $2I-D^{-1/2}AD^{-1/2}$ 的特征值为正, 故矩阵
\[
  2I-D^{-1/2}AD^{-1/2}
\]
正定. 由于
\[
  D^{-1/2}(2D-A)D^{-1/2}=2I-D^{-1/2}AD^{-1/2},
\]
可知 $2D-A$ 也正定.

反之, 由于
\[
  D^{1/2}(I-B_J)D^{-1/2}=D^{-1/2}AD^{-1/2},
\]
且 $A$ 正定, 推知 $I-B_J$ 的特征值为正, 即 $B_J$ 的特征值小于 $1$. 又因 $2D-A$ 正定, 且
\[
  D^{-1/2}(2D-A)D^{-1/2}=D^{1/2}\bigl(I+B_J\bigr)D^{-1/2},
\]
可推知 $I+B_J$ 的特征值为正, 即 $B_J$ 的特征值大于 $-1$. 于是 $\rho(B_J)<1$. 再次由定理 \ref{thm:5-2} 知, Jacobi 方法收敛.
\end{proof}

\begin{theorem}
\label{thm:5-6}
设 $A$ 为对称正定矩阵. 则 Gauss-Seidel 方法对任意初始近似 $\mathbf{x}_0$ 收敛.
\end{theorem}

\begin{proof}
设 $\lambda$ 是迭代矩阵 $B_{GS}$ 的特征值, $\mathbf{u}$ 是相应的特征向量. 则
\[
  (D-L)^{-1}U\mathbf{u}=\lambda\mathbf{u}.
\]
由于 $A$ 对称正定, 有 $U=L^T$, 且
\[
  \lambda(D-L)\mathbf{u}=L^T\mathbf{u}.
\]
因此
\[
  \lambda\mathbf{u}^*(D-L)\mathbf{u}=\mathbf{u}^*L^T\mathbf{u}.
\]
令
\[
  \mathbf{u}^*D\mathbf{u}=\delta,\qquad \mathbf{u}^*L\mathbf{u}=\alpha+\ii\beta.
\]
于是有
\[
  \mathbf{u}^*L^T\mathbf{u}=(L\mathbf{u})^*\mathbf{u}=\overline{\mathbf{u}^*L\mathbf{u}}=\alpha-\ii\beta.
\]
因此
\[
  \lambda\bigl[\delta-(\alpha+\ii\beta)\bigr]=\alpha-\ii\beta.
\]
两边取模, 得
\[
  |\lambda|^2=\frac{\alpha^2+\beta^2}{(\delta-\alpha)^2+\beta^2}.
\]
另一方面,
\[
  0<\mathbf{u}^*A\mathbf{u}=\mathbf{u}^*(D-L-L^T)\mathbf{u}=\delta-2\alpha.
\]
因此
\[
  \begin{aligned}
    (\delta-\alpha)^2+\beta^2
    &=\delta^2+\alpha^2+\beta^2-2\delta\alpha\\
    &=\delta(\delta-2\alpha)+\alpha^2+\beta^2\\
    &>\alpha^2+\beta^2.
  \end{aligned}
\]
于是得到 $|\lambda|<1$. 由定理 \ref{thm:5-2} 知, Gauss-Seidel 方法收敛.
\end{proof}

\begin{definition}
\label{def:5-2}
若不存在置换矩阵 $P$ 使得
\[
  PAP^T=
  \begin{bmatrix}
    A_{11} & A_{12}\\
    \mathbf{0} & A_{22}
  \end{bmatrix},
\]
其中 $A_{11}$ 与 $A_{22}$ 为方阵, 则称矩阵 $A$ 是\newterm{irreducible}{不可约的}{irreducible}.
\end{definition}

我们有下述引理 \cite{src:85}.

\begin{lemma}
\label{lem:5-1}
若矩阵 $A$ 严格对角占优; 或者 $A$ 不可约且弱对角占优, 则 $A$ 非奇异.
\end{lemma}

此外,

\begin{theorem}
\label{thm:5-7}
我们有
\begin{romanlist}
  \item 若 $A$ 严格对角占优, 则 Jacobi 方法与 Gauss-Seidel 方法都收敛. 事实上,
  \[
    \norm{B_{GS}}_\infty\leqslant\norm{B_J}_\infty<1.
  \]
  \item 若 $A$ 不可约且弱对角占优, 则 Jacobi 方法与 Gauss-Seidel 方法都收敛. 此外,
  \[
    \rho(B_{GS})<\rho(B_J)<1.
  \]
\end{romanlist}
\end{theorem}

本定理的证明参见 \cite{src:30,src:85}.

\section{收敛速度}
\label{sec:5-3}

我们考虑线性定常迭代法
\[
  \mathbf{x}_{k+1}=M\mathbf{x}_k+\mathbf{g},\qquad k=0,1,\ldots,
\]
其中 $M$ 是 $n\times n$ 矩阵. 若 $I-M$ 非奇异, 则
\[
  (I-M)\mathbf{x}_*=\mathbf{g}
\]
有唯一解. 误差向量 $\mathbf{y}_k$ 定义为 $\mathbf{y}_k=\mathbf{x}_k-\mathbf{x}_*$, $k=0,1,\ldots$, 于是
\[
  \mathbf{y}_k=M\mathbf{y}_{k-1}=\cdots=M^k\mathbf{y}_0.
\]
利用矩阵范数与向量范数, 有
\[
  \norm{\mathbf{y}_k}\leqslant\norm{M^k}\,\norm{\mathbf{y}_0},
\]
且对每个 $k$, 都可能对某个向量 $\mathbf{y}_0$ 取等号. 因此, 若 $\mathbf{y}_0$ 不是零向量, 则 $\norm{M^k}$ 给出了比值 $\norm{\mathbf{y}_k}/\norm{\mathbf{y}_0}$ 的一个精确上界估计. 由于实际问题中初始向量 $\mathbf{y}_0$ 未知, 故 $\norm{M^k}$ 可作为比较不同迭代法的一种度量.

\begin{definition}
\label{def:5-3}
设 $M$ 为 $n\times n$ 迭代矩阵. 若对某个正整数 $k$ 有 $\norm{M^k}<1$, 则量
\[
  R_k(M)\equiv-\ln\bigl[(\norm{M^k})^{1/k}\bigr]=-\frac{\ln\norm{M^k}}{k}
\]
称为 $M$ 迭代 $k$ 次的\newterm{averagerate}{平均收敛速度}{average rate of convergence}.
\end{definition}

从实际计算的角度看, 平均收敛速度 $R_k(M)$ 的意义如下. 显然, 量
\[
  \sigma=\left(\frac{\norm{\mathbf{y}_k}}{\norm{\mathbf{y}_0}}\right)^{1/k}
\]
是误差范数每次迭代的平均缩减因子. 若 $\norm{M^k}<1$, 则由定义 \ref{def:5-3} 有
\[
  \sigma\leqslant(\norm{M^k})^{1/k}=e^{-R_k(M)},
\]
其中 $e$ 是自然对数的底.

若 $M$ 是对称的 (或 Hermitian 的, 或正规的, 即 $MM^*=M^*M$), 则利用迭代矩阵 $M$ 的谱半径, 有
\[
  \norm{M^k}_2=[\rho(M)]^k,
\]
从而
\[
  R_k(M)=-\ln\rho(M),
\]
它与 $k$ 无关.

下面考虑\newterm{asymptoticrate}{渐近收敛速度}{asymptotic convergence rate}
\[
  R_\infty(M)\equiv\lim_{k\to\infty}R_k(M).
\]

\begin{theorem}
\label{thm:5-8}
有
\[
  R_\infty(M)=-\ln\rho(M).
\]
\end{theorem}

\begin{proof}
只需证明
\[
  \lim_{k\to\infty}\norm{M^k}^{1/k}=\rho(M).
\]
由于
\[
  [\rho(M)]^k=\rho(M^k)\leqslant\norm{M^k},
\]
有
\[
  \rho(M)\leqslant\norm{M^k}^{1/k}.
\]
另一方面, 对任意 $\epsilon>0$, 考虑矩阵
\[
  B_\epsilon=\frac{1}{\rho(M)+\epsilon}M.
\]
显然 $\rho(B_\epsilon)<1$, 从而 $\lim_{k\to\infty}B_\epsilon^k=0$. 因此存在自然数 $K$, 当 $k\geqslant K$ 时有
\[
  \norm{B_\epsilon^k}\leqslant1,
\]
即
\[
  \norm{M^k}\leqslant[\rho(M)+\epsilon]^k.
\]
于是
\[
  \rho(M)\leqslant\norm{M^k}^{1/k}\leqslant\rho(M)+\epsilon,
\]
这意味着
\[
  \lim_{k\to\infty}\norm{M^k}^{1/k}=\rho(M).\qedhere
\]
\end{proof}

\section{SOR 方法}
\label{sec:5-4}

当 $\rho(B_{GS})$ 接近 $1$ 时, Gauss-Seidel 方法非常慢. 然而在某些情形下, 引入一个称为\newterm{relaxationparameter}{松弛因子}{relaxation parameter}的参数 $\omega$ 可以提高 Gauss-Seidel 迭代的收敛速度. 这种方法称为\newterm{successiveoverrelaxation}{逐次超松弛方法}{successive overrelaxation (SOR) method}.

\subsection{迭代形式}
\label{subsec:5-4-1}

设 $A=[a_{ij}]\in\R^{n\times n}$, 且 $A=D-L-U$ 如 \ref{subsec:5-1-1} 小节所定义. 再次考虑 $A\mathbf{x}=\mathbf{b}$ 的求解. SOR 方法的动机是对 $x_i^{(k)}$ 与 $x_i^{(k+1)}$ 取适当的加权平均来改进 Gauss-Seidel 迭代, 即
\begin{equation}
  x_i^{(k+1)}=(1-\omega)x_i^{(k)}
  +\omega\Biggl(\sum_{j=1}^{i-1}c_{ij}x_j^{(k+1)}
  +\sum_{j=i+1}^{n}c_{ij}x_j^{(k)}+g_i\Biggr).
  \label{eq:5-8}
\end{equation}
这里 $D^{-1}(L+U)=[c_{ij}]$, $\mathbf{x}_k=\bigl[x_1^{(k)},x_2^{(k)},\ldots,x_n^{(k)}\bigr]^T$, 且 $\mathbf{g}=D^{-1}\mathbf{b}=[g_1,g_2,\ldots,g_n]^T$. 写成矩阵形式, 有
\[
  \mathbf{x}_{k+1}=L_\omega\mathbf{x}_k+\omega(D-\omega L)^{-1}\mathbf{b},
\]
其中
\[
  L_\omega=(D-\omega L)^{-1}\bigl[(1-\omega)D+\omega U\bigr]
\]
称为 SOR 方法的迭代矩阵, $\omega$ 为松弛因子.

按 $\omega$ 的取值有三种情形:
\begin{parenlist}
  \item 若 $\omega=1$, 则 \eqref{eq:5-8} 等价于 Gauss-Seidel 方法;
  \item 若 $\omega<1$, 则 \eqref{eq:5-8} 称为\newterm{underrelaxation}{低松弛}{underrelaxation};
  \item 若 $\omega>1$, 则 \eqref{eq:5-8} 称为\newterm{overrelaxation}{超松弛}{overrelaxation}.
\end{parenlist}

\subsection{收敛准则}
\label{subsec:5-4-2}

自然会问: 对 $\omega$ 的哪些取值范围, SOR 迭代收敛? 为此, 我们先证明下述重要结果 \cite{src:30,src:85}.

\begin{theorem}
\label{thm:5-9}
若 $\omega$ 落在区间 $(0,2)$ 之外, 则 SOR 迭代对任何初始近似都不收敛.
\end{theorem}

\begin{proof}
回忆 SOR 迭代矩阵 $L_\omega$ 由
\[
  L_\omega=(D-\omega L)^{-1}\bigl[(1-\omega)D+\omega U\bigr]
\]
给出, 其中 $A=[a_{ij}]=D-L-U$.

矩阵 $(D-\omega L)^{-1}$ 是以 $1/a_{ii}$ ($i=1,2,\ldots,n$) 为对角元的下三角矩阵, 而矩阵
\[
  (1-\omega)D+\omega U
\]
是以 $(1-\omega)a_{ii}$ ($i=1,2,\ldots,n$) 为对角元的上三角矩阵. 因此
\[
  \det(L_\omega)=(1-\omega)^n.
\]
由于矩阵的行列式等于其特征值之积, 我们断定
\[
  \rho(L_\omega)\geqslant|1-\omega|,
\]
其中 $\rho(L_\omega)$ 是 $L_\omega$ 的谱半径. 由定理 \ref{thm:5-2} 知, 迭代矩阵的谱半径应小于 $1$ 才能收敛. 于是断定: SOR 方法的收敛要求 $0<\omega<2$.
\end{proof}

\begin{theorem}
\label{thm:5-10}
若 $A$ 对称正定, 则 $\rho(L_\omega)<1$ 对一切 $0<\omega<2$ 成立.
\end{theorem}

\begin{proof}
设 $\lambda$ 是 SOR 迭代矩阵
\[
  L_\omega=(D-\omega L)^{-1}\bigl[(1-\omega)D+\omega L^T\bigr]
\]
的任一特征值, $\mathbf{x}$ 是相应的特征向量. 有
\[
  \bigl[(1-\omega)D+\omega L^T\bigr]\mathbf{x}=\lambda(D-\omega L)\mathbf{x},
\]
或
\[
  \mathbf{x}^*\bigl[(1-\omega)D+\omega L^T\bigr]\mathbf{x}
  =\lambda\mathbf{x}^*(D-\omega L)\mathbf{x}.
\]
令
\[
  \mathbf{x}^*D\mathbf{x}=\delta,\qquad \mathbf{x}^*L\mathbf{x}=\alpha+\ii\beta.
\]
因此 $\mathbf{x}^*L^T\mathbf{x}=\alpha-\ii\beta$, 从而
\[
  (1-\omega)\delta+\omega(\alpha-\ii\beta)
  =\lambda\bigl[\delta-\omega(\alpha+\ii\beta)\bigr].
\]
两边取模, 得
\begin{equation}
  |\lambda|^2=
  \frac{\bigl[(1-\omega)\delta+\omega\alpha\bigr]^2+\omega^2\beta^2}
       {(\delta-\omega\alpha)^2+\omega^2\beta^2}.
  \label{eq:5-9}
\end{equation}
注意
\[
  \begin{aligned}
    &\bigl[(1-\omega)\delta+\omega\alpha\bigr]^2+\omega^2\beta^2
     -(\delta-\omega\alpha)^2-\omega^2\beta^2\\
    &\quad=\bigl[\delta-\omega(\delta-\alpha)\bigr]^2-(\delta-\omega\alpha)^2\\
    &\quad=\omega\delta(\delta-2\alpha)(\omega-2).
  \end{aligned}
\]
由于 $A$ 对称正定, 有
\[
  \delta>0,\qquad \delta-2\alpha>0.
\]
因此, 若 $0<\omega<2$, 则
\begin{equation}
  \bigl[(1-\omega)\delta+\omega\alpha\bigr]^2+\omega^2\beta^2
  <(\delta-\omega\alpha)^2+\omega^2\beta^2.
  \label{eq:5-10}
\end{equation}
于是, 对 $0<\omega<2$, 由 \eqref{eq:5-9} 与 \eqref{eq:5-10} 得
\[
  |\lambda|^2<1,
\]
即 SOR 方法收敛.
\end{proof}

\subsection{SOR 迭代中的最优 $\omega$}
\label{subsec:5-4-3}

为了进一步比较 Jacobi, Gauss-Seidel 与 SOR 方法, 我们对矩阵再施加一个条件. 这一条件使我们能用 $\rho(B_J)$ 显式地算出 $\rho(B_{GS})$ 与 $\rho(L_\omega)$.

\begin{definition}
\label{def:5-4}
若存在置换 $P$ 使得
\[
  PMP^T=
  \begin{bmatrix}
    D_{11} & M_{12}\\
    M_{21} & D_{22}
  \end{bmatrix},
\]
其中 $D_{11}$ 与 $D_{22}$ 为对角矩阵, 则称矩阵 $M$ 具有\newterm{propertya}{性质 A}{property A}.
\end{definition}

若 $M$ 具有性质 A, 则可写
\[
  PMP^T=\widehat{D}-\widehat{L}-\widehat{U},
\]
其中
\[
  \widehat{D}=
  \begin{bmatrix}
    D_{11} & \mathbf{0}\\
    \mathbf{0} & D_{22}
  \end{bmatrix},
  \qquad
  \widehat{L}=
  \begin{bmatrix}
    \mathbf{0} & \mathbf{0}\\
    -M_{21} & \mathbf{0}
  \end{bmatrix},
  \qquad
  \widehat{U}=
  \begin{bmatrix}
    \mathbf{0} & -M_{12}\\
    \mathbf{0} & \mathbf{0}
  \end{bmatrix}.
\]
令
\[
  \widehat{B_J}(\alpha)\equiv\alpha\widehat{D}^{-1}\widehat{L}
  +\frac{1}{\alpha}\widehat{D}^{-1}\widehat{U}.
\]
我们有

\begin{theorem}
\label{thm:5-11}
$\widehat{B_J}(\alpha)$ 的特征值与 $\alpha$ 无关.
\end{theorem}

\begin{proof}
只需注意矩阵
\[
  \widehat{B_J}(\alpha)=-
  \begin{bmatrix}
    \mathbf{0} & \dfrac{1}{\alpha}D_{11}^{-1}M_{12}\\[8pt]
    \alpha D_{22}^{-1}M_{21} & \mathbf{0}
  \end{bmatrix}
\]
与矩阵
\[
  \begin{bmatrix}
    I & \mathbf{0}\\
    \mathbf{0} & \alpha I
  \end{bmatrix}^{-1}
  \widehat{B_J}(\alpha)
  \begin{bmatrix}
    I & \mathbf{0}\\
    \mathbf{0} & \alpha I
  \end{bmatrix}
  =-
  \begin{bmatrix}
    \mathbf{0} & D_{11}^{-1}M_{12}\\
    D_{22}^{-1}M_{21} & \mathbf{0}
  \end{bmatrix}
  =\widehat{B_J}(1).\qedhere
\]
\end{proof}

\begin{definition}
\label{def:5-5}
设 $M=D-L-U$, 且
\[
  B_J(\alpha)=\alpha D^{-1}L+\frac{1}{\alpha}D^{-1}U.
\]
若 $B_J(\alpha)$ 的特征值与 $\alpha$ 无关, 则称 $M$ 具有\newterm{consistentordering}{相容次序}{consistent ordering}(或称 $M$ 是相容次序的).
\end{definition}

注意 $B_J(1)=B_J$ 是 Jacobi 迭代矩阵. 由定理 \ref{thm:5-11} 知, 若 $M$ 具有性质 A, 则 $PMP^T$ 是相容次序的, 其中 $P$ 是使
\[
  PMP^T=
  \begin{bmatrix}
    D_{11} & M_{12}\\
    M_{21} & D_{22}
  \end{bmatrix}
\]
的置换矩阵, 而 $D_{11}$ 与 $D_{22}$ 为对角矩阵. 相容次序并不蕴涵性质 A.

\begin{example}
\label{ex:5-1}
当 $D_i$ 为对角矩阵时, 任何\newterm{blocktridiagonal}{块三对角矩阵}{block tridiagonal matrix}
\[
  \begin{bmatrix}
    D_1 & A_1 & & \\
    B_1 & \ddots & \ddots & \\
      & \ddots & \ddots & A_{n-1}\\
      & & B_{n-1} & D_n
  \end{bmatrix}
\]
都是相容次序的.
\end{example}

下述定理给出 $B_J$ 的特征值与 $L_\omega$ 的特征值之间的关系 \cite{src:30,src:85}.

\begin{theorem}
\label{thm:5-12}
若 $A$ 是相容次序的, 且 $\omega\neq0$, 则
\begin{romanlist}
  \item $B_J$ 的特征值以 $\pm$ 对的形式出现.
  \item 若 $\mu$ 是 $B_J$ 的特征值, 且 $\lambda$ 满足
  \begin{equation}
    (\lambda+\omega-1)^2=\lambda\omega^2\mu^2,
    \label{eq:5-11}
  \end{equation}
  则 $\lambda$ 是 $L_\omega$ 的特征值.
  \item 若 $\lambda\neq0$ 是 $L_\omega$ 的特征值, 且 $\mu$ 满足 \eqref{eq:5-11}, 则 $\mu$ 是 $B_J$ 的特征值.
\end{romanlist}
\end{theorem}

\begin{corollary}
\label{cor:5-1}
若 $A$ 是相容次序的, 则
\[
  \rho(B_{GS})=[\rho(B_J)]^2.
\]
这意味着 Gauss-Seidel 方法的收敛速度是 Jacobi 方法的两倍.
\end{corollary}

\begin{proof}
SOR 方法中取 $\omega=1$ 等价于 Gauss-Seidel 方法. 因此由 \eqref{eq:5-11} 有 $\lambda^2=\lambda\mu^2$, 从而
\[
  \lambda=\mu^2.\qedhere
\]
\end{proof}

为了从超松弛中获得最大收益, 我们希望找一个最优的 $\omega$, 记为 $\omega_{opt}$, 使 $\rho(L_\omega)$ 达到最小. 我们有下述定理 \cite{src:30,src:85}.

\begin{theorem}
\label{thm:5-13}
设 $A$ 是相容次序的, 且 $B_J$ 的特征值为实数, $\mu=\rho(B_J)<1$. 则
\[
  \omega_{opt}=\frac{2}{1+\sqrt{1-\mu^2}},
\]
\[
  \rho(L_{\omega_{opt}})=\frac{\mu^2}{\bigl(1+\sqrt{1-\mu^2}\bigr)^2},
\]
以及
\[
  \rho(L_\omega)=
  \begin{cases}
    \omega-1, & \omega_{opt}\leqslant\omega<2,\\[6pt]
    1-\omega+\dfrac{1}{2}\omega^2\mu^2
    +\omega\mu\sqrt{1-\omega+\dfrac{1}{4}\omega^2\mu^2},
    & 0<\omega\leqslant\omega_{opt}.
  \end{cases}
\]
\end{theorem}

关于经典迭代法更多的收敛性分析, 可参见 \cite{src:14,src:74}.

\begin{exercises}
  \item 判断下列例子中 Jacobi 方法与 Gauss-Seidel 方法的收敛性:
  \[
    A_1=
    \begin{bmatrix}
      1 & 2 & -2\\
      1 & 1 & 1\\
      2 & 2 & 1
    \end{bmatrix},
    \qquad
    A_2=
    \begin{bmatrix}
      2 & -1 & 1\\
      1 & 1 & 1\\
      1 & 1 & -2
    \end{bmatrix}.
  \]

  \item 证明对 $2\times2$ 对称正定方程组, Jacobi 方法收敛.

  \item 证明: 若 $A=M-N$ 奇异, 则即使 $M$ 非奇异, 也绝不可能有 $\rho(M^{-1}N)<1$.

  \item 考虑 $A\mathbf{x}=\mathbf{b}$, 其中
  \[
    A=
    \begin{bmatrix}
      1 & 0 & \alpha\\
      0 & 1 & 0\\
      \alpha & 0 & 1
    \end{bmatrix}.
  \]
  \subquest 对哪些 $\alpha$, $A$ 正定?
  \subquest 对哪些 $\alpha$, Jacobi 方法收敛?
  \subquest 对哪些 $\alpha$, Gauss-Seidel 方法收敛?

  \item 设 $A\in\R^{n\times n}$ 非奇异. 证明存在置换矩阵 $P$, 使得 $PA$ 的对角元均非零.

  \item 证明
  \[
    A=
    \begin{bmatrix}
      4 & -1 & -1 & 0\\
      -1 & 4 & 0 & -1\\
      -1 & 0 & 4 & -1\\
      0 & -1 & -1 & 4
    \end{bmatrix}
  \]
  是相容次序的.

  \item 设 $A\in\R^{n\times n}$. 则 $\rho(A)<1$ 当且仅当 $I-A$ 非奇异, 且 $(I-A)^{-1}(I+A)$ 的每个特征值都有正实部.

  \item 设 $B\in\R^{n\times n}$ 满足 $\rho(B)=0$. 证明: 对任意 $\mathbf{g},\mathbf{x}_0\in\R^n$, 迭代公式
  \[
    \mathbf{x}_{k+1}=B\mathbf{x}_k+\mathbf{g},\qquad k=0,1,\ldots
  \]
  至多经过 $n$ 次迭代就收敛到 $\mathbf{x}=B\mathbf{x}+\mathbf{g}$ 的精确解.

  \item 若 $A=[a_{ij}]\in\R^{n\times n}$ 不可约, $a_{ij}\geqslant0$, 而 $B=[b_{ij}]\in\R^{n\times n}$, $b_{ij}\geqslant0$. 证明 $A+B$ 不可约.

  \item 设
  \[
    A=
    \begin{bmatrix}
      0 & a & 0\\
      0 & 0 & b\\
      c & 0 & 0
    \end{bmatrix}.
  \]
  $A^k$ 是否不可约 ($k=1,2,3$)?

  \item 设
  \[
    A=
    \begin{bmatrix}
      2 & -1\\
      1 & 0
    \end{bmatrix}.
  \]
  证明
  \subquest
  \[
    A^k=k
    \begin{bmatrix}
      1+1/k & -1\\
      1 & 1/k-1
    \end{bmatrix},
    \qquad k=1,2,\ldots
  \]
  \subquest
  \[
    \lim_{k\to\infty}\frac{\norm{A^k}_p}{k}=2,
    \qquad p=1,2,\infty.
  \]
  \subquest $\rho(A^k)=1$.

  \item 设
  \[
    B_k=B_{k-1}+B_{k-1}(I-AB_{k-1}),\qquad k=1,2,\ldots
  \]
  证明: 若 $\norm{I-AB_0}=c<1$, 则
  \[
    \lim_{k\to\infty}B_k=A^{-1}
  \]
  且
  \[
    \norm{A^{-1}-B_k}\leqslant\frac{c^{2^k}}{1-c}\norm{B_0}.
  \]

  \item 证明定理 \ref{thm:5-12}.

  \item 证明定理 \ref{thm:5-13}.
\end{exercises}

% chapters/ch06.tex —— 第 6 章：Krylov 子空间方法
% 源书：Chapter 6 Krylov Subspace Methods，书页 76–102（物理页 90–116）
% 公式编号 (6.1)–(6.23)；定理 6.1–6.14；引理 6.1–6.2；推论 6.1；习题 1–11
\chapter{Krylov 子空间方法}
\label{ch:krylov}

本章介绍一类称为 \newterm{krylovmethod}{Krylov 子空间方法}{Krylov subspace method}的迭代法. 在为大型稀疏问题发展起来的各类 Krylov 子空间方法中, 我们主要研究两种: 共轭梯度 (CG) 法与广义极小残量 (GMRES) 法. CG 法由 Hestenes 与 Stiefel 于 1952 年提出, 是求解对称正定线性方程组最著名的迭代法之一 \cite{src:37}; GMRES 法由 Saad 与 Schultz 于 1986 年针对非对称线性方程组提出 \cite{src:71}. 与往常一样, 我们先从最速下降法谈起.

\section{最速下降法}
\label{sec:6-1}

考虑线性方程组
\[
  A\mathbf{x}=\mathbf{b},
\]
其中 $A=[a_{ij}]\in\R^{n\times n}$ 为对称正定矩阵, $\mathbf{x}=[x_1,x_2,\ldots,x_n]^T\in\R^n$ 为未知向量, $\mathbf{b}=[b_1,b_2,\ldots,b_n]^T\in\R^n$ 为已知向量. 定义如下二次函数:
\begin{equation}
  \phi(\mathbf{x})\equiv\mathbf{x}^TA\mathbf{x}-2\mathbf{b}^T\mathbf{x}.
  \label{eq:6-1}
\end{equation}

\begin{theorem}
\label{thm:6-1}
设 $A\in\R^{n\times n}$ 为对称正定矩阵. 则求 $A\mathbf{x}=\mathbf{b}$ 的解等价于求函数 \eqref{eq:6-1} 的极小值点.
\end{theorem}

\begin{proof}
注意
\[
  \frac{\partial\phi}{\partial x_i}
  =2(a_{i1}x_1+\cdots+a_{in}x_n)-2b_i,
  \qquad i=1,2,\ldots,n.
\]
于是
\begin{equation}
  \operatorname{grad}\phi(\mathbf{x})=2(A\mathbf{x}-\mathbf{b})=-2\mathbf{r},
  \label{eq:6-2}
\end{equation}
其中 $\operatorname{grad}\phi(\mathbf{x})$ 表示 $\phi(\mathbf{x})$ 的梯度, $\mathbf{r}=\mathbf{b}-A\mathbf{x}$. 若 $\phi(\mathbf{x})$ 在点 $\mathbf{x}_*$ 处达到极小, 则
\[
  \operatorname{grad}\phi(\mathbf{x}_*)=\mathbf{0},
\]
即 $A\mathbf{x}_*=\mathbf{b}$, 这说明 $\mathbf{x}_*$ 是方程组的解.

反之, 若 $\mathbf{x}_*$ 是方程组的解, 则对任意向量 $\mathbf{y}$, 有
\[
  \begin{aligned}
    \phi(\mathbf{x}_*+\mathbf{y})
    &=(\mathbf{x}_*+\mathbf{y})^TA(\mathbf{x}_*+\mathbf{y})
      -2\mathbf{b}^T(\mathbf{x}_*+\mathbf{y})\\
    &=\mathbf{x}_*^TA\mathbf{x}_*-2\mathbf{b}^T\mathbf{x}_*+\mathbf{y}^TA\mathbf{y}
     =\phi(\mathbf{x}_*)+\mathbf{y}^TA\mathbf{y}.
  \end{aligned}
\]
由于 $A$ 对称正定, 故 $\mathbf{y}^TA\mathbf{y}\geqslant0$. 因此
\[
  \phi(\mathbf{x}_*+\mathbf{y})\geqslant\phi(\mathbf{x}_*),
\]
即 $\phi(\mathbf{x})$ 在点 $\mathbf{x}_*$ 处达到极小.
\end{proof}

那么如何求 \eqref{eq:6-1} 的极小值? 通常, 对给定的初始向量 $\mathbf{x}_0$, 我们先选取一个方向 $\mathbf{p}_0$, 然后在这条直线 $\mathbf{x}=\mathbf{x}_0+\alpha\mathbf{p}_0$ 上寻找一点
\[
  \mathbf{x}_1=\mathbf{x}_0+\alpha_0\mathbf{p}_0,
\]
使得
\[
  \phi(\mathbf{x}_1)=\phi(\mathbf{x}_0+\alpha_0\mathbf{p}_0)
  \leqslant\phi(\mathbf{x}_0+\alpha\mathbf{p}_0).
\]
这表示沿这条直线, $\phi(\mathbf{x})$ 在点 $\mathbf{x}_1$ 处达到极小. 随后, 从 $\mathbf{x}_1$ 出发再选取另一个方向 $\mathbf{p}_1$, 并在直线 $\mathbf{x}=\mathbf{x}_1+\alpha\mathbf{p}_1$ 上寻找一点
\[
  \mathbf{x}_2=\mathbf{x}_1+\alpha_1\mathbf{p}_1,
\]
使得
\[
  \phi(\mathbf{x}_2)=\phi(\mathbf{x}_1+\alpha_1\mathbf{p}_1)
  \leqslant\phi(\mathbf{x}_1+\alpha\mathbf{p}_1),
\]
即沿这条直线, $\phi(\mathbf{x})$ 在点 $\mathbf{x}_2$ 处达到极小. 如此一步一步地进行下去, 便得到
\[
  \mathbf{p}_0,\mathbf{p}_1,\ldots,
  \qquad\text{以及}\qquad
  \alpha_0,\alpha_1,\ldots,
\]
其中 $\{\mathbf{p}_k\}$ 称为\newterm{linesearch}{线搜索方向}{line search direction}, $\{\alpha_k\}$ 称为\newterm{stepsize}{步长}{step size}. 一般地, 在点 $\mathbf{x}_k$ 处选取一个方向 $\mathbf{p}_k$, 再沿直线 $\mathbf{x}=\mathbf{x}_k+\alpha\mathbf{p}_k$ 确定步长 $\alpha_k$, 使得
\[
  \phi(\mathbf{x}_k+\alpha_k\mathbf{p}_k)
  \leqslant\phi(\mathbf{x}_k+\alpha\mathbf{p}_k).
\]
于是得到 $\mathbf{x}_{k+1}=\mathbf{x}_k+\alpha_k\mathbf{p}_k$. 我们指出, 选取线搜索方向与步长的不同方式, 就对应着求解 \eqref{eq:6-1} 的不同算法.

下面先考虑如何确定步长 $\alpha_k$. 从点 $\mathbf{x}_k$ 出发沿方向 $\mathbf{p}_k$, 我们希望在直线 $\mathbf{x}=\mathbf{x}_k+\alpha\mathbf{p}_k$ 上找到步长 $\alpha_k$, 使得
\[
  \phi(\mathbf{x}_k+\alpha_k\mathbf{p}_k)
  \leqslant\phi(\mathbf{x}_k+\alpha\mathbf{p}_k).
\]
令
\[
  \begin{aligned}
    f(\alpha)&=\phi(\mathbf{x}_k+\alpha\mathbf{p}_k)\\
    &=(\mathbf{x}_k+\alpha\mathbf{p}_k)^TA(\mathbf{x}_k+\alpha\mathbf{p}_k)
      -2\mathbf{b}^T(\mathbf{x}_k+\alpha\mathbf{p}_k)\\
    &=\alpha^2\mathbf{p}_k^TA\mathbf{p}_k-2\alpha\mathbf{r}_k^T\mathbf{p}_k
      +\phi(\mathbf{x}_k),
  \end{aligned}
\]
其中 $\mathbf{r}_k=\mathbf{b}-A\mathbf{x}_k$. 由
\[
  \frac{df}{d\alpha}=2\alpha\mathbf{p}_k^TA\mathbf{p}_k-2\mathbf{r}_k^T\mathbf{p}_k=0
\]
即得
\begin{equation}
  \alpha_k=\frac{\mathbf{r}_k^T\mathbf{p}_k}{\mathbf{p}_k^TA\mathbf{p}_k}.
  \label{eq:6-3}
\end{equation}
一旦求出 $\alpha_k$, 便可计算
\[
  \mathbf{x}_{k+1}=\mathbf{x}_k+\alpha_k\mathbf{p}_k.
\]
是否有 $\phi(\mathbf{x}_{k+1})\leqslant\phi(\mathbf{x}_k)$? 考虑
\[
  \begin{aligned}
    \phi(\mathbf{x}_{k+1})-\phi(\mathbf{x}_k)
    &=\phi(\mathbf{x}_k+\alpha_k\mathbf{p}_k)-\phi(\mathbf{x}_k)\\
    &=\alpha_k^2\mathbf{p}_k^TA\mathbf{p}_k-2\alpha_k\mathbf{r}_k^T\mathbf{p}_k\\
    &=-\frac{(\mathbf{r}_k^T\mathbf{p}_k)^2}{\mathbf{p}_k^TA\mathbf{p}_k}\leqslant0.
  \end{aligned}
\]
若 $\mathbf{r}_k^T\mathbf{p}_k\neq0$, 则 $\phi(\mathbf{x}_{k+1})<\phi(\mathbf{x}_k)$.

现在考虑如何选取方向 $\mathbf{p}_k$. 众所周知, $\phi(\mathbf{x})$ 的最速下降方向是梯度的负方向, 即由 \eqref{eq:6-2} 得 $\mathbf{p}_k=\mathbf{r}_k$. 由此便得到\newterm{steepestdescent}{最速下降法}{steepest descent method}. 为了讨论该方法的收敛速度, 先引入如下引理.

\begin{lemma}
\label{lem:6-1}
设 $0<\lambda_1\leqslant\cdots\leqslant\lambda_n$ 为对称正定矩阵 $A$ 的特征值, $P(t)$ 为 $t$ 的实系数多项式. 则
\[
  \norm{P(A)\mathbf{x}}_A
  \leqslant\max_{1\leqslant j\leqslant n}|P(\lambda_j)|\cdot\norm{\mathbf{x}}_A,
  \qquad\mathbf{x}\in\R^n,
\]
其中 $\norm{\mathbf{x}}_A\equiv\sqrt{\mathbf{x}^TA\mathbf{x}}$ 称为\newterm{energynorm}{能量范数}{energy norm}.
\end{lemma}

\begin{proof}
设 $\mathbf{y}_1,\mathbf{y}_2,\ldots,\mathbf{y}_n$ 依次为 $A$ 对应于特征值 $\lambda_1,\lambda_2,\ldots,\lambda_n$ 的特征向量, 并设 $\mathbf{y}_1,\mathbf{y}_2,\ldots,\mathbf{y}_n$ 同时构成 $\R^n$ 的一组标准正交基. 于是对任意 $\mathbf{x}\in\R^n$, 有 $\mathbf{x}=\sum_{i=1}^n\beta_i\mathbf{y}_i$, 进而
\[
  \begin{aligned}
    \mathbf{x}^TP(A)AP(A)\mathbf{x}
    &=\left(\sum_{i=1}^n\beta_iP(\lambda_i)\mathbf{y}_i\right)^T
      A\left(\sum_{i=1}^n\beta_iP(\lambda_i)\mathbf{y}_i\right)\\
    &=\sum_{i=1}^n\lambda_i\beta_i^2P^2(\lambda_i)
     \leqslant\max_{1\leqslant j\leqslant n}P^2(\lambda_j)
      \sum_{i=1}^n\lambda_i\beta_i^2\\
    &=\max_{1\leqslant j\leqslant n}P^2(\lambda_j)\,\mathbf{x}^TA\mathbf{x}.
  \end{aligned}
\]
因此
\[
  \norm{P(A)\mathbf{x}}_A
  \leqslant\max_{1\leqslant j\leqslant n}|P(\lambda_j)|\cdot\norm{\mathbf{x}}_A.
  \qedhere
\]
\end{proof}

对最速下降法, 有如下收敛性定理.

\begin{theorem}
\label{thm:6-2}
设 $0<\lambda_1\leqslant\cdots\leqslant\lambda_n$ 为对称正定矩阵 $A$ 的特征值. 则由最速下降法产生的序列 $\{\mathbf{x}_k\}$ 满足
\[
  \norm{\mathbf{x}_k-\mathbf{x}_*}_A
  \leqslant\left(\frac{\lambda_n-\lambda_1}{\lambda_n+\lambda_1}\right)^k
  \norm{\mathbf{x}_0-\mathbf{x}_*}_A,
\]
其中 $\mathbf{x}_*$ 为 $A\mathbf{x}=\mathbf{b}$ 的精确解.
\end{theorem}

\begin{proof}
最速下降法产生的 $\mathbf{x}_k$ 满足
\[
  \phi(\mathbf{x}_k)\leqslant\phi(\mathbf{x}_{k-1}+\alpha\mathbf{r}_{k-1}),
  \qquad\alpha\in\R.
\]
注意到
\[
  \phi(\mathbf{x})+\mathbf{x}_*^TA\mathbf{x}_*=(\mathbf{x}-\mathbf{x}_*)^TA(\mathbf{x}-\mathbf{x}_*),
\]
故对任意 $\alpha\in\R$ 有
\begin{equation}
  \begin{aligned}
    (\mathbf{x}_k-\mathbf{x}_*)^TA(\mathbf{x}_k-\mathbf{x}_*)
    &\leqslant(\mathbf{x}_{k-1}+\alpha\mathbf{r}_{k-1}-\mathbf{x}_*)^T
      A(\mathbf{x}_{k-1}+\alpha\mathbf{r}_{k-1}-\mathbf{x}_*)\\
    &=[(I-\alpha A)(\mathbf{x}_{k-1}-\mathbf{x}_*)]^T
      A[(I-\alpha A)(\mathbf{x}_{k-1}-\mathbf{x}_*)].
  \end{aligned}
  \label{eq:6-4}
\end{equation}
令 $P_\alpha(t)=1-\alpha t$. 由 \eqref{eq:6-4} 与引理 \ref{lem:6-1}, 得
\begin{equation}
  \begin{aligned}
    \norm{\mathbf{x}_k-\mathbf{x}_*}_A
    &\leqslant\norm{P_\alpha(A)(\mathbf{x}_{k-1}-\mathbf{x}_*)}_A\\
    &\leqslant\max_{1\leqslant j\leqslant n}|P_\alpha(\lambda_j)|
      \cdot\norm{\mathbf{x}_{k-1}-\mathbf{x}_*}_A
  \end{aligned}
  \label{eq:6-5}
\end{equation}
对任意 $\alpha\in\R$ 成立. 利用 Chebyshev 逼近的性质 \cite{src:70}, 有
\begin{equation}
  \min_\alpha\max_{\lambda_1\leqslant t\leqslant\lambda_n}|1-\alpha t|
  =\frac{\lambda_n-\lambda_1}{\lambda_n+\lambda_1}.
  \label{eq:6-6}
\end{equation}
把 \eqref{eq:6-6} 代入 \eqref{eq:6-5}, 即得
\[
  \norm{\mathbf{x}_k-\mathbf{x}_*}_A
  \leqslant\frac{\lambda_n-\lambda_1}{\lambda_n+\lambda_1}
  \norm{\mathbf{x}_{k-1}-\mathbf{x}_*}_A.
\]
从而
\[
  \norm{\mathbf{x}_k-\mathbf{x}_*}_A
  \leqslant\left(\frac{\lambda_n-\lambda_1}{\lambda_n+\lambda_1}\right)^k
  \norm{\mathbf{x}_0-\mathbf{x}_*}_A.
  \qedhere
\]
\end{proof}

\section{共轭梯度法}
\label{sec:6-2}

本节介绍\newterm{conjugategradient}{共轭梯度法}{conjugate gradient (CG) method}, 它是 Krylov 子空间方法中最重要的方法之一.

\subsection{共轭梯度法}
\label{subsec:6-2-1}

CG 法的基本思想如下. 对给定的初始向量 $\mathbf{x}_0$, 第一步仍取负梯度方向, 即 $\mathbf{p}_0=\mathbf{r}_0$. 于是
\[
  \alpha_0=\frac{\mathbf{r}_0^T\mathbf{p}_0}{\mathbf{p}_0^TA\mathbf{p}_0},
  \qquad
  \mathbf{x}_1=\mathbf{x}_0+\alpha_0\mathbf{p}_0,
  \qquad
  \mathbf{r}_1=\mathbf{b}-A\mathbf{x}_1.
\]
此后, 在第 $(k+1)$ 步 ($k\geqslant1$), 我们希望在平面
\[
  \pi_2=\left\{\mathbf{x}=\mathbf{x}_k+\xi\mathbf{r}_k+\eta\mathbf{p}_{k-1}:
  \ \xi,\eta\in\R\right\}
\]
上选取一个方向 $\mathbf{p}_k$, 使 $\phi$ 下降得最快. 考虑 $\phi$ 在 $\pi_2$ 上的限制:
\[
  \begin{aligned}
    \psi(\xi,\eta)&=\phi(\mathbf{x}_k+\xi\mathbf{r}_k+\eta\mathbf{p}_{k-1})\\
    &=(\mathbf{x}_k+\xi\mathbf{r}_k+\eta\mathbf{p}_{k-1})^T
      A(\mathbf{x}_k+\xi\mathbf{r}_k+\eta\mathbf{p}_{k-1})\\
    &\quad-2\mathbf{b}^T(\mathbf{x}_k+\xi\mathbf{r}_k+\eta\mathbf{p}_{k-1}).
  \end{aligned}
\]
直接计算可得
\[
  \frac{\partial\psi}{\partial\xi}
  =2\left(\xi\mathbf{r}_k^TA\mathbf{r}_k+\eta\mathbf{r}_k^TA\mathbf{p}_{k-1}
  -\mathbf{r}_k^T\mathbf{r}_k\right),
\]
\[
  \frac{\partial\psi}{\partial\eta}
  =2\left(\xi\mathbf{r}_k^TA\mathbf{p}_{k-1}
  +\eta\mathbf{p}_{k-1}^TA\mathbf{p}_{k-1}\right),
\]
其中用到了 $\mathbf{r}_k^T\mathbf{p}_{k-1}=0$ (见后面的定理 \ref{thm:6-3}). 令
\[
  \frac{\partial\psi}{\partial\xi}=\frac{\partial\psi}{\partial\eta}=0,
\]
便得到 $\phi$ 在平面 $\pi_2$ 上的唯一极小值点
\[
  \tilde{\mathbf{x}}=\mathbf{x}_k+\xi_0\mathbf{r}_k+\eta_0\mathbf{p}_{k-1},
\]
其中 $\xi_0$ 与 $\eta_0$ 满足
\begin{equation}
  \begin{aligned}
    \xi_0\mathbf{r}_k^TA\mathbf{r}_k+\eta_0\mathbf{r}_k^TA\mathbf{p}_{k-1}
    &=\mathbf{r}_k^T\mathbf{r}_k,\\
    \xi_0\mathbf{r}_k^TA\mathbf{p}_{k-1}
    +\eta_0\mathbf{p}_{k-1}^TA\mathbf{p}_{k-1}&=0.
  \end{aligned}
  \label{eq:6-7}
\end{equation}
由 \eqref{eq:6-7} 可见, 若 $\mathbf{r}_k\neq\mathbf{0}$, 则 $\xi_0\neq0$. 因此可以取
\[
  \mathbf{p}_k=\frac{1}{\xi_0}(\tilde{\mathbf{x}}-\mathbf{x}_k)
  =\mathbf{r}_k+\frac{\eta_0}{\xi_0}\mathbf{p}_{k-1}
\]
作为新的方向, 它是使 $\phi$ 在平面 $\pi_2$ 上达到极小的最优方向. 令 $\beta_{k-1}=\eta_0/\xi_0$, 则由 \eqref{eq:6-7} 的第二式得
\[
  \beta_{k-1}=-\frac{\mathbf{r}_k^TA\mathbf{p}_{k-1}}
  {\mathbf{p}_{k-1}^TA\mathbf{p}_{k-1}}.
\]
注意 $\mathbf{p}_k$ 满足 $\mathbf{p}_k^TA\mathbf{p}_{k-1}=0$ (见后面的定理 \ref{thm:6-3}), 即 $\mathbf{p}_k$ 与 $\mathbf{p}_{k-1}$ 关于 $A$ 共轭.

一旦求出 $\mathbf{p}_k$, 便可利用 \eqref{eq:6-3} 确定 $\alpha_k$, 进而计算
\[
  \mathbf{x}_{k+1}=\mathbf{x}_k+\alpha_k\mathbf{p}_k.
\]
归纳起来, 我们有如下公式:
\[
  \alpha_k=\frac{\mathbf{r}_k^T\mathbf{p}_k}{\mathbf{p}_k^TA\mathbf{p}_k},
  \qquad
  \mathbf{x}_{k+1}=\mathbf{x}_k+\alpha_k\mathbf{p}_k,
  \qquad
  \mathbf{r}_{k+1}=\mathbf{b}-A\mathbf{x}_{k+1},
\]
\[
  \beta_k=-\frac{\mathbf{r}_{k+1}^TA\mathbf{p}_k}{\mathbf{p}_k^TA\mathbf{p}_k},
  \qquad
  \mathbf{p}_{k+1}=\mathbf{r}_{k+1}+\beta_k\mathbf{p}_k.
\]
经过几步初等运算后, 可得
\[
  \alpha_k=\frac{\mathbf{r}_k^T\mathbf{r}_k}{\mathbf{p}_k^TA\mathbf{p}_k},
  \qquad
  \beta_k=\frac{\mathbf{r}_{k+1}^T\mathbf{r}_{k+1}}{\mathbf{r}_k^T\mathbf{r}_k}.
\]
于是, 求解对称正定方程组 $A\mathbf{x}=\mathbf{b}$ 的最常用, 最成功的迭代法之一——CG 法的计算格式如下. 在初始化步, 取 $k=0$, 选定 $\mathbf{x}_0$ 后计算
\[
  \mathbf{r}_0=\mathbf{b}-A\mathbf{x}_0.
\]
当 $\mathbf{r}_k\neq\mathbf{0}$ 时, 迭代步为
\[
  \left\{
  \begin{array}{l}
    k=k+1\\
    \textbf{if}\ k=1\\
    \quad \mathbf{p}_0=\mathbf{r}_0\\
    \textbf{else}\\
    \quad \beta_{k-2}=\mathbf{r}_{k-1}^T\mathbf{r}_{k-1}
      /\mathbf{r}_{k-2}^T\mathbf{r}_{k-2}\\
    \quad \mathbf{p}_{k-1}=\mathbf{r}_{k-1}+\beta_{k-2}\mathbf{p}_{k-2}\\
    \textbf{end}\\
    \alpha_{k-1}=\mathbf{r}_{k-1}^T\mathbf{r}_{k-1}
      /\mathbf{p}_{k-1}^TA\mathbf{p}_{k-1}\\
    \mathbf{x}_k=\mathbf{x}_{k-1}+\alpha_{k-1}\mathbf{p}_{k-1}\\
    \mathbf{r}_k=\mathbf{r}_{k-1}-\alpha_{k-1}A\mathbf{p}_{k-1}
  \end{array}
  \right.
\]
其中 $\mathbf{r}_k$ 与 $\mathbf{p}_k$ 为向量, $\alpha_k$ 与 $\beta_k$ 为标量, $k=0,1,\ldots$. 这里的 $\mathbf{x}_k$ 是第 $k$ 次迭代后对精确解的近似. 当 $\mathbf{r}_k=\mathbf{0}$ 时, 解即为 $\mathbf{x}=\mathbf{x}_k$.

\subsection{基本性质}
\label{subsec:6-2-2}

\begin{theorem}
\label{thm:6-3}
向量组 $\{\mathbf{r}_i\}$ 与 $\{\mathbf{p}_i\}$ 满足如下性质:
\begin{parenlist}
  \item $\mathbf{p}_i^T\mathbf{r}_j=0$, \quad $0\leqslant i<j\leqslant k$;
  \item $\mathbf{r}_i^T\mathbf{r}_j=0$, \quad $i\neq j$, \ $0\leqslant i,j\leqslant k$;
  \item $\mathbf{p}_i^TA\mathbf{p}_j=0$, \quad $i\neq j$, \ $0\leqslant i,j\leqslant k$;
  \item $\spn\{\mathbf{r}_0,\ldots,\mathbf{r}_k\}
        =\spn\{\mathbf{p}_0,\ldots,\mathbf{p}_k\}
        =\mathcal{K}(A,\mathbf{r}_0,k+1)$, 其中
    \[
      \mathcal{K}(A,\mathbf{r}_0,k+1)
      \equiv\spn\left\{\mathbf{r}_0,A\mathbf{r}_0,\ldots,A^k\mathbf{r}_0\right\}
    \]
    称为 \newterm{krylovsubspace}{Krylov 子空间}{Krylov subspace}.
\end{parenlist}
\end{theorem}

\begin{proof}
对 $k=1$, 用归纳法, 有
\[
  \mathbf{p}_0=\mathbf{r}_0,\qquad
  \mathbf{r}_1=\mathbf{r}_0-\alpha_0A\mathbf{p}_0,\qquad
  \mathbf{p}_1=\mathbf{r}_1+\beta_0\mathbf{p}_0.
\]
于是, 当 $\alpha_0=\mathbf{r}_0^T\mathbf{r}_0/\mathbf{p}_0^TA\mathbf{p}_0$ 时,
\[
  \begin{aligned}
    \mathbf{p}_0^T\mathbf{r}_1&=\mathbf{r}_0^T\mathbf{r}_1
    =\mathbf{r}_0^T(\mathbf{r}_0-\alpha_0A\mathbf{p}_0)\\
    &=\mathbf{r}_0^T\mathbf{r}_0-\alpha_0\mathbf{p}_0^TA\mathbf{p}_0=0,
  \end{aligned}
\]
并且
\[
  \begin{aligned}
    \mathbf{p}_1^TA\mathbf{p}_0
    &=(\mathbf{r}_1+\beta_0\mathbf{r}_0)^TA\mathbf{r}_0\\
    &=\mathbf{r}_1^TA\mathbf{r}_0
      -\frac{\mathbf{r}_1^TA\mathbf{r}_0}{\mathbf{r}_0^TA\mathbf{r}_0}
      \mathbf{r}_0^TA\mathbf{r}_0=0.
  \end{aligned}
\]
现假定定理对 $k$ 成立, 证明它对 $k+1$ 也成立.

对 (1), 利用归纳假设与 $\mathbf{r}_{k+1}=\mathbf{r}_k-\alpha_kA\mathbf{p}_k$, 有
\[
  \mathbf{p}_i^T\mathbf{r}_{k+1}
  =\mathbf{p}_i^T\mathbf{r}_k-\alpha_k\mathbf{p}_i^TA\mathbf{p}_k=0,
  \qquad 0\leqslant i\leqslant k-1,
\]
同时
\[
  \mathbf{p}_k^T\mathbf{r}_{k+1}
  =\mathbf{p}_k^T\mathbf{r}_k
  -\frac{\mathbf{p}_k^T\mathbf{r}_k}{\mathbf{p}_k^TA\mathbf{p}_k}
  \mathbf{p}_k^TA\mathbf{p}_k=0.
\]
因此 (1) 对 $k+1$ 成立.

对 (2), 由归纳假设有
\[
  \spn\{\mathbf{r}_0,\ldots,\mathbf{r}_k\}
  =\spn\{\mathbf{p}_0,\ldots,\mathbf{p}_k\}.
\]
由 (1) 知 $\mathbf{r}_{k+1}$ 与该子空间正交, 因此 (2) 对 $k+1$ 成立.

对 (3), 利用归纳假设, (2) 以及
\[
  \mathbf{p}_{k+1}=\mathbf{r}_{k+1}+\beta_k\mathbf{p}_k,
  \qquad
  \mathbf{r}_{i+1}=\mathbf{r}_i-\alpha_iA\mathbf{p}_i,
\]
得
\[
  \mathbf{p}_{k+1}^TA\mathbf{p}_i
  =\frac{1}{\alpha_i}\mathbf{r}_{k+1}^T(\mathbf{r}_i-\mathbf{r}_{i+1})
  +\beta_k\mathbf{p}_k^TA\mathbf{p}_i=0,
  \qquad i=0,1,\ldots,k-1.
\]
又由 $\beta_k$ 的定义,
\[
  \begin{aligned}
    \mathbf{p}_{k+1}^TA\mathbf{p}_k
    &=(\mathbf{r}_{k+1}+\beta_k\mathbf{p}_k)^TA\mathbf{p}_k\\
    &=\mathbf{r}_{k+1}^TA\mathbf{p}_k
      -\frac{\mathbf{r}_{k+1}^TA\mathbf{p}_k}{\mathbf{p}_k^TA\mathbf{p}_k}
      \mathbf{p}_k^TA\mathbf{p}_k=0.
  \end{aligned}
\]
于是 (3) 对 $k+1$ 成立.

对 (4), 由归纳假设知
\[
  \mathbf{r}_k,\mathbf{p}_k\in\mathcal{K}(A,\mathbf{r}_0,k+1)
  =\spn\left\{\mathbf{r}_0,A\mathbf{r}_0,\ldots,A^k\mathbf{r}_0\right\}.
\]
因此
\[
  \mathbf{r}_{k+1}=\mathbf{r}_k-\alpha_kA\mathbf{p}_k
  \in\mathcal{K}(A,\mathbf{r}_0,k+2)
  =\spn\left\{\mathbf{r}_0,A\mathbf{r}_0,\ldots,A^{k+1}\mathbf{r}_0\right\},
\]
且
\[
  \mathbf{p}_{k+1}=\mathbf{r}_{k+1}+\beta_k\mathbf{p}_k
  \in\mathcal{K}(A,\mathbf{r}_0,k+2)
  =\spn\left\{\mathbf{r}_0,A\mathbf{r}_0,\ldots,A^{k+1}\mathbf{r}_0\right\}.
\]
由 (2) 与 (3) 可知, 向量组 $\mathbf{r}_0,\ldots,\mathbf{r}_{k+1}$ 与 $\mathbf{p}_0,\ldots,\mathbf{p}_{k+1}$ 均线性无关. 因此 (4) 对 $k+1$ 成立.
\end{proof}

我们指出, 由定理 \ref{thm:6-3} 可知, 用 CG 法至多 $n$ 步即可得到 $n$ 阶对称正定方程组的精确解. 因此从理论上看, CG 法是一种直接法.

\begin{theorem}
\label{thm:6-4}
由 CG 法得到的 $\mathbf{x}_k$ 满足
\begin{equation}
  \phi(\mathbf{x}_k)=\min\left\{\phi(\mathbf{x}):
  \ \mathbf{x}\in\mathbf{x}_0+\mathcal{K}(A,\mathbf{r}_0,k)\right\},
  \label{eq:6-8}
\end{equation}
或
\begin{equation}
  \norm{\mathbf{x}_k-\mathbf{x}_*}_A
  =\min\left\{\norm{\mathbf{x}-\mathbf{x}_*}_A:
  \ \mathbf{x}\in\mathbf{x}_0+\mathcal{K}(A,\mathbf{r}_0,k)\right\},
  \label{eq:6-9}
\end{equation}
其中 $\norm{\mathbf{x}}_A=\sqrt{\mathbf{x}^TA\mathbf{x}}$, $\mathbf{x}_*$ 为 $A\mathbf{x}=\mathbf{b}$ 的精确解.
\end{theorem}

\begin{proof}
由于 \eqref{eq:6-8} 与 \eqref{eq:6-9} 等价, 只需证明 \eqref{eq:6-9}. 设 CG 法进行到第 $l$ 步时 $\mathbf{r}_l=\mathbf{0}$, 则
\[
  \begin{aligned}
    \mathbf{x}_*&=\mathbf{x}_l=\mathbf{x}_{l-1}+\alpha_{l-1}\mathbf{p}_{l-1}\\
    &=\mathbf{x}_{l-2}+\alpha_{l-2}\mathbf{p}_{l-2}+\alpha_{l-1}\mathbf{p}_{l-1}\\
    &=\cdots\cdots\\
    &=\mathbf{x}_0+\alpha_0\mathbf{p}_0+\cdots+\alpha_{l-1}\mathbf{p}_{l-1}.
  \end{aligned}
\]
对 $k<l$, 有
\[
  \mathbf{x}_k=\mathbf{x}_0+\alpha_0\mathbf{p}_0+\cdots+\alpha_{k-1}\mathbf{p}_{k-1}
  \in\mathbf{x}_0+\mathcal{K}(A,\mathbf{r}_0,k).
\]
设 $\mathbf{x}$ 为 $\mathbf{x}_0+\mathcal{K}(A,\mathbf{r}_0,k)$ 中的任一向量. 则由定理 \ref{thm:6-3} (4) 知
\[
  \mathbf{x}=\mathbf{x}_0+\gamma_0\mathbf{p}_0+\cdots+\gamma_{k-1}\mathbf{p}_{k-1}.
\]
进而
\[
  \mathbf{x}_*-\mathbf{x}
  =(\alpha_0-\gamma_0)\mathbf{p}_0+\cdots
  +(\alpha_{k-1}-\gamma_{k-1})\mathbf{p}_{k-1}
  +\alpha_k\mathbf{p}_k+\cdots+\alpha_{l-1}\mathbf{p}_{l-1}.
\]
由于
\[
  \mathbf{x}_*-\mathbf{x}_k=\alpha_k\mathbf{p}_k+\cdots+\alpha_{l-1}\mathbf{p}_{l-1},
\]
利用定理 \ref{thm:6-3} (3) 得
\[
  \begin{aligned}
    \norm{\mathbf{x}_*-\mathbf{x}}_A^2
    &=\norm{(\alpha_0-\gamma_0)\mathbf{p}_0+\cdots
      +(\alpha_{k-1}-\gamma_{k-1})\mathbf{p}_{k-1}}_A^2\\
    &\quad+\norm{\alpha_k\mathbf{p}_k+\cdots
      +\alpha_{l-1}\mathbf{p}_{l-1}}_A^2\\
    &\geqslant\norm{\alpha_k\mathbf{p}_k+\cdots
      +\alpha_{l-1}\mathbf{p}_{l-1}}_A^2
     =\norm{\mathbf{x}_*-\mathbf{x}_k}_A^2.
    \qedhere
  \end{aligned}
\]
\end{proof}

\section{实用的 CG 方法与收敛性分析}
\label{sec:6-3}

本节给出 CG 法的一个实用算法, 并分析其收敛速度.

\subsection{实用的 CG 方法}
\label{subsec:6-3-1}

由定理 \ref{thm:6-3} 可知, 在精确算术下 CG 法经 $n$ 步即可得到精确解, 其中 $n$ 为方程组的阶数. 换言之, CG 法可以被看作直接法而非迭代法. 当 $n$ 很大时, 实践中我们只把 CG 法当作迭代法使用, 并在下列情形之一发生时停止迭代:
\begin{romanlist}
  \item $\norm{\mathbf{r}_k}$ 小于 $\epsilon$, 其中 $\mathbf{r}_k=\mathbf{b}-A\mathbf{x}_k$, $\epsilon$ 为事先给定的误差界; 或
  \item 迭代次数达到 $k_{\max}$, 即事先给定的最大迭代次数, 其中 $k_{\max}\ll n$.
\end{romanlist}
于是得到求解对称正定方程组 $A\mathbf{x}=\mathbf{b}$ 的如下实用算法. 在初始化步 $k=0$, 选定初始向量 $\mathbf{x}$ 并计算
\[
  \mathbf{r}=\mathbf{b}-A\mathbf{x},\qquad \rho=\mathbf{r}^T\mathbf{r}.
\]
当 $\sqrt{\rho}>\epsilon\norm{\mathbf{b}}_2$ 且 $k<k_{\max}$ 时, 迭代步为
\[
  \left\{
  \begin{array}{l}
    k=k+1\\
    \textbf{if}\ k=1\\
    \quad \mathbf{p}=\mathbf{r}\\
    \textbf{else}\\
    \quad \beta=\rho/\tilde{\rho};\qquad \mathbf{p}=\mathbf{r}+\beta\mathbf{p}\\
    \textbf{end}\\
    \mathbf{w}=A\mathbf{p};\qquad
    \alpha=\rho/\mathbf{p}^T\mathbf{w};\qquad
    \mathbf{x}=\mathbf{x}+\alpha\mathbf{p}\\
    \mathbf{r}=\mathbf{r}-\alpha\mathbf{w};\qquad
    \tilde{\rho}=\rho;\qquad
    \rho=\mathbf{r}^T\mathbf{r}
  \end{array}
  \right.
\]
其中 $\mathbf{r},\mathbf{p},\mathbf{w}$ 为向量, $\rho,\alpha,\beta$ 为标量.

我们指出:
\begin{parenlist}
  \item 算法中只出现矩阵与向量的乘法. 通常计算一次这样的乘法需要 $\mathcal{O}(n^2)$ 的运算量. 若矩阵是稀疏的或具有特殊结构, 则可用某些稀疏求解器或快速求解器高效地完成这些乘法.
  \item 与 SOR 方法不同, 算法中不需要估计任何参数.
  \item 每次迭代中, 矩阵与向量的运算都可以并行执行.
\end{parenlist}

下面简要讨论如何用 CG 法求解一般线性方程组 $A\mathbf{x}=\mathbf{b}$. 由于不能直接把 CG 法用于该方程组, 我们可以改为用 CG 法求解正规方程组
\[
  A^TA\mathbf{x}=A^T\mathbf{b}.
\]
当方程组良态时, 这种正规化 CG 方法是适用的. 但若方程组病态, 则由 $\kappa_2(A^TA)=(\kappa_2(A))^2$ 可知, 正规方程组的条件数可能变得非常大, 因此正规化 CG 方法不适用于病态方程组.

\subsection{收敛性分析}
\label{subsec:6-3-2}

关于 CG 法的收敛速度, 有如下定理.

\begin{theorem}
\label{thm:6-5}
若 $A=I+B$, 其中 $I$ 为单位矩阵且 $\rank(B)=p$, 则 CG 法至多经 $p+1$ 次迭代即可得到 $A\mathbf{x}=\mathbf{b}$ 的精确解.
\end{theorem}

\begin{proof}
由 $A=I+B$ 易证
\[
  \spn\left\{\mathbf{r}_0,A\mathbf{r}_0,\ldots,A^k\mathbf{r}_0\right\}
  =\spn\left\{\mathbf{r}_0,B\mathbf{r}_0,\ldots,B^k\mathbf{r}_0\right\}.
\]
注意到 $\rank(B)=p$, 于是
\[
  \dim\left(\spn\{\mathbf{r}_0,B\mathbf{r}_0,\ldots,B^k\mathbf{r}_0\}\right)\leqslant p+1.
\]
由定理 \ref{thm:6-3} 知
\[
  \dim\left(\spn\{\mathbf{r}_0,\ldots,\mathbf{r}_p,\mathbf{r}_{p+1}\}\right)\leqslant p+1,
\]
并且
\[
  \mathbf{r}_i^T\mathbf{r}_{p+1}=0,\qquad i=0,1,\ldots,p.
\]
因此 $\mathbf{r}_{p+1}=\mathbf{0}$, 即 $A\mathbf{x}_{p+1}=\mathbf{b}$.
\end{proof}

\begin{theorem}
\label{thm:6-6}
设 $A\in\R^{n\times n}$ 对称正定, $\mathbf{x}_*$ 为 $A\mathbf{x}=\mathbf{b}$ 的精确解. 则
\[
  \norm{\mathbf{x}_*-\mathbf{x}_k}_A
  \leqslant2\left(\frac{\sqrt{\kappa_2}-1}{\sqrt{\kappa_2}+1}\right)^k
  \norm{\mathbf{x}_*-\mathbf{x}_0}_A,
\]
其中 $\mathbf{x}_k$ 由 CG 法产生,
\[
  \kappa_2=\kappa_2(A)=\norm{A}_2\norm{A^{-1}}_2.
\]
\end{theorem}

\begin{proof}
由定理 \ref{thm:6-3} 知, 对任意 $\mathbf{x}\in\mathbf{x}_0+\mathcal{K}(A,\mathbf{r}_0,k)$, 有
\[
  \begin{aligned}
    \mathbf{x}_*-\mathbf{x}
    &=\mathbf{x}_*-\mathbf{x}_0+a_{k1}\mathbf{r}_0+a_{k2}A\mathbf{r}_0
      +\cdots+a_{kk}A^{k-1}\mathbf{r}_0\\
    &=A^{-1}\left(\mathbf{r}_0+a_{k1}A\mathbf{r}_0+a_{k2}A^2\mathbf{r}_0
      +\cdots+a_{kk}A^k\mathbf{r}_0\right)\\
    &=A^{-1}P_k(A)\mathbf{r}_0,
  \end{aligned}
\]
其中 $P_k(\lambda)=\sum_{j=0}^ka_{kj}\lambda^j$ 且 $P_k(0)=1$. 设 $\mathcal{P}_k$ 为次数不超过 $k$ 且满足 $P_k(0)=1$ 的一切多项式 $P_k$ 的集合. 由定理 \ref{thm:6-4} 与引理 \ref{lem:6-1}, 有
\begin{equation}
  \begin{aligned}
    \norm{\mathbf{x}_k-\mathbf{x}_*}_A
    &=\min\left\{\norm{\mathbf{x}-\mathbf{x}_*}_A:
      \ \mathbf{x}\in\mathbf{x}_0+\mathcal{K}(A,\mathbf{r}_0,k)\right\}\\
    &=\min_{P_k\in\mathcal{P}_k}\norm{A^{-1}P_k(A)\mathbf{r}_0}_A
     =\min_{P_k\in\mathcal{P}_k}\norm{P_k(A)A^{-1}\mathbf{r}_0}_A\\
    &\leqslant\min_{P_k\in\mathcal{P}_k}\max_{1\leqslant i\leqslant n}
      |P_k(\lambda_i)|\cdot\norm{A^{-1}\mathbf{r}_0}_A\\
    &\leqslant\min_{P_k\in\mathcal{P}_k}\max_{a_1\leqslant\lambda\leqslant a_2}
      |P_k(\lambda)|\cdot\norm{\mathbf{x}_*-\mathbf{x}_0}_A,
  \end{aligned}
  \label{eq:6-10}
\end{equation}
其中
\[
  0<a_1=\lambda_1\leqslant\cdots\leqslant\lambda_n=a_2
\]
为 $A$ 的特征值. 由著名的 Chebyshev 多项式逼近定理 \cite{src:70} 可知, 最优问题
\[
  \min_{P_k\in\mathcal{P}_k}\max_{a_1\leqslant\lambda\leqslant a_2}|P_k(\lambda)|
\]
存在唯一解
\[
  \widetilde{P}_k(\lambda)
  =\frac{T_k\left(\dfrac{a_2+a_1-2\lambda}{a_2-a_1}\right)}
  {T_k\left(\dfrac{a_2+a_1}{a_2-a_1}\right)}.
\]
这里 $T_k(z)$ 为第 $k$ 个 Chebyshev 多项式, 由递推式
\[
  T_k(z)=2zT_{k-1}(z)-T_{k-2}(z)
\]
定义, 且 $T_0(z)=1$, $T_1(z)=z$. 再次利用 Chebyshev 多项式的性质 \cite{src:70}, 可知
\begin{equation}
  \max_{a_1\leqslant\lambda\leqslant a_2}|\widetilde{P}_k(\lambda)|
  =\frac{1}{T_k\left(\dfrac{a_2+a_1}{a_2-a_1}\right)}
  \leqslant2\left(\frac{\sqrt{\kappa_2}-1}{\sqrt{\kappa_2}+1}\right)^k,
  \label{eq:6-11}
\end{equation}
其中 $\kappa_2=\kappa_2(A)=\norm{A}_2\norm{A^{-1}}_2$. 因此
\[
  \norm{\mathbf{x}_*-\mathbf{x}_k}_A
  \leqslant2\left(\frac{\sqrt{\kappa_2}-1}{\sqrt{\kappa_2}+1}\right)^k
  \norm{\mathbf{x}_*-\mathbf{x}_0}_A.
  \qedhere
\]
\end{proof}

进一步, 还有如下定理 \cite{src:43}.

\begin{theorem}
\label{thm:6-7}
若对称正定矩阵 $A$ 的特征值 $\lambda_j$ 满足如下排序:
\[
  0<\lambda_1\leqslant\cdots\leqslant\lambda_p\leqslant b_1
  \leqslant\lambda_{p+1}\leqslant\cdots\leqslant\lambda_{n-q}\leqslant b_2
  \leqslant\lambda_{n-q+1}\leqslant\cdots\leqslant\lambda_n,
\]
其中 $b_1$ 与 $b_2$ 为两个常数, 则
\[
  \frac{\norm{\mathbf{x}_*-\mathbf{x}_k}_A}{\norm{\mathbf{x}_*-\mathbf{x}_0}_A}
  \leqslant2\left(\frac{\alpha-1}{\alpha+1}\right)^{k-p-q}
  \cdot\max_{\lambda\in[b_1,b_2]}\prod_{j=1}^p
  \left(\frac{\lambda-\lambda_j}{\lambda_j}\right),
\]
其中 $\alpha\equiv(b_2/b_1)^{1/2}\geqslant1$.
\end{theorem}

\begin{proof}
令 $\mathbf{e}_k=\mathbf{x}_*-\mathbf{x}_k$. 设 $\mathcal{P}_k$ 为次数不超过 $k$ 且满足 $P_k(0)=1$ 的一切多项式 $P_k$ 的集合. 对任意 $P_k\in\mathcal{P}_k$, 由 \eqref{eq:6-10} 有
\begin{equation}
  \norm{\mathbf{e}_k}_A
  \leqslant\min_{P_k\in\mathcal{P}_k}\max_{1\leqslant j\leqslant n}
  |P_k(\lambda_j)|\cdot\norm{\mathbf{e}_0}_A
  \leqslant\max_{1\leqslant i\leqslant n}|P_k(\lambda_i)|
  \cdot\norm{\mathbf{e}_0}_A.
  \label{eq:6-12}
\end{equation}
给定多项式 $P_{k-p-q}\in\mathcal{P}_{k-p-q}$, 构造如下特殊多项式:
\[
  \widehat{P}_k\equiv P_{k-p-q}\cdot(-1)^p
  \cdot\prod_{j=1}^p\left(\frac{\lambda-\lambda_j}{\lambda_j}\right)
  \cdot\prod_{j=n-q+1}^n\left(\frac{\lambda_j-\lambda}{\lambda_j}\right).
\]
对这样定义的 $\widehat{P}_k$, 当 $i=1,2,\ldots,p,n-q+1,\ldots,n$ 时, 有
\begin{equation}
  \widehat{P}_k(\lambda_i)=(-1)^pP_{k-p-q}(\lambda_i)
  \cdot\prod_{j=1}^p\left(\frac{\lambda_i-\lambda_j}{\lambda_j}\right)
  \cdot\prod_{j=n-q+1}^n\left(\frac{\lambda_j-\lambda_i}{\lambda_j}\right)=0.
  \label{eq:6-13}
\end{equation}
这里可以看出
\begin{equation}
  \widehat{P}_k\in\mathcal{P}_k.
  \label{eq:6-14}
\end{equation}
根据 \eqref{eq:6-12}, \eqref{eq:6-13} 与 \eqref{eq:6-14}, 有
\[
  \begin{aligned}
    \norm{\mathbf{e}_k}_A
    &\leqslant\max_{1\leqslant i\leqslant n}|\widehat{P}_k(\lambda_i)|
      \cdot\norm{\mathbf{e}_0}_A\\
    &\leqslant\max_{p+1\leqslant i\leqslant n-q}
      \left|P_{k-p-q}(\lambda_i)\right|
      \cdot\max_{p+1\leqslant i\leqslant n-q}\left|\prod_{j=1}^p
      \left(\frac{\lambda_i-\lambda_j}{\lambda_j}\right)\right.\\
    &\qquad\cdot\left.\prod_{j=n-q+1}^n
      \left(\frac{\lambda_j-\lambda_i}{\lambda_j}\right)\right|
      \norm{\mathbf{e}_0}_A\\
    &\leqslant\max_{b_1\leqslant\lambda\leqslant b_2}
      \left|P_{k-p-q}(\lambda)\right|
      \cdot\max_{b_1\leqslant\lambda\leqslant b_2}\left|\prod_{j=1}^p
      \left(\frac{\lambda-\lambda_j}{\lambda_j}\right)\right.\\
    &\qquad\cdot\left.\prod_{j=n-q+1}^n
      \left(\frac{\lambda_j-\lambda}{\lambda_j}\right)\right|
      \norm{\mathbf{e}_0}_A.
  \end{aligned}
\]
由于上述不等式对任意 $P_{k-p-q}\in\mathcal{P}_{k-p-q}$ 成立, 故由 \eqref{eq:6-11} 得
\begin{equation}
  \max_{b_1\leqslant\lambda\leqslant b_2}\left|P_{k-p-q}(\lambda)\right|
  \leqslant2\left(\frac{\alpha-1}{\alpha+1}\right)^{k-p-q},
  \qquad\alpha=\sqrt{b_2/b_1}.
  \label{eq:6-15}
\end{equation}
此外, 当 $b_1\leqslant\lambda\leqslant b_2$ 时,
\begin{equation}
  0\leqslant\frac{\lambda_j-\lambda}{\lambda_j}\leqslant1.
  \label{eq:6-16}
\end{equation}
最后由 \eqref{eq:6-15} 与 \eqref{eq:6-16} 得
\[
  \norm{\mathbf{e}_k}_A
  \leqslant2\left(\frac{\alpha-1}{\alpha+1}\right)^{k-p-q}
  \cdot\max_{\lambda\in[b_1,b_2]}\prod_{j=1}^p
  \left(\frac{\lambda-\lambda_j}{\lambda_j}\right)\norm{\mathbf{e}_0}_A.
  \qedhere
\]
\end{proof}

由定理 \ref{thm:6-7} 可见, 当 $n$ 增大时, 若 $p$, $q$ 为不依赖于 $n$ 的常数且 $\lambda_1$ 一致地有正下界, 则收敛速度是线性的, 即迭代次数与 $n$ 无关. 我们还注意到, 特征值聚得越紧, 收敛越快.

\begin{corollary}
\label{cor:6-1}
若对称正定矩阵 $A$ 的特征值 $\lambda_j$ 满足如下排序:
\[
  0<\delta<\lambda_1\leqslant\cdots\leqslant\lambda_p\leqslant1-\epsilon
  \leqslant\lambda_{p+1}\leqslant\cdots\leqslant\lambda_{n-q}\leqslant1+\epsilon
  \leqslant\lambda_{n-q+1}\leqslant\cdots\leqslant\lambda_n,
\]
其中 $0<\epsilon<1$, 则
\[
  \frac{\norm{\mathbf{x}_*-\mathbf{x}_k}_A}{\norm{\mathbf{x}_*-\mathbf{x}_0}_A}
  \leqslant2\left(\frac{1+\epsilon}{\delta}\right)^p\epsilon^{k-p-q},
\]
其中 $k\geqslant p+q$.
\end{corollary}

\begin{proof}
对定理 \ref{thm:6-7} 中的 $\alpha$, 有
\[
  \alpha\equiv\left(\frac{b_2}{b_1}\right)^{\frac{1}{2}}
  =\left(\frac{1+\epsilon}{1-\epsilon}\right)^{\frac{1}{2}}.
\]
因此
\[
  \frac{\alpha-1}{\alpha+1}=\frac{1-\sqrt{1-\epsilon^2}}{\epsilon}<\epsilon.
\]
对 $1\leqslant j\leqslant p$ 与 $\lambda\in[1-\epsilon,1+\epsilon]$, 有
\[
  0\leqslant\frac{\lambda-\lambda_j}{\lambda_j}\leqslant\frac{1+\epsilon}{\delta}.
\]
于是利用定理 \ref{thm:6-7} 即得
\[
  \begin{aligned}
    \frac{\norm{\mathbf{x}_*-\mathbf{x}_k}_A}{\norm{\mathbf{x}_*-\mathbf{x}_0}_A}
    &\leqslant2\left(\frac{\alpha-1}{\alpha+1}\right)^{k-p-q}
      \cdot\max_{\lambda\in[b_1,b_2]}\prod_{j=1}^p
      \left(\frac{\lambda-\lambda_j}{\lambda_j}\right)\\
    &\leqslant2\left(\frac{1+\epsilon}{\delta}\right)^p\epsilon^{k-p-q}.
  \end{aligned}
  \qedhere
\]
\end{proof}

\section{预处理}
\label{sec:6-4}

由第 3 节可知, 若方程组
\[
  A\mathbf{x}=\mathbf{b}
\]
的系数矩阵 $A$ 是良态的或其谱是聚簇的, 则 CG 法的收敛速度很快. 因此, 为了加快收敛速度, 通常对方程组作\newterm{preconditioning}{预处理}{preconditioning}, 即不直接求解原方程组, 而求解如下预处理后的方程组
\begin{equation}
  \widetilde{A}\tilde{\mathbf{x}}=\tilde{\mathbf{b}},
  \label{eq:6-17}
\end{equation}
其中
\[
  \widetilde{A}=C^{-1}AC^{-1},\qquad
  \tilde{\mathbf{x}}=C\mathbf{x},\qquad
  \tilde{\mathbf{b}}=C^{-1}\mathbf{b},
\]
且 $C$ 对称正定. 我们希望预处理后的矩阵 $\widetilde{A}$ 具有比 $A$ 更好的谱性质.

对 \eqref{eq:6-17} 使用 CG 法, 得
\begin{equation}
  \left\{
  \begin{array}{l}
    \alpha_k=\dfrac{\tilde{\mathbf{r}}_k^T\tilde{\mathbf{r}}_k}
      {\tilde{\mathbf{p}}_k^T\widetilde{A}\tilde{\mathbf{p}}_k},\\[1mm]
    \tilde{\mathbf{x}}_{k+1}=\tilde{\mathbf{x}}_k+\alpha_k\tilde{\mathbf{p}}_k,\\[1mm]
    \tilde{\mathbf{r}}_{k+1}=\tilde{\mathbf{r}}_k
      -\alpha_k\widetilde{A}\tilde{\mathbf{p}}_k,\\[1mm]
    \beta_k=\dfrac{\tilde{\mathbf{r}}_{k+1}^T\tilde{\mathbf{r}}_{k+1}}
      {\tilde{\mathbf{r}}_k^T\tilde{\mathbf{r}}_k},\\[1mm]
    \tilde{\mathbf{p}}_{k+1}=\tilde{\mathbf{r}}_{k+1}
      +\beta_k\tilde{\mathbf{p}}_k,
  \end{array}
  \right.
  \label{eq:6-18}
\end{equation}
其中 $\tilde{\mathbf{x}}_0$ 为任意给定的初始向量, $\tilde{\mathbf{r}}_0=\tilde{\mathbf{b}}-\widetilde{A}\tilde{\mathbf{x}}_0$, $\tilde{\mathbf{p}}_0=\tilde{\mathbf{r}}_0$. 令
\[
  \tilde{\mathbf{x}}_k=C\mathbf{x}_k,\qquad
  \tilde{\mathbf{r}}_k=C^{-1}\mathbf{r}_k,\qquad
  \tilde{\mathbf{p}}_k=C\mathbf{p}_k,\qquad
  M=C^2.
\]
把它们代入 \eqref{eq:6-18}, 实际得到
\[
  \begin{array}{ll}
    \mathbf{w}_k=A\mathbf{p}_k, &
    \alpha_k=\rho_k/\mathbf{p}_k^T\mathbf{w}_k,\\[1mm]
    \mathbf{x}_{k+1}=\mathbf{x}_k+\alpha_k\mathbf{p}_k, &
    \mathbf{r}_{k+1}=\mathbf{r}_k-\alpha_k\mathbf{w}_k,\\[1mm]
    \mathbf{z}_{k+1}=M^{-1}\mathbf{r}_{k+1}, &
    \rho_{k+1}=\mathbf{r}_{k+1}^T\mathbf{z}_{k+1},\\[1mm]
    \beta_k=\rho_{k+1}/\rho_k, &
    \mathbf{p}_{k+1}=\mathbf{z}_{k+1}+\beta_k\mathbf{p}_k,
  \end{array}
\]
其中 $\mathbf{x}_0$ 为任意给定的初始向量, $\mathbf{r}_0=\mathbf{b}-A\mathbf{x}_0$, $\mathbf{z}_0=M^{-1}\mathbf{r}_0$, $\rho_0=\mathbf{r}_0^T\mathbf{z}_0$, $\mathbf{p}_0=\mathbf{z}_0$.

于是得到如下预处理算法. 在初始化步 $k=0$, 选定初始向量 $\mathbf{x}$ 并计算 $\mathbf{r}=\mathbf{b}-A\mathbf{x}$. 当 $\sqrt{\mathbf{r}^T\mathbf{r}}>\epsilon\norm{\mathbf{b}}_2$ 且 $k<k_{\max}$ 时, 迭代步为
\[
  \left\{
  \begin{array}{l}
    \text{解}\ M\mathbf{z}=\mathbf{r}\ \text{得}\ \mathbf{z}\\
    k=k+1\\
    \textbf{if}\ k=1\\
    \quad \mathbf{p}=\mathbf{z}\\
    \textbf{else}\\
    \quad \beta=\rho/\tilde{\rho};\qquad \mathbf{p}=\mathbf{z}+\beta\mathbf{p}\\
    \textbf{end}\\
    \mathbf{w}=A\mathbf{p};\qquad
    \alpha=\rho/\mathbf{p}^T\mathbf{w};\qquad
    \mathbf{x}=\mathbf{x}+\alpha\mathbf{p}\\
    \mathbf{r}=\mathbf{r}-\alpha\mathbf{w};\qquad
    \tilde{\rho}=\rho;\qquad
    \rho=\mathbf{r}^T\mathbf{z}
  \end{array}
  \right.
\]
其中 $\mathbf{z},\mathbf{r},\mathbf{p},\mathbf{w}$ 为向量, $\rho,\alpha,\beta$ 为标量. 该算法称为\newterm{pcg}{预处理共轭梯度 (PCG) 法}{preconditioned conjugate gradient (PCG) method}. 注意 PCG 法具有如下性质:
\begin{romanlist}
  \item 当 $i\neq j$ 时, $\mathbf{r}_i^TM^{-1}\mathbf{r}_j=0$;
  \item 当 $i\neq j$ 时, $\mathbf{p}_i^TA\mathbf{p}_j=0$;
  \item 近似解 $\mathbf{x}_k$ 满足
    \[
      \norm{\mathbf{x}_*-\mathbf{x}_k}_{\widetilde{A}}
      \leqslant2\left(\frac{\sqrt{\kappa}-1}{\sqrt{\kappa}+1}\right)^k
      \norm{\mathbf{x}_*-\mathbf{x}_0}_{\widetilde{A}},
    \]
    其中 $\kappa=\lambda_n/\lambda_1$, $\lambda_n$ 与 $\lambda_1$ 分别为 $M^{-1}A$ 的最大与最小特征值.
\end{romanlist}

选取一个好的\newterm{preconditioner}{预条件子}{preconditioner} $M=C^2$ 需考虑三条准则 \cite{src:2,src:16,src:32,src:42,src:43}:
\begin{parenlist}
  \item $M$ 易于构造;
  \item $M\mathbf{z}=\mathbf{d}$ 易于求解;
  \item $M^{-1}A$ 的谱是聚簇的, 且 (或)$M^{-1}A$ 相对于 $A$ 是良态的. 注意 $\widetilde{A}$ 相似于 $M^{-1}A$.
\end{parenlist}
对预处理后的方程组 $\widetilde{A}\tilde{\mathbf{x}}=\tilde{\mathbf{b}}$ 实施 CG 法, 其主要工作量在于每次迭代都要对某个向量 $\mathbf{v}$ 计算矩阵与向量的乘积 $M^{-1}A\mathbf{v}$.

通常很难选出同时满足这三条准则的预条件子. 下面简要讨论三类预条件子.
\begin{parenlist}
  \item \textbf{对角预条件子}. 若系数矩阵 $A$ 的对角元彼此相差很大, 可以用矩阵
    \[
      M=\diag(a_{11},\ldots,a_{nn})
    \]
    作为预条件子, 以期加快 CG 法的收敛速度. 对于分块矩阵
    \[
      A=\begin{bmatrix}
        A_{11} & \cdots & A_{1k}\\
        \vdots & \ddots & \vdots\\
        A_{k1} & \cdots & A_{kk}
      \end{bmatrix},
    \]
    若 $A_{ii}^{-1}$ 易于求得, 则可以用分块对角矩阵
    \[
      M=\diag(A_{11},\ldots,A_{kk})
    \]
    作为预条件子.

  \item \textbf{基于不完全 Cholesky 分解的预条件子}. 若先计算\newterm{incompletecholesky}{不完全 Cholesky 分解}{incomplete Cholesky factorization} \cite{src:32}
    \[
      A=LL^T+R,
    \]
    则可以用矩阵 $M=LL^T$ 作为预条件子. 我们可以要求矩阵 $L$ 与矩阵 $A$ 具有相同的稀疏结构, 且 $LL^T\approx A$.

  \item \textbf{最优 (循环) 预条件子}. 这类预条件子在文献 \cite{src:16,src:18,src:21,src:42,src:43,src:77} 中研究过. \newterm{circulant}{循环矩阵}{circulant matrix}定义如下:
    \[
      C_n=\begin{bmatrix}
        c_0 & c_{-1} & \cdots & c_{2-n} & c_{1-n}\\
        c_1 & c_0 & c_{-1} & \cdots & c_{2-n}\\
        \vdots & c_1 & c_0 & \ddots & \vdots\\
        c_{n-2} & \cdots & \ddots & \ddots & c_{-1}\\
        c_{n-1} & c_{n-2} & \cdots & c_1 & c_0
      \end{bmatrix},
    \]
    其中 $c_{-k}=c_{n-k}$, $1\leqslant k\leqslant n-1$. 众所周知, 循环矩阵可被 \newterm{fouriermatrix}{Fourier 矩阵}{Fourier matrix} $F_n$ 对角化 \cite{src:29}, 即
    \begin{equation}
      C_n=F_n^*\Lambda_nF_n,
      \label{eq:6-19}
    \end{equation}
    其中 $F_n$ 的元素为
    \[
      (F_n)_{j,k}=\frac{1}{\sqrt{n}}e^{2\pi\ii(j-1)(k-1)/n},
      \qquad \ii\equiv\sqrt{-1},
    \]
    $1\leqslant j,k\leqslant n$, 而 $\Lambda_n$ 为以 $C_n$ 的特征值为对角元的对角矩阵. 我们指出, 对 $C_n$ 的第一列作\newterm{fft}{快速 Fourier 变换 (FFT)}{fast Fourier transform (FFT)}, 即可用 $\mathcal{O}(n\log n)$ 次运算得到 $\Lambda_n$. 一旦得到 $\Lambda_n$, 对任意向量 $\mathbf{y}$, 乘积 $C_n\mathbf{y}$ 与 $C_n^{-1}\mathbf{y}$ 都可用 FFT 在 $\mathcal{O}(n\log n)$ 次运算内算出. 对 Fourier 矩阵 $F_n$, 在不引起混淆时我们记其为 $F$.
\end{parenlist}

下面研究一类称为\newterm{optimalprecond}{最优预条件子}{optimal preconditioner}的预条件子 \cite{src:16,src:18,src:21,src:42,src:43}. 给定任意酉矩阵 $U\in\C^{n\times n}$, 设 $\mathcal{M}_U$ 为所有可被 $U$ 同时对角化的矩阵之集, 即
\begin{equation}
  \mathcal{M}_U=\left\{U^*\Lambda U:
  \ \Lambda\ \text{为}\ n\times n\ \text{对角矩阵}\right\}.
  \label{eq:6-20}
\end{equation}
我们指出, 在 \eqref{eq:6-20} 中取 $U=F$(Fourier 矩阵) 时, $\mathcal{M}_F$ 就是全体循环矩阵之集. 记 $\delta(A)$ 为对角元与矩阵 $A$ 的对角元相同的对角矩阵. 有如下引理 \cite{src:16,src:42,src:43,src:51,src:82}.

\begin{lemma}
\label{lem:6-2}
对任意 $A=[a_{pq}]\in\C^{n\times n}$, 设 $c_U(A)$ 为 $\norm{W-A}_F$ 在一切 $W\in\mathcal{M}_U$ 上的极小点. 则
\begin{romanlist}
  \item $c_U(A)$ 由 $A$ 唯一确定, 且
    \[
      c_U(A)=U^*\delta(UAU^*)U.
    \]
  \item 若 $A$ 为 Hermitian 矩阵, 则 $c_U(A)$ 也是 Hermitian 矩阵. 进而, 若 $\lambda_{\min}(\cdot)$ 与 $\lambda_{\max}(\cdot)$ 分别表示最小与最大特征值, 则
    \[
      \lambda_{\min}(A)\leqslant\lambda_{\min}(c_U(A))
      \leqslant\lambda_{\max}(c_U(A))\leqslant\lambda_{\max}(A).
    \]
    特别地, 若 $A$ 正定, 则 $c_U(A)$ 也正定.
  \item 若 $A$ 正规且稳定, 即 $A^*A=AA^*$ 且 $A$ 的所有特征值实部为负, 则 $c_U(A)$ 也稳定.
  \item 当 $U$ 为 Fourier 矩阵 $F$ 时,
    \[
      c_F(A)=\sum_{j=0}^{n-1}
      \left(\frac{1}{n}\sum_{p-q\equiv j\ (\mathrm{mod}\ n)}a_{pq}\right)Q^j,
    \]
    其中 $Q$ 为 $n\times n$ 循环矩阵
    \[
      Q\equiv\begin{bmatrix}
        0 & & & & 1\\
        1 & 0 & & & \\
        0 & 1 & \ddots & & \\
        \vdots & \ddots & \ddots & \ddots & \\
        0 & \cdots & 0 & 1 & 0
      \end{bmatrix}.
    \]
\end{romanlist}
\end{lemma}

矩阵 $c_U(A)$ 称为 $A$ 的最优预条件子, 矩阵 $c_F(A)$ 称为 $A$ 的\newterm{optimalcirculant}{最优循环预条件子}{optimal circulant preconditioner}. 我们指出, 对一大类结构型线性方程组 $A\mathbf{x}=\mathbf{b}$ (例如 Toeplitz 方程组, Hankel 方程组等), $c_F(A)$ 是一个良好的预条件子, 参见 \cite{src:16,src:18,src:21,src:25,src:42,src:43,src:49,src:50,src:52}. 其他预条件子的应用可参阅 \cite{src:4,src:67,src:97} 以及第 9, 10 章.

\section{GMRES 方法}
\label{sec:6-5}

本节介绍求解一般方程组
\[
  A\mathbf{x}=\mathbf{b}
\]
的\newterm{gmres}{广义极小残量 (GMRES) 法}{generalized minimum residual (GMRES) method}, 其中 $A\in\R^{n\times n}$ 非奇异. GMRES 法由 Saad 与 Schultz 于 1986 年提出, 是求解非对称方程组最重要的 Krylov 子空间方法之一 \cite{src:70,src:71}.

\subsection{GMRES 方法的基本性质}
\label{subsec:6-5-1}

对 GMRES 法, 在第 $k$ 次迭代中, 我们要求解如下最小二乘问题的解 $\mathbf{x}_k$:
\[
  \min_{\mathbf{x}\in\mathbf{x}_0+\mathcal{K}(A,\mathbf{r}_0,k)}
  \norm{\mathbf{b}-A\mathbf{x}}_2,
\]
其中
\[
  \mathcal{K}(A,\mathbf{r}_0,k)
  \equiv\spn\left\{\mathbf{r}_0,A\mathbf{r}_0,\ldots,A^{k-1}\mathbf{r}_0\right\},
\]
$\mathbf{r}_0=\mathbf{b}-A\mathbf{x}_0$. 设 $\mathbf{x}\in\mathbf{x}_0+\mathcal{K}(A,\mathbf{r}_0,k)$, 则有
\[
  \mathbf{x}=\mathbf{x}_0+\sum_{j=0}^{k-1}\gamma_jA^j\mathbf{r}_0,
\]
进而
\[
  \mathbf{r}=\mathbf{b}-A\mathbf{x}
  =\mathbf{b}-A\mathbf{x}_0-\sum_{j=0}^{k-1}\gamma_jA^{j+1}\mathbf{r}_0
  =\mathbf{r}_0-\sum_{j=1}^{k}\gamma_{j-1}A^j\mathbf{r}_0.
\]
因此
\[
  \mathbf{r}=\bar{P}_k(A)\mathbf{r}_0,
\]
其中 $\bar{P}_k\in\mathcal{P}_k$, $\mathcal{P}_k$ 为次数不超过 $k$ 且满足 $P_k(0)=1$ 的一切多项式 $P_k$ 的集合. 于是有如下定理.

\begin{theorem}
\label{thm:6-8}
设 $\mathbf{x}_k$ 为 GMRES 法第 $k$ 次迭代后的解. 则
\[
  \norm{\mathbf{r}_k}_2=\min_{P\in\mathcal{P}_k}\norm{P(A)\mathbf{r}_0}_2
  \leqslant\norm{\bar{P}_k(A)\mathbf{r}_0}_2.
\]
进而
\[
  \frac{\norm{\mathbf{r}_k}_2}{\norm{\mathbf{r}_0}_2}
  \leqslant\norm{\bar{P}_k(A)}_2.
\]
\end{theorem}

此外还有如下定理.

\begin{theorem}
\label{thm:6-9}
设 $A\in\R^{n\times n}$, 则 GMRES 法至多 $n$ 次迭代即可得到 $A\mathbf{x}=\mathbf{b}$ 的精确解.
\end{theorem}

\begin{proof}
$A$ 的特征多项式为
\[
  p_A(z)=\det(zI-A),
\]
其阶为 $n$, 且 $p_A(0)=(-1)^n\det(A)\neq0$. 于是
\[
  \bar{P}_n(z)=\frac{p_A(z)}{p_A(0)}\in\mathcal{P}_n.
\]
由 Hamilton-Cayley 定理 \cite{src:39} 知
\[
  \bar{P}_n(A)=p_A(A)=\mathbf{0}.
\]
由定理 \ref{thm:6-8} 得
\[
  \mathbf{r}_n=\mathbf{b}-A\mathbf{x}_n=\mathbf{0}.
\]
因此 $\mathbf{x}_n$ 为 $A\mathbf{x}=\mathbf{b}$ 的精确解.
\end{proof}

若 $A$ 可对角化, 即 $A=V\Lambda V^{-1}$, 其中 $\Lambda$ 为对角矩阵, 则
\[
  P(A)=VP(\Lambda)V^{-1}.
\]
若 $V$ 为正交矩阵, 则 $A$ 为\newterm{normalmatrix}{正规矩阵}{normal matrix}.

\begin{theorem}
\label{thm:6-10}
设 $A=V\Lambda V^{-1}$. 则
\[
  \frac{\norm{\mathbf{r}_k}_2}{\norm{\mathbf{r}_0}_2}
  \leqslant\kappa_2(V)\min_{P\in\mathcal{P}_k}
  \max_{\lambda_i}|P(\lambda_i)|,
\]
其中 $\lambda_i$ ($i=1,2,\ldots,n$) 为矩阵 $A$ 的特征值.
\end{theorem}

\begin{proof}
由定理 \ref{thm:6-8}, 有
\[
  \begin{aligned}
    \norm{\mathbf{r}_k}_2
    &=\min_{P\in\mathcal{P}_k}\norm{P(A)\mathbf{r}_0}_2
     =\min_{P\in\mathcal{P}_k}\norm{VP(\Lambda)V^{-1}\mathbf{r}_0}_2\\
    &\leqslant\min_{P\in\mathcal{P}_k}\norm{V}_2\norm{V^{-1}}_2
      \norm{P(\Lambda)}_2\norm{\mathbf{r}_0}_2\\
    &=\kappa_2(V)\min_{P\in\mathcal{P}_k}\norm{P(\Lambda)}_2
      \norm{\mathbf{r}_0}_2\\
    &=\kappa_2(V)\min_{P\in\mathcal{P}_k}\max_{\lambda_i}|P(\lambda_i)|
      \cdot\norm{\mathbf{r}_0}_2.
  \end{aligned}
\]
\end{proof}

我们指出, 当 $A$ 为正规矩阵时, $\kappa_2(V)=1$. 因此
\[
  \frac{\norm{\mathbf{r}_k}_2}{\norm{\mathbf{r}_0}_2}
  \leqslant\min_{P\in\mathcal{P}_k}\max_{\lambda_i}|P(\lambda_i)|.
\]

\begin{theorem}
\label{thm:6-11}
若 $A$ 可对角化且恰有 $k$ 个互不相同的特征值, 则 GMRES 法至多 $k$ 次迭代即终止.
\end{theorem}

\begin{proof}
构造多项式
\[
  P(z)=\prod_{i=1}^{k}\frac{\lambda_i-z}{\lambda_i}\in\mathcal{P}_k,
\]
其中 $\lambda_i$ 为 $A$ 的特征值. 由定理 \ref{thm:6-10} 知 $\mathbf{r}_k=\mathbf{0}$, 即 $A\mathbf{x}_k=\mathbf{b}$.
\end{proof}

应当强调, 一般而言 GMRES 法的表现不能仅由特征值决定. 事实上, 文献 \cite{src:35} 指出: 对某个非正规矩阵问题, GMRES 法可以具有任意预先指定的非增收敛速度, 而且该问题的特征值可以具有任意想要的分布. 因此, 例如紧聚在 1 附近的特征值对非正规矩阵未必是好事, 而在正规矩阵情形下却是好事. 不过, 我们仍有如下两个定理.

\begin{theorem}
\label{thm:6-12}
若 $\mathbf{b}$ 是 $A$ 的 $k$ 个特征向量的线性组合, 即
\[
  \mathbf{b}=\sum_{l=1}^{k}\gamma_l\mathbf{u}_{i_l},
\]
则 GMRES 法至多 $k$ 次迭代即终止.
\end{theorem}

\begin{proof}
首先把特征向量组 $\{\mathbf{u}_{i_l}\}$ 扩充为 $\R^n$ 的一组基, 即向量
\[
  \mathbf{u}_{i_1},\mathbf{u}_{i_2},\ldots,\mathbf{u}_{i_k},
  \mathbf{v}_1,\ldots,\mathbf{v}_{n-k}
\]
构成 $\R^n$ 的一组基. 设 $\mathbf{x}_*$ 为 $A\mathbf{x}=\mathbf{b}$ 的精确解, 则
\[
  \mathbf{x}_*=\sum_{l=1}^{k}\alpha_l\mathbf{u}_{i_l}
  +\sum_{j=1}^{n-k}\beta_j\mathbf{v}_j.
\]
进而
\[
  \begin{aligned}
    A\mathbf{x}_*&=\sum_{l=1}^{k}\alpha_lA\mathbf{u}_{i_l}
      +\sum_{j=1}^{n-k}\beta_jA\mathbf{v}_j\\
    &=\sum_{l=1}^{k}\alpha_l\lambda_{i_l}\mathbf{u}_{i_l}
      +\sum_{j=1}^{n-k}\beta_jA\mathbf{v}_j
     =\mathbf{b}=\sum_{l=1}^{k}\gamma_l\mathbf{u}_{i_l}.
  \end{aligned}
\]
因此
\[
  \left\{
  \begin{array}{ll}
    \beta_j=0, & j=1,2,\ldots,n-k;\\
    \alpha_l=\gamma_l/\lambda_{i_l}, & l=1,2,\ldots,k.
  \end{array}
  \right.
\]
于是
\[
  \mathbf{x}_*=\sum_{l=1}^{k}(\gamma_l/\lambda_{i_l})\mathbf{u}_{i_l}.
\]
令
\[
  \bar{P}_k(z)=\prod_{l=1}^{k}\frac{\lambda_{i_l}-z}{\lambda_{i_l}}
  \in\mathcal{P}_k.
\]
注意当 $1\leqslant l\leqslant k$ 时 $\bar{P}_k(\lambda_{i_l})=0$, 且
\[
  \bar{P}_k(A)\mathbf{x}_*=\sum_{l=1}^{k}\bar{P}_k(\lambda_{i_l})
  (\gamma_l/\lambda_{i_l})\mathbf{u}_{i_l}=\mathbf{0}.
\]
于是, 取初始向量 $\mathbf{x}_0=\mathbf{0}$, 得
\[
  \begin{aligned}
    \norm{\mathbf{r}_k}_2
    &\leqslant\norm{\bar{P}_k(A)\mathbf{r}_0}_2
     =\norm{\bar{P}_k(A)\mathbf{b}}_2\\
    &=\norm{\bar{P}_k(A)A\mathbf{x}_*}_2
     =\norm{A\bar{P}_k(A)\mathbf{x}_*}_2=0.
  \end{aligned}
\]
\end{proof}

\begin{theorem}
\label{thm:6-13}
用 GMRES 法求解线性方程组 $A\mathbf{x}=\mathbf{b}$, 其中 $A=I+L$, 则该方法至多 $\rank(L)+1$ 次迭代即收敛.
\end{theorem}

\begin{proof}
先回忆: $\mathbf{r}_0$ 关于 $A$ 的\newterm{minimalpoly}{极小多项式}{minimal polynomial}是指使 $p(A)\mathbf{r}_0=\mathbf{0}$ 成立的最低次首一非零多项式 $p$ \cite{src:70}. 由定理 \ref{thm:6-8}, GMRES 法必在 $\nu$ 次迭代内收敛, 其中 $\nu$ 为残差 $\mathbf{r}_0$ 关于 $A=I+L$ 的极小多项式的次数. 设 $\mu$ 为 $\mathbf{r}_0$ 关于 $L$ 的极小多项式的次数. 由于
\[
  \sum_{i=0}^{\mu}\alpha_i(I+L)^i\mathbf{r}_0=\mathbf{0}
\]
蕴含
\[
  \sum_{i=0}^{\mu}\beta_iL^i\mathbf{r}_0=\mathbf{0}
\]
(其中 $\beta_i$ 为某些常数), 反之亦然, 故 $\nu=\mu$. 此外, 由 $\mu$ 的定义, 集合
\[
  \left\{\mathbf{r}_0,L\mathbf{r}_0,\ldots,L^{\mu-1}\mathbf{r}_0\right\}
\]
线性无关. 设 $B$ 为 $L$ 的列向量空间, 则 $B$ 的维数等于 $L$ 的秩. 由于当 $i\geqslant1$ 时 $L^i\mathbf{r}_0\in B$, 故
\[
  \left\{L\mathbf{r}_0,\ldots,L^{\mu-1}\mathbf{r}_0\right\}\subseteq B.
\]
因此 $\mu-1\leqslant\rank(L)$, 即 $\mu\leqslant\rank(L)+1$. 于是 GMRES 法在
\[
  \nu=\mu\leqslant\rank(L)+1
\]
次迭代内收敛.
\end{proof}

\subsection{GMRES 方法的实现}
\label{subsec:6-5-2}

回忆 GMRES 法第 $k$ 次迭代所要求解的最小二乘问题为
\[
  \min_{\mathbf{x}\in\mathbf{x}_0+\mathcal{K}(A,\mathbf{r}_0,k)}
  \norm{\mathbf{b}-A\mathbf{x}}_2.
\]
设已有矩阵
\[
  V_k=\left[\mathbf{v}_1^k\mid\mathbf{v}_2^k\mid\cdots\mid\mathbf{v}_k^k\right],
\]
其列构成 $\mathcal{K}(A,\mathbf{r}_0,k)$ 的一组标准正交基. 则对任意 $\mathbf{z}\in\mathcal{K}(A,\mathbf{r}_0,k)$, 可写为
\[
  \mathbf{z}=\sum_{l=1}^{k}u_l\mathbf{v}_l^k=V_k\mathbf{u},
\]
其中 $\mathbf{u}=[u_1,u_2,\ldots,u_k]^T\in\R^k$. 因此, 一旦求出 $V_k$, 就可以把 Krylov 子空间中的原最小二乘问题转化为 $\R^k$ 中的最小二乘问题如下. 设 $\mathbf{x}_k$ 为第 $k$ 次迭代后的解, 则
\[
  \mathbf{x}_k=\mathbf{x}_0+V_k\mathbf{y}_k,
\]
其中向量 $\mathbf{y}_k$ 极小化
\[
  \min_{\mathbf{y}\in\R^k}\norm{\mathbf{b}-A(\mathbf{x}_0+V_k\mathbf{y})}_2
  =\min_{\mathbf{y}\in\R^k}\norm{\mathbf{r}_0-AV_k\mathbf{y}}_2.
\]
这是一个标准线性最小二乘问题, 可用 QR 分解求解.

可以用 Gram-Schmidt 正交化求 $\mathcal{K}(A,\mathbf{r}_0,k)$ 的一组标准正交基. 算法如下:
\begin{parenlist}
  \item 定义 $\mathbf{r}_0=\mathbf{b}-A\mathbf{x}_0$ 与
    $\mathbf{v}_1=\mathbf{r}_0/\norm{\mathbf{r}_0}_2$.
  \item 对 $i=1,2,\ldots,k-1$, 计算
    \[
      \mathbf{v}_{i+1}
      =\frac{A\mathbf{v}_i-\sum_{j=1}^{i}
      ((A\mathbf{v}_i)^T\mathbf{v}_j)\mathbf{v}_j}
      {\norm{A\mathbf{v}_i-\sum_{j=1}^{i}
      ((A\mathbf{v}_i)^T\mathbf{v}_j)\mathbf{v}_j}_2}.
    \]
\end{parenlist}
该算法产生的矩阵 $V_k$ 的列也构成 $\mathcal{K}(A,\mathbf{r}_0,k)$ 的一组标准正交基. 我们指出, 当出现除以零时, 就发生\newterm{breakdown}{中断}{breakdown}. 关于中断的发生, 有如下定理.

\begin{theorem}
\label{thm:6-14}
设 $A$ 非奇异, 向量 $\mathbf{v}_j$ 由上述算法产生, $i$ 为使
\begin{equation}
  A\mathbf{v}_i-\sum_{j=1}^{i}((A\mathbf{v}_i)^T\mathbf{v}_j)\mathbf{v}_j
  =\mathbf{0}
  \label{eq:6-21}
\end{equation}
成立的最小整数. 则 $\mathbf{x}=A^{-1}\mathbf{b}
\in\mathbf{x}_0+\mathcal{K}(A,\mathbf{r}_0,i)$.
\end{theorem}

\begin{proof}
由 \eqref{eq:6-21} 有
\[
  A\mathbf{v}_i=\sum_{j=1}^{i}((A\mathbf{v}_i)^T\mathbf{v}_j)\mathbf{v}_j
  \in\mathcal{K}(A,\mathbf{r}_0,i),
\]
故
\[
  A\mathcal{K}(A,\mathbf{r}_0,i)\subset\mathcal{K}(A,\mathbf{r}_0,i).
\]
注意 $V_i=[\mathbf{v}_1\mid\mathbf{v}_2\mid\cdots\mid\mathbf{v}_i]$ 的列构成 $\mathcal{K}(A,\mathbf{r}_0,i)$ 的一组标准正交基, 即
\[
  \mathcal{K}(A,\mathbf{r}_0,i)=\spn\{\mathbf{v}_1,\mathbf{v}_2,\ldots,\mathbf{v}_i\},
\]
于是
\begin{equation}
  AV_i=V_iH,
  \label{eq:6-22}
\end{equation}
其中 $H\in\R^{i\times i}$ 非奇异 (因 $A$ 非奇异). 由于 $\mathbf{x}_i-\mathbf{x}_0
\in\mathcal{K}(A,\mathbf{r}_0,i)$, 存在向量 $\mathbf{y}\in\R^i$ 使 $\mathbf{x}_i-\mathbf{x}_0=V_i\mathbf{y}$. 因此
\begin{equation}
  \norm{\mathbf{r}_i}_2=\norm{\mathbf{b}-A\mathbf{x}_i}_2
  =\norm{\mathbf{r}_0-A(\mathbf{x}_i-\mathbf{x}_0)}_2
  =\norm{\mathbf{r}_0-AV_i\mathbf{y}}_2.
  \label{eq:6-23}
\end{equation}
令 $\beta=\norm{\mathbf{r}_0}_2$, $\mathbf{e}_1=[1,0,\ldots,0]^T\in\R^i$, 则 $\mathbf{r}_0=\beta V_i\mathbf{e}_1$. 由于 $V_i$ 的列标准正交, 由 \eqref{eq:6-22} 与 \eqref{eq:6-23} 得
\[
  \norm{\mathbf{r}_i}_2=\norm{V_i(\beta\mathbf{e}_1-H\mathbf{y})}_2
  =\norm{\beta\mathbf{e}_1-H\mathbf{y}}_2.
\]
取 $\mathbf{y}=\beta H^{-1}\mathbf{e}_1$, 即得 $\mathbf{r}_i=\mathbf{0}$, 亦即
\[
  \mathbf{x}_i=A^{-1}\mathbf{b}\in\mathbf{x}_0+\mathcal{K}(A,\mathbf{r}_0,i).
  \qedhere
\]
\end{proof}

若 Gram-Schmidt 过程不中断, 就可用如下高效方式实现 GMRES 法. 令 $h_{ij}=(A\mathbf{v}_j)^T\mathbf{v}_i$. 由 Gram-Schmidt 算法得到 $(k+1)\times k$ 矩阵 $H_k$, 它是上 Hessenberg 矩阵, 即当 $i>j+1$ 时 $h_{ij}=0$. 该过程产生一列列标准正交的矩阵序列 $\{V_k\}$, 使得
\[
  AV_k=V_{k+1}H_k.
\]
因此
\[
  \begin{aligned}
    \mathbf{r}_k&=\mathbf{b}-A\mathbf{x}_k
    =\mathbf{r}_0-A(\mathbf{x}_k-\mathbf{x}_0)\\
    &=\beta V_{k+1}\mathbf{e}_1-AV_k\mathbf{y}_k
    =V_{k+1}(\beta\mathbf{e}_1-H_k\mathbf{y}_k),
  \end{aligned}
\]
其中 $\mathbf{y}_k$ 为如下问题的解:
\[
  \min_{\mathbf{y}\in\R^k}\norm{\beta\mathbf{e}_1-H_k\mathbf{y}}_2.
\]
于是
\[
  \mathbf{x}_k=\mathbf{x}_0+V_k\mathbf{y}_k.
\]

当找到向量 $\mathbf{x}$, 使对给定的 $\epsilon$ 有
\[
  \norm{\mathbf{r}}_2=\norm{\mathbf{b}-A\mathbf{x}}_2
  \leqslant\epsilon\norm{\mathbf{b}}_2
\]
时, 就终止 GMRES 迭代. 于是得到如下求解
\[
  A\mathbf{x}=\mathbf{b}
\]
的 GMRES 算法 \cite{src:70}, 其中 $A\in\R^{n\times n}$ 与 $\mathbf{b}\in\R^n$ 已知. 在初始化步, 令
\[
  \mathbf{r}_0=\mathbf{b}-A\mathbf{x}_0,\qquad
  \beta=\norm{\mathbf{r}_0}_2,\qquad
  \mathbf{v}_1=\mathbf{r}_0/\beta.
\]
在迭代步, 有
\[
  \left\{
  \begin{array}{l}
    \textbf{for}\ j=1{:}m\\
    \quad \mathbf{w}_j=A\mathbf{v}_j\\
    \quad \textbf{for}\ i=1{:}j\\
    \quad\quad h_{ij}=\mathbf{w}_j^T\mathbf{v}_i\\
    \quad\quad \mathbf{w}_j=\mathbf{w}_j-h_{ij}\mathbf{v}_i\\
    \quad \textbf{end}\\
    \quad h_{j+1,j}=\norm{\mathbf{w}_j}_2;\ \text{若}\ h_{j+1,j}=0,\
      \text{则置}\ m=j\ \text{并转至}\ (*)\\
    \quad \mathbf{v}_{j+1}=\mathbf{w}_j/h_{j+1,j}\\
    \textbf{end}\\
    (*)\ \text{计算}\ \mathbf{y}_m,\ \text{使其极小化}\
      \norm{\beta\mathbf{e}_1-\bar{H}_m\mathbf{y}}_2\\
    \mathbf{x}_m=\mathbf{x}_0+V_m\mathbf{y}_m
  \end{array}
  \right.
\]
其中 $\bar{H}_m=(h_{ij})_{1\leqslant i\leqslant m+1,\,1\leqslant j\leqslant m}$. 算法中的矩阵
\[
  V_m=[\mathbf{v}_1\mid\mathbf{v}_2\mid\cdots\mid\mathbf{v}_m]\in\R^{n\times m}
  \quad(m\leqslant n)
\]
是列标准正交的矩阵.

\needspace{6\baselineskip}
\begin{exercises}
  \item 设 $\mathbf{x}_k$ 由最速下降法产生. 证明
    \[
      \phi(\mathbf{x}_k)\leqslant\left(1-\frac{1}{\kappa_2(A)}\right)
      \phi(\mathbf{x}_{k-1}),
    \]
    其中 $\kappa_2(A)=\norm{A}_2\norm{A^{-1}}_2$.

  \item 设 $A$ 为对称正定矩阵, 定义 $\norm{\mathbf{x}}_A\equiv\sqrt{\mathbf{x}^TA\mathbf{x}}$. 证明它是一个向量范数.

  \item 设 $A\in\R^{n\times n}$ 对称正定, $\mathbf{p}_1,\mathbf{p}_2,\ldots,\mathbf{p}_k\in\R^n$ 两两关于 $A$ 共轭, 即当 $i\neq j$ 时 $\mathbf{p}_i^TA\mathbf{p}_j=0$. 证明 $\{\mathbf{p}_1,\mathbf{p}_2,\ldots,\mathbf{p}_k\}$ 线性无关.

  \item 设 $\mathbf{x}_k$ 由 CG 法产生. 证明
    \[
      \norm{\mathbf{x}_k-\mathbf{x}_*}_2
      \leqslant2\sqrt{\kappa_2}
      \left(\frac{\sqrt{\kappa_2}-1}{\sqrt{\kappa_2}+1}\right)^k
      \norm{\mathbf{x}_0-\mathbf{x}_*}_2,
    \]
    其中 $\kappa_2=\kappa_2(A)=\norm{A}_2\norm{A^{-1}}_2$.

  \item 设 $\mathbf{z}_k$ 与 $\mathbf{r}_k$ 由 PCG 法产生. 证明: 若 $\mathbf{r}_k\neq\mathbf{0}$, 则 $\mathbf{z}_k^T\mathbf{r}_k>0$.

  \item 给出用 CG 法求解 $A^TA\mathbf{x}=A^T\mathbf{b}$ 的一个高效算法.

  \item 证明: 若 $A\in\R^{n\times n}$ 对称正定且恰有 $k$ 个互不相同的特征值, 则 CG 法至多 $k$ 次迭代即终止.

  \item 取初始向量 $\mathbf{x}_0=\mathbf{0}$. 当用 GMRES 法求解线性方程组 $A\mathbf{x}=\mathbf{b}$ 时, 其中
    \[
      A=\begin{bmatrix}
        0 & & & & 1\\
        1 & 0 & & & \\
        & 1 & 0 & & \\
        & & 1 & 0 & \\
        & & & 1 & 0
      \end{bmatrix},
      \qquad
      \mathbf{b}=[1,0,0,0,0]^T,
    \]
    收敛速度是多少?

  \item 设
    \[
      A=\begin{bmatrix}
        I & Y\\
        0 & I
      \end{bmatrix}.
    \]
    当用 GMRES 法求解 $A\mathbf{x}=\mathbf{b}$ 时, 收敛所需的最大迭代次数是多少?

  \item 证明 GMRES 法中的最小二乘问题具有列满秩.

  \item 证明引理 \ref{lem:6-2}.
\end{exercises}

% chapters/ch07.tex —— 第 7 章：非对称特征值问题
% 源书：Chapter 7 Nonsymmetric Eigenvalue Problems，书页 103–120（物理页 117–134）
% 公式编号 (7.1)–(7.31)；定理 7.1–7.5；算法 7.1；习题 1–13
\chapter{非对称特征值问题}
\label{ch:nonsym-eig}

特征值问题在数值线性代数中格外引人注目. 本章研究非对称特征值问题, 讨论幂法, 反幂法与 QR 方法等几种著名方法.

\section{基本性质}
\label{sec:7-1}

设 $A\in\C^{n\times n}$. 回忆: 若存在非零向量 $\mathbf{x}\in\C^n$ 使
\[
  A\mathbf{x}=\lambda\mathbf{x},
\]
则称复数 $\lambda$ 为 $A$ 的特征值. 这里 $\mathbf{x}$ 称为 $A$ 的与 $\lambda$ 相伴的特征向量. 众所周知, $\lambda$ 是 $A$ 的特征值当且仅当
\[
  \det(\lambda I-A)=0.
\]
令
\[
  p_A(\lambda)=\det(\lambda I-A)
\]
为 $A$ 的特征多项式. 由代数基本定理知, $p_A(\lambda)$ 在 $\C$ 中有 $n$ 个根, 即 $A$ 有 $n$ 个特征值.

现设 $p_A(\lambda)$ 有如下分解:
\[
  p_A(\lambda)=(\lambda-\lambda_1)^{n_1}(\lambda-\lambda_2)^{n_2}\cdots(\lambda-\lambda_p)^{n_p},
\]
其中 $n_1+n_2+\cdots+n_p=n$, 且当 $i\neq j$ 时 $\lambda_i\neq\lambda_j$. 称 $n_i$ 为 $\lambda_i$ 的\newterm{algebraicmultiplicity}{代数重数}{algebraic multiplicity}, 而数
\[
  m_i=n-\rank(\lambda_i I-A)
\]
称为 $\lambda_i$ 的\newterm{geometricmultiplicity}{几何重数}{geometric multiplicity}. 实际上, $m_i$ 就是 $\lambda_i$ 的\newterm{eigenspace}{特征子空间}{eigenspace}的维数. 我们指出, $\lambda_i$ 的特征子空间是 $(\lambda_i I-A)\mathbf{x}=\mathbf{0}$ 的解空间. 容易看出 $m_i\leqslant n_i$ ($i=1,2,\ldots,p$). 若 $n_i=1$, 则称 $\lambda_i$ 为\newterm{simpleeigenvalue}{单特征值}{simple eigenvalue}. 若对某个特征值 $\lambda_i$ 有 $m_i<n_i$, 则称 $A$ 为\newterm{defective}{亏损的}{defective}. 若 $A$ 的每个特征值的几何重数都等于其代数重数, 则称 $A$ 是\newterm{nondefective}{非亏损的}{nondefective}的. 注意 $A$ 可对角化当且仅当 $A$ 非亏损.

设 $A,B\in\C^{n\times n}$. 回忆: 若存在非奇异矩阵 $X\in\C^{n\times n}$ 使
\[
  B=XAX^{-1},
\]
则称 $A$ 与 $B$ 为\newterm{similarmatrix}{相似矩阵}{similar matrices}. 变换
\[
  A\rightarrow B=XAX^{-1}
\]
称为用相似矩阵 $X$ 作的\newterm{similaritytransformation}{相似变换}{similarity transformation}. 相似矩阵有相同的特征值. 若 $\mathbf{x}$ 是 $A$ 的特征向量, 则 $\mathbf{y}=X\mathbf{x}$ 是 $B$ 的特征向量. 由 Jordan 分解定理 (定理 \ref{thm:ch1-jordan}) 知, 任一 $n$ 阶矩阵都相似于其 Jordan 标准形. 若要求相似矩阵为酉矩阵, 便得到数值线性代数中或许最基本, 最有用的下述定理 \cite{src:39}.

\begin{theorem}[Schur 分解定理]
\label{thm:7-1}
设 $A\in\C^{n\times n}$, 其特征值为 $\lambda_1,\ldots,\lambda_n$, 可按任意指定的次序排列. 则存在酉矩阵 $U\in\C^{n\times n}$ 使
\[
  U^*AU=T=[t_{ij}],
\]
其中 $T$ 为上三角矩阵, 对角元 $t_{ii}=\lambda_i$ ($i=1,\ldots,n$). 此外, 若 $A\in\R^{n\times n}$ 且 $A$ 的所有特征值都是实数, 则 $U$ 可取为实正交矩阵.
\end{theorem}

\begin{proof}
设 $\mathbf{x}_1\in\C^n$ 是 $A$ 的与特征值 $\lambda_1$ 相伴的标准化特征向量. 向量 $\mathbf{x}_1$ 可扩充为 $\C^n$ 的一组基:
\[
  \mathbf{x}_1,\mathbf{y}_2,\ldots,\mathbf{y}_n.
\]
对此基施行 Gram-Schmidt 正交规范化, 便得到 $\C^n$ 的一组标准正交基:
\[
  \mathbf{x}_1,\mathbf{z}_2,\ldots,\mathbf{z}_n.
\]
令
\[
  U_1=[\mathbf{x}_1\mid\mathbf{z}_2\mid\cdots\mid\mathbf{z}_n]
\]
为酉矩阵. 经简单计算得
\[
  U_1^*AU_1=\begin{bmatrix}
    \lambda_1 & *\\
    \mathbf{0} & A_1
  \end{bmatrix}.
\]

矩阵 $A_1\in\C^{(n-1)\times(n-1)}$ 的特征值为 $\lambda_2,\ldots,\lambda_n$. 设 $\mathbf{x}_2\in\C^{n-1}$ 是 $A_1$ 的与 $\lambda_2$ 相伴的标准化特征向量, 照此重做一遍. 确定酉矩阵 $V_2\in\C^{(n-1)\times(n-1)}$ 使
\[
  V_2^*A_1V_2=\begin{bmatrix}
    \lambda_2 & *\\
    \mathbf{0} & A_2
  \end{bmatrix}.
\]
令
\[
  U_2=\begin{bmatrix}
    1 & \mathbf{0}\\
    \mathbf{0} & V_2
  \end{bmatrix}.
\]
则矩阵 $U_2$ 与 $U_1U_2$ 都是酉矩阵, 且
\[
  U_2^*U_1^*AU_1U_2=\begin{bmatrix}
    \lambda_1 &          & *\\
              & \lambda_2 & \\
    \mathbf{0} &          & A_2
  \end{bmatrix}.
\]
如此继续, 可作出酉矩阵 $U_1,U_2,\ldots,U_{n-1}$, 使矩阵
\[
  U=U_1U_2\cdots U_{n-1}
\]
为酉矩阵, 且 $U^*AU$ 具有所需的形式.
\end{proof}

\begin{theorem}[实 Schur 分解定理 \cite{src:39}]
\label{thm:7-2}
设 $A\in\R^{n\times n}$. 则存在正交矩阵 $Q\in\R^{n\times n}$ 使
\[
  Q^TAQ=\begin{bmatrix}
    R_{11} & R_{12} & \cdots & R_{1m}\\
           & R_{22} & \cdots & R_{2m}\\
           &        & \ddots & \vdots\\
    \mathbf{0} &   &        & R_{mm}
  \end{bmatrix},
\]
其中 $R_{ii}$ 或为实数, 或为具有一对共轭复特征值的 $2\times2$ 矩阵.
\end{theorem}

一般地, 不能指望用正交相似变换把实矩阵约化为严格上三角形式, 因为此时对角元就是特征值, 而它们 (若为复数) 不可能是实数.

\section{幂法}
\label{sec:7-2}

幂法是计算绝对值最大的特征值及其对应特征向量的迭代算法. 下面介绍该方法的基本思想. 为简单起见, 先设矩阵 $A\in\C^{n\times n}$ 可对角化, 即 $A$ 有如下 Jordan 分解:
\[
  A=X\Lambda X^{-1},
\]
其中 $\Lambda=\diag(\lambda_1,\ldots,\lambda_n)$, 且其次序为
\[
  \abs{\lambda_1}>\abs{\lambda_2}\geqslant\cdots\geqslant\abs{\lambda_n},
\]
而 $X=[\mathbf{x}_1\mid\mathbf{x}_2\mid\cdots\mid\mathbf{x}_n]\in\C^{n\times n}$. 设 $\mathbf{u}_0\in\C^n$ 为任意向量. 由于列向量 $\mathbf{x}_1,\ldots,\mathbf{x}_n$ 构成 $\C^n$ 的一组基, 故有
\[
  \mathbf{u}_0=\sum_{j=1}^{n}\alpha_j\mathbf{x}_j,
\]
其中 $\alpha_j\in\C$. 于是
\[
  \begin{aligned}
    A^k\mathbf{u}_0
    &=\sum_{j=1}^{n}\alpha_jA^k\mathbf{x}_j
     =\sum_{j=1}^{n}\alpha_j\lambda_j^k\mathbf{x}_j\\
    &=\lambda_1^k\left(\alpha_1\mathbf{x}_1
      +\sum_{j=2}^{n}\alpha_j\left(\frac{\lambda_j}{\lambda_1}\right)^k\mathbf{x}_j\right),
  \end{aligned}
\]
从而
\[
  \lim_{k\rightarrow\infty}\frac{A^k\mathbf{u}_0}{\lambda_1^k}=\alpha_1\mathbf{x}_1.
\]
当 $\alpha_1\neq0$ 且 $k$ 充分大时, 可知向量
\begin{equation}
  \mathbf{u}_k=\frac{A^k\mathbf{u}_0}{\lambda_1^k}
  \label{eq:7-1}
\end{equation}
是 $A$ 的一个好的近似特征向量.

实际计算中不能直接用 \eqref{eq:7-1} 求近似特征向量, 因为事先并不知道特征值 $\lambda_1$, 且 $k$ 很大时 $A^k$ 的运算量极大. 因此提出如下迭代算法:
\begin{equation}
  \left\{
  \begin{array}{l}
    \mathbf{y}_k=A\mathbf{u}_{k-1},\\
    \mu_k=\zeta_k,\\
    \mathbf{u}_k=\mathbf{y}_k/\mu_k,
  \end{array}
  \right.
  \label{eq:7-2}
\end{equation}
其中 $\mathbf{u}_0\in\C^n$ 为任意给定的初始向量, 通常取 $\norm{\mathbf{u}_0}_\infty=1$, 而 $\zeta_k$ 是 $\mathbf{y}_k$ 各分量绝对值的最大者. 该迭代算法称为\newterm{powermethod}{幂法}{power method}. 关于幂法的收敛性, 有如下定理.

\begin{theorem}
\label{thm:7-3}
设 $A\in\C^{n\times n}$ 有 $p$ 个互不相同的特征值, 满足
\[
  \abs{\lambda_1}>\abs{\lambda_2}\geqslant\cdots\geqslant\abs{\lambda_p}.
\]
若 $\lambda_1$ 的几何重数等于其代数重数, 且初始向量 $\mathbf{u}_0$ 在 $\lambda_1$ 的特征子空间上的投影非零, 则由 \eqref{eq:7-2} 产生的 $\{\mathbf{u}_k\}$ 收敛到 $A$ 的与 $\lambda_1$ 相伴的一个特征向量 $\mathbf{x}_1$; 同时由 \eqref{eq:7-2} 产生的 $\{\mu_k\}$ 收敛到 $\lambda_1$.
\end{theorem}

\begin{proof}
注意 $A$ 有如下 Jordan 分解:
\begin{equation}
  A=X\diag(J_1,\ldots,J_p)X^{-1},
  \label{eq:7-3}
\end{equation}
其中 $X\in\C^{n\times n}$, $J_i\in\C^{n_i\times n_i}$ 是 $\lambda_i$ 对应的 Jordan 块 ($i=1,\ldots,p$), 且
\[
  n_1+n_2+\cdots+n_p=n.
\]
由于 $\lambda_1$ 的几何重数与其代数重数相同, 故
\[
  J_1=\lambda_1I_{n_1},
\]
其中 $I_{n_1}\in\R^{n_1\times n_1}$ 是单位矩阵. 令 $\mathbf{y}=X^{-1}\mathbf{u}_0$, 并把 $\mathbf{y}$ 与 $X$ 作如下分解:
\[
  \mathbf{y}=[\mathbf{y}_1^T,\mathbf{y}_2^T,\ldots,\mathbf{y}_p^T]^T,
  \qquad
  X=[X_1\mid X_2\mid\cdots\mid X_p],
\]
其中 $\mathbf{y}_i\in\C^{n_i}$, $X_i\in\C^{n\times n_i}$ ($i=1,\ldots,p$). 利用 \eqref{eq:7-3} 得
\[
  \begin{aligned}
    A^k\mathbf{u}_0
    &=X\diag(J_1^k,\ldots,J_p^k)X^{-1}\mathbf{u}_0\\
    &=X_1J_1^k\mathbf{y}_1+X_2J_2^k\mathbf{y}_2+\cdots+X_pJ_p^k\mathbf{y}_p\\
    &=\lambda_1^kX_1\mathbf{y}_1+X_2J_2^k\mathbf{y}_2+\cdots+X_pJ_p^k\mathbf{y}_p\\
    &=\lambda_1^k\left(X_1\mathbf{y}_1
      +X_2(\lambda_1^{-1}J_2)^k\mathbf{y}_2
      +\cdots+X_p(\lambda_1^{-1}J_p)^k\mathbf{y}_p\right).
  \end{aligned}
\]
注意 $\lambda_1^{-1}J_i$ 的谱半径满足
\[
  \rho(\lambda_1^{-1}J_i)=\abs{\lambda_i}/\abs{\lambda_1}<1,
\]
其中 $i=2,3,\ldots,p$. 因此
\begin{equation}
  \lim_{k\rightarrow\infty}\frac{1}{\lambda_1^k}A^k\mathbf{u}_0=X_1\mathbf{y}_1.
  \label{eq:7-4}
\end{equation}
由于初始向量 $\mathbf{u}_0$ 在 $\lambda_1$ 的特征子空间上的投影非零, 故 $X_1\mathbf{y}_1\neq\mathbf{0}$. 令
\[
  \mathbf{x}_1=\zeta^{-1}X_1\mathbf{y}_1,
\]
其中 $\zeta$ 是 $X_1\mathbf{y}_1$ 各分量绝对值的最大者. 显然, $\mathbf{x}_1$ 是 $A$ 的与 $\lambda_1$ 相伴的特征向量. 设 $\zeta_k$ 是 $A^k\mathbf{u}_0$ 各分量绝对值的最大者, 则 $\lambda_1^{-k}A^k\mathbf{u}_0$ 各分量绝对值的最大者为 $\zeta_k\lambda_1^{-k}$. 由 \eqref{eq:7-2} 有
\[
  \mathbf{u}_k=\frac{A\mathbf{u}_{k-1}}{\mu_k}
  =\frac{A^k\mathbf{u}_0}{\mu_k\mu_{k-1}\cdots\mu_1}
  =\frac{A^k\mathbf{u}_0}{\zeta_k}
  =\frac{A^k\mathbf{u}_0/\lambda_1^k}{\zeta_k/\lambda_1^k}.
\]
利用 \eqref{eq:7-4} 知 $\{\mathbf{u}_k\}$ 收敛, 且
\[
  \lim_{k\rightarrow\infty}\mathbf{u}_k
  =\frac{\lim_{k\rightarrow\infty}(A^k\mathbf{u}_0/\lambda_1^k)}
        {\lim_{k\rightarrow\infty}(\zeta_k/\lambda_1^k)}
  =\frac{X_1\mathbf{y}_1}{\zeta}=\mathbf{x}_1.
\]
利用 $A\mathbf{u}_{k-1}=\mu_k\mathbf{u}_k$ 以及 $\{\mathbf{u}_k\}$ 收敛到 $\lambda_1$ 的, 分量绝对值最大者为 $1$ 的特征向量这一事实, 立即可知 $\{\mu_k\}$ 收敛到 $\lambda_1$.
\end{proof}

我们指出, 由定理 \ref{thm:7-3} 的证明可知, 幂法的收敛速度由 $\abs{\lambda_2}/\abs{\lambda_1}$ 的值决定. 在定理条件下有
\[
  \frac{\abs{\lambda_2}}{\abs{\lambda_1}}<1.
\]
$\abs{\lambda_2}/\abs{\lambda_1}$ 越小, 收敛越快; 当 $\abs{\lambda_2}/\abs{\lambda_1}$ 接近 $1$ 时, 收敛很慢. 为加快幂法的收敛, 可对 $A-\mu I$ 使用该方法, 其中 $\mu$ 称为位移. 可选取 $\mu$ 使绝对值最大的特征值与其他特征值之间的距离变大, 从而提高幂法的收敛速度.

\section{反幂法}
\label{sec:7-3}

若对 $A^{-1}$ 使用幂法, 以求 $A$ 的绝对值最小的特征值及其对应特征向量, 则该算法称为\newterm{inversepowermethod}{反幂法}{inverse power method}. 其基本迭代格式如下:
\[
  \left\{
  \begin{array}{l}
    A\mathbf{y}_k=\mathbf{z}_{k-1},\\
    \mu_k=\zeta_k,\\
    \mathbf{z}_k=\mathbf{y}_k/\mu_k,
  \end{array}
  \right.
\]
其中 $\mathbf{z}_0\in\C^n$ 为任意给定的初始向量, $\zeta_k$ 是 $\mathbf{y}_k$ 各分量绝对值的最大者. 由定理 \ref{thm:7-3} 知, 若 $A$ 的特征值满足
\[
  \abs{\lambda_n}<\abs{\lambda_{n-1}}\leqslant\cdots\leqslant\abs{\lambda_1},
\]
则 $\{\mathbf{z}_k\}$ 收敛到与 $\lambda_n$ 相伴的特征向量, 而 $\{\mu_k\}$ 收敛到 $\lambda_n^{-1}$. 反幂法的收敛速度由 $\abs{\lambda_n}/\abs{\lambda_{n-1}}$ 决定.

实际计算中, 通常对矩阵
\[
  A-\mu I
\]
使用反幂法: 当预先得到 $A$ 的某个互异特征值 $\lambda_i$ 的近似值 $\mu$ 后, 用它来计算近似特征向量. 于是有如下带\newterm{shift}{位移}{shift} $\mu$ 的反幂法:
\begin{equation}
  \left\{
  \begin{array}{l}
    (A-\mu I)\mathbf{v}_k=\mathbf{z}_{k-1},\\
    \mathbf{z}_k=\mathbf{v}_k/\norm{\mathbf{v}_k}_2.
  \end{array}
  \right.
  \label{eq:7-5}
\end{equation}

由 \eqref{eq:7-5} 知, 反幂法每迭代一次都要解一个线性方程组, 故其运算量远大于幂法. 不过, 可预先作部分主元 $LU$ 分解, 此后每次迭代只需解两个三角形方程组.

设 $A-\mu I$ 的特征值按如下次序排列:
\[
  0<\abs{\lambda_1-\mu}<\abs{\lambda_2-\mu}\leqslant\abs{\lambda_3-\mu}\leqslant\cdots\leqslant\abs{\lambda_n-\mu}.
\]
再次由定理 \ref{thm:7-3} 知, 由 \eqref{eq:7-5} 产生的序列 $\{\mathbf{z}_k\}$ 收敛到与 $\lambda_1$ 相伴的特征向量, 收敛速度由 $\abs{\lambda_1-\mu}/\abs{\lambda_2-\mu}$ 的值决定. $\mu$ 越接近 $\lambda_1$, 收敛越快. 但若 $\mu$ 接近 $A$ 的某个特征值, 则 $A-\mu I$ 接近奇异矩阵, 因此反幂法每迭代一次都要解一个病态的线性方程组. 然而由实际计算可知, 方程组的病态性并不影响该方法的收敛速度. 通常, 若 $\mu$ 接近 $A$ 的某个特征值, 只迭代一次就能得到 $A$ 的一个好的近似特征向量.

\section{QR 方法}
\label{sec:7-4}

本节介绍著名的 QR 方法, 它是矩阵计算中几个最重要的进展之一. 对任意给定的 $A_0=A\in\C^{n\times n}$, QR 算法的基本迭代格式如下:
\begin{equation}
  \left\{
  \begin{array}{l}
    A_{m-1}=Q_mR_m,\\
    A_m=R_mQ_m,
  \end{array}
  \right.
  \label{eq:7-6}
\end{equation}
其中 $m=1,2,\ldots$, $Q_m$ 为酉矩阵, $R_m$ 为上三角矩阵. 为便于后文分析, 要求 $R_m$ 的对角元非负. 由 \eqref{eq:7-6} 容易得到
\begin{equation}
  A_m=Q_m^*A_{m-1}Q_m,
  \label{eq:7-7}
\end{equation}
即序列 $\{A_m\}$ 中每个矩阵都与矩阵 $A$ 相似. 反复使用 \eqref{eq:7-7} 得
\begin{equation}
  A_m=\widetilde{Q}_m^*A\widetilde{Q}_m,
  \label{eq:7-8}
\end{equation}
其中 $\widetilde{Q}_m=Q_1Q_2\cdots Q_m$. 把 $A_m=Q_{m+1}R_{m+1}$ 代入 \eqref{eq:7-8} 得
\[
  \widetilde{Q}_mQ_{m+1}R_{m+1}=A\widetilde{Q}_m.
\]
因此
\[
  \widetilde{Q}_mQ_{m+1}R_{m+1}R_m\cdots R_1=A\widetilde{Q}_mR_m\cdots R_1,
\]
即
\[
  \widetilde{Q}_{m+1}\widetilde{R}_{m+1}=A\widetilde{Q}_m\widetilde{R}_m,
\]
其中 $\widetilde{R}_k=R_kR_{k-1}\cdots R_1$ ($k=m,m+1$). 此外还有
\begin{equation}
  A^m=\widetilde{Q}_m\widetilde{R}_m.
  \label{eq:7-9}
\end{equation}

\begin{theorem}
\label{thm:7-4}
设 $A$ 的特征值满足
\[
  \abs{\lambda_1}>\abs{\lambda_2}>\cdots>\abs{\lambda_n}>0.
\]
设 $Y$ 为矩阵, $\mathbf{y}_i^T$ 表示 $Y$ 的第 $i$ 行, 满足
\[
  \mathbf{y}_i^TA=\lambda_i\mathbf{y}_i^T.
\]
若 $Y$ 有 $LU$ 分解, 则由 \eqref{eq:7-6} 产生的矩阵 $A_m=[\alpha_{ij}^{(m)}]$ 中位于对角线以下的元素当 $m\rightarrow\infty$ 时趋于零. 同时
\[
  \alpha_{ii}^{(m)}\rightarrow\lambda_i,
\]
其中 $i=1,2,\ldots,n$.
\end{theorem}

\begin{proof}
令
\[
  X=Y^{-1},\qquad \Lambda=\diag(\lambda_1,\ldots,\lambda_n).
\]
则 $A=X\Lambda Y$. 由假设, $Y$ 有 $LU$ 分解
\[
  Y=LU,
\]
其中 $L$ 为单位下三角矩阵, $U$ 为上三角矩阵. 因此
\begin{equation}
  \begin{aligned}
    A^m=X\Lambda^mY=X\Lambda^mLU
    &=X(\Lambda^mL\Lambda^{-m})\Lambda^mU\\
    &=X(I+E_m)\Lambda^mU,
  \end{aligned}
  \label{eq:7-10}
\end{equation}
其中 $I+E_m=\Lambda^mL\Lambda^{-m}$. 由于 $L$ 为单位下三角矩阵, 且当 $i>j$ 时 $\abs{\lambda_i}<\abs{\lambda_j}$, 故有
\begin{equation}
  \lim_{m\rightarrow\infty}E_m=\mathbf{0}.
  \label{eq:7-11}
\end{equation}
令 $X=QR$, 其中 $Q$ 为酉矩阵, $R$ 为上三角矩阵. 由于 $X$ 非奇异, 可要求 $R$ 的对角元为正. 把 $X=QR$ 代入 \eqref{eq:7-10} 得
\begin{equation}
  A^m=QR(I+E_m)\Lambda^mU
  =Q(I+RE_mR^{-1})R\Lambda^mU.
  \label{eq:7-12}
\end{equation}

当 $m$ 充分大时, $I+RE_mR^{-1}$ 非奇异, 且有如下 QR 分解:
\begin{equation}
  I+RE_mR^{-1}=\widehat{Q}_m\widehat{R}_m,
  \label{eq:7-13}
\end{equation}
其中 $\widehat{R}_m$ 的对角元为正. 利用 \eqref{eq:7-11} 与 \eqref{eq:7-13} 容易证明
\begin{equation}
  \lim_{m\rightarrow\infty}\widehat{Q}_m
  =\lim_{m\rightarrow\infty}\widehat{R}_m=I.
  \label{eq:7-14}
\end{equation}
把 \eqref{eq:7-13} 代入 \eqref{eq:7-12} 得
\[
  A^m=(Q\widehat{Q}_m)(\widehat{R}_mR\Lambda^mU),
\]
这是 $A^m$ 的一个 QR 分解. 为保证上三角矩阵的所有对角元为正, 定义
\[
  D_1=\diag\left(\frac{\lambda_1}{\abs{\lambda_1}},\ldots,\frac{\lambda_n}{\abs{\lambda_n}}\right),
\]
以及
\[
  D_2=\diag\left(\frac{u_{11}}{\abs{u_{11}}},\ldots,\frac{u_{nn}}{\abs{u_{nn}}}\right),
\]
其中 $u_{ii}$ 是 $U$ 的第 $i$ 个对角元. 于是
\[
  A^m=(Q\widehat{Q}_mD_1^mD_2)
  (D_2^{-1}D_1^{-m}\widehat{R}_mR\Lambda^mU).
\]
与 \eqref{eq:7-9} 比较, 并注意 QR 分解的唯一性, 得
\[
  \widetilde{Q}_m=Q\widehat{Q}_mD_1^mD_2,
  \qquad
  \widetilde{R}_m=D_2^{-1}D_1^{-m}\widehat{R}_mR\Lambda^mU.
\]
把它们代入 \eqref{eq:7-8} 得
\[
  A_m=D_2^*(D_1^*)^m\widehat{Q}_m^*Q^*AQ\widehat{Q}_mD_1^mD_2.
\]
注意
\[
  A=X\Lambda Y=X\Lambda X^{-1}=QR\Lambda R^{-1}Q^*.
\]
最终得到
\[
  A_m=D_2^*(D_1^*)^m\widehat{Q}_m^*R\Lambda R^{-1}\widehat{Q}_mD_1^mD_2.
\]
当 $m\rightarrow\infty$ 时, 由 \eqref{eq:7-14} 知, 由 \eqref{eq:7-6} 产生的矩阵 $A_m$ 中位于对角线以下的元素趋于零. 同时
\[
  \alpha_{ii}^{(m)}\rightarrow\lambda_i,
\]
其中 $i=1,2,\ldots,n$.
\end{proof}

由定理 \ref{thm:7-4} 知, 由 \eqref{eq:7-6} 产生的序列 $\{A_m\}$ 收敛到 $A$ 的 Schur 分解.

\section{QR 算法的实形式}
\label{sec:7-5}

若 $A\in\R^{n\times n}$, 我们希望找到一种只用实数运算的快速而有效的 QR 算法. 我们集中发展 \eqref{eq:7-6} 的实形式如下: 令 $A_1=A$, 然后构造迭代算法
\begin{equation}
  \left\{
  \begin{array}{l}
    A_m=Q_mR_m,\\
    A_{m+1}=R_mQ_m,
  \end{array}
  \right.
  \label{eq:7-15}
\end{equation}
其中 $m=1,2,\ldots$, $Q_m\in\R^{n\times n}$ 为正交矩阵, $R_m\in\R^{n\times n}$ 为上三角矩阵. 与 \eqref{eq:7-15} 相关的一个困难在于: 当 $A$ 有复特征值时, $A_m$ 绝不可能收敛到严格上三角形式. 由定理 \ref{thm:7-2} 可以期望 \eqref{eq:7-15} 收敛到 $A$ 的实 Schur 分解.

我们指出, 实际上 \eqref{eq:7-15} 并不是一个好的迭代法, 原因有二:
\begin{romanlist}
  \item 每次迭代的运算量太大;
  \item 收敛速度太慢.
\end{romanlist}
因此需要减少每次迭代的运算量并提高方法的收敛速度. 于是引入上 Hessenberg 化与位移策略.

\subsection{上 Hessenberg 化}
\label{sec:7-5-1}

构造正交矩阵 $Q_0$ 使
\begin{equation}
  Q_0^TAQ_0=H
  \label{eq:7-16}
\end{equation}
具有某种特殊结构 (Hessenberg 型), 含有大量零元. 此后即可对 \eqref{eq:7-16} 中的矩阵 $H$ 应用 QR 算法 \eqref{eq:7-15}, 每次迭代的运算量便可大幅减少.

对 $A=[\alpha_{ij}]\in\R^{n\times n}$, 第一步可选取 Householder 变换 $H_1$ 使 $H_1A$ 的第一列含有大量零元 (至多 $n-1$ 个零元). 为保持相似变换, 还需再作一次列变换:
\[
  H_1AH_1.
\]
因此 $H_1$ 可取如下形式
\begin{equation}
  H_1=\begin{bmatrix}
    1 & \mathbf{0}\\
    \mathbf{0} & \widetilde{H}_1
  \end{bmatrix}
  \label{eq:7-17}
\end{equation}
以使第一列中的零元保持不变. 利用 \eqref{eq:7-17} 定义的 $H_1$, 有
\begin{equation}
  H_1AH_1=\begin{bmatrix}
    \alpha_{11} & \mathbf{a}_2^T\widetilde{H}_1\\
    \widetilde{H}_1\mathbf{a}_1 & \widetilde{H}_1A_{22}\widetilde{H}_1
  \end{bmatrix},
  \label{eq:7-18}
\end{equation}
其中 $\mathbf{a}_1^T=[\alpha_{21},\alpha_{31},\ldots,\alpha_{n1}]$, $\mathbf{a}_2^T=[\alpha_{12},\alpha_{13},\ldots,\alpha_{1n}]$, 而 $A_{22}$ 是 $A$ 右下角的 $(n-1)$ 阶主子矩阵. 由 \eqref{eq:7-18} 知, Householder 变换 $\widetilde{H}_1$ 的最佳选择应是
\begin{equation}
  \widetilde{H}_1\mathbf{a}_1=p\mathbf{e}_1,
  \label{eq:7-19}
\end{equation}
其中 $p\in\R$, $\mathbf{e}_1\in\R^{n-1}$ 是第一个单位向量. 因此, 利用 \eqref{eq:7-17} 与 \eqref{eq:7-19} 定义的 $H_1$, \eqref{eq:7-18} 中矩阵的第一列有 $n-2$ 个零元.

此后, 对 $\widetilde{A}_{22}=\widetilde{H}_1A_{22}\widetilde{H}_1$, 可找到 Householder 变换
\[
  \widetilde{H}_2=\begin{bmatrix}
    1 & \mathbf{0}\\
    \mathbf{0} & \widehat{H}_2
  \end{bmatrix}
\]
使
\[
  (\widetilde{H}_2\widetilde{A}_{22}\widetilde{H}_2)\mathbf{e}_1
  =[*,*,0,\ldots,0]^T.
\]
令
\[
  H_2=\begin{bmatrix}
    1 & \mathbf{0}\\
    \mathbf{0} & \widetilde{H}_2
  \end{bmatrix}.
\]
于是有
\[
  H_2H_1AH_1H_2=\begin{bmatrix}
    h_{11} & h_{12} & \cdots & *\\
    h_{21} & h_{22} & \vdots & \vdots\\
    0 & h_{32} & \vdots & *\\
    0 & 0 & \vdots & *\\
    \vdots & \vdots & \vdots & \vdots\\
    0 & 0 & \cdots & *
  \end{bmatrix}.
\]
经 $n-2$ 步后, 求得 $n-2$ 个 Householder 变换 $H_1,H_2,\ldots,H_{n-2}$, 使
\[
  H_{n-2}\cdots H_2H_1AH_1H_2\cdots H_{n-2}=H,
\]
其中
\[
  H=\begin{bmatrix}
    h_{11} & h_{12} & h_{13} & \cdots & h_{1,n-1} & h_{1n}\\
    h_{21} & h_{22} & h_{23} & \cdots & h_{2,n-1} & h_{2n}\\
    0 & h_{32} & h_{33} & \cdots & h_{3,n-1} & h_{3n}\\
    \vdots & 0 & h_{43} & & \vdots & \vdots\\
    \vdots & \vdots & & \ddots & \vdots & \vdots\\
    0 & 0 & 0 & \cdots & h_{n,n-1} & h_{nn}
  \end{bmatrix}
\]
满足当 $i>j+1$ 时 $h_{ij}=0$, 称为\newterm{hessenberg}{上 Hessenberg 矩阵}{upper Hessenberg matrix}. 令
\[
  Q_0=H_1H_2\cdots H_{n-2},
\]
则
\[
  Q_0^TAQ_0=H,
\]
它称为 $A$ 的\newterm{hessenbergdecomp}{上 Hessenberg 分解}{upper Hessenberg decomposition}. 用 Householder 变换可得如下算法.

\begin{algorithm}{上 Hessenberg 分解}
  \label{alg:7-1}
  \[
    \left\{
    \begin{array}{ll}
      \textbf{for} & k=1:n-2\\
        & [\mathbf{v},\beta]=\mathrm{house}(A(k+1:n,k))\\
        & A(k+1:n,k:n)=(I-\beta\mathbf{v}\mathbf{v}^T)A(k+1:n,k:n)\\
        & A(1:n,k+1:n)=A(1:n,k+1:n)(I-\beta\mathbf{v}\mathbf{v}^T)\\
      \textbf{end}
    \end{array}
    \right.
  \]
\end{algorithm}

用 Householder 变换作 Hessenberg 分解的运算量为 $\mathcal{O}(n^3)$. 当然, 也可用 Givens 旋转得到 Hessenberg 分解, 但其运算量是用 Householder 变换时的两倍. 虽然矩阵的 Hessenberg 分解不唯一, 但有下述定理.

\begin{theorem}
\label{thm:7-5}
设 $A\in\R^{n\times n}$ 有如下两个上 Hessenberg 分解:
\begin{equation}
  U^TAU=H,\qquad V^TAV=G,
  \label{eq:7-20}
\end{equation}
其中
\[
  U=[\mathbf{u}_1\mid\mathbf{u}_2\mid\cdots\mid\mathbf{u}_n],
  \qquad
  V=[\mathbf{v}_1\mid\mathbf{v}_2\mid\cdots\mid\mathbf{v}_n]
\]
为 $n$ 阶正交矩阵, 而 $H=[h_{ij}]$, $G=[g_{ij}]$ 为上 Hessenberg 矩阵. 若 $\mathbf{u}_1=\mathbf{v}_1$ 且所有元素 $h_{i+1,i}$ 均不为零, 则存在对角元为 $1$ 或 $-1$ 的对角矩阵 $D$ 使
\[
  U=VD,\qquad H=DGD.
\]
\end{theorem}

\begin{proof}
设已对某个 $m$ ($1\leqslant m<n$) 证明了
\begin{equation}
  \mathbf{u}_j=\epsilon_j\mathbf{v}_j,\qquad j=1,2,\ldots,m,
  \label{eq:7-21}
\end{equation}
其中 $\epsilon_1=1$, $\epsilon_j=1$ 或 $-1$ ($j=2,3,\ldots,m$). 现在要证明存在 $\epsilon_{m+1}=1$ 或 $-1$ 使
\[
  \mathbf{u}_{m+1}=\epsilon_{m+1}\mathbf{v}_{m+1}.
\]

由 \eqref{eq:7-20} 有
\[
  AU=UH,\qquad AV=VG.
\]
分别比较上述矩阵等式的第 $m$ 列, 得
\begin{equation}
  A\mathbf{u}_m=h_{1m}\mathbf{u}_1+\cdots+h_{mm}\mathbf{u}_m+h_{m+1,m}\mathbf{u}_{m+1},
  \label{eq:7-22}
\end{equation}
以及
\begin{equation}
  A\mathbf{v}_m=g_{1m}\mathbf{v}_1+\cdots+g_{mm}\mathbf{v}_m+g_{m+1,m}\mathbf{v}_{m+1}.
  \label{eq:7-23}
\end{equation}
用 $\mathbf{u}_i^T$ 与 $\mathbf{v}_i^T$ 分别乘 \eqref{eq:7-22} 与 \eqref{eq:7-23}, 得
\[
  h_{im}=\mathbf{u}_i^TA\mathbf{u}_m,\qquad
  g_{im}=\mathbf{v}_i^TA\mathbf{v}_m,\qquad
  i=1,2,\ldots,m.
\]
因此由 \eqref{eq:7-21} 得
\begin{equation}
  h_{im}=\epsilon_i\epsilon_mg_{im},\qquad i=1,2,\ldots,m.
  \label{eq:7-24}
\end{equation}
把 \eqref{eq:7-24} 代入 \eqref{eq:7-22}, 并利用 \eqref{eq:7-21} 与 \eqref{eq:7-23}, 得
\begin{equation}
  \begin{aligned}
    h_{m+1,m}\mathbf{u}_{m+1}
    &=\epsilon_m(A\mathbf{v}_m-\epsilon_1^2g_{1m}\mathbf{v}_1-\cdots-\epsilon_m^2g_{mm}\mathbf{v}_m)\\
    &=\epsilon_m(A\mathbf{v}_m-g_{1m}\mathbf{v}_1-\cdots-g_{mm}\mathbf{v}_m)\\
    &=\epsilon_mg_{m+1,m}\mathbf{v}_{m+1}.
  \end{aligned}
  \label{eq:7-25}
\end{equation}
于是
\[
  \abs{h_{m+1,m}}=\abs{g_{m+1,m}}.
\]
由于 $h_{m+1,m}\neq0$, 由 \eqref{eq:7-25} 知
\[
  \mathbf{u}_{m+1}=\epsilon_{m+1}\mathbf{v}_{m+1},
\]
其中 $\epsilon_{m+1}=1$ 或 $-1$. 由归纳法, 证明完成.
\end{proof}

我们指出, 对上 Hessenberg 矩阵 $H=[h_{ij}]$, 若 $h_{i+1,i}\neq0$ ($i=1,2,\ldots,n-1$), 则它是不可约的. 定理 \ref{thm:7-5} 表明: 若 $Q^TAQ=H$ 是不可约上 Hessenberg 矩阵, 其中 $Q$ 为正交矩阵, 则 $Q$ 与 $H$ 完全由 $Q$ 的第一列确定 (相差 $\pm1$).

现设 $H\in\R^{n\times n}$ 为上 Hessenberg 矩阵, 要对 $H$ 作一次 QR 迭代. 首先需构造 $H$ 的 QR 分解. 由于 $H$ 具有特殊结构, 可用 $n-1$ 个 Givens 旋转得到 QR 分解. 为简单起见, 下面考虑 $n=5$ 的情形. 设已求得两个 Givens 旋转 $P_{12}$ 与 $P_{23}$ 使
\[
  P_{23}P_{12}H=\begin{bmatrix}
    \times & \times & \times & \times & \times\\
    0 & \times & \times & \times & \times\\
    0 & 0 & h_{33} & \times & \times\\
    0 & 0 & h_{43} & \times & \times\\
    0 & 0 & 0 & \times & \times
  \end{bmatrix}.
\]
然后构造 Givens 旋转 $P_{34}=G(3,4,\theta_3)$, 使 $\theta_3$ 满足
\[
  \begin{bmatrix}
    \cos\theta_3 & \sin\theta_3\\
    -\sin\theta_3 & \cos\theta_3
  \end{bmatrix}
  \begin{bmatrix}
    h_{33}\\ h_{43}
  \end{bmatrix}
  =\begin{bmatrix}
    \times\\ 0
  \end{bmatrix}.
\]
于是
\[
  P_{34}P_{23}P_{12}H=\begin{bmatrix}
    \times & \times & \times & \times & \times\\
    0 & \times & \times & \times & \times\\
    0 & 0 & \times & \times & \times\\
    0 & 0 & 0 & \times & \times\\
    0 & 0 & 0 & \times & \times
  \end{bmatrix}.
\]
因此容易看出: 对 $n$ 阶上 Hessenberg 矩阵 $H$, 可构造 $n-1$ 个 Givens 旋转 $P_{12},P_{23},\ldots,P_{n-1,n}$ 使
\[
  P_{n-1,n}P_{n-2,n-1}\cdots P_{12}H=R
\]
为上三角矩阵. 令
\[
  Q=(P_{n-1,n}P_{n-2,n-1}\cdots P_{12})^T.
\]
则有 $H=QR$. 为完成一次 QR 迭代, 还须计算
\[
  \widetilde{H}=RQ=RP_{12}^TP_{23}^T\cdots P_{n-1,n}^T.
\]
注意 $RP_{12}^T$ 与 $R$ 只在头两列不同. 由于 $R$ 为上三角矩阵, $RP_{12}^T$ 应具有如下形式 ($n=5$):
\[
  RP_{12}^T=\begin{bmatrix}
    \times & \times & \times & \times & \times\\
    \times & \times & \times & \times & \times\\
    0 & 0 & \times & \times & \times\\
    0 & 0 & 0 & \times & \times\\
    0 & 0 & 0 & 0 & \times
  \end{bmatrix}.
\]
类似地, $RP_{12}^TP_{23}^T$ 与 $RP_{12}^T$ 只在第二列与第三列不同. 因此
\[
  RP_{12}^TP_{23}^T=\begin{bmatrix}
    \times & \times & \times & \times & \times\\
    \times & \times & \times & \times & \times\\
    0 & \times & \times & \times & \times\\
    0 & 0 & 0 & \times & \times\\
    0 & 0 & 0 & 0 & \times
  \end{bmatrix}.
\]
如此继续, 最终得到矩阵 $\widetilde{H}$, 它也是上 Hessenberg 矩阵. 容易知道, 对上 Hessenberg 矩阵作一次 QR 迭代的运算量为 $\mathcal{O}(n^2)$; 而对一般矩阵作一次 QR 迭代的运算量为 $\mathcal{O}(n^3)$.

\subsection{单步位移 QR 迭代}
\label{sec:7-5-2}

现在只需讨论 Hessenberg 形式. 由定理 \ref{thm:7-4} 知, 基本 QR 算法的收敛速度是线性的, 且取决于特征值之间的距离. 为加快收敛速度, 引入位移策略. 带单步位移的 QR 迭代如下:
\begin{equation}
  \left\{
  \begin{array}{l}
    H_m-\mu_mI=Q_mR_m,\\
    H_{m+1}=R_mQ_m+\mu_mI,
  \end{array}
  \right.
  \label{eq:7-26}
\end{equation}
其中 $m=1,2,\ldots$, $H_1=H\in\R^{n\times n}$ 是给定的上 Hessenberg 矩阵, 满足定理 \ref{thm:7-4} 的条件, 而 $\mu_m\in\R$. 若 $H$ 的所有特征值都假定为实数, 我们来考虑如何选取位移. 由于 $H_m$ 是上 Hessenberg 矩阵, 其最后一行只有两个非零元 $h_{n,n-1}^{(m)}$ 与 $h_{nn}^{(m)}$ ($n=5$):
\[
  H_m=\begin{bmatrix}
    \times & \times & \times & \times & \times\\
    \times & \times & \times & \times & \times\\
    0 & \times & \times & \times & \times\\
    0 & 0 & \times & \times & \times\\
    0 & 0 & 0 & h_{n,n-1}^{(m)} & h_{nn}^{(m)}
  \end{bmatrix}.
\]
当 QR 算法收敛时, $h_{n,n-1}^{(m)}$ 很小, 而 $h_{nn}^{(m)}$ 趋于 $H$ 的某个特征值. 因此可取位移 $\mu_m=h_{nn}^{(m)}$. 事实上可以证明: 若 $h_{n,n-1}^{(m)}=\epsilon$ 很小, 则 QR 算法迭代一次后有
\begin{equation}
  h_{n,n-1}^{(m+1)}=\mathcal{O}(\epsilon^2).
  \label{eq:7-27}
\end{equation}

由上述讨论知, 把 $H_m-h_{nn}^{(m)}I$ 化为上三角矩阵需要 $n-1$ 步. 设前 $n-2$ 步已完成:
\[
  \widehat{H}=\begin{bmatrix}
    \times & \times & \times & \times & \times\\
    0 & \times & \times & \times & \times\\
    0 & 0 & \times & \times & \times\\
    0 & 0 & 0 & \alpha & \beta\\
    0 & 0 & 0 & \epsilon & 0
  \end{bmatrix}.
\]
实际上只需研究矩阵 $\widehat{H}$ 右下角的 $2\times2$ 子矩阵
\[
  H_{m2}=\begin{bmatrix}
    \alpha & \beta\\
    \epsilon & 0
  \end{bmatrix}.
\]
在第 $(n-1)$ 步约化中, 要确定 $c=\cos\theta$ 与 $s=\sin\theta$ 使
\[
  \begin{bmatrix}
    c & s\\
    -s & c
  \end{bmatrix}
  \begin{bmatrix}
    \alpha\\ \epsilon
  \end{bmatrix}
  =\begin{bmatrix}
    \sigma\\ 0
  \end{bmatrix}.
\]
容易看出
\[
  c=\alpha/\sigma,\qquad s=\epsilon/\sigma,\qquad
  \sigma=\sqrt{\alpha^2+\epsilon^2}.
\]
于是, 对矩阵 $H_{m+1}=R_mQ_m+h_{nn}^{(m)}I$ 右下角的 $2\times2$ 子矩阵作若干简单计算后得
\[
  h_{n,n-1}^{(m+1)}=-s^2\beta=-\beta\epsilon^2/\sigma^2=\mathcal{O}(\epsilon^2),
\]
即 \eqref{eq:7-27} 成立. 通过一次位移, QR 迭代的收敛速度可望达到二次. 若 $H$ 有复特征值, 则可用双步位移策略加快 QR 迭代的收敛速度.

\subsection{双步位移 QR 迭代}
\label{sec:7-5-3}

注意若 $H_m$ 右下角的子矩阵
\[
  G=\begin{bmatrix}
    h_{pp}^{(m)} & h_{pn}^{(m)}\\
    h_{np}^{(m)} & h_{nn}^{(m)}
  \end{bmatrix},
  \qquad p=n-1
\]
有一对共轭复特征值 $\mu_1$ 与 $\mu_2$, 则对 \eqref{eq:7-26} 可以预料会出现困难: 不能指望 $h_{nn}^{(m)}$ 趋于 $A$ 的某个特征值. 绕开这一困难的办法是施行如下带双步位移的 QR 算法:
\begin{equation}
  \left\{
  \begin{array}{l}
    H-\mu_1I=Q_1R_1,\\
    H_1=R_1Q_1+\mu_1I,\\
    H_1-\mu_2I=Q_2R_2,\\
    H_2=R_2Q_2+\mu_2I,
  \end{array}
  \right.
  \label{eq:7-28}
\end{equation}
其中 $H=H_m$. 令
\begin{equation}
  M\equiv(H-\mu_1I)(H-\mu_2I).
  \label{eq:7-29}
\end{equation}
经若干简单计算, 有
\begin{equation}
  M=QR,
  \label{eq:7-30}
\end{equation}
以及
\begin{equation}
  H_2=Q^THQ,
  \label{eq:7-31}
\end{equation}
其中
\[
  Q=Q_1Q_2,\qquad R=R_2R_1.
\]
由 \eqref{eq:7-29} 得
\[
  M=H^2-sH+tI,
\]
其中
\[
  s=\mu_1+\mu_2=h_{pp}^{(m)}+h_{nn}^{(m)}\in\R,
\]
以及
\[
  t=\mu_1\mu_2=\det(G)\in\R.
\]
因此 $M$ 也是实矩阵. 若 $\mu_1$, $\mu_2$ 不是 $H$ 的特征值, 且每次迭代中 $R_1$, $R_2$ 的对角元取为正, 则由 \eqref{eq:7-30} 知 $Q$ 也是实矩阵; 再由 \eqref{eq:7-31} 便知 $H_2$ 是实矩阵. 因此, 在不计舍入误差的假设下, 由 \eqref{eq:7-28} 得到的 $H_2$ 仍是实上 Hessenberg 矩阵. 实际计算中, 由于舍入误差, $H_2$ 通常不是实矩阵. 为保持 $H_2$ 的实性, 利用 \eqref{eq:7-30} 与 \eqref{eq:7-31} 提出如下计算 $H_2$ 的过程:
\begin{parenlist}
  \item 计算 $M=H^2-sH+tI$.
  \item 计算 $M$ 的 QR 分解: $M=QR$.
  \item 计算 $H_2=Q^THQ$.
\end{parenlist}
注意第 (1) 步中形成矩阵 $M$ 的运算量为 $\mathcal{O}(n^3)$. 我们指出, 实际上不必显式形成矩阵 $M$ \cite{src:30}. 关于带位移策略的 QR 算法的实用实现, 参见 \cite{src:30,src:32,src:98}.

\needspace{6\baselineskip}
\begin{exercises}
  \item 证明: 若 $T\in\C^{n\times n}$ 为上三角矩阵且正规, 则 $T$ 是对角矩阵.

  \item 设 $A,B\in\C^{n\times n}$. 证明 $AB$ 的谱等于 $BA$ 的谱.

  \item 设 $A\in\C^{n\times n}$, $\mathbf{x}\in\C^n$, $X=[\mathbf{x}\mid A\mathbf{x}\mid\cdots\mid A^{n-1}\mathbf{x}]$. 证明: 若 $X$ 非奇异, 则 $X^{-1}AX$ 是上 Hessenberg 矩阵.

  \item 设 $A\in\C^{n\times n}$ 有互不相同的特征值. 证明: 若 $Q^*AQ=T$ 是 Schur 分解, 且 $AB=BA$, 则 $Q^*BQ$ 是上三角矩阵.

  \item 设 $A\in\R^{n\times n}$, $\mathbf{z}\in\R^n$. 给出一个详细算法, 用以计算正交矩阵 $Q$, 使 $Q^TAQ$ 为上 Hessenberg 矩阵, 且 $Q^T\mathbf{z}$ 是 $\mathbf{e}_1$ 的倍数, 其中 $\mathbf{e}_1$ 是第一个单位向量.

  \item 设 $W,Y\in\R^{n\times n}$, 定义矩阵 $C,B$ 为
    \[
      C=W+\ii Y,\qquad
      B=\begin{bmatrix}
        W & -Y\\
        Y & W
      \end{bmatrix}.
    \]
    证明: 若 $\lambda\in\R$ 是 $C$ 的特征值, 则 $\lambda$ 也是 $B$ 的特征值. 两个对应的特征向量之间有何关系?

  \item 设
    \[
      A=\begin{bmatrix}
        w & x\\
        y & z
      \end{bmatrix}\in\R^{2\times2}
    \]
    有特征值 $\lambda\pm\ii\mu$, 其中 $\mu\neq0$. 求一个稳定地确定 $c=\cos\theta$ 与 $s=\sin\theta$ 的算法, 使
    \[
      \begin{bmatrix}
        c & s\\
        -s & c
      \end{bmatrix}^T
      \begin{bmatrix}
        w & x\\
        y & z
      \end{bmatrix}
      \begin{bmatrix}
        c & s\\
        -s & c
      \end{bmatrix}
      =\begin{bmatrix}
        \lambda & \beta\\
        \alpha & \lambda
      \end{bmatrix},
    \]
    其中 $\alpha\beta=-\mu^2$.

  \item 求一个 $2\times2$ 对角矩阵 $D$, 使 $\norm{D^{-1}AD}_F$ 达到极小, 其中
    \[
      A=\begin{bmatrix}
        w & x\\
        y & z
      \end{bmatrix}.
    \]

  \item 设 $H=H_1$ 为给定矩阵. 按下式生成矩阵 $H_k$:
    \[
      H_k-\mu_kI=Q_kR_k,\qquad
      H_{k+1}=R_kQ_k+\mu_kI.
    \]
    证明
    \[
      (Q_1\cdots Q_j)(R_j\cdots R_1)=(H-\mu_1I)\cdots(H-\mu_jI).
    \]

  \item 证明: 若
    \[
      Y=\begin{bmatrix}
        I & Z\\
        \mathbf{0} & I
      \end{bmatrix},
    \]
    则
    \[
      \kappa_2(Y)\equiv\norm{Y}_2\norm{Y^{-1}}_2
      =\frac{1}{2}\left(2+\sigma^2+\sqrt{4\sigma^2+\sigma^4}\right),
    \]
    其中 $\sigma=\norm{Z}_2$.

  \item 设 $A$ 的对角元为实数. 证明
    \[
      \abs{\operatorname{Im}(\lambda)}
      \leqslant\max_i\sum_{j\neq i}\abs{a_{ij}},
    \]
    其中 $\lambda$ 是 $A$ 的任一特征值, $\operatorname{Im}(\cdot)$ 表示复数的虚部.

  \item 证明: 若
    \[
      \frac{1}{2}(A+A^T)
    \]
    正定, 则 $\operatorname{Re}(\lambda)>0$, 其中 $\lambda$ 是 $A$ 的任一特征值, $\operatorname{Re}(\cdot)$ 表示复数的实部.

  \item 设 $B$ 满足 $\norm{B}_2<1$. 证明 $I-B$ 非奇异, 且 $I+2(B-I)^{-1}$ 的特征值有负实部.
\end{exercises}

% chapters/ch08.tex —— 第 8 章：对称特征值问题
% 源书：Chapter 8 Symmetric Eigenvalue Problems，书页 121–143（物理页 135–157）
% 公式 (8.1)–(8.7)；定理 8.1–8.9；引理 8.1；推论 8.1；算法 8.1–8.7；习题 1–12
\chapter{对称特征值问题}
\label{ch:sym-eig}

对称特征值问题性质优良, 数学理论丰富, 是数值线性代数中最令人愉悦的课题之一. 本章研究对称特征值问题. 我们先介绍对称矩阵的一些基本谱性质, 然后讨论对称 QR 方法, Jacobi 方法与对分法, 最后介绍一种分治算法.

\section{基本谱性质}
\label{sec:8-1}

我们先介绍任一对称矩阵 $A\in\R^{n\times n}$ 的特征值与特征向量的一些基本性质. 众所周知, 任一对称矩阵 $A$ 的特征值都是实数, 且 $A$ 的特征向量构成 $\R^n$ 的一组标准正交基.

\begin{theorem}[谱分解定理 \cite{src:39}]
\label{thm:8-1}
设 $A\in\R^{n\times n}$ 为对称矩阵, 则存在正交矩阵 $Q\in\R^{n\times n}$ 使
\[
  Q^TAQ=\diag(\lambda_1,\ldots,\lambda_n).
\]
\end{theorem}

对称矩阵的特征值具有基于所谓 \newterm{rayleighquot}{Rayleigh 商}{Rayleigh quotient}的极小极大性质:
\[
  \frac{\mathbf{x}^TA\mathbf{x}}{\mathbf{x}^T\mathbf{x}}.
\]
我们有下述定理, 其证明可在 \cite{src:32} 中找到.

\begin{theorem}[Courant--Fischer 极小极大定理]
\label{thm:8-2}
设 $A\in\R^{n\times n}$ 为对称矩阵, 其特征值按如下次序排列:
\[
  \lambda_1\geqslant\cdots\geqslant\lambda_n.
\]
则
\[
  \lambda_i=\max_{\dim(S)=i}\min_{\mathbf{0}\neq\mathbf{u}\in S}
  \frac{\mathbf{u}^TA\mathbf{u}}{\mathbf{u}^T\mathbf{u}}
  =\min_{\dim(S)=n-i+1}\max_{\mathbf{0}\neq\mathbf{u}\in S}
  \frac{\mathbf{u}^TA\mathbf{u}}{\mathbf{u}^T\mathbf{u}},
\]
其中 $S$ 是 $\R^n$ 的任一子空间.
\end{theorem}

下一个定理 \cite{src:95} 刻画对称矩阵特征值的敏感性.

\begin{theorem}[Weyl--Wielandt--Hoffman 定理]
\label{thm:8-3}
设 $A,B\in\R^{n\times n}$ 为对称矩阵, 其特征值分别按如下次序排列:
\[
  \lambda_1(A)\geqslant\cdots\geqslant\lambda_n(A),
  \qquad
  \mu_1(B)\geqslant\cdots\geqslant\mu_n(B),
\]
则
\[
  \abs{\lambda_i(A)-\mu_i(B)}\leqslant\norm{A-B}_2,
  \qquad i=1,2,\ldots,n,
\]
并且
\[
  \sum_{i=1}^{n}\bigl(\lambda_i(A)-\mu_i(B)\bigr)^2\leqslant\norm{A-B}_F^2.
\]
\end{theorem}

定理 \ref{thm:8-3} 告诉我们, 任一对称矩阵的特征值都是良态的, 即 $A$ 的元素受到的微小扰动只会引起其特征值的微小变化.

\begin{theorem}[Cauchy 交错定理 \cite{src:39,src:95}]
\label{thm:8-4}
设 $A\in\R^{n\times n}$ 对称, $A_r$ 表示 $A$ 的 $r$ 阶顺序主子矩阵, 则
\[
  \lambda_{r+1}(A_{r+1})\leqslant\lambda_r(A_r)\leqslant\lambda_r(A_{r+1})
  \leqslant\cdots\leqslant\lambda_2(A_{r+1})\leqslant\lambda_1(A_r)
  \leqslant\lambda_1(A_{r+1}),
\]
其中 $r=1,2,\ldots,n-1$.
\end{theorem}

至于特征向量的敏感性, 我们有下述定理 \cite{src:95}.

\begin{theorem}
\label{thm:8-5}
设 $A,A+E\in\R^{n\times n}$ 为对称矩阵,
\[
  Q=[\mathbf{q}_1\mid Q_2]\in\R^{n\times n}
\]
为正交矩阵, 其中 $\mathbf{q}_1$ 是 $A$ 的单位特征向量. 把矩阵 $Q^TAQ$ 与 $Q^TEQ$ 分块如下:
\[
  Q^TAQ=\begin{bmatrix}
    \lambda & \mathbf{0}\\
    \mathbf{0} & D_{22}
  \end{bmatrix},
  \qquad
  Q^TEQ=\begin{bmatrix}
    \epsilon & \mathbf{e}^T\\
    \mathbf{e} & E_{22}
  \end{bmatrix},
\]
其中 $D_{22},E_{22}\in\R^{(n-1)\times(n-1)}$. 若
\[
  d=\min_{\mu}\abs{\lambda-\mu}>0,
  \qquad
  \norm{E}_2\leqslant d/4,
\]
其中 $\mu$ 是 $D_{22}$ 的任一特征值, 则存在单位特征向量 $\tilde{\mathbf{q}}_1$ 使
\[
  \sin\theta=\sqrt{1-\abs{\mathbf{q}_1^T\tilde{\mathbf{q}}_1}^2}
  \leqslant\frac{4}{d}\norm{\mathbf{e}}_2\leqslant\frac{4}{d}\norm{E}_2,
\]
其中 $\theta=\arccos\abs{\mathbf{q}_1^T\tilde{\mathbf{q}}_1}$.
\end{theorem}

看来 $\theta$ 能够很好地度量 $\mathbf{q}_1$ 与 $\tilde{\mathbf{q}}_1$ 之间的接近程度. 可以看出, 单个特征向量对扰动的敏感程度取决于其对应特征值与其余特征值的分离程度.

任一对称矩阵的特征值同该矩阵的奇异值密切相关. \newterm{svdecomp}{奇异值分解}{singular value decomposition} \cite{src:32} 在数值线性代数中至关重要.

\begin{theorem}[奇异值分解定理]
\label{thm:8-6}
设 $A\in\R^{m\times n}$ 且 $\rank(A)=r$. 则存在正交矩阵
\[
  U=[\mathbf{u}_1\mid\mathbf{u}_2\mid\cdots\mid\mathbf{u}_m]\in\R^{m\times m},
  \qquad
  V=[\mathbf{v}_1\mid\mathbf{v}_2\mid\cdots\mid\mathbf{v}_n]\in\R^{n\times n}
\]
使
\[
  U^TAV=\begin{bmatrix}
    \Sigma_r & \mathbf{0}\\
    \mathbf{0} & \mathbf{0}
  \end{bmatrix},
\]
其中
\[
  \Sigma_r=\diag(\sigma_1,\ldots,\sigma_r)
\]
满足 $\sigma_1\geqslant\sigma_2\geqslant\cdots\geqslant\sigma_r>0$.
\end{theorem}

称 $\sigma_i$ ($i=1,2,\ldots,n$) 为 $A$ 的奇异值. 向量 $\mathbf{u}_i$ 与 $\mathbf{v}_i$ 分别称为第 $i$ 个左奇异向量与第 $i$ 个右奇异向量. 下一个推论就是关于奇异值的 Weyl--Wielandt--Hoffman 定理.

\begin{corollary}
\label{cor:8-1}
设 $A,B\in\R^{n\times n}$, 其奇异值分别按如下次序排列:
\[
  \sigma_1(A)\geqslant\cdots\geqslant\sigma_n(A),
  \qquad
  \tau_1(B)\geqslant\cdots\geqslant\tau_n(B),
\]
则
\[
  \abs{\sigma_i(A)-\tau_i(B)}\leqslant\norm{A-B}_2,
  \qquad i=1,2,\ldots,n,
\]
并且
\[
  \sum_{i=1}^{n}\bigl(\sigma_i(A)-\tau_i(B)\bigr)^2\leqslant\norm{A-B}_F^2.
\]
\end{corollary}

该推论表明, 任一实矩阵的奇异值也是良态的, 即 $A$ 的元素受到的微小扰动只会引起其奇异值的微小变化.

\section{对称 QR 方法}
\label{sec:8-2}

\newterm{symqrmethod}{对称 QR 方法}{symmetric QR method}是求解对称特征值问题的一种 QR 迭代. 它利用对称性对对称矩阵 $A$ 施行 QR 算法. 为构造高效的 QR 方法, 我们先把对称矩阵 $A$ 约化为三对角矩阵 $T$, 再对矩阵 $T$ 施行 QR 迭代.

\subsection{三对角 QR 迭代}
\label{sec:8-2-1}

设 $A$ 对称, 且假定 $A$ 可分解为
\[
  Q^TAQ=T,
\]
其中 $Q$ 为正交矩阵, $T$ 为上 Hessenberg 矩阵. 此时 $T$ 必为对称三对角矩阵. 于是我们只需处理一个对称三对角矩阵 $T$ 的特征值问题.

把 $A$ 分块如下:
\[
  A=\begin{bmatrix}
    \alpha_1 & \mathbf{v}_0^T\\
    \mathbf{v}_0 & A_0
  \end{bmatrix},
\]
其中 $A_0\in\R^{(n-1)\times(n-1)}$. 利用 Householder 变换, 可按如下方式把矩阵 $A$ 约化为对称三对角矩阵: 在第 $k$ 步,
\begin{parenlist}
  \item 计算 Householder 变换 $\tilde{H}_k\in\R^{(n-k)\times(n-k)}$ 使
    \[
      \tilde{H}_k\mathbf{v}_{k-1}=\beta_k\mathbf{e}_1,
      \qquad \beta_k\in\R.
    \]
  \item 计算
    \[
      \begin{bmatrix}
        \alpha_{k+1} & \mathbf{v}_k^T\\
        \mathbf{v}_k & A_k
      \end{bmatrix}
      =\tilde{H}_kA_{k-1}\tilde{H}_k,
    \]
    其中 $A_k\in\R^{(n-k-1)\times(n-k-1)}$.
\end{parenlist}

若用上述约化产生的 $\alpha_k$, $\beta_k$ 与 $\tilde{H}_k$ 定义
\[
  T=\begin{bmatrix}
    \alpha_1 & \beta_1 & & \mathbf{0}\\
    \beta_1 & \alpha_2 & \ddots & \\
     & \ddots & \ddots & \beta_{n-1}\\
    \mathbf{0} & & \beta_{n-1} & \alpha_n
  \end{bmatrix},
\]
\[
  H_k=\diag(I_k,\tilde{H}_k)\in\R^{n\times n},
  \qquad
  Q=H_1H_2\cdots H_{n-2},
\]
其中
\[
  \begin{bmatrix}
    \alpha_{n-1} & \beta_{n-1}\\
    \beta_{n-1} & \alpha_n
  \end{bmatrix}
  =\tilde{H}_{n-2}A_{n-3}\tilde{H}_{n-2},
\]
则有
\[
  Q^TAQ=T.
\]
由上述约化容易看出, 第 $k$ 步的主要运算量在于计算 $\tilde{H}_kA_{k-1}\tilde{H}_k$. 令
\[
  \tilde{H}_k=I-\beta\mathbf{v}\mathbf{v}^T,
  \qquad \mathbf{v}\in\R^{n-k}.
\]

利用 $A_{k-1}$ 的对称性, 可得
\[
  \tilde{H}_kA_{k-1}\tilde{H}_k=A_{k-1}-\mathbf{v}\mathbf{w}^T-\mathbf{w}\mathbf{v}^T,
\]
其中
\[
  \mathbf{w}=\mathbf{u}-\frac{1}{2}\beta(\mathbf{v}^T\mathbf{u})\mathbf{v},
  \qquad
  \mathbf{u}=\beta A_{k-1}\mathbf{v}.
\]
由于只需计算该矩阵的上三角部分, 可见从 $A_{k-1}$ 到 $A_k$ 的过渡只需 $4(n-k)^2$ 次运算. 给定对称矩阵 $A\in\R^{n\times n}$, 下述算法用 $T=Q^TAQ$ 覆盖 $A$, 其中 $T$ 为三对角矩阵, $Q$ 为若干 Householder 变换之积.

\needspace{10\baselineskip}
\begin{algorithm}{用 Householder 变换作三对角化}
  \label{alg:8-1}
  \[
    \left\{
    \begin{array}{ll}
      \textbf{for} & k=1:n-2\\
        & [\mathbf{v},\beta]=\mathrm{house}(A(k+1:n,k))\\
        & \mathbf{u}=\beta A(k+1:n,k+1:n)\mathbf{v}\\
        & \mathbf{w}=\mathbf{u}-(\beta\mathbf{u}^T\mathbf{v}/2)\mathbf{v}\\
        & A(k+1,k)=\norm{A(k+1:n,k)}_2\\
        & A(k,k+1)=A(k+1,k)\\
        & A(k+1:n,k+1:n)=A(k+1:n,k+1:n)-\mathbf{v}\mathbf{w}^T-\mathbf{w}\mathbf{v}^T\\
      \textbf{end}
    \end{array}
    \right.
  \]
\end{algorithm}

该算法需要 $4n^3/3$ 次运算. 若需要显式求出 $Q$, 则还要再花 $4n^3/3$ 次运算.

把对称矩阵 $A$ 约化为对称三对角矩阵 $T$ 之后, 我们的目标便转为为 QR 迭代选取合适的位移. 考虑如下带单步位移 $\mu_k$ 的 QR 迭代:
\begin{equation}
  \left\{
  \begin{array}{l}
    T_k-\mu_kI=Q_kR_k,\\
    T_{k+1}=R_kQ_k+\mu_kI,
  \end{array}
  \right.
  \label{eq:8-1}
\end{equation}
其中 $k=1,2,\ldots$, $T_1=T$ 为对称三对角矩阵, \eqref{eq:8-1} 中的每个矩阵 $T_k$ 也都是对称三对角矩阵. 与非对称情形的 QR 算法一样, 这里假定 $T_k$ 不可约, 即次对角元非零. 下面讨论如何选取位移 $\mu_k$. 可以在每次迭代中取 $\mu_k=T_k(n,n)$, 即 $(n,n)$ 元作为位移. 然而, 更好的选法是取
\[
  \mu_k=\alpha_n+\delta-\sign(\delta)\sqrt{\delta^2+\beta_{n-1}^2},
\]
其中 $\delta=(\alpha_{n-1}-\alpha_n)/2$. 这就是著名的 \newterm{wilkinsonshift}{Wilkinson 位移}{Wilkinson shift} \cite{src:95}. 注意 $\mu_k$ 恰是矩阵
\[
  T_k(n-1:n,n-1:n)=\begin{bmatrix}
    \alpha_{n-1} & \beta_{n-1}\\
    \beta_{n-1} & \alpha_n
  \end{bmatrix}
\]
的两个特征值中较接近 $\alpha_n$ 的那一个.

\subsection{隐式对称 QR 迭代}
\label{sec:8-2-2}

对显式 QR 算法 \eqref{eq:8-1}, 可以不必显式形成矩阵 $T-\mu I$ 而完成从
\[
  T-\mu I=QR
\]
到
\[
  \tilde{T}=RQ+\mu I
\]
的过渡. \eqref{eq:8-1} 的实质是用正交相似变换把 $T$ 变成 $\tilde{T}$. 由定理 \ref{thm:7-5} 知, $\tilde{T}$ 可由 $Q$ 的第一列完全确定. 由用 Givens 旋转作 $T-\mu I$ 的 QR 分解的过程可知
\[
  Q\mathbf{e}_1=G_1\mathbf{e}_1,
\]
其中 $G_1=G(1,2,\theta_1)$ 是使 $T-\mu I$ 第一列中第二个元素变为零的旋转. $\theta_1$ 可由
\[
  \begin{bmatrix}
    \cos\theta_1 & \sin\theta_1\\
    -\sin\theta_1 & \cos\theta_1
  \end{bmatrix}
  \begin{bmatrix}
    \alpha_1-\mu\\
    \beta_1
  \end{bmatrix}
  =\begin{bmatrix}
    *\\
    0
  \end{bmatrix}
\]
算出. 令
\[
  B=G_1TG_1^T.
\]
则 $B$ 具有如下形式 ($n=4$):
\[
  B=\begin{bmatrix}
    * & * & + & 0\\
    * & * & * & 0\\
    + & * & * & *\\
    0 & 0 & * & *
  \end{bmatrix}.
\]
令 $G_i=G(i,i+1,\theta_i)$, $i=2,3$. 则这种形式的 $G_i$ 可用来把矩阵 $B$ 中不需要的非零元``+''追赶出去, 如下所示:
\[
  B\ \xrightarrow{\ G_2\ }\
  \begin{bmatrix}
    * & * & 0 & 0\\
    * & * & * & +\\
    0 & * & * & *\\
    0 & + & * & *
  \end{bmatrix}
  \ \xrightarrow{\ G_3\ }\
  \begin{bmatrix}
    * & * & 0 & 0\\
    * & * & * & 0\\
    0 & * & * & *\\
    0 & 0 & * & *
  \end{bmatrix}.
\]
一般地, 若 $Z=G_1G_2\cdots G_{n-1}$, 则
\[
  Z\mathbf{e}_1=G_1\mathbf{e}_1=Q\mathbf{e}_1
\]
且 $ZTZ^T$ 三对角. 于是由定理 \ref{thm:7-5} 知, 用这种``追赶零元''技术得到的三对角矩阵 $ZTZ^T$ 与用显式方法 \eqref{eq:8-1} 得到的三对角矩阵 $T$ 本质上是同一个矩阵. 总之, 我们得到

\needspace{10\baselineskip}
\begin{algorithm}{带 Wilkinson 位移的隐式对称 QR 迭代}
  \label{alg:8-2}
  {\footnotesize
  \[
    \left\{
    \begin{array}{l}
      d=(T(n-1,n-1)-T(n,n))/2\\
      \mu=T(n,n)-T(n,n-1)^2/\bigl(d+\sign(d)\sqrt{d^2+T(n,n-1)^2}\bigr)\\
      x=T(1,1)-\mu;\quad z=T(2,1)\\
      \textbf{for}\ k=1:n-1\\
      \quad [c,s]=\mathrm{givens}(x,z)\\
      \quad T=G_kTG_k^T,\quad \text{其中}\ G_k=G(k,k+1,\theta_k)\\
      \quad \textbf{if}\ k<n-1\\
      \quad\quad x=T(k+1,k);\quad z=T(k+2,k)\\
      \quad \textbf{end}\\
      \textbf{end}
    \end{array}
    \right.
  \]}
\end{algorithm}

该算法约需 $30n$ 次运算与 $n$ 次开平方. 当然, 在实际实现中三对角矩阵 $T$ 应存放在一对 $n$ 维向量中.

\subsection{隐式对称 QR 算法}
\label{sec:8-2-3}

算法 \ref{alg:8-2} 是对称 QR 算法的基础——后者是计算任一对称矩阵\newterm{spectraldecomp}{谱分解}{spectral decomposition}的标准手段. 应用算法 \ref{alg:8-2}, 我们得到下述算法.

\needspace{10\baselineskip}
\begin{algorithm}{隐式对称 QR 算法}
  \label{alg:8-3}
  \renewcommand{\labelenumii}{(\roman{enumii})}
  \begin{parenlist}
    \item 输入 $A$ (实对称矩阵).
    \item 三对角化: 用算法 \ref{alg:8-1} 计算 $A$ 的三对角化,
      \[
        T=U_0^TAU_0.
      \]
      置 $Q=U_0$.
    \item 收敛性判别:
      \begin{parenlist}
        \item 对 $i=1,\ldots,n-1$, 若
          \[
            \abs{t_{i+1,i}}=\abs{t_{i,i+1}}
            \leqslant\bigl(\abs{t_{i,i}}+\abs{t_{i+1,i+1}}\bigr)\mu_\epsilon,
          \]
          则令 $t_{i+1,i}$ 与 $t_{i,i+1}$ 为零, 其中 $\mu_\epsilon$ 是机器精度.
        \item 求最大的整数 $m\geqslant0$ 与最小的整数 $l\geqslant0$ 使
          \[
            T=\begin{bmatrix}
              T_{11} & \mathbf{0} & \mathbf{0}\\
              \mathbf{0} & T_{22} & \mathbf{0}\\
              \mathbf{0} & \mathbf{0} & T_{33}
            \end{bmatrix},
          \]
          其中 $T_{11}\in\R^{l\times l}$, $T_{22}\in\R^{(n-l-m)\times(n-l-m)}$ 为不可约三对角矩阵, 而 $T_{33}\in\R^{m\times m}$ 为对角矩阵.
        \item 若 $m=n$, 则输出; 否则转下一步.
      \end{parenlist}
    \item QR 迭代: 对 $T_{22}$ 应用算法 \ref{alg:8-2}:
      \[
        T_{22}=GT_{22}G^T,
        \qquad
        G=G_1G_2\cdots G_{n-l-m-1}.
      \]
    \item 置 $Q=Q\diag(I_l,G,I_m)$, 然后转到第 (3) 步.
  \end{parenlist}
\end{algorithm}

若只需计算特征值, 该算法约需 $4n^3/3$ 次运算. 若同时需要特征值与特征向量, 则约需 $9n^3$ 次运算. 可以证明 \cite{src:95}, 由算法 \ref{alg:8-3} 得到的计算特征值 $\tilde{\lambda}_i$ ($i=1,2,\ldots,n$) 满足
\[
  Q^T(A+E)Q=\diag(\tilde{\lambda}_1,\ldots,\tilde{\lambda}_n),
\]
其中 $Q\in\R^{n\times n}$ 正交, 且 $\norm{E}_2\approx\norm{A}_2\mu_\epsilon$, $\mu_\epsilon$ 为机器精度. 利用定理 \ref{thm:8-3}, 有
\[
  \abs{\lambda_i-\tilde{\lambda}_i}\approx\norm{A}_2\mu_\epsilon,
  \qquad i=1,2,\ldots,n,
\]
其中每个 $\lambda_i$ 都是 $A$ 的特征值. 每个 $\tilde{\lambda}_i$ 的绝对误差都很小, 且相对误差小于机器精度 $\mu_\epsilon$. 若记计算所得标准正交特征向量构成的矩阵为 $\tilde{Q}=[\tilde{\mathbf{q}}_1\mid\tilde{\mathbf{q}}_2\mid\cdots\mid\tilde{\mathbf{q}}_n]$, 则每个 $\tilde{\mathbf{q}}_i$ 的精度取决于 $\lambda_i$ 与其余特征值的分离程度.

\section{Jacobi 方法}
\label{sec:8-3}

Jacobi 方法是计算对称特征值问题最早的方法之一. 它由 Jacobi 于 1846 年提出 (见 \cite{src:32}). 众所周知, 实对称矩阵可用正交相似变换约化为对角矩阵. Jacobi 方法利用矩阵的对称性, 选取适当的旋转把对称矩阵约化为对角形式. Jacobi 方法通常比对称 QR 算法慢得多. 然而, Jacobi 方法至今仍令人感兴趣, 因为该方法易于编程, 而且近年人们认识到它具有内在的并行性.

\subsection{基本思想}
\label{sec:8-3-1}

设 $A=[a_{ij}]\in\R^{n\times n}$ 为对称矩阵,
\[
  \mathrm{off}(A)\equiv\Bigl(\norm{A}_F^2-\sum_{i=1}^{n}a_{ii}^2\Bigr)^{1/2}
  =\Bigl(\sum_{\substack{i=1\\ j\neq i}}^{n}a_{ij}^2\Bigr)^{1/2}.
\]
Jacobi 方法的思想是系统地把 $\mathrm{off}(A)$ 约化为零. 实现这一点的基本工具称为 Jacobi 旋转, 其定义如下:
\[
  J(p,q,\theta)=I+\sin\theta\bigl(\mathbf{e}_p\mathbf{e}_q^T-\mathbf{e}_q\mathbf{e}_p^T\bigr)
  +(\cos\theta-1)\bigl(\mathbf{e}_p\mathbf{e}_p^T+\mathbf{e}_q\mathbf{e}_q^T\bigr),
\]
其中 $p<q$, $\mathbf{e}_k$ 表示第 $k$ 个单位向量. 注意 Jacobi 旋转与 Givens 旋转并无区别 (见第 4.2.2 小节). 本节改用这一名称是为了纪念其发明者. Jacobi 过程的每一步包含:
\begin{parenlist}
  \item 选取 $p$ 与 $q$ 作旋转, 其中 $1\leqslant p<q\leqslant n$.
  \item 计算旋转角 $\theta$ 使
    \begin{equation}
      \begin{bmatrix}
        b_{pp} & b_{pq}\\
        b_{qp} & b_{qq}
      \end{bmatrix}
      =\begin{bmatrix}
        c & s\\
        -s & c
      \end{bmatrix}^T
      \begin{bmatrix}
        a_{pp} & a_{pq}\\
        a_{qp} & a_{qq}
      \end{bmatrix}
      \begin{bmatrix}
        c & s\\
        -s & c
      \end{bmatrix}
      \label{eq:8-2}
    \end{equation}
    为对角矩阵, 即 $b_{pq}=b_{qp}=0$, 其中 $c=\cos\theta$, $s=\sin\theta$.
  \item 用 $B=[b_{ij}]=J^TAJ$ 覆盖 $A$, 其中 $J=J(p,q,\theta)$.
\end{parenlist}

注意矩阵 $B$ 除第 $p$ 行 (列) 与第 $q$ 行 (列) 外与矩阵 $A$ 相同. 具体关系为:
\[
  b_{ip}=b_{pi}=ca_{ip}-sa_{iq},
  \qquad i\neq p,q,
\]
\[
  b_{iq}=b_{qi}=sa_{ip}+ca_{iq},
  \qquad i\neq p,q,
\]
\[
  b_{pp}=c^2a_{pp}-2sca_{pq}+s^2a_{qq},
\]
\[
  b_{qq}=s^2a_{pp}+2sca_{pq}+c^2a_{qq},
\]
\[
  b_{pq}=b_{qp}=(c^2-s^2)a_{pq}+sc(a_{pp}-a_{qq}).
\]

我们先考虑如何实际计算 $s=\sin\theta$ 与 $c=\cos\theta$, 使 \eqref{eq:8-2} 中 $b_{pq}=b_{qp}=0$. 这实际上等价于
\begin{equation}
  a_{pq}(c^2-s^2)+(a_{pp}-a_{qq})cs=0.
  \label{eq:8-3}
\end{equation}
若 $a_{pq}=0$, 则直接取 $c=1$, $s=0$. 否则定义
\[
  \tau=\frac{a_{qq}-a_{pp}}{2a_{pq}},
  \qquad
  t=\tan\theta=\frac{s}{c},
\]
并由 \eqref{eq:8-3} 推知 $t$ 满足二次方程
\[
  t^2+2\tau t-1=0.
\]
于是
\[
  t=-\tau\pm\sqrt{1+\tau^2}.
\]
我们取 $t$ 为两个根中较小者, 这既保证 $\abs{\theta}\leqslant\pi/4$, 又使 $\norm{B-A}_F^2$ 达到极小, 因为
\[
  \norm{B-A}_F^2
  =4(1-c)\sum_{\substack{i=1\\ i\neq p,q}}^{n}\bigl(a_{ip}^2+a_{iq}^2\bigr)
  +\frac{2a_{pq}^2}{c^2}.
\]
$t$ 确定后, 便可由公式
\[
  c=\frac{1}{\sqrt{1+t^2}},
  \qquad
  s=tc
\]
得到 $c$ 与 $s$.

下面总结 Jacobi 旋转 $J(p,q,\theta)$ 的计算. 给定对称矩阵 $A\in\R^{n\times n}$ 与满足 $1\leqslant p<q\leqslant n$ 的下标 $p,q$, 下述算法计算一对余弦–正弦值, 使 $b_{pq}=b_{qp}=0$, 其中 $b_{jk}$ 是矩阵 $B=J(p,q,\theta)^TAJ(p,q,\theta)$ 的第 $(j,k)$ 元.

\par\medskip
\needspace{8\baselineskip}
\refstepcounter{algorithm}\label{alg:8-4}
\noindent{\headfont 算法\ \textbf{\thealgorithm}}\par
\smallskip
  \[
    \left\{
    \begin{array}{ll}
      \textbf{function}: & [c,s]=\mathrm{sym}(A,p,q)\\
        & \textbf{if}\ A(p,q)\neq0\\
        & \quad \tau=(A(q,q)-A(p,p))/(2A(p,q))\\
        & \quad \textbf{if}\ \tau\geqslant0\\
        & \qquad t=1/(\tau+\sqrt{1+\tau^2})\\
        & \quad \textbf{else}\\
        & \qquad t=-1/(-\tau+\sqrt{1+\tau^2})\\
        & \quad \textbf{end}\\
        & \quad c=1/\sqrt{1+t^2}\\
        & \quad s=tc\\
        & \textbf{else}\\
        & \quad c=1\\
        & \quad s=0\\
        & \textbf{end}\\
      \textbf{end}
    \end{array}
    \right.
  \]
\par\medskip

$J(p,q,\theta)$ 确定后, 更新
\[
  A\leftarrow J(p,q,\theta)^TAJ(p,q,\theta)
\]
可用 $6n$ 次运算完成.

那么如何选取整数 $p$ 与 $q$? 由于正交变换保持 Frobenius 范数 (见第 \ref{sec:3-1} 节), 有 $\norm{B}_F=\norm{A}_F$. 注意
\[
  a_{pp}^2+a_{qq}^2+2a_{pq}^2=b_{pp}^2+b_{qq}^2+2b_{pq}^2=b_{pp}^2+b_{qq}^2,
\]
于是
\begin{align}
    \mathrm{off}(B)^2
    &=\norm{B}_F^2-\sum_{i=1}^{n}b_{ii}^2\notag\\
    &=\norm{A}_F^2-\sum_{i=1}^{n}a_{ii}^2
      +\bigl(a_{pp}^2+a_{qq}^2-b_{pp}^2-b_{qq}^2\bigr)\notag\\
    &=\mathrm{off}(A)^2-2a_{pq}^2.
  \label{eq:8-4}
\end{align}
我们的目标是使 $\mathrm{off}(B)$ 达到极小, 故 $p,q$ 的最佳选取应为
\[
  \abs{a_{pq}}=\max_{1\leqslant i<j\leqslant n}\abs{a_{ij}}.
\]
这就是\newterm{classicaljacobi}{经典 Jacobi 方法}{classical Jacobi method}的基本思想.

\subsection{经典 Jacobi 方法}
\label{sec:8-3-2}

给定对称矩阵 $A\in\R^{n\times n}$ 与容差 $\epsilon>0$, 下述算法用 $U^TAU$ 覆盖 $A$, 其中 $U$ 为正交矩阵, 且 $\mathrm{off}(U^TAU)\leqslant\epsilon\norm{A}_F$.

\needspace{10\baselineskip}
\begin{algorithm}{经典 Jacobi 方法}
  \label{alg:8-5}
  \[
    \left\{
    \begin{array}{l}
      U=I_n;\quad \mathrm{eps}=\epsilon\norm{A}_F\\
      \textbf{while}\ \mathrm{off}(A)>\mathrm{eps}\\
      \quad \text{选取}\ p,q\ \text{使}\ \abs{a_{pq}}=\max\limits_{i\neq j}\abs{a_{ij}}\\
      \quad [c,s]=\mathrm{sym}(A,p,q)\\
      \quad A=J(p,q,\theta)^TAJ(p,q,\theta)\\
      \quad U=UJ(p,q,\theta)\\
      \textbf{end}
    \end{array}
    \right.
  \]
\end{algorithm}

由于 $\abs{a_{pq}}$ 是最大的非对角元, 有
\[
  \mathrm{off}(A)^2\leqslant N\bigl(a_{pq}^2+a_{qp}^2\bigr),
\]
其中 $N=n(n-1)/2$. 由 \eqref{eq:8-4} 知
\[
  \mathrm{off}(B)^2\leqslant\Bigl(1-\frac{1}{N}\Bigr)\mathrm{off}(A)^2.
\]
若 $A_k$ 表示第 $k$ 次 Jacobi 迭代后的矩阵, 且 $A_0=A$, 则由归纳法得
\[
  \mathrm{off}(A_k)^2\leqslant\Bigl(1-\frac{1}{N}\Bigr)^{k}\mathrm{off}(A_0)^2.
\]
这意味着经典 Jacobi 方法线性收敛. 然而实际上, Jacobi 方法的渐近收敛速度是二次的. 我们可以证明, 当 $k$ 充分大时, 存在常数 $c$ 使
\[
  \mathrm{off}(A_{k+N})\leqslant c\cdot\mathrm{off}(A_k)^2.
\]
参见 \cite{src:32} 及其中的参考文献. 因此, 经过足够多次迭代后, 非对角``范数''以二次速度趋于零.

Jacobi 方法的另一个优点是容易计算特征向量. 若迭代在第 $k$ 次旋转后停止, 则有
\[
  A_k=J_k^TJ_{k-1}^T\cdots J_1^TAJ_1J_2\cdots J_k.
\]
记
\[
  Q_k=J_1J_2\cdots J_k.
\]
于是
\[
  AQ_k=Q_kA_k.
\]
由于 $A_k$ 的非对角元很小, $A_k$ 的对角元便是 $A$ 的特征值的良好近似, 上述恒等式表明 $Q_k$ 的各列是 $A$ 的特征向量的良好近似, 且所有近似特征向量都是标准正交的. 我们可在 Jacobi 迭代过程中同时得到近似特征向量 $Q_k$.

\subsection{并行 Jacobi 方法}
\label{sec:8-3-3}

用 Jacobi 方法求解对称特征值问题本身具有并行性. 为说明这一点, 设 $A\in\R^{8\times8}$ 对称. 若有一台含 4 个处理器的并行计算机, 则可把 28 个子问题分成如下 7 组旋转:
\[
  \begin{array}{ll}
    \text{group (1)}: & (1,2),(3,4),(5,6),(7,8);\\
    \text{group (2)}: & (1,3),(2,4),(5,7),(6,8);\\
    \text{group (3)}: & (1,4),(2,3),(5,8),(6,7);\\
    \text{group (4)}: & (1,5),(2,6),(3,7),(4,8);\\
    \text{group (5)}: & (1,6),(2,5),(3,8),(4,7);\\
    \text{group (6)}: & (1,7),(2,8),(3,5),(4,6);\\
    \text{group (7)}: & (1,8),(2,7),(3,6),(4,5).
  \end{array}
\]
注意每一组内的 4 个旋转互不冲突. 例如, 第一组中的子问题 $J(2i-1,2i,\theta_i)$ ($i=1,2,3,4$) 可并行完成. 当我们计算 $J(1,2,\theta_1)^TAJ(1,2,\theta_1)$ 时, 它不影响旋转 $(3,4)$, $(5,6)$ 与 $(7,8)$. 于是, 计算
\[
  A=AJ(1,2,\theta_1),\qquad A=AJ(3,4,\theta_2),
\]
\[
  A=AJ(5,6,\theta_3),\qquad A=AJ(7,8,\theta_4)
\]
可由 4 个处理器并行执行. 类似地, 计算
\[
  A=J(1,2,\theta_1)^TA,\qquad A=J(3,4,\theta_2)^TA,
\]
\[
  A=J(5,6,\theta_3)^TA,\qquad A=J(7,8,\theta_4)^TA
\]
也可由 4 个处理器并行完成. 对上述例子, 只需单处理器计算机 $1/4$ 的计算时间. 并行 Jacobi 算法可在 \cite{src:32} 中找到.

\section{对分法}
\label{sec:8-4}

本节介绍求解对称三对角特征值问题的\newterm{bisection}{对分法}{bisection method}. 把对分法与三对角化技巧结合起来, 便可得到计算对称矩阵某个指定特征值及其对应特征向量的一种数值方法. 对给定的三对角矩阵
\begin{equation}
  T=\begin{bmatrix}
    a_1 & b_1 & & \mathbf{0}\\
    b_1 & a_2 & \ddots & \\
     & \ddots & \ddots & b_{n-1}\\
    \mathbf{0} & & b_{n-1} & a_n
  \end{bmatrix}\in\R^{n\times n},
  \label{eq:8-5}
\end{equation}
我们考虑 $T$ 的特征值的计算. 不失一般性, 假设 $b_i\neq0$ ($i=1,2,\ldots,n-1$), 即 $T$ 是不可约对称三对角矩阵. 否则, $T$ 可分成若干较小的不可约对称三对角矩阵.

对每个 $i$, 令 $p_i(\lambda)$ 为 $T$ 的 $i$ 阶顺序主子矩阵 $T_i$ 的特征多项式. 则这些多项式满足三项递推:
\begin{equation}
  \begin{aligned}
    p_0(\lambda)&\equiv1,\qquad p_1(\lambda)=a_1-\lambda,\\
    p_i(\lambda)&=(a_i-\lambda)p_{i-1}(\lambda)-b_{i-1}^2p_{i-2}(\lambda),
    \qquad i=2,3,\ldots,n.
  \end{aligned}
  \label{eq:8-6}
\end{equation}
由于 $T$ 对称, 多项式 $p_i(\lambda)$ ($i=1,2,\ldots,n$) 的根都是实数. 下述\newterm{interlace}{交错性质}{interlacing property}非常重要.

\begin{theorem}[Sturm 序列性质]
\label{thm:8-7}
设 \eqref{eq:8-5} 中的对称三对角矩阵 $T$ 不可约. 则 $T_{i-1}$ 的特征值严格分离 $T_i$ 的特征值:
\[
  \lambda_i(T_i)<\lambda_{i-1}(T_{i-1})<\lambda_{i-1}(T_i)
  <\cdots<\lambda_2(T_i)<\lambda_1(T_{i-1})<\lambda_1(T_i).
\]
此外, 若 $s_n(\lambda)$ 表示序列
\[
  \{p_0(\lambda),p_1(\lambda),\ldots,p_n(\lambda)\}
\]
中变号的次数, 则 $s_n(\lambda)$ 等于 $T$ 的小于 $\lambda$ 的特征值的个数, 其中 $p_i(\lambda)$ 由 \eqref{eq:8-6} 定义. 若 $p_i(\lambda)=0$, 则 $p_{i-1}(\lambda)p_{i+1}(\lambda)<0$.
\end{theorem}

\begin{proof}
由定理 \ref{thm:8-4} 知, $T_{i-1}$ 的特征值弱分离 $T_i$ 的特征值. 下面证明这种分离必为严格分离. 设对某个 $i$ 与 $\mu$ 有 $p_i(\mu)=p_{i-1}(\mu)=0$. 由于 $T$ 不可约, 由 \eqref{eq:8-6} 注意到
\[
  p_0(\mu)=p_1(\mu)=\cdots=p_i(\mu)=0,
\]
这与 $p_0(\lambda)\equiv1$ 矛盾. 故为严格分离. 关于 $s_n(\lambda)$ 的论断在 \cite{src:95} 中有详细论述.
\end{proof}

计算 $T$ 的某个指定特征值的对分法可陈述如下.

\needspace{10\baselineskip}
\begin{algorithm}{对分算法}
  \label{alg:8-6}
  设 $\lambda_1<\lambda_2<\cdots<\lambda_n$ 为 $T$ 的特征值, 即 $p_n(\lambda)$ 的根, $\epsilon$ 为容差. 设对给定的 $m\leqslant n$, 所求特征值为 $\lambda_m$. 则
  \begin{parenlist}
    \item 求包含 $\lambda_m$ 的区间 $[l_0,u_0]$. 由于 $\abs{\lambda_1}\leqslant\rho(T)\leqslant\norm{T}_\infty$, 可取 $l_0=-\norm{T}_\infty$, $u_0=\norm{T}_\infty$.
    \item 计算 $r_1=\dfrac{l_0+u_0}{2}$ 与 $s_n(r_1)$.
    \item 若 $s_n(r_1)\geqslant m$, 则 $\lambda_m\in[l_0,r_1]$, 置 $l_1=l_0$, $u_1=r_1$; 否则 $\lambda_m\in[r_1,u_0]$, 置 $l_1=r_1$, $u_1=u_0$.
    \item 若 $\abs{l_1-u_1}<\epsilon$, 取 $r_2=\dfrac{l_1+u_1}{2}$ 作为 $\lambda_m$ 的近似值. 否则转到第 (2) 步.
  \end{parenlist}
\end{algorithm}

由上述算法可见, 主要运算量在于计算 $s_n(\mu)$. 然而在实际计算中, 不能通过计算 $p_i(\mu)$ 的值来得到 $s_n(\mu)$, 因为高次多项式的求值很困难. 为避免这一困难, 定义
\[
  q_i(\lambda)=\frac{p_i(\lambda)}{p_{i-1}(\lambda)},
  \qquad i=1,2,\ldots,n.
\]
由 \eqref{eq:8-6} 得
\[
  q_1(\lambda)=p_1(\lambda)=a_1-\lambda,
  \qquad
  q_i(\lambda)=a_i-\lambda-\frac{b_{i-1}^2}{q_{i-1}(\lambda)},
  \quad i=2,3,\ldots,n.
\]
容易验证, $s_n(\mu)$ 恰好是序列 $q_1(\mu),\ldots,q_n(\mu)$ 中负值的个数. 下述算法是计算 $s_n(\mu)$ 的实用算法.

\needspace{10\baselineskip}
\begin{algorithm}{计算变号次数}
  \label{alg:8-7}
  \[
    \left\{
    \begin{array}{l}
      \mathbf{x}=[a_1,a_2,\ldots,a_n]\\
      \mathbf{y}=[0,b_1,\ldots,b_{n-1}]\\
      s=0;\quad q=\mathbf{x}(1)-\mu\\
      \textbf{for}\ k=1:n\\
      \quad \textbf{if}\ q<0\\
      \quad\quad s=s+1\\
      \quad \textbf{end}\\
      \quad \textbf{if}\ k<n\\
      \quad\quad \textbf{if}\ q=0\\
      \quad\quad\quad q=\abs{\mathbf{y}(k+1)}\mu_\epsilon\\
      \quad\quad \textbf{end}\\
      \quad\quad q=\mathbf{x}(k+1)-\mu-\mathbf{y}(k+1)^2/q\\
      \quad \textbf{end}\\
      \textbf{end}
    \end{array}
    \right.
  \]
\end{algorithm}

其中 $\mu_\epsilon$ 是机器精度.

当 $q_i(\mu)=0$ 时, 我们把它当作正数处理. 因此可在算法 \ref{alg:8-7} 中用一个小正数 $\abs{b_i}\mu_\epsilon$ 代替 $q_i(\mu)$. 若事先存好 $b_i^2$, 则算法 \ref{alg:8-7} 需要 $3n$ 次运算. 若用对分法平均 $m$ 次算出一个特征值, 则运算量为 $3nm$. 因此, 用对分法计算任一对称三对角矩阵的特征值是高效的. 另一方面, 舍入误差分析表明对分法是数值稳定的 \cite{src:95}.

\section{分治方法}
\label{sec:8-5}

\newterm{divideconquer}{分治方法}{divide-and-conquer method}是 Dongarra 与 Sorensen 于 1987 年为计算对称三对角矩阵的全部特征值与特征向量而提出的一种数值方法 \cite{src:32}. 其基本思想是把原来的对称三对角矩阵撕裂成 $2^k$ 个规模较小的对称三对角矩阵, 然后计算每个较小对称三对角矩阵的谱分解. 得到这些较小的谱分解之后, 再把它们合并起来, 形成原矩阵的谱分解. 因此, 该方法适于并行计算.

\subsection{撕裂}
\label{sec:8-5-1}

设 $T\in\R^{n\times n}$ 给定如下:
\[
  T=\begin{bmatrix}
    a_1 & b_1 & & \mathbf{0}\\
    b_1 & a_2 & \ddots & \\
     & \ddots & \ddots & b_{n-1}\\
    \mathbf{0} & & b_{n-1} & a_n
  \end{bmatrix}.
\]
不失一般性, 设 $n=2m$. 令
\[
  \mathbf{v}=[\underbrace{0,\ldots,0}_{m-1},1,\theta,
  \underbrace{0,\ldots,0}_{m-1}]^T\in\R^n.
\]
考虑矩阵
\[
  \tilde{T}=T-\rho\mathbf{v}\mathbf{v}^T,
\]
其中 $\theta,\rho\in\R$ 待定. 容易看出, $\tilde{T}$ 与 $T$ 除中间的 4 个元素外完全相同:
\[
  \begin{bmatrix}
    a_m-\rho & b_m-\rho\theta\\
    b_m-\rho\theta & a_{m+1}-\rho\theta^2
  \end{bmatrix}.
\]
若取 $\rho\theta=b_m$, 则
\[
  T=\begin{bmatrix}
    T_1 & \mathbf{0}\\
    \mathbf{0} & T_2
  \end{bmatrix}+\rho\mathbf{v}\mathbf{v}^T,
\]
其中
\[
  T_1=\begin{bmatrix}
    a_1 & b_1 & & & \mathbf{0}\\
    b_1 & a_2 & b_2 & & \\
     & \ddots & \ddots & \ddots & \\
     & & b_{m-2} & a_{m-1} & b_{m-1}\\
    \mathbf{0} & & & b_{m-1} & \tilde{a}_m
  \end{bmatrix}
\]
以及
\[
  T_2=\begin{bmatrix}
    \tilde{a}_{m+1} & b_{m+1} & & & \mathbf{0}\\
    b_{m+1} & a_{m+2} & b_{m+2} & & \\
     & \ddots & \ddots & \ddots & \\
     & & b_{n-2} & a_{n-1} & b_{n-1}\\
    \mathbf{0} & & & b_{n-1} & a_n
  \end{bmatrix},
\]
且
\[
  \tilde{a}_m=a_m-\rho,
  \qquad
  \tilde{a}_{m+1}=a_{m+1}-\rho\theta^2.
\]
因此, $T$ 被分成一个分块矩阵与一个秩一矩阵之和. 若对 $T_1$ 与 $T_2$ 反复施行这样的分裂, 最终可把 $T$ 分成 $2^k$ 块.

\subsection{合并}
\label{sec:8-5-2}

一旦得到 $T_1$ 与 $T_2$ 的谱分解:
\[
  Q_1^TT_1Q_1=D_1,
  \qquad
  Q_2^TT_2Q_2=D_2,
\]
其中 $Q_1,Q_2\in\R^{m\times m}$ 为正交矩阵, $D_1,D_2\in\R^{m\times m}$ 为对角矩阵, 我们的目标便是由 $T_1$ 与 $T_2$ 的谱分解计算 $T$ 的谱分解, 即求一个正交矩阵 $V$ 使
\[
  V^TTV=\diag(\lambda_1,\ldots,\lambda_n).
\]
令
\[
  U=\begin{bmatrix}
    Q_1 & \mathbf{0}\\
    \mathbf{0} & Q_2
  \end{bmatrix},
\]
则
\[
  \begin{aligned}
    U^TTU
    &=\begin{bmatrix}
        Q_1 & \mathbf{0}\\
        \mathbf{0} & Q_2
      \end{bmatrix}^T
      \left(\begin{bmatrix}
        T_1 & \mathbf{0}\\
        \mathbf{0} & T_2
      \end{bmatrix}+\rho\mathbf{v}\mathbf{v}^T\right)
      \begin{bmatrix}
        Q_1 & \mathbf{0}\\
        \mathbf{0} & Q_2
      \end{bmatrix}\\
    &=D+\rho\mathbf{z}\mathbf{z}^T,
  \end{aligned}
\]
其中
\[
  D=\diag(D_1,D_2),
  \qquad
  \mathbf{z}=U^T\mathbf{v}.
\]
于是, 求 $T$ 的谱分解的问题便归结为计算 $D+\rho\mathbf{z}\mathbf{z}^T$ 的谱分解. 下面考虑如何快速而稳定地计算 $D+\rho\mathbf{z}\mathbf{z}^T$ 的谱分解.

\begin{lemma}
\label{lem:8-1}
设 $D=\diag(d_1,\ldots,d_n)\in\R^{n\times n}$ 且 $d_1>d_2>\cdots>d_n$. 假定 $0\neq\rho\in\R$, $\mathbf{z}=[z_1,z_2,\ldots,z_n]^T\in\R^n$ 且对所有 $i$ 有 $z_i\neq0$. 设 $\mathbf{u}\in\R^n$ 与 $\lambda\in\R$ 满足
\[
  (D+\rho\mathbf{z}\mathbf{z}^T)\mathbf{u}=\lambda\mathbf{u},
  \qquad \mathbf{u}\neq\mathbf{0}.
\]
则 $\mathbf{z}^T\mathbf{u}\neq0$ 且 $D-\lambda I$ 非奇异.
\end{lemma}

\begin{proof}
若 $\mathbf{z}^T\mathbf{u}=0$, 则 $D\mathbf{u}=\lambda\mathbf{u}$ 且 $\mathbf{u}\neq\mathbf{0}$, 即 $\lambda$ 是 $D$ 的特征值, $\mathbf{u}$ 是与 $\lambda$ 相伴的特征向量. 由于 $D$ 是对角元互不相同的对角矩阵, 必存在某个 $i$ 使 $d_i=\lambda$ 且 $\mathbf{u}=\alpha\mathbf{e}_i$, 其中 $\alpha\neq0$, $\mathbf{e}_i$ 是第 $i$ 个单位向量. 于是
\[
  0=\mathbf{z}^T\mathbf{u}=\alpha\mathbf{z}^T\mathbf{e}_i=\alpha z_i,
\]
推出 $z_i=0$, 矛盾. 因此 $\mathbf{z}^T\mathbf{u}\neq0$.

另一方面, 若 $D-\lambda I$ 奇异, 则存在某个 $i$ 使 $\mathbf{e}_i^T(D-\lambda I)=0$, 于是
\[
  0=\mathbf{e}_i^T(D-\lambda I)\mathbf{u}=-\rho\mathbf{z}^T\mathbf{u}\mathbf{e}_i^T\mathbf{z}.
\]
由于 $\rho\mathbf{z}^T\mathbf{u}\neq0$, 得 $\mathbf{e}_i^T\mathbf{z}=z_i=0$, 矛盾. 故 $D-\lambda I$ 非奇异.
\end{proof}

\begin{theorem}
\label{thm:8-8}
设 $D=\diag(d_1,\ldots,d_n)\in\R^{n\times n}$ 且 $d_1>d_2>\cdots>d_n$. 假定 $0\neq\rho\in\R$, $\mathbf{z}=[z_1,z_2,\ldots,z_n]^T\in\R^n$ 且对所有 $i$ 有 $z_i\neq0$. 设 $D+\rho\mathbf{z}\mathbf{z}^T$ 的谱分解为
\[
  V^T(D+\rho\mathbf{z}\mathbf{z}^T)V=\diag(\lambda_1,\ldots,\lambda_n),
\]
其中 $V=[\mathbf{v}_1\mid\mathbf{v}_2\mid\cdots\mid\mathbf{v}_n]$ 为正交矩阵, 且 $\lambda_1\geqslant\cdots\geqslant\lambda_n$. 则
\begin{romanlist}
  \item $\lambda_1,\ldots,\lambda_n$ 是函数
    \[
      f(\lambda)=1+\rho\mathbf{z}^T(D-\lambda I)^{-1}\mathbf{z}
    \]
    的 $n$ 个根.
  \item 若 $\rho>0$, 则
    \[
      \lambda_1>d_1>\lambda_2>\cdots>\lambda_n>d_n;
    \]
    若 $\rho<0$, 则
    \[
      d_1>\lambda_1>d_2>\cdots>d_n>\lambda_n.
    \]
  \item 存在常数 $\alpha_i\neq0$ 使
    \[
      \mathbf{v}_i=\alpha_i(D-\lambda_iI)^{-1}\mathbf{z},
      \qquad i=1,2,\ldots,n.
    \]
\end{romanlist}
\end{theorem}

\begin{proof}
由假定有
\[
  (D+\rho\mathbf{z}\mathbf{z}^T)\mathbf{v}_i=\lambda_i\mathbf{v}_i,
  \qquad \norm{\mathbf{v}_i}_2=1.
\]
由引理 \ref{lem:8-1} 知 $D-\lambda_iI$ 非奇异. 于是
\begin{equation}
  \mathbf{v}_i=-\rho\mathbf{z}^T\mathbf{v}_i(D-\lambda_iI)^{-1}\mathbf{z},
  \qquad i=1,2,\ldots,n,
  \label{eq:8-7}
\end{equation}
这就证明了 (iii). 注意 $D+\rho\mathbf{z}\mathbf{z}^T$ 的特征值互不相同. 否则, 若 $\lambda_i=\lambda_j$, 则 $\mathbf{v}_i$ 与 $\mathbf{v}_j$ 线性相关, 这与 $\mathbf{v}_i$ 和 $\mathbf{v}_j$ 的正交性矛盾.

用 $\mathbf{z}^T$ 乘 \eqref{eq:8-7} 两端, 并注意 $\mathbf{z}^T\mathbf{v}_i\neq0$, 得
\[
  1=-\rho\mathbf{z}^T(D-\lambda_iI)^{-1}\mathbf{z},
\]
即
\[
  f(\lambda_i)=0,
  \qquad i=1,2,\ldots,n.
\]
于是, $\lambda_i$ ($i=1,2,\ldots,n$) 都是 $f(\lambda)$ 的根. 下面证明 $f(\lambda)$ 恰有 $n$ 个零点. 注意
\[
  f(\lambda)=1+\rho\left(\frac{z_1^2}{d_1-\lambda}+\cdots+\frac{z_n^2}{d_n-\lambda}\right),
\]
并且
\[
  f'(\lambda)=\rho\left(\frac{z_1^2}{(d_1-\lambda)^2}+\cdots+\frac{z_n^2}{(d_n-\lambda)^2}\right).
\]
因此, $f(\lambda)$ 在相邻极点 $d_i$ 与 $d_{i+1}$ 之间严格单调. 若 $\rho>0$, $f(\lambda)$ 严格递增; 若 $\rho<0$, $f(\lambda)$ 严格递减. 于是容易看出 $f(\lambda)$ 恰有 $n$ 个根: 当 $\rho>0$ 时, 区间
\[
  (d_n,d_{n-1}),\ldots,(d_2,d_1),(d_1,\infty)
\]
中各有一个; 当 $\rho<0$ 时, 区间
\[
  (-\infty,d_n),(d_n,d_{n-1}),\ldots,(d_2,d_1)
\]
中各有一个. 故 (i) 与 (ii) 得证.
\end{proof}

由定理 \ref{thm:8-8}, 可按如下两步高效地计算 $D+\rho\mathbf{z}\mathbf{z}^T$ 的谱分解:
\begin{parenlist}
  \item 求 $f(\lambda)$ 的根 $\lambda_1,\lambda_2,\ldots,\lambda_n$. $f(\lambda)$ 在每个区间 $(d_{i+1},d_i)$ 中恰有一个根, 且在该区间内严格单调. 因此这一步可用类 Newton 方法 \cite{src:30} 快速而稳定地实现.
  \item 计算
    \[
      \mathbf{v}_i=\frac{(D-\lambda_iI)^{-1}\mathbf{z}}
      {\norm{(D-\lambda_iI)^{-1}\mathbf{z}}_2},
      \qquad i=1,2,\ldots,n.
    \]
\end{parenlist}

一般 $D+\rho\mathbf{z}\mathbf{z}^T$ 的谱分解问题可化归为定理 \ref{thm:8-8} 的情形. 为此, 我们可以构造性地证明下述定理.

\begin{theorem}
\label{thm:8-9}
设 $D=\diag(d_1,\ldots,d_n)\in\R^{n\times n}$, $\mathbf{z}\in\R^n$. 则存在正交矩阵 $V$ 以及 $\{1,2,\ldots,n\}$ 的一个排列 $\pi$ 使
\begin{romanlist}
  \item $V^T\mathbf{z}=[\eta_1,\ldots,\eta_r,0,\ldots,0]^T$, 其中 $\eta_i\neq0$ ($i=1,2,\ldots,r$).
  \item $V^TDV=\diag(d_{\pi(1)},\ldots,d_{\pi(n)})$, 其中 $d_{\pi(1)}>d_{\pi(2)}>\cdots>d_{\pi(r)}$.
\end{romanlist}
\end{theorem}

\begin{proof}
设两个下标 $i<j$ 满足 $d_i=d_j$. 则可取旋转 $P_{ij}=G(i,j,\theta)$ 使 $P_{ij}\mathbf{z}$ 的第 $j$ 个分量为零. 容易证明 $P_{ij}^TDP_{ij}=D$. 经若干步后, 可找到正交矩阵 $V_1$, 它是若干旋转之积, 使
\[
  V_1^TDV_1=D,
  \qquad
  V_1^T\mathbf{z}=[\xi_1,\ldots,\xi_n]^T,
\]
且具有性质: 若 $\xi_i\xi_j\neq0$ ($i\neq j$), 则 $d_i\neq d_j$.

若 $\xi_i=0$, $\xi_j\neq0$ ($i<j$), 则可找到一个置换矩阵交换第 $i$ 列与第 $j$ 列. 如此便可找到一个置换矩阵 $P_1$, 把所有非零的 $\xi_j$ 都移到前面, 即
\[
  P_1^TV_1^T\mathbf{z}=[\xi_{\pi_1(1)},\ldots,\xi_{\pi_1(n)}]^T,
\]
满足
\[
  \xi_{\pi_1(i)}\neq0,\quad i=1,2,\ldots,r,
  \qquad
  \xi_{\pi_1(i)}=0,\quad i=r+1,\ldots,n.
\]
这里 $\pi_1$ 是 $\{1,2,\ldots,n\}$ 的一个排列. 由 $P_1$ 的构造知
\[
  P_1^TV_1^TDV_1P_1=P_1^TDP_1=\diag\bigl(d_{\pi_1(1)},\ldots,d_{\pi_1(n)}\bigr),
\]
其中前 $r$ 个对角元 $d_{\pi_1(1)},\ldots,d_{\pi_1(r)}$ 互不相同.

最后, 可找到另一个 $r$ 阶置换矩阵 $P_2$ 使
\[
  P_2^T\diag\bigl(d_{\pi_1(1)},\ldots,d_{\pi_1(r)}\bigr)P_2
  =\diag(\mu_1,\ldots,\mu_r),
\]
其中 $\mu_1>\mu_2>\cdots>\mu_r$. 令
\[
  V=V_1P_1\diag(P_2,I_{n-r}),
\]
并令 $\pi$ 为 $\{1,2,\ldots,n\}$ 的由 $P_1$ 与 $P_2$ 确定的排列. 则
\[
  V^T\mathbf{z}=[\xi_{\pi(1)},\ldots,\xi_{\pi(r)},0,\ldots,0]^T
  =[\eta_1,\ldots,\eta_r,0,\ldots,0]^T,
\]
其中 $\eta_i\neq0$ ($i=1,2,\ldots,r$), 且
\[
  V^TDV=\diag\bigl(d_{\pi(1)},\ldots,d_{\pi(n)}\bigr),
\]
满足
\[
  d_{\pi(1)}>d_{\pi(2)}>\cdots>d_{\pi(r)}.
\]
证明完成.
\end{proof}

对任一 $D=\diag(d_1,\ldots,d_n)\in\R^{n\times n}$ 与 $\mathbf{z}\in\R^n$, 由定理 \ref{thm:8-9} 可构造正交矩阵 $V$ 使
\[
  V^T(D+\rho\mathbf{z}\mathbf{z}^T)V=
  \begin{bmatrix}
    D_1+\rho\mathbf{w}\mathbf{w}^T & \mathbf{0}\\
    \mathbf{0} & D_2
  \end{bmatrix},
\]
其中
\[
  D_1=\diag(d_{\pi(1)},\ldots,d_{\pi(r)})\in\R^{r\times r},
  \qquad
  d_{\pi(1)}>\cdots>d_{\pi(r)};
\]
\[
  D_2=\diag(d_{\pi(r+1)},\ldots,d_{\pi(n)})\in\R^{(n-r)\times(n-r)};
\]
并且
\[
  \mathbf{w}=[\eta_1,\ldots,\eta_r]^T,
  \qquad
  \eta_i\neq0,\quad i=1,2,\ldots,r.
\]
于是我们只需计算 $D_1+\rho\mathbf{w}\mathbf{w}^T$ 的谱分解, 而不必计算 $D+\rho\mathbf{z}\mathbf{z}^T$ 的谱分解.

最后, 我们简要介绍分治方法的并行计算. 为简单起见, 用一台含 4 个处理器的并行计算机计算一个 $4N$ 阶对称三对角矩阵 $T$ 的谱分解. 它可归纳为如下 4 步:
\begin{parenlist}
  \item 撕裂
    \[
      T=\begin{bmatrix}
        T_1 & \mathbf{0}\\
        \mathbf{0} & T_2
      \end{bmatrix}+\rho\mathbf{v}\mathbf{v}^T,
      \qquad
      T_1\in\R^{2N\times2N},
      \qquad
      \mathbf{v}\in\R^{4N};
    \]
    以及
    \[
      T_i=\begin{bmatrix}
        T_{i1} & \mathbf{0}\\
        \mathbf{0} & T_{i2}
      \end{bmatrix}+\rho_i\mathbf{w}_i\mathbf{w}_i^T,
    \]
    其中 $T_{ij}\in\R^{N\times N}$, $\mathbf{w}_i\in\R^{2N}$ ($i=1,2$).
  \item 由 4 个处理器并行计算 $T_{11},T_{12},T_{21},T_{22}$ 的谱分解.
  \item 合并 $T_{11}$, $T_{12}$ 的谱分解, 得到 $T_1$ 的谱分解; 合并 $T_{21}$, $T_{22}$ 的谱分解, 得到 $T_2$ 的谱分解. 这两件事可由 4 个处理器同时完成.
  \item 由 $T_1$ 与 $T_2$ 的谱分解合并得到 $T$ 的谱分解. 这可由 4 个处理器并行完成.
\end{parenlist}

由上述讨论可知, 分治方法可用于并行计算任一大型对称三对角矩阵的全部特征值与特征向量.

\needspace{6\baselineskip}
\begin{exercises}
  \item 求
    \[
      A=\begin{bmatrix}
        1 & 2\\
        2 & 3
      \end{bmatrix}
    \]
    的 Schur 分解.

  \item 证明: 若 $X\in\R^{n\times r}$, $r\leqslant n$, 且 $\norm{X^TX-I}_2=\tau<1$, 则
    \[
      \sigma_{\min}(X)\geqslant1-\tau,
    \]
    其中 $\sigma_{\min}$ 表示最小奇异值.

  \item 证明: $A=B+\ii C$ 为 Hermite 矩阵当且仅当
    \[
      M=\begin{bmatrix}
        B & -C\\
        C & B
      \end{bmatrix}
    \]
    是对称矩阵. 并给出 $A$ 的特征值, 特征向量与 $M$ 的特征值, 特征向量之间的关系.

  \item 把 $A=B+\ii C$ 的奇异值与奇异向量同
    \[
      \begin{bmatrix}
        B & -C\\
        C & B
      \end{bmatrix}
    \]
    的奇异值与奇异向量联系起来, 其中 $B,C\in\R^{m\times n}$.

  \item 利用奇异值分解证明: 若 $A\in\R^{m\times n}$ 且 $m\geqslant n$, 则存在矩阵 $Q\in\R^{m\times n}$ ($Q^TQ=I$) 与半正定矩阵 $P\in\R^{n\times n}$ 使 $A=QP$.

  \item 设
    \[
      A=\begin{bmatrix}
        I & B\\
        B^* & I
      \end{bmatrix},
    \]
    其中 $\norm{B}_2<1$. 证明
    \[
      \norm{A}_2\norm{A^{-1}}_2=\frac{1+\norm{B}_2}{1-\norm{B}_2}.
    \]

  \item 设
    \[
      A=\begin{bmatrix}
        a_1 & b_1 & & & \mathbf{0}\\
        c_1 & a_2 & b_2 & & \\
         & \ddots & \ddots & \ddots & \\
         & & \ddots & \ddots & b_{n-1}\\
        \mathbf{0} & & & c_{n-1} & a_n
      \end{bmatrix},
    \]
    其中 $b_ic_i>0$. 证明存在对角矩阵 $D$ 使 $D^{-1}AD$ 为对称三对角矩阵.

  \item 设
    \[
      T=\begin{bmatrix}
        2 & -1 & 0 & 0\\
        -1 & 2 & -1 & 0\\
        0 & -1 & 2 & -1\\
        0 & 0 & -1 & 2
      \end{bmatrix}.
    \]
    \subquest $T$ 是否正定?
    \subquest $T$ 在区间 $[0,2]$ 内有多少个特征值?

  \item 设 $A,E\in\R^{n\times n}$ 为两个对称矩阵. 证明: 若 $A$ 正定且 $\norm{A^{-1}}_2\norm{E}_2<1$, 则 $A+E$ 也正定.

  \item 设 $A\in\R^{m\times n}$ 且 $m\geqslant n$, 并假定 $A$ 的奇异值按如下次序排列:
    \[
      \sigma_1\geqslant\sigma_2\geqslant\cdots\geqslant\sigma_n.
    \]
    证明
    \[
      \sigma_i=\max_{\dim(S)=i}\min_{\mathbf{0}\neq\mathbf{u}\in S}
      \frac{\norm{A\mathbf{u}}_2}{\norm{\mathbf{u}}_2}
      =\min_{\dim(S)=n-i+1}\max_{\mathbf{0}\neq\mathbf{u}\in S}
      \frac{\norm{A\mathbf{u}}_2}{\norm{\mathbf{u}}_2},
    \]
    其中 $S$ 是 $\R^n$ 的任一子空间.

  \item 设 $A=[a_{ij}]\in\R^{n\times n}$ 对称, 且满足
    \subquest $a_{ii}>0$, $i=1,2,\ldots,n$.
    \subquest $a_{ij}\leqslant0$, $i\neq j$.
    \subquest $\displaystyle\sum_{i=1}^{n}a_{i1}>0$.
    \subquest $\displaystyle\sum_{i=1}^{n}a_{ij}=0$, $j=2,3,\ldots,n$.

    \hspace*{2em}证明 $A$ 的特征值非负.

  \item 设 $A$ 对称且带宽为 $p$. 证明: 若作带位移的 QR 迭代 $A-\mu I=QR$, $\tilde{A}=RQ+\mu I$, 则 $\tilde{A}$ 的带宽为 $p$.
\end{exercises}

% chapters/ch09.tex —— 第 9 章：在常微分方程与时滞微分方程中的应用
% 源书：Chapter 9 Applications in ODEs and DDEs，书页 144–169（物理页 158–183）
% 公式编号 (9.1)–(9.31)；定义 9.1；引理 9.1–9.3；定理 9.1–9.7；推论 9.1；
% 例 9.1–9.4；表 9.1–9.4；本章无习题
\chapter{在常微分方程与时滞微分方程中的应用}
\label{ch:ode-dde}

本章简要研究利用\newterm{bvm}{边值法}{boundary value method (BVM)}求解常微分方程组初值问题的一些新近进展. 这类方法需要求解一个或多个非对称, 大型且稀疏的线性方程组. 因此, 我们用第 6 章研究的 GMRES 方法并配以某些预条件子来求解这些线性方程组. 主要结果之一是: 若用 \newterm{astable}{$A_{\nu_1,\nu_2}$-稳定}{$A_{\nu_1,\nu_2}$-stable}的 BVM 求解 $n$ 阶常微分方程组, 则预处理矩阵可分解为 $I+L$, 其中 $I$ 是单位矩阵, $L$ 的秩至多为 $2n(\nu_1+\nu_2)$. 当对预处理后的方程组应用 GMRES 方法时, 该方法至多 $2n(\nu_1+\nu_2)+1$ 次迭代即收敛. 本章还给出对不同类型\newterm{dde}{时滞微分方程}{delay differential equation (DDE)}的应用. 关于用 BVMs 求解 ODEs 与 DDEs 的文献, 参见 \cite{src:5,src:8,src:10,src:17,src:19,src:44,src:45,src:47,src:48,src:55,src:58}.

\section{引言}
\label{sec:9-1}

我们从下述初值问题开始:
\begin{equation}
  \left\{
  \begin{array}{l}
    \mathbf{y}'(t)=J_n\mathbf{y}(t)+\mathbf{g}(t),\quad t\in(t_0,T],\\
    \mathbf{y}(t_0)=\mathbf{z},
  \end{array}
  \right.
  \label{eq:9-1}
\end{equation}
其中 $\mathbf{y}(t),\mathbf{g}(t):\R\to\R^n$, $\mathbf{z}\in\R^n$, $J_n\in\R^{n\times n}$. \newterm{ivm}{初值法}{initial value method (IVM)}, 如 Runge-Kutta 方法, 是求解 \eqref{eq:9-1} 的著名方法 \cite{src:83}. 最近, \cite{src:10} 中提出了另一类称为边值法 (BVMs) 的方法. 用 BVMs 离散 \eqref{eq:9-1}, 得到线性方程组
\[
  M\mathbf{u}=\mathbf{b}.
\]
使用 BVMs 的优点是方法更为稳定, 因而所得的线性方程组 $M\mathbf{u}=\mathbf{b}$ 的条件数也更好. 然而, 该方程组通常是大型且稀疏的 (具有带状结构), 求解它是应用 BVMs 时的主要问题.

第 6 章研究的 GMRES 方法用于求解 $M\mathbf{u}=\mathbf{b}$. 为加快 GMRES 迭代的收敛速度, 用一个称为 \newterm{strangblock}{Strang 型块循环预条件子}{Strang-type block-circulant preconditioner}\cite{src:19} 的预条件子 $S$ 对离散方程组作预处理. Strang 型预条件子的优点是: 若用 $A_{\nu_1,\nu_2}$-稳定的 BVM 求解 \eqref{eq:9-1}, 则 $S$ 可逆, 且预处理矩阵可分解为
\[
  S^{-1}M=I+L,
\]
其中 $L$ 的秩至多为 $2n(\nu_1+\nu_2)$, 它与积分步长无关. 由此可知, 在精确算术下, 对预处理后的方程组应用 GMRES 方法至多 $2n(\nu_1+\nu_2)+1$ 次迭代即收敛.

本章的安排如下. 9.2 节给出关于线性多步公式 (LMFs) 与 BVMs 的一些背景知识. 然后, 9.3 节研究常微分方程的 Strang 型块循环预条件子的性质, 并连同数值例子给出该方法的收敛性与运算量分析. 最后, 9.4–9.6 节讨论用 BVMs 的 Strang 型预条件子求解不同类型时滞微分方程 (DDEs) 的应用.

\section{边值法的背景}
\label{sec:9-2}

本节给出关于 LMFs 与 BVMs 的一些背景知识.

\subsection{线性多步公式}
\label{sec:9-2-1}

考虑初值问题
\[
  \left\{
  \begin{array}{ll}
    y'=f(t,y), & t\in(t_0,T],\\
    y(t_0)=y_0,
  \end{array}
  \right.
\]
其中 $y(t):\R\to\R$, $f(t,y):\R^2\to\R$. 步长为 $h$ 的均匀网格上的 $\mu$ 步\newterm{lmf}{线性多步公式}{linear multistep formula (LMF)}定义如下:
\begin{equation}
  \sum_{j=0}^{\mu}\alpha_jy_{m+j}=h\sum_{j=0}^{\mu}\beta_jf_{m+j},\quad m=0,1,\ldots,
  \label{eq:9-2}
\end{equation}
其中 $y_m$ 是 $y(t_m)$ 的离散逼近, $f_m$ 表示 $f(t_m,y_m)$.

为得到 \eqref{eq:9-2} 的解, 需要 $\mu$ 个初值
\[
  y_0,y_1,\ldots,y_{\mu-1}.
\]
由于原问题只提供 $y_0$, 我们不得不为其余的值
\[
  y_1,y_2,\ldots,y_{\mu-1}
\]
补充条件. 带 $\mu-1$ 个附加条件的方程 \eqref{eq:9-2} 称为初值法 (IVMs). 若 $\beta_{\mu}\neq0$, 则称 IVM 为隐式的; 若 $\beta_{\mu}=0$, 则称为显式的. 若把 IVM 用于区间 $[t_0,t_{N+\mu-1}]$ 上的初值问题, 则得到如下离散问题
\begin{equation}
  A_N\mathbf{y}=hB_N\mathbf{f}-\left[
  \begin{array}{c}
    \sum\limits_{i=0}^{\mu-1}(\alpha_iy_i-h\beta_if_i)\\
    \vdots\\
    \alpha_0y_{\mu-1}-h\beta_0f_{\mu-1}\\
    0\\
    \vdots\\
    0
  \end{array}\right],
  \label{eq:9-3}
\end{equation}
其中
\[
  \mathbf{y}=[y_{\mu},y_{\mu+1},\ldots,y_{N+\mu-1}]^T,\qquad
  \mathbf{f}=[f_{\mu},f_{\mu+1},\ldots,f_{N+\mu-1}]^T,
\]
\[
  A_N=\begin{bmatrix}
    \alpha_{\mu} & & & & \mathbf{0}\\
    \vdots & \ddots & & & \\
    \alpha_0 & & \ddots & & \\
    & \ddots & & \ddots & \\
    \mathbf{0} & & \alpha_0 & \cdots & \alpha_{\mu}
  \end{bmatrix},\qquad
  B_N=\begin{bmatrix}
    \beta_{\mu} & & & & \mathbf{0}\\
    \vdots & \ddots & & & \\
    \beta_0 & & \ddots & & \\
    & \ddots & & \ddots & \\
    \mathbf{0} & & \beta_0 & \cdots & \beta_{\mu}
  \end{bmatrix}.
\]
注意矩阵 $A_N,B_N\in\R^{N\times N}$ 是下带宽为 $\mu$ 的下三角带状 \newterm{toeplitzmatrix}{Toeplitz 矩阵}{Toeplitz matrix}. 我们回忆: 若矩阵的元素沿各对角线为常数, 则称之为 Toeplitz 矩阵. 此外, 线性方程组 \eqref{eq:9-3} 容易用前推递推求解. IVM 的一个经典例子是二阶\newterm{bdf}{向后差分公式}{backward differentiation formula (BDF)}
\[
  3y_{m+2}-4y_{m+1}+y_m=2hf_{m+1},
\]
这是一个两步方法, 其中 $\alpha_0=1$, $\alpha_1=-4$, $\alpha_2=3$, $\beta_1=2$.

除了用带 $\mu$ 个初值的 IVM 求解 \eqref{eq:9-1}, 我们也可以使用所谓的边值法 (BVMs). 给定 $\nu_1,\nu_2\geqslant0$ 使 $\nu_1+\nu_2=\mu$, 则相应的 BVM 需要 $\nu_1$ 个初始附加条件
\[
  y_0,y_1,\ldots,y_{\nu_1-1},
\]
以及 $\nu_2$ 个末端附加条件
\[
  y_N,y_{N+1},\ldots,y_{N+\nu_2-1},
\]
它们称为 $(\nu_1,\nu_2)$-边界条件. 注意 BVMs 类包含 IVMs 类 (即 $\nu_1=\mu$, $\nu_2=0$).

由带 $(\nu_1,\nu_2)$-边界条件的 $\mu$ 步 BVM 生成的离散问题可写成如下矩阵形式
\[
  A\mathbf{y}=hB\mathbf{f}-\left[
  \begin{array}{c}
    \sum\limits_{i=0}^{\nu_1-1}(\alpha_iy_i-h\beta_if_i)\\
    \vdots\\
    \alpha_0y_{\nu_1-1}-h\beta_0f_{\nu_1-1}\\
    0\\
    \vdots\\
    0\\
    \alpha_{\mu}y_N-h\beta_{\mu}f_N\\
    \vdots\\
    \sum\limits_{i=1}^{\nu_2}(\alpha_{\nu_1+i}y_{N-1+i}-h\beta_{\nu_1+i}f_{N-1+i})
  \end{array}\right],
\]
其中
\[
  \mathbf{y}=[y_{\nu_1},y_{\nu_1+1},\ldots,y_{N-1}]^T,\qquad
  \mathbf{f}=[f_{\nu_1},f_{\nu_1+1},\ldots,f_{N-1}]^T,
\]
而 $A$ 和 $B\in\R^{(N-\nu_1)\times(N-\nu_1)}$ 定义如下,
\begin{equation}
  A=\begin{bmatrix}
    \alpha_{\nu_1} & \cdots & \alpha_{\mu} & & \mathbf{0}\\
    \vdots & \ddots & \ddots & \ddots & \\
    \alpha_0 & \ddots & \ddots & \ddots & \alpha_{\mu}\\
    & \ddots & & \ddots & \vdots\\
    \mathbf{0} & & \alpha_0 & \cdots & \alpha_{\nu_1}
  \end{bmatrix},\qquad
  B=\begin{bmatrix}
    \beta_{\nu_1} & \cdots & \beta_{\mu} & & \mathbf{0}\\
    \vdots & \ddots & \ddots & \ddots & \\
    \beta_0 & \ddots & \ddots & \ddots & \beta_{\mu}\\
    & \ddots & & \ddots & \vdots\\
    \mathbf{0} & & \beta_0 & \cdots & \beta_{\nu_1}
  \end{bmatrix}.
  \label{eq:9-4}
\end{equation}
注意系数矩阵是下带宽为 $\nu_1$, 上带宽为 $\nu_2$ 的带状 Toeplitz 矩阵. BVMs 的一个例子是三阶\newterm{gbdf}{广义向后差分公式}{generalized backward differentiation formula (GBDF)}
\[
  2y_{m+1}+3y_m-6y_{m-1}+y_{m-2}=6hf_m,
\]
这是一个带 $(2,1)$-边界条件的三步方法, 其中
\[
  \alpha_0=1,\quad\alpha_1=-6,\quad\alpha_2=3,\quad\alpha_3=2,\quad\beta_2=6.
\]

虽然 IVMs 比 BVMs 更有效 (后者不能用前推递推求解), 但使用 BVMs 相对于 IVMs 的优点来自其稳定性性质. 例如, 通常的 BDF 对 $\mu>2$ 不是 $A$-稳定的, 而 GBDF 对任意 $\mu\geqslant1$ 都是 $A_{\nu_1,\nu_2}$-稳定的. 参见例如 \cite{src:1} 与 \cite{src:10} 第 79 页及图 5.1–5.3.

\subsection{块-BVMs 及其矩阵形式}
\label{sec:9-2-2}

令 $\mu=\nu_1+\nu_2$. 用基于 LMF 的 $\mu$ 步块-BVM 在均匀网格 $h=(T-t_0)/s$ 上求解 \eqref{eq:9-1}, 得
\begin{equation}
  \sum_{i=-\nu_1}^{\nu_2}\alpha_{i+\nu_1}\mathbf{y}_{m+i}
  =h\sum_{i=-\nu_1}^{\nu_2}\beta_{i+\nu_1}\mathbf{f}_{m+i},
  \quad m=\nu_1,\ldots,s-\nu_2.
  \label{eq:9-5}
\end{equation}
这里, $\mathbf{y}_m$ 是 $\mathbf{y}(t_m)$ 的离散逼近,
\[
  \mathbf{f}_m=J_n\mathbf{y}_m+\mathbf{g}_m,\qquad
  \mathbf{g}_m=\mathbf{g}(t_m).
\]
此外, \eqref{eq:9-5} 需要 $\nu_1$ 个初始条件与 $\nu_2$ 个末端条件, 它们由下述 $\mu-1$ 个附加方程提供:
\begin{equation}
  \sum_{i=0}^{\mu}\alpha_i^{(j)}\mathbf{y}_i
  =h\sum_{i=0}^{\mu}\beta_i^{(j)}\mathbf{f}_i,
  \quad j=1,\ldots,\nu_1-1,
  \label{eq:9-6}
\end{equation}
以及
\begin{equation}
  \sum_{i=0}^{\mu}\alpha_{\mu-i}^{(j)}\mathbf{y}_{s-i}
  =h\sum_{i=0}^{\mu}\beta_{\mu-i}^{(j)}\mathbf{f}_{s-i},
  \quad j=s-\nu_2+1,\ldots,s.
  \label{eq:9-7}
\end{equation}
\eqref{eq:9-6} 与 \eqref{eq:9-7} 中的系数 $\{\alpha^{(j)}\}$, $\{\beta^{(j)}\}$ 应选得使这些初始条件与末端条件的截断误差与 \eqref{eq:9-5} 中的同阶. 把 \eqref{eq:9-5}, \eqref{eq:9-6}, \eqref{eq:9-7} 与初始条件 $\mathbf{y}(t_0)=\mathbf{y}_0=\mathbf{z}$ 结合起来, \eqref{eq:9-1} 的离散方程组由如下块形式给出
\begin{equation}
  M\mathbf{y}\equiv(\widetilde{A}\otimes I_n-h\widetilde{B}\otimes J_n)\mathbf{y}
  =\mathbf{e}_1\otimes\mathbf{z}+h(\widetilde{B}\otimes I_n)\mathbf{g}.
  \label{eq:9-8}
\end{equation}
这里
\[
  \mathbf{e}_1=[1,0,\ldots,0]^T\in\R^{s+1},\qquad
  \mathbf{y}=[\mathbf{y}_0^T,\ldots,\mathbf{y}_s^T]^T\in\R^{(s+1)n},
\]
\[
  \mathbf{g}=[\mathbf{g}_0^T,\ldots,\mathbf{g}_s^T]^T\in\R^{(s+1)n},
\]
而 $\widetilde{A},\widetilde{B}\in\R^{(s+1)\times(s+1)}$ 由下式给出:
\[
  \widetilde{A}=\begin{bmatrix}
    1 & \cdots & 0 & & & & \\
    \alpha_0^{(1)} & \cdots & \alpha_{\mu}^{(1)} & & & & \\
    \vdots & \vdots & \vdots & & & & \\
    \alpha_0^{(\nu_1-1)} & \cdots & \alpha_{\mu}^{(\nu_1-1)} & & & & \\
    \alpha_0 & \cdots & \alpha_{\mu} & & & \mathbf{0} & \\
    & \alpha_0 & \cdots & \alpha_{\mu} & & & \\
    & & \ddots & \ddots & \ddots & & \\
    & & & \ddots & \ddots & \ddots & \\
    & & & & \alpha_0 & \cdots & \alpha_{\mu}\\
    & & & & \alpha_0^{(s-\nu_2+1)} & \cdots & \alpha_{\mu}^{(s-\nu_2+1)}\\
    & & \mathbf{0} & & \vdots & & \vdots\\
    & & & & \alpha_0^{(s)} & \cdots & \alpha_{\mu}^{(s)}
  \end{bmatrix},
\]
\[
  \widetilde{B}=\begin{bmatrix}
    0 & \cdots & 0 & & & & \\
    \beta_0^{(1)} & \cdots & \beta_{\mu}^{(1)} & & & & \\
    \vdots & \vdots & \vdots & & & & \\
    \beta_0^{(\nu_1-1)} & \cdots & \beta_{\mu}^{(\nu_1-1)} & & & & \\
    \beta_0 & \cdots & \beta_{\mu} & & & \mathbf{0} & \\
    & \beta_0 & \cdots & \beta_{\mu} & & & \\
    & & \ddots & \ddots & \ddots & & \\
    & & & \ddots & \ddots & \ddots & \\
    & & & & \beta_0 & \cdots & \beta_{\mu}\\
    & & & & \beta_0^{(s-\nu_2+1)} & \cdots & \beta_{\mu}^{(s-\nu_2+1)}\\
    & & \mathbf{0} & & \vdots & & \vdots\\
    & & & & \beta_0^{(s)} & \cdots & \beta_{\mu}^{(s)}
  \end{bmatrix},
\]
而 ``$\otimes$'' 是\newterm{tensorproduct}{张量积}{tensor product}.

我们回忆: $A=[a_{ij}]\in\R^{m\times n}$ 与 $B\in\R^{p\times q}$ 的张量积定义如下:
\[
  A\otimes B\equiv\begin{bmatrix}
    a_{11}B & a_{12}B & \cdots & a_{1n}B\\
    a_{21}B & a_{22}B & \cdots & a_{2n}B\\
    \vdots & \vdots & & \vdots\\
    a_{m1}B & a_{m2}B & \cdots & a_{mn}B
  \end{bmatrix},
\]
它是一个 $mp$ 行乘 $nq$ 列的矩阵. 张量积的基本性质可在 \cite{src:32,src:42} 中找到.

我们指出, 通常线性方程组 \eqref{eq:9-8} 是大型且稀疏的 (具有带状结构), 求解它是应用 BVMs 时的主要问题. 我们用第 6 章的 GMRES 方法求解 \eqref{eq:9-8}. 为加快 GMRES 迭代的收敛速度, 我们使用一个称为 Strang 型块循环预条件子的预条件子 $S$.

\section{常微分方程的 Strang 型预条件子}
\label{sec:9-3}

现在, 我们构造用于求解 \eqref{eq:9-8} 的 Strang 型块循环预条件子 $S$. 我们将说明 Strang 型预条件子的主要优点:

\begin{parenlist}
  \item 若使用 $A_{\nu_1,\nu_2}$-稳定的 BVM, 则 $S$ 可逆.
  \item 预处理后系统的谱是聚集的.
  \item 预处理 GMRES 方法每次迭代的运算量小于直接解法的运算量.
\end{parenlist}

\subsection{预条件子的构造}
\label{sec:9-3-1}

我们首先回忆 Strang 关于 Toeplitz 矩阵的\newterm{strangcirc}{循环预条件子}{circulant preconditioner}的定义. 给定 Toeplitz 矩阵
\[
  T_l=[t_{i-j}]_{i,j=1}^{l}=[t_q],
\]
Strang 预条件子 $s(T_l)$ 是一个循环矩阵, 其各对角线由下式给出:
\[
  [s(T_l)]_q=\begin{cases}
    t_q, & 0\leqslant q\leqslant\lfloor l/2\rfloor,\\
    t_{q-l}, & \lfloor l/2\rfloor<q<l,\\
    [s(T_l)]_{l+q}, & 0<-q<l.
  \end{cases}
\]
参见 \cite{src:16,src:18,src:42,src:43,src:77}.

忽略 $\widetilde{A}$ 与 $\widetilde{B}$ 的扰动, 我们提出下述预条件子 $S\in\R^{(s+1)n\times(s+1)n}$, 它称为 \eqref{eq:9-8} 中 $M$ 的 Strang 型块循环预条件子:
\begin{equation}
  S=s(A)\otimes I_n-hs(B)\otimes J_n,
  \label{eq:9-9}
\end{equation}
其中
\[
  s(A)=\begin{bmatrix}
    \alpha_{\nu_1} & \cdots & \alpha_{\mu} & & & \alpha_0 & \cdots & \alpha_{\nu_1-1}\\
    \vdots & \ddots & & \ddots & & & \ddots & \vdots\\
    \alpha_0 & & \ddots & & \ddots & & & \alpha_0\\
    & \ddots & & \ddots & & \mathbf{0} & & \\
    & & \ddots & & \ddots & & \ddots & \\
    \alpha_{\mu} & & \ddots & & \ddots & & & \alpha_{\mu}\\
    \vdots & \mathbf{0} & & \ddots & & \ddots & & \vdots\\
    \alpha_{\nu_1+1} & \cdots & \alpha_{\mu} & & & \alpha_0 & \cdots & \alpha_{\nu_1}
  \end{bmatrix},
\]
而 $s(B)$ 通过把 $s(A)$ 中的 $\{\alpha_i\}_{i=0}^{\mu}$ 换成 $\{\beta_i\}_{i=0}^{\mu}$ 类似地定义. 这里的 $\{\alpha_i\}_{i=0}^{\mu}$ 与 $\{\beta_i\}_{i=0}^{\mu}$ 是 \eqref{eq:9-5} 中给出的系数. 我们指出, $s(A)$, $s(B)$ 实际上只是 Toeplitz 矩阵 $A$, $B$ 的 Strang 循环预条件子, 其中 $A$, $B$ 由 \eqref{eq:9-4} 给出.

我们证明: 只要给定的 BVM 是 $A_{\nu_1,\nu_2}$-稳定的, 且 $J_n$ 的特征值属于
\[
  \C^-=\{q\in\C:\ \operatorname{Re}(q)<0\},
\]
则预条件子 $S$ 可逆. 其中 $\operatorname{Re}(\cdot)$ 表示复数的实部. BVM 的稳定性与两个 $\mu=\nu_1+\nu_2$ 次特征多项式密切相关, 它们定义如下:
\begin{equation}
  \rho(z)\equiv\sum_{j=-\nu_1}^{\nu_2}\alpha_{j+\nu_1}z^{j+\nu_1}
  \qquad\text{与}\qquad
  \sigma(z)\equiv\sum_{j=-\nu_1}^{\nu_2}\beta_{j+\nu_1}z^{j+\nu_1}.
  \label{eq:9-10}
\end{equation}
$A_{\nu_1,\nu_2}$-稳定多项式定义为
\begin{equation}
  \pi(z,q)\equiv\rho(z)-q\sigma(z),
  \label{eq:9-11}
\end{equation}
其中 $z,q\in\C$.

现在考虑方程 $\pi(z,q)=0$. 它定义了复 $z$ 平面与复 $q$ 平面之间的一个映射. 对每个使 $\pi(z,q)$ 有定义的根 $z\in\C$, \eqref{eq:9-11} 给出
\[
  q=q(z)=\frac{\rho(z)}{\sigma(z)}.
\]
令
\begin{equation}
  \Gamma\equiv\left\{q\in\C:\ q=\frac{\rho(e^{\ii\theta})}{\sigma(e^{\ii\theta})},\ 0\leqslant\theta<2\pi\right\}.
  \label{eq:9-12}
\end{equation}
$\Gamma$ 是对应于单位圆周上根的点集, 称为\newterm{boundarylocus}{边界轨迹}{boundary locus}. 我们有下述定义与引理 \cite{src:10}.

\begin{definition}[$A_{\nu_1,\nu_2}$-稳定区域]
\label{def:9-1}
考虑具有由 \eqref{eq:9-10} 定义的 $A_{\nu_1,\nu_2}$-稳定多项式 $\pi(z,q)$ 的 BVM. 区域
\[
  \mathcal{D}_{\nu_1,\nu_2}
  =\bigl\{q\in\C:\ \pi(z,q)\ \text{在}\ \abs{z}=1\ \text{内有}\ \nu_1\ \text{个零点,
  在}\ \abs{z}=1\ \text{外有}\ \nu_2\ \text{个零点}\bigr\}
\]
称为该 BVM 的 $A_{\nu_1,\nu_2}$-稳定区域. 此外, 若 $\C^-\subseteq\mathcal{D}_{\nu_1,\nu_2}$, 则称该 BVM 是 $A_{\nu_1,\nu_2}$-稳定的.
\end{definition}

\begin{lemma}
\label{lem:9-1}
若 BVM 是 $A_{\nu_1,\nu_2}$-稳定的, 且 $\Gamma$ 由 \eqref{eq:9-12} 定义, 则对所有 $q\in\Gamma$ 有 $\operatorname{Re}(q)\geqslant0$.
\end{lemma}

现在, 我们要证明在稳定性条件下预条件子 $S$ 可逆.

\begin{theorem}
\label{thm:9-1}
若用于 \eqref{eq:9-1} 的 BVM 是 $A_{\nu_1,\nu_2}$-稳定的, 且 $h\lambda_k(J_n)\in\mathcal{D}_{\nu_1,\nu_2}$, 其中 $\lambda_k(J_n)$ ($k=1,\ldots,n$) 是 $J_n$ 的特征值, 则由 \eqref{eq:9-9} 定义的预条件子 $S$ 可逆.
\end{theorem}

\begin{proof}
由于 $s(A)$ 与 $s(B)$ 是循环矩阵, 它们的特征值由
\[
  g_A(z)=\alpha_{\mu}z^{\nu_2}+\cdots+\alpha_{\nu_1}
  +\alpha_{\nu_1-1}\frac{1}{z}+\cdots+\alpha_0\frac{1}{z^{\nu_1}}
  =\frac{\rho(z)}{z^{\nu_1}}
\]
与
\[
  g_B(z)=\beta_{\mu}z^{\nu_2}+\cdots+\beta_{\nu_1}
  +\beta_{\nu_1-1}\frac{1}{z}+\cdots+\beta_0\frac{1}{z^{\nu_1}}
  =\frac{\sigma(z)}{z^{\nu_1}}
\]
给出, 取值点为 $\omega_j=e^{\frac{2\pi\ii j}{s+1}}$, 其中 $\ii\equiv\sqrt{-1}$, $j=0,\ldots,s$. 参见 \cite{src:18,src:29}. $S$ 的特征值 $\lambda_{jk}(S)$ 因而由
\[
  \lambda_{jk}(S)=g_A(\omega_j)-h\lambda_k(J_n)g_B(\omega_j),
  \qquad j=0,\ldots,s,\quad k=1,\ldots,n
\]
给出. 由于 BVM 是 $A_{\nu_1,\nu_2}$-稳定的, $\mu$ 次多项式
\[
  \pi[z,h\lambda_k(J_n)]=\rho(z)-h\lambda_k(J_n)\sigma(z)
\]
在单位圆 $\abs{z}=1$ 上没有根. 于是, 对所有 $k=1,\ldots,n$ 与任意 $\abs{z}=1$, 有
\[
  g_A(z)-h\lambda_k(J_n)g_B(z)
  =\frac{1}{z^{\nu_1}}\pi[z,h\lambda_k(J_n)]\neq0.
\]
由此推出
\[
  \lambda_{jk}(S)\neq0,\qquad j=0,\ldots,s,\quad k=1,\ldots,n.
\]
故 $S$ 可逆.
\end{proof}

特别地, 我们有

\begin{corollary}
\label{cor:9-1}
若 BVM 是 $A_{\nu_1,\nu_2}$-稳定的, 且 $\lambda_k(J_n)\in\C^-$, 则预条件子 $S$ 可逆.
\end{corollary}

\subsection{收敛率与运算量}
\label{sec:9-3-2}

关于收敛率, 我们有下述定理.

\begin{theorem}
\label{thm:9-2}
有
\[
  S^{-1}M=I_{n(s+1)}+L,
\]
其中
\[
  \rank(L)\leqslant2n\mu.
\]
\end{theorem}

\begin{proof}
令 $E=M-S$. 由 \eqref{eq:9-8} 与 \eqref{eq:9-9},
\[
  E=\bigl(\widetilde{A}-s(A)\bigr)\otimes I_n
  -h\bigl(\widetilde{B}-s(B)\bigr)\otimes J_n
  =L_A\otimes I_n-hL_B\otimes J_n.
\]
容易验证, $L_A$ 与 $L_B$ 是 $(s+1)$ 阶矩阵, 其非零元仅出现在以下四个角: 左上角一个 $\nu_1$ 行乘 $(\mu+1)$ 列的块; 右上角一个 $\nu_1$ 阶块; 右下角一个 $\nu_2$ 行乘 $(\mu+1)$ 列的块; 左下角一个 $\nu_2$ 阶块. 注意 $\mu=\nu_1+\nu_2$, 于是有
\[
  \rank(L_A)\leqslant\mu,\qquad \rank(L_B)\leqslant\mu.
\]
因此,
\[
  \rank(L_A\otimes I_n)=\rank(L_A)\cdot n\leqslant n\mu
\]
与
\[
  \rank(L_B\otimes J_n)=\rank(L_B)\cdot n\leqslant n\mu.
\]
故
\[
  S^{-1}M=I_{n(s+1)}+S^{-1}E=I_{n(s+1)}+L,
\]
其中 $L$ 的秩至多为 $2n\mu$.
\end{proof}

因此, 当把 GMRES 方法用于
\[
  S^{-1}M\mathbf{y}=S^{-1}\mathbf{b}
\]
时, 由定理 \ref{thm:6-13} 与 \ref{thm:9-2} 知, 在精确算术下该方法至多 $2n\mu+1$ 次迭代即收敛.

关于每次迭代的运算量, GMRES 方法每次迭代的主要工作是矩阵-向量乘法
\[
  S^{-1}M\mathbf{z}
  =(s(A)\otimes I_n-hs(B)\otimes J_n)^{-1}
   (\widetilde{A}\otimes I_n-h\widetilde{B}\otimes J_n)\mathbf{z}
\]
(见 \ref{sec:6-5} 节). 由于 $\widetilde{A}$, $\widetilde{B}$ 是带状矩阵, 且假定 $J_n$ 稀疏, 矩阵-向量乘法
\[
  M\mathbf{z}=(\widetilde{A}\otimes I_n-h\widetilde{B}\otimes J_n)\mathbf{z}
\]
可以很快完成.

现在我们计算 $S^{-1}(M\mathbf{z})$. 注意任何循环矩阵都可由 Fourier 矩阵 $F$ 对角化 (见 \ref{sec:6-4} 节与 \cite{src:29}). 由于 $s(A)$ 与 $s(B)$ 是循环矩阵, 由 \eqref{eq:6-19} 有分解
\[
  s(A)=F\Lambda_AF^*,\qquad s(B)=F\Lambda_BF^*,
\]
其中 $\Lambda_A$, $\Lambda_B$ 分别是含 $s(A)$, $s(B)$ 特征值的对角矩阵. 由此可得
\[
  S^{-1}(M\mathbf{z})
  =(F^*\otimes I_n)(\Lambda_A\otimes I_n-h\Lambda_B\otimes J_n)^{-1}
   (F\otimes I_n)(M\mathbf{z}).
\]
该乘积可用 FFT 与求解 $s+1$ 个 $n$ 阶线性方程组得到. 由于 $J_n$ 稀疏, 矩阵
\[
  \Lambda_A\otimes I_n-h\Lambda_B\otimes J_n
\]
也稀疏. 因此 $S^{-1}(M\mathbf{z})$ 可通过求解 $s+1$ 个 $n$ 阶稀疏线性方程组得到. 于是每次迭代的总运算量为
\[
  \gamma_1n(s+1)\log(s+1)+\gamma_2(s+1)nq,
\]
其中 $q$ 是 $J_n$ 的非零元个数, $\gamma_1$ 与 $\gamma_2$ 是某些正常数. 关于本方法的运算量与直接解法 (对线性方程组 \eqref{eq:9-8}) 的比较, 参见 \cite{src:19}.

\subsection{数值结果}
\label{sec:9-3-3}

现在我们通过求解 \cite{src:8} 中给出的一个测试问题来说明预条件子 $S$ 的效率. 实验在 MATLAB 中进行. 我们用 MATLAB 提供的 M 文件 ``gmres'' 求解预处理后的方程组. 我们应当强调, 本章所有试验中, 零向量都取作初始猜测, 停止准则为
\[
  \frac{\norm{\mathbf{r}_q}_2}{\norm{\mathbf{r}_0}_2}<10^{-6},
\]
其中 $\mathbf{r}_q$ 是 $q$ 次迭代后的残差. 我们使用的 BVM 是三阶\newterm{gam}{广义 Adams 方法}{generalized Adams method (GAM)}. 其公式与初始, 末端附加条件可在 \cite{src:10} 中找到.

\begin{example}[热传导方程]
\label{ex:9-1}
\[
  \left\{
  \begin{array}{ll}
    \dfrac{\partial u}{\partial t}=\dfrac{\partial^2u}{\partial x^2}, & \\[2ex]
    u(0,t)=\dfrac{\partial u}{\partial x}(\pi,t)=0, & t\in[0,2\pi],\\[1ex]
    u(x,0)=x, & x\in[0,\pi].
  \end{array}
  \right.
\]
\end{example}

我们用中心差分, 步长等于 $\pi/(n+1)$ 离散偏微分算子 $\partial^2/\partial x^2$. 所得的常微分方程组为
\[
  \left\{
  \begin{array}{ll}
    \mathbf{y}'(t)=T_n\mathbf{y}(t), & t\in[0,2\pi],\\
    \mathbf{y}(0)=[x_1,x_2,\ldots,x_n]^T,
  \end{array}
  \right.
\]
其中 $T_n$ 是缩放后的离散 Laplacian 矩阵
\[
  T_n=\frac{(n+1)^2}{\pi^2}\begin{bmatrix}
    -2 & 1 & & & \mathbf{0}\\
    1 & \ddots & \ddots & & \\
    & \ddots & \ddots & \ddots & \\
    & & \ddots & -2 & 1\\
    \mathbf{0} & & & 1 & -1
  \end{bmatrix}.
\]

表 \ref{tab:9-1} 列出了不同 $n$ 与 $s$ 时 GMRES 方法收敛所需的迭代次数. 表中 $I$ 表示不使用预条件子, $S$ 表示由 \eqref{eq:9-9} 定义的 Strang 型块循环预条件子. 我们看到, 使用 $S$ 时收敛所需的迭代次数远少于不使用预条件子时的迭代次数. $S$ 列下的数在 $s$ 与 $n$ 增大时几乎保持不变.

\begin{table}[htbp]
  \centering
  \caption{收敛所需的迭代次数}
  \label{tab:9-1}
  \begin{tabular}{cccc}
    \toprule
    $n$ & $s$ & $I$ & $S$\\ \midrule
    24 & 6  & 19  & 4\\
       & 12 & 70  & 4\\
       & 24 & 152 & 4\\
       & 48 & 227 & 3\\
       & 96 & 314 & 3\\ \bottomrule
  \end{tabular}
  \hspace{2em}
  \begin{tabular}{cccc}
    \toprule
    $n$ & $s$ & $I$ & $S$\\ \midrule
    48 & 6  & 47   & 4\\
       & 12 & 167  & 4\\
       & 24 & 359  & 4\\
       & 48 & $>\!400$ & 3\\
       & 96 & $>\!400$ & 3\\ \bottomrule
  \end{tabular}
\end{table}

\section{时滞微分方程的 Strang 型预条件子}
\label{sec:9-4}

现在, 我们研究时滞微分方程 (DDEs).

\subsection{多时滞微分方程}
\label{sec:9-4-1}

考虑多时滞微分方程的解:
\begin{equation}
  \left\{
  \begin{array}{ll}
    \mathbf{y}'(t)=J_n\mathbf{y}(t)+D_n^{(1)}\mathbf{y}(t-\tau_1)
      +\cdots+D_n^{(s)}\mathbf{y}(t-\tau_s)+\mathbf{f}(t), & t\geqslant t_0,\\
    \mathbf{y}(t)=\boldsymbol{\phi}(t), & t\leqslant t_0,
  \end{array}
  \right.
  \label{eq:9-13}
\end{equation}
其中 $\mathbf{y}(t),\mathbf{f}(t),\boldsymbol{\phi}(t):\R\to\R^n$; $J_n,D_n^{(1)},\ldots,D_n^{(s)}\in\R^{n\times n}$; 而 $\tau_1,\ldots,\tau_s>0$ 是某些有理数.

为求得合理的数值解, 我们要求 \eqref{eq:9-13} 的解是渐近稳定的. 我们有下述引理 \cite{src:63,src:64,src:91}.

\begin{lemma}
\label{lem:9-2}
对任意 $s\geqslant1$, 若
\[
  \eta(J_n)\equiv\frac{1}{2}\lambda_{\max}(J_n+J_n^T)<0
\]
且
\begin{equation}
  \eta(J_n)+\sum_{j=1}^{s}\norm{D_n^{(j)}}_2<0,
  \label{eq:9-14}
\end{equation}
则 \eqref{eq:9-13} 的解是渐近稳定的.
\end{lemma}

下面为简单起见, 我们只考虑 \eqref{eq:9-13} 中 $s=2$ 的情形; 对任意 $s$ 的推广是直接的. 令
\[
  h=\tau_1/m_1=\tau_2/m_2
\]
为步长, 其中 $m_1$ 与 $m_2$ 是正整数且 $m_2>m_1$ ($\tau_2>\tau_1$). 对 \eqref{eq:9-13}, 在区间 $[t_0,t_0+r_1h]$ 的均匀网格
\[
  t_j=t_0+jh,\qquad j=0,\ldots,r_1
\]
上使用带 $(\nu_1,\nu_2)$-边界条件的 BVM, 得
\begin{equation}
  \begin{aligned}
    \sum_{i=0}^{\mu}\alpha_i\mathbf{y}_{p+i-\nu_1}
    ={}&h\sum_{i=0}^{\mu}\beta_i\Bigl(J_n\mathbf{y}_{p+i-\nu_1}
      +D_n^{(1)}\mathbf{y}_{p+i-\nu_1-m_1}\\
    &\quad+D_n^{(2)}\mathbf{y}_{p+i-\nu_1-m_2}
      +\mathbf{f}_{p+i-\nu_1}\Bigr),
  \end{aligned}
  \label{eq:9-15}
\end{equation}
其中 $p=\nu_1,\ldots,r_1-1$, $\mu=\nu_1+\nu_2$. 通过提供值
\begin{equation}
  \mathbf{y}_{-m_2},\ldots,\mathbf{y}_{-m_1},\ldots,\mathbf{y}_0,\qquad
  \mathbf{y}_1,\ldots,\mathbf{y}_{\nu_1-1},\qquad
  \mathbf{y}_{r_1},\ldots,\mathbf{y}_{r_1+\nu_2-1},
  \label{eq:9-16}
\end{equation}
\eqref{eq:9-15} 可写成矩阵形式
\[
  R\mathbf{y}=\mathbf{b},
\]
其中
\begin{equation}
  \begin{aligned}
    R&\equiv A\otimes I_n-hB\otimes J_n-hC^{(1)}\otimes D_n^{(1)}
      -hC^{(2)}\otimes D_n^{(2)},\\
    \mathbf{y}&=[\mathbf{y}_{\nu_1}^T,\mathbf{y}_{\nu_1+1}^T,\ldots,
      \mathbf{y}_{r_1-1}^T]^T\in\R^{n(r_1-\nu_1)},
  \end{aligned}
  \label{eq:9-17}
\end{equation}
$\mathbf{b}\in\R^{n(r_1-\nu_1)}$ 依赖于 $\mathbf{f}$, 边界值以及方法的系数. 矩阵 $A$, $B\in\R^{(r_1-\nu_1)\times(r_1-\nu_1)}$ 按 \eqref{eq:9-4} 定义, 而 $C^{(1)}$, $C^{(2)}\in\R^{(r_1-\nu_1)\times(r_1-\nu_1)}$ 定义如下:
\[
  C^{(1)}=\begin{bmatrix}
    \mathbf{0} & & & & & \mathbf{0}\\
    \beta_{\mu} & \ddots & & & & \\
    \vdots & \ddots & \ddots & & & \\
    \beta_0 & \cdots & \beta_{\mu} & \ddots & & \\
    & \ddots & & \ddots & \beta_{\mu} & \\
    \mathbf{0} & & \beta_0 & \cdots & \beta_{\mu} & \mathbf{0}
  \end{bmatrix},\qquad
  C^{(2)}=\begin{bmatrix}
    \mathbf{0} & & & & & \mathbf{0}\\
    \beta_{\mu} & \ddots & & & & \\
    \vdots & \ddots & \ddots & & & \\
    \beta_0 & \cdots & \beta_{\mu} & \ddots & & \\
    & \ddots & & \ddots & \beta_{\mu} & \\
    \mathbf{0} & & \beta_0 & \cdots & \beta_{\mu} & \mathbf{0}
  \end{bmatrix}.
\]
我们指出, $C^{(1)}$ 的第一列由下式给出:
\[
  \Bigl[\,\underbrace{0,\ldots,0}_{m_1-\nu_2},\beta_{\mu},\ldots,\beta_0,
  \underbrace{0,\ldots,0}_{r_1-m_1-2\nu_1-1}\,\Bigr]^T,
\]
而 $C^{(2)}$ 的第一列由下式给出:
\[
  \Bigl[\,\underbrace{0,\ldots,0}_{m_2-\nu_2},\beta_{\mu},\ldots,\beta_0,
  \underbrace{0,\ldots,0}_{r_1-m_2-2\nu_1-1}\,\Bigr]^T.
\]

\subsection{预条件子的构造}
\label{sec:9-4-2}

\eqref{eq:9-17} 的 Strang 型块循环预条件子定义如下:
\begin{equation}
  \widetilde{S}\equiv s(A)\otimes I_n-hs(B)\otimes J_n
  -hs(C^{(1)})\otimes D_n^{(1)}-hs(C^{(2)})\otimes D_n^{(2)},
  \label{eq:9-18}
\end{equation}
其中 $s(E)$ 是 Toeplitz 矩阵 $E$ 的 Strang 循环预条件子, $E$ 分别取 $A$, $B$, $C^{(1)}$, $C^{(2)}$.

现在讨论由 \eqref{eq:9-18} 定义的 Strang 型预条件子 $\widetilde{S}$ 的可逆性. 由于任何循环矩阵都可由 Fourier 矩阵 $F$ 对角化, 由 \eqref{eq:6-19} 有
\[
  s(E)=F^*\Lambda_EF,
\]
其中 $\Lambda_E$ 是含 $s(E)$ 特征值的对角矩阵, $E$ 分别取 $A$, $B$, $C^{(1)}$, $C^{(2)}$. 因此得到
\[
  \widetilde{S}=(F^*\otimes I_n)\bigl(\Lambda_A\otimes I_n-h\Lambda_B\otimes J_n
  -h\Lambda_{C^{(1)}}\otimes D_n^{(1)}-h\Lambda_{C^{(2)}}\otimes D_n^{(2)}\bigr)
  (F\otimes I_n).
\]

注意
\[
  \Lambda_A\otimes I_n-h\Lambda_B\otimes J_n
  -h\Lambda_{C^{(1)}}\otimes D_n^{(1)}-h\Lambda_{C^{(2)}}\otimes D_n^{(2)}
\]
的第 $j$ 块由
\[
  \widetilde{S}_j=[\Lambda_A]_{jj}I_n-h[\Lambda_B]_{jj}J_n
  -h[\Lambda_{C^{(1)}}]_{jj}D_n^{(1)}-h[\Lambda_{C^{(2)}}]_{jj}D_n^{(2)}
\]
给出, $j=1,2,\ldots,r_1-\nu_1$. 令
\[
  w_j=e^{-\frac{2\pi\ii j}{r_1-\nu_1}}.
\]
我们有
\[
  [\Lambda_A]_{jj}=\frac{\rho(w_j)}{w_j^{\nu_1}},\qquad
  [\Lambda_B]_{jj}=\frac{\sigma(w_j)}{w_j^{\nu_1}},
\]
\[
  [\Lambda_{C^{(1)}}]_{jj}
  =\beta_{\mu}w_j^{-m_1+\nu_2}+\cdots+\beta_0w_j^{-m_1-\nu_1}
  =\frac{\sigma(w_j)}{w_j^{m_1+\nu_1}},
\]
以及
\[
  [\Lambda_{C^{(2)}}]_{jj}
  =\beta_{\mu}w_j^{-m_2+\nu_2}+\cdots+\beta_0w_j^{-m_2-\nu_1}
  =\frac{\sigma(w_j)}{w_j^{m_2+\nu_1}},
\]
其中 $\rho(z)$ 与 $\sigma(z)$ 如 \eqref{eq:9-10} 定义. 因此,
\[
  \widetilde{S}_j=\frac{1}{w_j^{m_2+\nu_1}}
  \Bigl[w_j^{m_2}\bigl(\rho(w_j)I_n-h\sigma(w_j)J_n
  -hw_j^{-m_1}\sigma(w_j)D_n^{(1)}\bigr)
  -h\sigma(w_j)D_n^{(2)}\Bigr].
\]
因此我们只需证明 $\widetilde{S}_j$ ($j=1,2,\ldots,r_1-\nu_1$) 可逆. 我们有下述定理.

\begin{theorem}
\label{thm:9-3}
若带 $(\nu_1,\nu_2)$-边界条件的 BVM 是 $A_{\nu_1,\nu_2}$-稳定的, 且 \eqref{eq:9-14} 成立, 则对任意 $\theta\in\R$, 矩阵
\[
  e^{\ii m_2\theta}\bigl(\rho(e^{\ii\theta})I_n-h\sigma(e^{\ii\theta})J_n
  -he^{-\ii m_1\theta}\sigma(e^{\ii\theta})D_n^{(1)}\bigr)
  -h\sigma(e^{\ii\theta})D_n^{(2)}
\]
可逆. 由此可知, 由 \eqref{eq:9-18} 定义的 Strang 型预条件子 $\widetilde{S}$ 也可逆.
\end{theorem}

\begin{proof}
假设存在 $\mathbf{x}\in\C^n$ 且 $\norm{\mathbf{x}}_2=1$ 以及 $\theta\in\R$ 使
\[
  \Bigl[e^{\ii m_2\theta}\bigl(\rho(e^{\ii\theta})I_n-h\sigma(e^{\ii\theta})J_n
  -he^{-\ii m_1\theta}\sigma(e^{\ii\theta})D_n^{(1)}\bigr)
  -h\sigma(e^{\ii\theta})D_n^{(2)}\Bigr]\mathbf{x}=\mathbf{0}.
\]
则
\[
  \mathbf{x}^*\Bigl[\rho(e^{\ii\theta})I_n-h\sigma(e^{\ii\theta})J_n
  -he^{-\ii m_1\theta}\sigma(e^{\ii\theta})D_n^{(1)}
  -he^{-\ii m_2\theta}\sigma(e^{\ii\theta})D_n^{(2)}\Bigr]\mathbf{x}=0,
\]
即
\[
  \rho(e^{\ii\theta})-h\sigma(e^{\ii\theta})\mathbf{x}^*J_n\mathbf{x}
  -he^{-\ii m_1\theta}\sigma(e^{\ii\theta})\mathbf{x}^*D_n^{(1)}\mathbf{x}
  -he^{-\ii m_2\theta}\sigma(e^{\ii\theta})\mathbf{x}^*D_n^{(2)}\mathbf{x}=0.
\]
于是有
\[
  \begin{aligned}
    &\rho(e^{\ii\theta})
    -\bigl(h\mathbf{x}^*J_n\mathbf{x}
      +he^{-\ii m_1\theta}\mathbf{x}^*D_n^{(1)}\mathbf{x}
      +he^{-\ii m_2\theta}\mathbf{x}^*D_n^{(2)}\mathbf{x}\bigr)\sigma(e^{\ii\theta})\\
    &\qquad=\pi\Bigl(e^{\ii\theta},
      h\bigl(\mathbf{x}^*J_n\mathbf{x}+e^{-\ii m_1\theta}\mathbf{x}^*D_n^{(1)}\mathbf{x}
      +e^{-\ii m_2\theta}\mathbf{x}^*D_n^{(2)}\mathbf{x}\bigr)\Bigr)=0,
  \end{aligned}
\]
其中 $\pi(z,q)$ 由 \eqref{eq:9-11} 给出. 因此,
\[
  h\bigl(\mathbf{x}^*J_n\mathbf{x}+e^{-\ii m_1\theta}\mathbf{x}^*D_n^{(1)}\mathbf{x}
  +e^{-\ii m_2\theta}\mathbf{x}^*D_n^{(2)}\mathbf{x}\bigr)\in\Gamma,
\]
其中 $\Gamma$ 是由 \eqref{eq:9-12} 定义的边界轨迹. 由于 BVM 是 $A_{\nu_1,\nu_2}$-稳定的, 由引理 \ref{lem:9-1} 知
\[
  \operatorname{Re}\bigl(\mathbf{x}^*J_n\mathbf{x}
  +e^{-\ii m_1\theta}\mathbf{x}^*D_n^{(1)}\mathbf{x}
  +e^{-\ii m_2\theta}\mathbf{x}^*D_n^{(2)}\mathbf{x}\bigr)\geqslant0.
\]
由 Cauchy-Schwarz 不等式, 有
\[
  \operatorname{Re}\bigl(e^{-\ii m_1\theta}\mathbf{x}^*D_n^{(1)}\mathbf{x}\bigr)
  \leqslant\abs{\mathbf{x}^*D_n^{(1)}\mathbf{x}}
  \leqslant\norm{\mathbf{x}}_2\norm{D_n^{(1)}\mathbf{x}}_2
  \leqslant\norm{D_n^{(1)}}_2\norm{\mathbf{x}}_2^2
  =\norm{D_n^{(1)}}_2,
\]
类似地有
\[
  \operatorname{Re}\bigl(e^{-\ii m_2\theta}\mathbf{x}^*D_n^{(2)}\mathbf{x}\bigr)
  \leqslant\norm{D_n^{(2)}}_2.
\]
注意
\[
  \eta(J_n)=\max_{\norm{\mathbf{x}}_2=1}\operatorname{Re}(\mathbf{x}^*J_n\mathbf{x})
  \geqslant\operatorname{Re}(\mathbf{x}^*J_n\mathbf{x}).
\]
于是有
\[
  \eta(J_n)+\norm{D_n^{(1)}}_2+\norm{D_n^{(2)}}_2\geqslant0,
\]
这与 \eqref{eq:9-14} 矛盾. 因此, 矩阵
\[
  e^{\ii m_2\theta}\bigl(\rho(e^{\ii\theta})I_n-h\sigma(e^{\ii\theta})J_n
  -he^{-\ii m_1\theta}\sigma(e^{\ii\theta})D_n^{(1)}\bigr)
  -h\sigma(e^{\ii\theta})D_n^{(2)}
\]
可逆, 从而 Strang 型预条件子 $\widetilde{S}$ 也可逆.
\end{proof}

\subsection{收敛率}
\label{sec:9-4-3}

现在, 我们讨论带 Strang 型块循环预条件子的预处理 GMRES 方法的收敛率. 关于预处理矩阵的谱, 我们有下述结果 \cite{src:44}.

\begin{theorem}
\label{thm:9-4}
设 $R$ 由 \eqref{eq:9-17} 给出, $\widetilde{S}$ 由 \eqref{eq:9-18} 给出. 则有
\[
  \widetilde{S}^{-1}R=I_{n(r_1-\nu_1)}+L,
\]
其中
\[
  \rank(L)\leqslant(2\mu+m_1+m_2+2\nu_1+2)n.
\]
\end{theorem}

由定理 \ref{thm:6-13} 与 \ref{thm:9-4}, 当把 GMRES 方法用于
\[
  \widetilde{S}^{-1}R\mathbf{y}=\widetilde{S}^{-1}\mathbf{b}
\]
时, 在精确算术下该方法至多 $(2\mu+m_1+m_2+2\nu_1+2)n+1$ 次迭代即收敛.

由定理 \ref{thm:9-4} 我们知道, 若步长 $h=\tau_1/m_1=\tau_2/m_2$ 固定, 则对预处理后方程组
\[
  \widetilde{S}^{-1}R\mathbf{y}=\widetilde{S}^{-1}\mathbf{b}
\]
应用 GMRES 方法时收敛所需的迭代次数与 $r_1$ 无关, 因而与我们考虑区间的长度无关. 我们应当强调, 9.4.4 小节的数值例子表明收敛速度比定理 \ref{thm:9-4} 所给估计预示的快得多. 关于本算法的运算量, 参见 \cite{src:44,src:58}.

\subsection{数值结果}
\label{sec:9-4-4}

我们通过求解下述例子来说明预条件子的效率. 我们使用的 BVM 在 $t\in[0,4]$ 上是三阶 GBDF.

\begin{example}
\label{ex:9-2}
考虑
\[
  \left\{
  \begin{array}{ll}
    \mathbf{y}'(t)=J_n\mathbf{y}(t)+D_n^{(1)}\mathbf{y}(t-0.5)
      +D_n^{(2)}\mathbf{y}(t-1), & t\geqslant0,\\
    \mathbf{y}(t)=[\sin t,1,\ldots,1]^T, & t\leqslant0,
  \end{array}
  \right.
\]
其中
\[
  J_n=\begin{bmatrix}
    -10 & 2 & & & \mathbf{0}\\
    2 & \ddots & \ddots & & \\
    1 & \ddots & \ddots & \ddots & \\
    & \ddots & \ddots & \ddots & 2\\
    \mathbf{0} & & 1 & 2 & -10
  \end{bmatrix},\qquad
  D_n^{(1)}=\frac{1}{n}\begin{bmatrix}
    2 & -1 & & & \mathbf{0}\\
    -1 & \ddots & \ddots & & \\
    & \ddots & \ddots & \ddots & \\
    & & \ddots & \ddots & -1\\
    \mathbf{0} & & & -1 & 2
  \end{bmatrix},
\]
而
\[
  D_n^{(2)}=\frac{1}{n}\begin{bmatrix}
    2 & 1 & & & \mathbf{0}\\
    1 & \ddots & \ddots & & \\
    & \ddots & \ddots & \ddots & \\
    & & \ddots & \ddots & 1\\
    \mathbf{0} & & & 1 & 2
  \end{bmatrix}.
\]
\end{example}

实际上, 我们并没有由 \eqref{eq:9-16} 提供的边界值
\[
  \mathbf{y}_1,\ldots,\mathbf{y}_{\nu_1-1},\qquad
  \mathbf{y}_{r_1},\ldots,\mathbf{y}_{r_1+\nu_2-1}.
\]
与 9.2 节的做法一样, 改为给出 $\nu_1-1$ 个初始附加方程与 $\nu_2$ 个末端附加方程. 我们指出, 引入附加方程之后, \eqref{eq:9-17} 中的矩阵 $A$, $B$, $C^{(1)}$ 与 $C^{(2)}$ 是带小秩扰动的 Toeplitz 矩阵. 忽略这些小秩扰动, 我们也可以构造 Strang 型预条件子 \eqref{eq:9-18}.

表 \ref{tab:9-2} 列出了不同矩阵规模 $n$ 与步长 $h$ 的组合下 GMRES 方法收敛所需的迭代次数. 表中 $I$ 表示不使用预条件子, $\widetilde{S}$ 表示由 \eqref{eq:9-18} 定义的 Strang 型块循环预条件子. 我们看到, 在 $\widetilde{S}$ 列下, 收敛所需的迭代次数随 $n$ 增大与 $h$ 减小而缓慢增加.

\begin{table}[htbp]
  \centering
  \caption{收敛所需的迭代次数}
  \label{tab:9-2}
  \begin{tabular}{cccc}
    \toprule
    $n$ & $h$ & $I$ & $\widetilde{S}$\\ \midrule
    24 & $1/10$ & 52  & 9\\
       & $1/20$ & 97  & 11\\
       & $1/40$ & 185 & 15\\
       & $1/80$ & 367 & 19\\ \bottomrule
  \end{tabular}
  \hspace{2em}
  \begin{tabular}{cccc}
    \toprule
    $n$ & $h$ & $I$ & $\widetilde{S}$\\ \midrule
    48 & $1/10$ & 53  & 12\\
       & $1/20$ & 98  & 14\\
       & $1/40$ & 189 & 14\\
       & $1/80$ & 378 & 17\\ \bottomrule
  \end{tabular}
\end{table}

\section{中立型时滞微分方程的 Strang 型预条件子}
\label{sec:9-5}

本节研究\newterm{ndde}{中立型时滞微分方程}{neutral delay differential equation (NDDE)}.

\subsection{中立型时滞微分方程}
\label{sec:9-5-1}

考虑 NDDE 的解:
\begin{equation}
  \left\{
  \begin{array}{ll}
    \mathbf{y}'(t)=L_n\mathbf{y}'(t-\tau)+M_n\mathbf{y}(t)
      +N_n\mathbf{y}(t-\tau), & t\geqslant t_0,\\
    \mathbf{y}(t)=\boldsymbol{\phi}(t), & t\leqslant t_0,
  \end{array}
  \right.
  \label{eq:9-19}
\end{equation}
其中 $\mathbf{y}(t),\boldsymbol{\phi}(t):\R\to\R^n$; $L_n,M_n,N_n\in\R^{n\times n}$; 而 $\tau>0$ 为常数.

与 9.4.1 小节一样, 我们要求 \eqref{eq:9-19} 的解渐近稳定. 我们有下述引理 \cite{src:41,src:54}.

\begin{lemma}
\label{lem:9-3}
设 $L_n$, $M_n$, $N_n$ 是任意满足 $\norm{L_n}_2<1$ 的矩阵. 则 \eqref{eq:9-19} 的解是渐近稳定的, 若 $\operatorname{Re}(\lambda_i)<0$, 其中 $\lambda_i$ ($i=1,2,\ldots,n$) 是矩阵
\[
  (I_n-\eta L_n)^{-1}(M_n+\eta N_n)
\]
的特征值, 且 $\abs{\eta}\leqslant1$.
\end{lemma}

令 $h=\tau/k_1$ 为步长, 其中 $k_1$ 是正整数. 对 \eqref{eq:9-19}, 在均匀网格
\[
  t_j=t_0+jh,\qquad j=0,\ldots,r_2
\]
上使用带 $(\nu_1,\nu_2)$-边界条件的 BVM, 得
\begin{equation}
  \begin{aligned}
    \sum_{i=0}^{\mu}\alpha_i\mathbf{y}_{p+i-\nu_1}
    ={}&\sum_{i=0}^{\mu}\alpha_iL_n\mathbf{y}_{p+i-\nu_1-k_1}\\
    &+h\sum_{i=0}^{\mu}\beta_i\bigl(M_n\mathbf{y}_{p+i-\nu_1}
      +N_n\mathbf{y}_{p+i-\nu_1-k_1}\bigr),
  \end{aligned}
  \label{eq:9-20}
\end{equation}
其中 $p=\nu_1,\ldots,r_2-1$, $\mu=\nu_1+\nu_2$. 通过提供值
\begin{equation}
  \mathbf{y}_{-k_1},\ldots,\mathbf{y}_0,\qquad
  \mathbf{y}_1,\ldots,\mathbf{y}_{\nu_1-1},\qquad
  \mathbf{y}_{r_2},\ldots,\mathbf{y}_{r_2+\nu_2-1},
  \label{eq:9-21}
\end{equation}
\eqref{eq:9-20} 可写成矩阵形式
\[
  H\mathbf{y}=\mathbf{b},
\]
其中
\begin{equation}
  \begin{aligned}
    H&\equiv A\otimes I_n-A^{(1)}\otimes L_n-hB\otimes M_n
      -hB^{(1)}\otimes N_n,\\
    \mathbf{y}&=[\mathbf{y}_{\nu_1}^T,\mathbf{y}_{\nu_1+1}^T,\ldots,
      \mathbf{y}_{r_2-1}^T]^T\in\R^{n(r_2-\nu_1)},
  \end{aligned}
  \label{eq:9-22}
\end{equation}
$\mathbf{b}\in\R^{n(r_2-\nu_1)}$ 依赖于边界值与方法的系数. 在 \eqref{eq:9-22} 中, 矩阵 $A$, $B\in\R^{(r_2-\nu_1)\times(r_2-\nu_1)}$ 按 \eqref{eq:9-4} 定义, 而 $A^{(1)}$, $B^{(1)}\in\R^{(r_2-\nu_1)\times(r_2-\nu_1)}$ 定义如下 \cite{src:5}:
\[
  A^{(1)}=\begin{bmatrix}
    \mathbf{0} & & & & & \mathbf{0}\\
    \alpha_{\mu} & \ddots & & & & \\
    \vdots & \ddots & \ddots & & & \\
    \alpha_0 & \cdots & \alpha_{\mu} & \ddots & & \\
    & \ddots & & \ddots & \alpha_{\mu} & \\
    \mathbf{0} & & \alpha_0 & \cdots & \alpha_{\mu} & \mathbf{0}
  \end{bmatrix},\qquad
  B^{(1)}=\begin{bmatrix}
    \mathbf{0} & & & & & \mathbf{0}\\
    \beta_{\mu} & \ddots & & & & \\
    \vdots & \ddots & \ddots & & & \\
    \beta_0 & \cdots & \beta_{\mu} & \ddots & & \\
    & \ddots & & \ddots & \beta_{\mu} & \\
    \mathbf{0} & & \beta_0 & \cdots & \beta_{\mu} & \mathbf{0}
  \end{bmatrix}.
\]
我们指出, $A^{(1)}$ 的第一列由下式给出:
\[
  \Bigl[\,\underbrace{0,\ldots,0}_{k_1-\nu_2},\alpha_{\mu},\ldots,\alpha_0,
  \underbrace{0,\ldots,0}_{r_2-k_1-2\nu_1-1}\,\Bigr]^T,
\]
而 $B^{(1)}$ 的第一列由下式给出:
\[
  \Bigl[\,\underbrace{0,\ldots,0}_{k_1-\nu_2},\beta_{\mu},\ldots,\beta_0,
  \underbrace{0,\ldots,0}_{r_2-k_1-2\nu_1-1}\,\Bigr]^T.
\]

\subsection{预条件子的构造}
\label{sec:9-5-2}

\eqref{eq:9-22} 的 Strang 型块循环预条件子定义如下:
\begin{equation}
  \widetilde{S}\equiv s(A)\otimes I_n-s(A^{(1)})\otimes L_n
  -hs(B)\otimes M_n-hs(B^{(1)})\otimes N_n,
  \label{eq:9-23}
\end{equation}
其中 $s(E)$ 是 Toeplitz 矩阵 $E$ 的 Strang 循环预条件子, $E$ 分别取 $A$, $B$, $A^{(1)}$, $B^{(1)}$. 由 \eqref{eq:6-19} 再次得到
\[
  \widetilde{S}=(F^*\otimes I_n)\bigl(\Lambda_A\otimes I_n-\Lambda_{A^{(1)}}\otimes L_n
  -h\Lambda_B\otimes M_n-h\Lambda_{B^{(1)}}\otimes N_n\bigr)(F\otimes I_n),
\]
其中 $\Lambda_E$ 是含 $s(E)$ 特征值的对角矩阵, $E$ 分别取 $A$, $B$, $A^{(1)}$, $B^{(1)}$.

现在讨论 Strang 型预条件子 $\widetilde{S}$ 的可逆性. 令
\[
  w_j=e^{-\frac{2\pi\ii j}{r_2-\nu_1}}.
\]
我们有
\[
  [\Lambda_A]_{jj}=\frac{\rho(w_j)}{w_j^{\nu_1}},\qquad
  [\Lambda_B]_{jj}=\frac{\sigma(w_j)}{w_j^{\nu_1}},
\]
\[
  [\Lambda_{A^{(1)}}]_{jj}
  =\alpha_{\mu}w_j^{-k_1+\nu_2}+\cdots+\alpha_0w_j^{-k_1-\nu_1}
  =\frac{\rho(w_j)}{w_j^{k_1+\nu_1}},
\]
以及
\[
  [\Lambda_{B^{(1)}}]_{jj}
  =\beta_{\mu}w_j^{-k_1+\nu_2}+\cdots+\beta_0w_j^{-k_1-\nu_1}
  =\frac{\sigma(w_j)}{w_j^{k_1+\nu_1}},
\]
其中 $\rho(z)$ 与 $\sigma(z)$ 如 \eqref{eq:9-10} 定义. 因此, $\widetilde{S}$ 中
\[
  \Lambda_A\otimes I_n-\Lambda_{A^{(1)}}\otimes L_n
  -h\Lambda_B\otimes M_n-h\Lambda_{B^{(1)}}\otimes N_n
\]
的第 $j$ 块由
\[
  \begin{aligned}
    \widetilde{S}_j
    &=[\Lambda_A]_{jj}I_n-[\Lambda_{A^{(1)}}]_{jj}L_n
      -h[\Lambda_B]_{jj}M_n-h[\Lambda_{B^{(1)}}]_{jj}N_n\\
    &=\frac{1}{w_j^{k_1+\nu_1}}
      \Bigl[w_j^{k_1}\bigl(\rho(w_j)I_n-h\sigma(w_j)M_n\bigr)
      -\rho(w_j)L_n-h\sigma(w_j)N_n\Bigr]
  \end{aligned}
\]
给出, $j=1,2,\ldots,r_2-\nu_1$. 为证明 $\widetilde{S}$ 可逆, 我们只需证明 $\widetilde{S}_j$ ($j=1,2,\ldots,r_2-\nu_1$) 可逆. 令
\[
  \begin{aligned}
    \Delta&\equiv e^{\ii k_1\theta}
      \bigl(\rho(e^{\ii\theta})I_n-h\sigma(e^{\ii\theta})M_n\bigr)
      -\rho(e^{\ii\theta})L_n-h\sigma(e^{\ii\theta})N_n\\
    &=e^{\ii k_1\theta}\bigl(I_n-e^{-\ii k_1\theta}L_n\bigr)D,
  \end{aligned}
\]
其中
\begin{equation}
  D\equiv\rho(e^{\ii\theta})I_n
  -h\bigl(I_n-e^{-\ii k_1\theta}L_n\bigr)^{-1}
  \bigl(M_n+e^{-\ii k_1\theta}N_n\bigr)\sigma(e^{\ii\theta}).
  \label{eq:9-24}
\end{equation}
因此, 为证明 $\widetilde{S}_j$ 可逆, 只需证明对任意 $\theta\in\R$, $\Delta$ 可逆. 假设 $\norm{L_n}_2<1$, 则 $I_n-e^{-\ii k_1\theta}L_n$ 对任意 $\theta\in\R$ 非奇异. 因此, 我们只需证明 $D$ 对任意 $\theta\in\R$ 可逆. 我们有下述定理 \cite{src:5}.

\begin{theorem}
\label{thm:9-5}
若带 $(\nu_1,\nu_2)$-边界条件的 BVM 是 $A_{\nu_1,\nu_2}$-稳定的, 且 $\operatorname{Re}(\lambda_i)<0$, 其中 $\lambda_i$ ($i=1,2,\ldots,n$) 是矩阵 $(I_n-\eta L_n)^{-1}(M_n+\eta N_n)$ 的特征值, $\abs{\eta}\leqslant1$, 则对任意 $\theta\in\R$, 由 \eqref{eq:9-24} 定义的矩阵 $D$ 可逆. 由此可知, 由 \eqref{eq:9-23} 定义的 Strang 型预条件子 $\widetilde{S}$ 也可逆.
\end{theorem}

\begin{proof}
令
\[
  U\equiv(I_n-e^{-\ii m\theta}L_n)^{-1}(M_n+e^{-\ii m\theta}N_n).
\]
则 $D$ 可写成
\[
  D=\rho(z)I_n-hU\sigma(z).
\]
注意 $D$ 的特征值由
\[
  \lambda_i(D)=\rho(z)-h\lambda_i(U)\sigma(z),\qquad i=1,\ldots,n
\]
给出, 其中 $\lambda_i(U)$ ($i=1,2,\ldots,n$) 表示 $U$ 的特征值. 由于已知
\[
  \operatorname{Re}[\lambda_i(U)]<0,\qquad i=1,\ldots,n,
\]
故 $h\lambda_i(U)\in\C^-$. 注意 BVM 是 $A_{\nu_1,\nu_2}$-稳定的, 于是有
\[
  h\lambda_i(U)\in\C^-\subseteq\mathcal{D}_{\nu_1,\nu_2}.
\]
因此, 由 \eqref{eq:9-11} 定义的 $A_{\nu_1,\nu_2}$-稳定多项式
\[
  \pi[z,h\lambda_i(U)]\equiv\rho(z)-h\lambda_i(U)\sigma(z)
\]
在单位圆 $\abs{z}=1$ 上没有根. 于是, 对任意 $\abs{z}=1$, 有
\[
  \lambda_i(D)=\rho(z)-h\lambda_i(U)\sigma(z)=\pi[z,h\lambda_i(U)]\neq0,
  \qquad i=1,\ldots,n.
\]
由此可知 $D$ 可逆. 因此, 由 \eqref{eq:9-23} 定义的 Strang 型预条件子 $\widetilde{S}$ 也可逆.
\end{proof}

\subsection{收敛率}
\label{sec:9-5-3}

关于预处理矩阵的谱, 我们有下述结果 \cite{src:5}.

\begin{theorem}
\label{thm:9-6}
设 $H$ 由 \eqref{eq:9-22} 给出, $\widetilde{S}$ 由 \eqref{eq:9-23} 给出. 则有
\[
  \widetilde{S}^{-1}H=I_{n(r_2-\nu_1)}+L,
\]
其中
\[
  \rank(L)\leqslant2(\mu+k_1+\nu_1+1)n.
\]
\end{theorem}

由定理 \ref{thm:6-13} 与 \ref{thm:9-6}, 当把 GMRES 方法用于
\[
  \widetilde{S}^{-1}H\mathbf{y}=\widetilde{S}^{-1}\mathbf{b}
\]
时, 在精确算术下该方法至多 $2(\mu+k_1+\nu_1+1)n+1$ 次迭代即收敛.

由定理 \ref{thm:9-6} 我们观察到, 若步长 $h=\tau/k_1$ 固定, 则对
\[
  \widetilde{S}^{-1}H\mathbf{y}=\widetilde{S}^{-1}\mathbf{b}
\]
应用 GMRES 方法时收敛所需的迭代次数与 $r_2$ 无关, 即与我们考虑区间的长度无关.

\subsection{数值结果}
\label{sec:9-5-4}

我们通过求解下述例子来说明预条件子的效率.

\begin{example}
\label{ex:9-3}
考虑
\[
  \left\{
  \begin{array}{ll}
    \mathbf{y}'(t)=L_n\mathbf{y}'(t-1)+M_n\mathbf{y}(t)+N_n\mathbf{y}(t-1),
      & t\geqslant0,\\
    \mathbf{y}(t)=[1,1,\ldots,1]^T, & t\leqslant0,
  \end{array}
  \right.
\]
其中
\[
  L_n=\frac{1}{n}\begin{bmatrix}
    2 & 1 & & & \mathbf{0}\\
    1 & \ddots & \ddots & & \\
    & \ddots & \ddots & \ddots & \\
    & & \ddots & \ddots & 1\\
    \mathbf{0} & & & 1 & 2
  \end{bmatrix},\qquad
  M_n=\begin{bmatrix}
    -8 & 2 & 1 & & \mathbf{0}\\
    2 & \ddots & \ddots & \ddots & \\
    1 & \ddots & \ddots & \ddots & 1\\
    & \ddots & \ddots & \ddots & 2\\
    \mathbf{0} & & 1 & 2 & -8
  \end{bmatrix},
\]
而
\[
  N_n=\frac{1}{n}\begin{bmatrix}
    2 & -1 & & & \mathbf{0}\\
    -1 & \ddots & \ddots & & \\
    & \ddots & \ddots & \ddots & \\
    & & \ddots & \ddots & -1\\
    \mathbf{0} & & & -1 & 2
  \end{bmatrix}.
\]
\end{example}

例 \ref{ex:9-3} 用五阶 GAM 在 $t\in[0,4]$ 上求解. 实际上, 我们并没有由 \eqref{eq:9-21} 提供的边界值
\[
  \mathbf{y}_1,\ldots,\mathbf{y}_{\nu_1-1},\qquad
  \mathbf{y}_{r_2},\ldots,\mathbf{y}_{r_2+\nu_2-1}.
\]
与 9.2 节一样, 改为给出 $\nu_1-1$ 个初始附加方程与 $\nu_2$ 个末端附加方程. 引入附加方程之后, \eqref{eq:9-22} 中的矩阵 $A$, $A^{(1)}$, $B$ 与 $B^{(1)}$ 是带小秩扰动的 Toeplitz 矩阵. 同样可以忽略这些小秩扰动来构造 Strang 型预条件子 \eqref{eq:9-23}.

表 \ref{tab:9-3} 列出了不同 $n$ 与 $k_1$ 时 GMRES 方法收敛所需的迭代次数. 表中 $I$ 表示不使用预条件子, $\widetilde{S}$ 表示由 \eqref{eq:9-23} 定义的 Strang 型块循环预条件子. 我们看到, 使用 $\widetilde{S}$ 时收敛所需的迭代次数远少于不使用预条件子时的迭代次数. 我们应当强调, 我们的数值例子表明收敛速度比定理 \ref{thm:9-6} 所给估计预示的快得多.

\begin{table}[htbp]
  \centering
  \caption{收敛所需的迭代次数 (``$*$'' 表示内存不足)}
  \label{tab:9-3}
  \begin{tabular}{cccc}
    \toprule
    $n$ & $k_1$ & $I$ & $\widetilde{S}$\\ \midrule
    24 & 10 & 43  & 7\\
       & 20 & 83  & 7\\
       & 40 & 161 & 7\\
       & 80 & $*$ & 7\\ \bottomrule
  \end{tabular}
  \hspace{2em}
  \begin{tabular}{cccc}
    \toprule
    $n$ & $k_1$ & $I$ & $\widetilde{S}$\\ \midrule
    48 & 10 & 44  & 6\\
       & 20 & 83  & 6\\
       & 40 & 163 & 6\\
       & 80 & $*$ & 6\\ \bottomrule
  \end{tabular}
\end{table}

\section{奇摄动时滞微分方程的 Strang 型预条件子}
\label{sec:9-6}

本节研究\newterm{spdde}{奇摄动时滞微分方程}{singular perturbation delay differential equation (SPDDE)}的求解.

\subsection{奇摄动时滞微分方程}
\label{sec:9-6-1}

考虑 SPDDE 的解:
\begin{equation}
  \left\{
  \begin{array}{ll}
    \mathbf{x}'(t)=V^{(1)}\mathbf{x}(t)+V^{(2)}\mathbf{x}(t-\tau)
      +C^{(1)}\mathbf{y}(t)+C^{(2)}\mathbf{y}(t-\tau), & t\geqslant t_0,\\
    \epsilon\mathbf{y}'(t)=F^{(1)}\mathbf{x}(t)+F^{(2)}\mathbf{x}(t-\tau)
      +G^{(1)}\mathbf{y}(t)+G^{(2)}\mathbf{y}(t-\tau), & t\geqslant t_0,\\
    \mathbf{x}(t)=\boldsymbol{\phi}(t), & t\leqslant t_0,\\
    \mathbf{y}(t)=\boldsymbol{\psi}(t), & t\leqslant t_0,
  \end{array}
  \right.
  \label{eq:9-25}
\end{equation}
其中
\[
  \mathbf{x}(t),\boldsymbol{\phi}(t):\R\to\R^m;\qquad
  \mathbf{y}(t),\boldsymbol{\psi}(t):\R\to\R^n;
\]
\[
  V^{(1)},V^{(2)}\in\R^{m\times m};\qquad
  C^{(1)},C^{(2)}\in\R^{m\times n};
\]
\[
  F^{(1)},F^{(2)}\in\R^{n\times m};\qquad
  G^{(1)},G^{(2)}\in\R^{n\times n};
\]
而 $\tau>0$, $0<\epsilon\ll1$ 为常数. 我们可以把 SPDDE \eqref{eq:9-25} 改写为下述初值问题 \cite{src:45}:
\begin{equation}
  \left\{
  \begin{array}{ll}
    \mathbf{z}'(t)=P\mathbf{z}(t)+Q\mathbf{z}(t-\tau), & t\geqslant t_0,\\
    \mathbf{z}(t)=\left[\begin{array}{c}
      \boldsymbol{\phi}(t)\\
      \boldsymbol{\psi}(t)
    \end{array}\right], & t\leqslant t_0,
  \end{array}
  \right.
  \label{eq:9-26}
\end{equation}
其中 $P$ 与 $Q\in\R^{(m+n)\times(m+n)}$ 定义如下,
\[
  P=\begin{bmatrix}
    V^{(1)} & C^{(1)}\\
    \epsilon^{-1}F^{(1)} & \epsilon^{-1}G^{(1)}
  \end{bmatrix},\qquad
  Q=\begin{bmatrix}
    V^{(2)} & C^{(2)}\\
    \epsilon^{-1}F^{(2)} & \epsilon^{-1}G^{(2)}
  \end{bmatrix}.
\]

令
\[
  h=\tau/k_2
\]
为步长, 其中 $k_2$ 是正整数. 对 \eqref{eq:9-26}, 在区间 $[t_0,t_0+vh]$ 的均匀网格
\[
  t_j=t_0+jh,\qquad j=0,\ldots,v
\]
上使用带 $(\nu_1,\nu_2)$-边界条件的 BVM, 得
\begin{equation}
  \sum_{i=0}^{\mu}\alpha_i\mathbf{z}_{p+i-\nu_1}
  =h\sum_{i=0}^{\mu}\beta_i\bigl(P\mathbf{z}_{p+i-\nu_1}
  +Q\mathbf{z}_{p+i-\nu_1-k_2}\bigr),
  \label{eq:9-27}
\end{equation}
其中 $p=\nu_1,\ldots,v-1$, $\mu=\nu_1+\nu_2$. 通过提供值
\begin{equation}
  \mathbf{z}_{-k_2},\ldots,\mathbf{z}_0,\qquad
  \mathbf{z}_1,\ldots,\mathbf{z}_{\nu_1-1},\qquad
  \mathbf{z}_{v},\ldots,\mathbf{z}_{v+\nu_2-1},
  \label{eq:9-28}
\end{equation}
\eqref{eq:9-27} 可写成矩阵形式
\begin{equation}
  K\mathbf{v}=\mathbf{b},
  \label{eq:9-29}
\end{equation}
其中
\begin{equation}
  K\equiv A\otimes I_{m+n}-hB\otimes P-hU\otimes Q.
  \label{eq:9-30}
\end{equation}
\eqref{eq:9-29} 中的向量 $\mathbf{v}$ 由
\[
  \mathbf{v}=[\mathbf{z}_{\nu_1}^T,\mathbf{z}_{\nu_1+1}^T,\ldots,
  \mathbf{z}_{v-1}^T]^T\in\R^{(m+n)(v-\nu_1)}
\]
定义.

\eqref{eq:9-29} 的右端 $\mathbf{b}\in\R^{(m+n)(v-\nu_1)}$ 依赖于边界值与方法的系数. 在 \eqref{eq:9-30} 中, 矩阵 $A$, $B\in\R^{(v-\nu_1)\times(v-\nu_1)}$ 按 \eqref{eq:9-4} 定义, 而 $U\in\R^{(v-\nu_1)\times(v-\nu_1)}$ 定义为 \eqref{eq:9-17} 中的矩阵 $C^{(1)}$.

\subsection{预条件子的构造}
\label{sec:9-6-2}

可以为求解 \eqref{eq:9-29} 构造 Strang 型块循环预条件子:
\begin{equation}
  \widehat{S}\equiv s(A)\otimes I_{m+n}-hs(B)\otimes P-hs(U)\otimes Q,
  \label{eq:9-31}
\end{equation}
其中 $s(E)$ 是矩阵 $E$ 的 Strang 循环预条件子, $E$ 分别取 $A$, $B$, $U$. 关于我们的预条件子的可逆性以及方法的收敛率, 我们有下述定理 \cite{src:45}.

\begin{theorem}
\label{thm:9-7}
若带 $(\nu_1,\nu_2)$-边界条件的 BVM 是 $A_{\nu_1,\nu_2}$-稳定的,
\[
  \eta(P)\equiv\frac{1}{2}\lambda_{\max}(P+P^T)<0,\qquad
  \eta(P)+\norm{Q}_2<0,
\]
则由 \eqref{eq:9-31} 定义的 Strang 型预条件子 $\widehat{S}$ 可逆. 此外, 当把 GMRES 方法用于
\[
  \widehat{S}^{-1}K\mathbf{v}=\widehat{S}^{-1}\mathbf{b}
\]
时, 在精确算术下该方法至多 $\mathcal{O}(m+n)$ 次迭代即收敛.
\end{theorem}

\subsection{数值结果}
\label{sec:9-6-3}

我们测试下述 SPDDE.

\begin{example}
\label{ex:9-4}
考虑
\[
  \left\{
  \begin{array}{ll}
    \mathbf{x}'(t)=A\mathbf{x}(t)+B\mathbf{x}(t-1)+C\mathbf{y}(t)
      +D\mathbf{y}(t-1), & t\geqslant0,\\
    \epsilon\mathbf{y}'(t)=E\mathbf{x}(t)+F\mathbf{x}(t-1)+G\mathbf{y}(t)
      +H\mathbf{y}(t-1), & t\geqslant0,\ \epsilon=0.001,\\
    \mathbf{x}(t)=[1,1,\ldots,1]^T, & t\leqslant0,\\
    \mathbf{y}(t)=[2,2,\ldots,2]^T, & t\leqslant0,
  \end{array}
  \right.
\]
其中
\[
  A=\begin{bmatrix}
    -10 & 2 & 1 & & \mathbf{0}\\
    2 & \ddots & \ddots & \ddots & \\
    1 & \ddots & \ddots & \ddots & 1\\
    & \ddots & \ddots & \ddots & 2\\
    \mathbf{0} & & 1 & 2 & -10
  \end{bmatrix}_{m\times m},\qquad
  B=\begin{bmatrix}
    -2 & 1 & & & \mathbf{0}\\
    1 & \ddots & \ddots & & \\
    & \ddots & \ddots & \ddots & \\
    & & \ddots & \ddots & 1\\
    \mathbf{0} & & & 1 & -2
  \end{bmatrix}_{m\times m},
\]
\[
  C=\begin{bmatrix}
    1 & & \mathbf{0}\\
    & \ddots & \\
    \mathbf{0} & & 1
  \end{bmatrix}_{m\times n},\qquad
  D=3C,\qquad
  E=\begin{bmatrix}
    2 & -1 & & & \mathbf{0}\\
    1 & \ddots & \ddots & & \\
    & \ddots & \ddots & \ddots & \\
    & & \ddots & \ddots & -1\\
    \mathbf{0} & & 1 & 2 & -1
  \end{bmatrix}_{n\times m},
\]
\[
  F=\begin{bmatrix}
    1 & -1 & & & \mathbf{0}\\
    & \ddots & \ddots & & \\
    & & \ddots & \ddots & \\
    & & & \ddots & -1\\
    \mathbf{0} & & & 1 & -1
  \end{bmatrix}_{n\times m},\qquad
  G=\begin{bmatrix}
    -2 & & & & \mathbf{0}\\
    1 & \ddots & & & \\
    & \ddots & \ddots & & \\
    & & \ddots & \ddots & \\
    \mathbf{0} & & & 1 & -2
  \end{bmatrix}_{n\times n},
\]
以及
\[
  H=\begin{bmatrix}
    -5 & -1 & & & \mathbf{0}\\
    2 & \ddots & \ddots & & \\
    & \ddots & \ddots & \ddots & \\
    & & \ddots & \ddots & -1\\
    \mathbf{0} & & & 2 & -5
  \end{bmatrix}_{n\times n}.
\]
\end{example}

例 \ref{ex:9-4} 用三阶 GAM 在 $t\in[0,4]$ 上求解. 表 \ref{tab:9-4} 列出了不同 $m$, $n$ 与 $k_2$ 时 GMRES 方法收敛所需的迭代次数. 表中 $\widehat{S}$ 表示由 \eqref{eq:9-31} 定义的 Strang 型块循环预条件子.

\begin{table}[htbp]
  \centering
  \caption{收敛所需的迭代次数}
  \label{tab:9-4}
  \begin{tabular}{ccccc}
    \toprule
    $k_2$ & $m$ & $n$ & $I$ & $\widehat{S}$\\ \midrule
    24 & 8  & 2 & 55  & 29\\
       & 16 & 4 & 83  & 57\\
       & 32 & 8 & 109 & 90\\ \bottomrule
  \end{tabular}
  \hspace{2em}
  \begin{tabular}{ccccc}
    \toprule
    $k_2$ & $m$ & $n$ & $I$ & $\widehat{S}$\\ \midrule
    48 & 8  & 2 & 100 & 34\\
       & 16 & 4 & 131 & 63\\
       & 32 & 8 & 177 & 90\\ \bottomrule
  \end{tabular}
\end{table}

% chapters/ch10.tex —— 第 10 章：在分数阶扩散方程中的应用
% 源书：Chapter 10 Applications in FDEs（正文首段明确 FDEs = fractional diffusion equations）
%       书页 170–180（物理页 184–194）
% 公式 (10.1)–(10.5)；引理 10.1；算法 10.1；例 10.1；表 10.1；本章无习题
\chapter{在分数阶扩散方程中的应用}
\label{ch:fde}

本章简要研究用于求解二维空间分数阶扩散方程的 Crank--Nicolson 加权与移位 Grünwald 差分 (CN-WSGD) 格式. 该方法需要求解一个或多个非对称, 大型且满的线性方程组. 我们使用第 6 章介绍的 GMRES 方法, 并结合重启策略与预处理技术来求解这些方程组. 主要结果之一是: 对预处理后的方程组施用重启 GMRES 方法时, 每个时间层的计算量为 $\mathcal{O}(N^{2}\log N)$, 其中 $N=N_{1}=N_{2}$ 为 $x$ 方向与 $y$ 方向的空间剖分数 (见 \ref{sec:10-2} 节). 关于分数阶扩散方程 (FDE) 的文献, 可参阅 \cite{src:7,src:31,src:56,src:60,src:68}.

\section{引言}
\label{sec:10-1}

考虑如下二维空间分数阶扩散方程的初边值问题 \cite{src:62,src:89}:
\begin{equation}
\left\{
\begin{array}{l}
\displaystyle \frac{\partial u(x,y,t)}{\partial t}
 =d_{+}(x,y,t)\,{}_{a}D_{x}^{\alpha}u(x,y,t)+d_{-}(x,y,t)\,{}_{x}D_{b}^{\alpha}u(x,y,t)\\[1.5mm]
\qquad{}+e_{+}(x,y,t)\,{}_{c}D_{y}^{\beta}u(x,y,t)+e_{-}(x,y,t)\,{}_{y}D_{d}^{\beta}u(x,y,t)+f(x,y,t),\\[1.5mm]
\qquad\qquad (x,y,t)\in\Omega\times(0,T],\\[1.5mm]
u(x,y,0)=u_{0}(x,y),\qquad (x,y)\in\overline{\Omega},\\[1.5mm]
u(x,y,t)=0,\qquad (x,y,t)\in\partial\Omega\times[0,T],
\end{array}
\right.
\label{eq:10-1}
\end{equation}
其中 $1<\alpha,\beta<2$, $d_{\pm}(x,y,t),e_{\pm}(x,y,t)\geqslant 0$, $\Omega=(a,b)\times(c,d)$. 此外, ${}_{a}D_{x}^{\alpha}$, ${}_{x}D_{b}^{\alpha}$, ${}_{c}D_{y}^{\beta}$ 与 ${}_{y}D_{d}^{\beta}$ 是如下定义的左, 右 \newterm{riemannliouville}{Riemann--Liouville 分数阶导数}{Riemann--Liouville fractional derivative}:
\[
\left\{
\begin{array}{l}
\displaystyle {}_{a}D_{x}^{\alpha}u(x,y,t)\equiv\frac{1}{\Gamma(2-\alpha)}\frac{\mathrm{d}^{2}}{\mathrm{d}x^{2}}\int_{a}^{x}\frac{u(\xi,y,t)}{(x-\xi)^{\alpha-1}}\,\mathrm{d}\xi,\\[4mm]
\displaystyle {}_{x}D_{b}^{\alpha}u(x,y,t)\equiv\frac{1}{\Gamma(2-\alpha)}\frac{\mathrm{d}^{2}}{\mathrm{d}x^{2}}\int_{x}^{b}\frac{u(\xi,y,t)}{(\xi-x)^{\alpha-1}}\,\mathrm{d}\xi,\\[4mm]
\displaystyle {}_{c}D_{y}^{\beta}u(x,y,t)\equiv\frac{1}{\Gamma(2-\beta)}\frac{\mathrm{d}^{2}}{\mathrm{d}y^{2}}\int_{c}^{y}\frac{u(x,\eta,t)}{(y-\eta)^{\beta-1}}\,\mathrm{d}\eta,\\[4mm]
\displaystyle {}_{y}D_{d}^{\beta}u(x,y,t)\equiv\frac{1}{\Gamma(2-\beta)}\frac{\mathrm{d}^{2}}{\mathrm{d}y^{2}}\int_{y}^{d}\frac{u(x,\eta,t)}{(\eta-y)^{\beta-1}}\,\mathrm{d}\eta.
\end{array}
\right.
\]
这里, \newterm{gammafunction}{Gamma 函数}{Gamma function} $\Gamma(z)$ 由如下积分定义 \cite{src:68}:
\[
  \Gamma(z)\equiv\int_{0}^{\infty}e^{-t}t^{z-1}\,\mathrm{d}t,\qquad\operatorname{Re}(z)>0,
\]
其中 $\operatorname{Re}(\cdot)$ 表示复数的实部.

为逼近含左分数阶导数的 FDE 与含双边 Riemann--Liouville 分数阶导数的 FDE, Meerschaert 与 Tadjeran 提出了 \newterm{shiftedgrunwald}{移位 Grünwald 差分格式}{shifted Grünwald difference scheme} \cite{src:61,src:62}. 他们的方法被证明是无条件稳定的. 2012 年, Tian, Zhou 与 Deng 提出了 Crank--Nicolson 加权与移位 Grünwald 差分 (CN-WSGD) 格式 \cite{src:79}. 用 CN-WSGD 格式离散 \eqref{eq:10-1}, 可得一列线性方程组
\[
  \underline{\mathbf{A}}^{(m+1)}\mathbf{u}^{(m+1)}=\mathbf{b}^{(m+1)},\qquad m=0,1,\ldots,M-1,
\]
其中 $M$ 为时间剖分数.

上述系数矩阵具有特殊的\newterm{toeplitzlike}{类 Toeplitz 结构}{Toeplitz-like structure}: 它可写成若干对角乘 Toeplitz 矩阵之和 \cite{src:46,src:79,src:90}. 这些线性方程组非对称且规模大, 故用重启 GMRES 方法求解. 数值结果表明, 当扩散系数不小且网格尺寸较小时, 离散后的方程组可能变得非常病态. 为加速重启 GMRES 方法的收敛, 通过推广 \cite{src:59} 中提出的\newterm{bandedprecond}{带状预条件子}{banded preconditioner}, 用系数矩阵的一个块-带状-带状-块近似的\newterm{ilufact}{不完全 LU 分解}{incomplete LU factorization}来构造预条件子. 该预处理迭代法的一个优点是: 利用系数矩阵的类 Toeplitz 结构, 它在每个时间层只需 $\mathcal{O}(N^{2}\log N)$ 次运算.

本章安排如下. \ref{sec:10-2} 节用 CN-WSGD 格式离散 FDE \eqref{eq:10-1}, 并给出该有限差分方法的矩阵形式. 随后在 \ref{sec:10-3} 节为 FDE \eqref{eq:10-1} 的离散线性方程组构造预条件子. 最后给出数值结果以说明预条件子的效率.

\section{CN-WSGD 格式}
\label{sec:10-2}

\subsection{分数阶扩散方程的离散}
\label{sec:10-2-1}

本小节用 CN-WSGD 格式离散 FDE \eqref{eq:10-1}. 将时间区间 $[0,T]$ 分成 $M$ 个长度为 $\tau=T/M$ 的子区间, 空间区域 $\Omega$ 以网格尺寸 $h_{x}=(b-a)/(N_{1}+1)$ 与 $h_{y}=(d-c)/(N_{2}+1)$ 均匀剖分, 其中 $M$, $N_{1}$, $N_{2}$ 为正整数. 令 $x_{i}=a+ih_{x}$ ($i=0,1,\ldots,N_{1}+1$), $y_{j}=c+jh_{y}$ ($j=0,1,\ldots,N_{2}+1$), $t^{(m)}=m\tau$ ($m=0,1,\ldots,M$). 令 $t^{(m+1/2)}=(t^{(m)}+t^{(m+1)})/2$ ($m=1,2,\ldots,M-1$).

下面考虑离散后线性方程组的数值解法, 这些方程组由带移位参数 $(p,q)=(1,0)$ 的二阶精度 CN-WSGD 格式 \cite{src:79} 得到:
\[
\left\{
\begin{array}{l}
\displaystyle {}_{L}\mathcal{D}_{h_{x}}^{\alpha}u(x_{i},y_{j},t^{(m)})=\frac{1}{h_{x}^{\alpha}}\sum_{k=0}^{i}w_{k}^{(\alpha)}u(x_{i-k+1},y_{j},t^{(m)}),\\[4mm]
\displaystyle {}_{R}\mathcal{D}_{h_{x}}^{\alpha}u(x_{i},y_{j},t^{(m)})=\frac{1}{h_{x}^{\alpha}}\sum_{k=0}^{N_{1}-i+1}w_{k}^{(\alpha)}u(x_{i+k-1},y_{j},t^{(m)}),\\[4mm]
\displaystyle {}_{L}\mathcal{D}_{h_{y}}^{\beta}u(x_{i},y_{j},t^{(m)})=\frac{1}{h_{y}^{\beta}}\sum_{k=0}^{j}w_{k}^{(\beta)}u(x_{i},y_{j-k+1},t^{(m)}),\\[4mm]
\displaystyle {}_{R}\mathcal{D}_{h_{y}}^{\beta}u(x_{i},y_{j},t^{(m)})=\frac{1}{h_{y}^{\beta}}\sum_{k=0}^{N_{2}-j+1}w_{k}^{(\beta)}u(x_{i},y_{j+k-1},t^{(m)}).
\end{array}
\right.
\]
这里, 系数 $w_{k}^{(\gamma)}$ 由
\[
  w_{0}^{(\gamma)}=\frac{\gamma}{2}g_{0}^{(\gamma)},\qquad
  w_{k}^{(\gamma)}=\frac{\gamma}{2}g_{k}^{(\gamma)}+\frac{2-\gamma}{2}g_{k-1}^{(\gamma)},\quad k\geqslant 1,
\]
给出, 其中 $g_{k}^{(\gamma)}=(-1)^{k}\binom{\gamma}{k}$, $\binom{\gamma}{k}$ 为分数阶二项式系数. 对 $\gamma=\alpha,\beta$, 分数阶二项式系数 $g_{j}^{(\gamma)}$ 具有如下性质 \cite{src:61,src:62,src:66,src:87}:
\begin{equation}
\left\{
\begin{array}{l}
g_{0}^{(\gamma)}=1,\quad g_{1}^{(\gamma)}=-\gamma,\quad g_{2}^{(\gamma)}>g_{3}^{(\gamma)}>\cdots>0,\\[2mm]
\displaystyle \sum_{j=0}^{\infty}g_{j}^{(\gamma)}=0,\quad \sum_{j=0}^{n}g_{j}^{(\gamma)}<0,\quad n\geqslant 1,\\[2mm]
g_{j}^{(\gamma)}=\mathcal{O}(j^{-(\gamma+1)}).
\end{array}
\right.
\label{eq:10-2}
\end{equation}

类似地, 系数 $w_{k}^{(\gamma)}$ 具有如下基本性质 \cite{src:79}.

\begin{lemma}
\label{lem:10-1}
对 $1<\gamma<2$, 系数 $w_{k}^{(\gamma)}$ 满足如下性质:
\[
\left\{
\begin{array}{l}
\displaystyle w_{0}^{(\gamma)}=\frac{\gamma}{2},\quad
w_{1}^{(\gamma)}=\frac{2-\gamma-\gamma^{2}}{2}<0,\quad
w_{2}^{(\gamma)}=\frac{\gamma(\gamma^{2}+\gamma-4)}{4},\\[3mm]
1\geqslant w_{0}^{(\gamma)}\geqslant w_{3}^{(\gamma)}\geqslant w_{4}^{(\gamma)}\geqslant\cdots\geqslant 0,\\[2mm]
\displaystyle \sum_{j=0}^{\infty}w_{j}^{(\gamma)}=0,\quad \sum_{j=0}^{n}w_{j}^{(\gamma)}<0,\quad n\geqslant 2.
\end{array}
\right.
\]
此外, 还有
\[
  w_{j}^{(\gamma)}=\mathcal{O}(j^{-(\gamma+1)}).
\]
\end{lemma}

令
\[
\left\{
\begin{array}{l}
\displaystyle \delta_{x}^{\alpha,(m)}=\frac{d_{i,j}^{+,(m)}\tau}{2}\,{}_{L}\mathcal{D}_{h_{x}}^{\alpha}+\frac{d_{i,j}^{-,(m)}\tau}{2}\,{}_{R}\mathcal{D}_{h_{x}}^{\alpha},\\[4mm]
\displaystyle \delta_{y}^{\beta,(m)}=\frac{e_{i,j}^{+,(m)}\tau}{2}\,{}_{L}\mathcal{D}_{h_{y}}^{\beta}+\frac{e_{i,j}^{-,(m)}\tau}{2}\,{}_{R}\mathcal{D}_{h_{y}}^{\beta},
\end{array}
\right.
\]
其中
\[
\begin{array}{ll}
d_{i,j}^{+,(m)}=d_{+}(x_{i},y_{j},t^{(m)}), & d_{i,j}^{-,(m)}=d_{-}(x_{i},y_{j},t^{(m)}),\\[1mm]
e_{i,j}^{+,(m)}=e_{+}(x_{i},y_{j},t^{(m)}), & e_{i,j}^{-,(m)}=e_{-}(x_{i},y_{j},t^{(m)}).
\end{array}
\]
二维 CN-WSGD 有限差分方法可表述如下:
\begin{equation}
  \bigl(1-\delta_{x}^{\alpha,(m+1)}-\delta_{y}^{\beta,(m+1)}\bigr)u_{i,j}^{(m+1)}
  =\bigl(1+\delta_{x}^{\alpha,(m)}+\delta_{y}^{\beta,(m)}\bigr)u_{i,j}^{(m)}
  +\tau f_{i,j}^{(m+1/2)},
\label{eq:10-3}
\end{equation}
其中 $u_{i,j}^{(m)}=u(x_{i},y_{j},t^{(m)})$, $f_{i,j}^{(m+1/2)}=f(x_{i},y_{j},t^{(m+1/2)})$.

\subsection{CN-WSGD 格式的矩阵形式}
\label{sec:10-2-2}

为给出格式 \eqref{eq:10-3} 的矩阵形式, 先定义如下记号.

(1) 令 $N=N_{1}=N_{2}$. $N^{2}$ 维向量 $\mathbf{u}^{(m)}$, $\mathbf{d}_{+}^{(m)}$, $\mathbf{d}_{-}^{(m)}$, $\mathbf{e}_{+}^{(m)}$, $\mathbf{e}_{-}^{(m)}$ 与 $\mathbf{f}^{(m+1/2)}$ 定义如下:
\[
\left\{
\begin{array}{l}
\mathbf{u}^{(m)}=\bigl[u_{1,1}^{(m)},\ldots,u_{N_{1},1}^{(m)},u_{1,2}^{(m)},\ldots,u_{N_{1},2}^{(m)},\ldots,u_{1,N_{2}}^{(m)},\ldots,u_{N_{1},N_{2}}^{(m)}\bigr]^{T},\\[2mm]
\mathbf{d}_{+}^{(m)}=\bigl[d_{1,1}^{+,(m)},\ldots,d_{N_{1},1}^{+,(m)},d_{1,2}^{+,(m)},\ldots,d_{N_{1},2}^{+,(m)},\ldots,d_{1,N_{2}}^{+,(m)},\ldots,d_{N_{1},N_{2}}^{+,(m)}\bigr]^{T},\\[2mm]
\mathbf{d}_{-}^{(m)}=\bigl[d_{1,1}^{-,(m)},\ldots,d_{N_{1},1}^{-,(m)},d_{1,2}^{-,(m)},\ldots,d_{N_{1},2}^{-,(m)},\ldots,d_{1,N_{2}}^{-,(m)},\ldots,d_{N_{1},N_{2}}^{-,(m)}\bigr]^{T},\\[2mm]
\mathbf{e}_{+}^{(m)}=\bigl[e_{1,1}^{+,(m)},\ldots,e_{N_{1},1}^{+,(m)},e_{1,2}^{+,(m)},\ldots,e_{N_{1},2}^{+,(m)},\ldots,e_{1,N_{2}}^{+,(m)},\ldots,e_{N_{1},N_{2}}^{+,(m)}\bigr]^{T},\\[2mm]
\mathbf{e}_{-}^{(m)}=\bigl[e_{1,1}^{-,(m)},\ldots,e_{N_{1},1}^{-,(m)},e_{1,2}^{-,(m)},\ldots,e_{N_{1},2}^{-,(m)},\ldots,e_{1,N_{2}}^{-,(m)},\ldots,e_{N_{1},N_{2}}^{-,(m)}\bigr]^{T},\\[2mm]
\begin{array}{l}
\mathbf{f}^{(m+1/2)}=\bigl[f_{1,1}^{(m+1/2)},\ldots,f_{N_{1},1}^{(m+1/2)},f_{1,2}^{(m+1/2)},\ldots,f_{N_{1},2}^{(m+1/2)},\ldots,f_{1,N_{2}}^{(m+1/2)},\\
\qquad\qquad\qquad \ldots,f_{N_{1},N_{2}}^{(m+1/2)}\bigr]^{T}.
\end{array}
\end{array}
\right.
\]

(2) 定义如下 $N^{2}\times N^{2}$ 对角矩阵:
\[
  D_{+}^{(m+1)}=\operatorname{diag}\bigl(r\mathbf{d}_{+}^{(m+1)}\bigr),\qquad
  D_{-}^{(m+1)}=\operatorname{diag}\bigl(r\mathbf{d}_{-}^{(m+1)}\bigr),
\]
\[
  E_{+}^{(m+1)}=\operatorname{diag}\bigl(s\mathbf{e}_{+}^{(m+1)}\bigr),\qquad
  E_{-}^{(m+1)}=\operatorname{diag}\bigl(s\mathbf{e}_{-}^{(m+1)}\bigr),
\]
其中 $r=\dfrac{\tau}{2h_{x}^{\alpha}}$, $s=\dfrac{\tau}{2h_{y}^{\beta}}$.

(3) 定义 Toeplitz 矩阵 $A_{Q}^{\gamma,L}$ 与 $A_{Q}^{\gamma,R}$:
\[
  A_{Q}^{\gamma,L}=
  \begin{bmatrix}
    w_{1}^{(\gamma)} & w_{0}^{(\gamma)} & 0 & \cdots & 0 & 0\\
    w_{2}^{(\gamma)} & w_{1}^{(\gamma)} & w_{0}^{(\gamma)} & \ddots & & 0\\
    \vdots & w_{2}^{(\gamma)} & w_{1}^{(\gamma)} & \ddots & \ddots & \vdots\\
    \vdots & \ddots & \ddots & \ddots & \ddots & 0\\
    w_{Q-1}^{(\gamma)} & & \ddots & \ddots & w_{1}^{(\gamma)} & w_{0}^{(\gamma)}\\
    w_{Q}^{(\gamma)} & w_{Q-1}^{(\gamma)} & \cdots & \cdots & w_{2}^{(\gamma)} & w_{1}^{(\gamma)}
  \end{bmatrix},
  \qquad
  A_{Q}^{\gamma,R}=\bigl(A_{Q}^{\gamma,L}\bigr)^{T}.
\]

(4) 定义 $N^{2}\times N^{2}$ 矩阵 $A_{x}^{(m+1)}$ 与 $A_{y}^{(m+1)}$:
\[
\left\{
\begin{array}{l}
A_{x}^{(m+1)}=D_{+}^{(m+1)}\bigl(I_{N_{2}}\otimes A_{N_{1}}^{\alpha,L}\bigr)+D_{-}^{(m+1)}\bigl(I_{N_{2}}\otimes A_{N_{1}}^{\alpha,R}\bigr),\\[2mm]
A_{y}^{(m+1)}=E_{+}^{(m+1)}\bigl(A_{N_{2}}^{\beta,L}\otimes I_{N_{1}}\bigr)+E_{-}^{(m+1)}\bigl(A_{N_{2}}^{\beta,R}\otimes I_{N_{1}}\bigr).
\end{array}
\right.
\]

(5) 定义 $A^{(m+1)}=A_{x}^{(m+1)}+A_{y}^{(m+1)}$.

于是 CN-WSGD 有限差分方法 \eqref{eq:10-3} 可写成如下矩阵形式
\[
  \bigl(I_{N^{2}}-A^{(m+1)}\bigr)\mathbf{u}^{(m+1)}
  =\bigl(I_{N^{2}}+A^{(m)}\bigr)\mathbf{u}^{(m)}+\tau\mathbf{f}^{(m+1/2)},
\]
其中 $I_{Q}$ 表示 $Q\times Q$ 单位矩阵. 令
\[
  \underline{\mathbf{A}}^{(m+1)}=I_{N^{2}}-A^{(m+1)},\qquad
  \mathbf{b}^{(m+1)}=\bigl(I_{N^{2}}+A^{(m)}\bigr)\mathbf{u}^{(m)}+\tau\mathbf{f}^{(m+1/2)}.
\]
上述线性方程组可写成
\begin{equation}
  \underline{\mathbf{A}}^{(m+1)}\mathbf{u}^{(m+1)}=\mathbf{b}^{(m+1)}.
\label{eq:10-4}
\end{equation}

令 $\gamma_{0}=(-1+\sqrt{17})/2\approx 1.562$. 由引理 \ref{lem:10-1} 可知, 若 $\gamma\in[\gamma_{0},2)$, 则 $w_{2}^{(\gamma)}\geqslant 0$, 且系数矩阵 $\underline{\mathbf{A}}^{(m+1)}$ 是严格对角占优的 \newterm{mmatrix}{M-矩阵}{M-matrix}. 此外, 矩阵 $(I_{N^{2}}-A^{(m+1)})^{-1}(I_{N^{2}}+A^{(m+1)})$ 的谱半径小于 $1$. 当扩散系数与时间无关时, CN-WSGD 格式 \cite{src:79} 是无条件稳定的, 且截断误差为 $\mathcal{O}(h_{x}^{2}+h_{y}^{2}+\tau^{2})$. 与 \cite{src:78} 中给出的方法相比, CN-WSGD 格式无需 Richardson 外推即为二阶精度.

\section{分数阶扩散方程的稀疏带状预条件子}
\label{sec:10-3}

离散线性方程组 \eqref{eq:10-4} 通常规模很大, 且具有特殊的类 Toeplitz 结构, 求解这类结构方程组是 CN-WSGD 格式应用中的核心问题. 我们用第 6 章介绍的重启 GMRES 方法求解 \eqref{eq:10-4}. 为加快重启 GMRES 方法的收敛速度, 使用由系数矩阵的一个块-带状-带状-块近似的不完全 LU 分解所定义的预条件子.

\subsection{预条件子的构造}
\label{sec:10-3-1}

注意 Toeplitz 矩阵 $A_{N_{1}}^{\alpha,L}$, $A_{N_{1}}^{\alpha,R}$, $A_{N_{2}}^{\beta,L}$ 与 $A_{N_{2}}^{\beta,R}$ 具有特殊的对角衰减性质: 上述每个矩阵的第 $(i,j)$ 个元素为 $\mathcal{O}(\abs{i-j}^{-\gamma})$ 阶, 其中 $\gamma=\alpha,\beta$; 参见 \eqref{eq:10-2}. 因此, 我们考虑基于系数矩阵带状近似来构造预条件子. 由于对角补偿是设计预条件子的一种有用技巧 \cite{src:2,src:59}, 我们也考虑使用这一技巧.

先考虑离散方程组 \eqref{eq:10-4} 的带状预条件子. 令
\[
  A_{Q,k}^{\gamma,L}=
  \begin{bmatrix}
    w_{1}^{(\gamma)} & w_{0}^{(\gamma)} & & & \\
    \vdots & w_{1}^{(\gamma)} & w_{0}^{(\gamma)} & & \\
    w_{k}^{(\gamma)} & & \ddots & \ddots & \\
    & \ddots & & \ddots & w_{0}^{(\gamma)}\\
    & & w_{k}^{(\gamma)} & \cdots & w_{1}^{(\gamma)}
  \end{bmatrix}_{Q\times Q},
  \qquad
  A_{Q,k}^{\gamma,R}=\bigl(A_{Q,k}^{\gamma,L}\bigr)^{T}.
\]
令 $C_{Q}^{(\gamma)}$ 为由
\[
  C_{Q}^{\gamma}=
  \begin{bmatrix}
    0 & & & & & \\
    & \ddots & & & & \\
    & & 0 & & & \\
    & & & w_{k+1}^{(\gamma)} & & \\
    & & & & \ddots & \\
    & & & & & \displaystyle\sum_{j=k+1}^{Q}w_{j}^{(\gamma)}
  \end{bmatrix}_{Q\times Q}
\]
定义的对角矩阵. 矩阵 $A_{x,k}^{(m+1)}$ 与 $A_{y,k}^{(m+1)}$ 定义为
\[
\left\{
\begin{array}{l}
A_{x,k}^{(m+1)}=D_{+}^{(m+1)}\Bigl(I_{N_{2}}\otimes\bigl(A_{N_{1},k}^{\alpha,L}+C_{N_{1}}^{\alpha}\bigr)\Bigr)
 +D_{-}^{(m+1)}\Bigl(I_{N_{2}}\otimes\bigl(A_{N_{1},k}^{\alpha,R}+C_{N_{1}}^{\alpha}\bigr)\Bigr),\\[2mm]
A_{y,k}^{(m+1)}=E_{+}^{(m+1)}\Bigl(\bigl(A_{N_{2},k}^{\beta,L}+C_{N_{2}}^{\beta}\bigr)\otimes I_{N_{1}}\Bigr)
 +E_{-}^{(m+1)}\Bigl(\bigl(A_{N_{2},k}^{\beta,R}+C_{N_{2}}^{\beta}\bigr)\otimes I_{N_{1}}\Bigr).
\end{array}
\right.
\]
于是矩阵 $A^{(m+1)}$ 可分裂为两部分
\[
  A^{(m+1)}=A_{k}^{(m+1)}+B_{k}^{(m+1)},
\]
其中
\[
  A_{k}^{(m+1)}=A_{x,k}^{(m+1)}+A_{y,k}^{(m+1)},\qquad
  B_{k}^{(m+1)}=A^{(m+1)}-A_{k}^{(m+1)}.
\]
矩阵 $A_{k}^{(m+1)}$ 是稀疏带状矩阵, 只含 $A^{(m+1)}$ 的 $4k-3$ 条非零对角线, 而 $I_{N^{2}}-A_{k}^{(m+1)}$ 的带宽为 $2N_{1}(k-1)+1$. 由于对 $\gamma=\alpha,\beta$ 有 $w_{k}^{(\gamma)}=\mathcal{O}(k^{-(\gamma+1)})$, 故
\[
  \frac{\norm{\bigl(I_{N^{2}}-A^{(m+1)}\bigr)-\bigl(I_{N^{2}}-A_{k}^{(m+1)}\bigr)}_{\infty}}
       {\norm{I_{N^{2}}-A^{(m+1)}}_{\infty}}
  =\frac{\norm{B_{k}^{(m+1)}}_{\infty}}{\norm{I_{N^{2}}-A^{(m+1)}}_{\infty}}
  =\mathcal{O}(k^{-\varsigma}),
\]
其中 $\varsigma=\min\{\alpha,\beta\}$. 若 $k$ 充分大, 则 $I_{N^{2}}-A_{k}^{(m+1)}$ 与 $I_{N^{2}}-A^{(m+1)}$ 的相对距离很小. 因此, 稀疏带状矩阵 $I_{N^{2}}-A_{k}^{(m+1)}$ 是系数矩阵 $\underline{\mathbf{A}}^{(m+1)}$ 的一个良好近似. 于是可考虑用稀疏带状矩阵 $I_{N^{2}}-A_{k}^{(m+1)}$ 的不完全 LU 分解来为 $\underline{\mathbf{A}}^{(m+1)}$ 构造预条件子.

由于带状矩阵 $I_{N^{2}}-A_{k}^{(m+1)}$ 只含 $4k-3$ 条非零对角线, 它非常稀疏. 然而 $I_{N^{2}}-A_{k}^{(m+1)}$ 的带宽为 $2N_{1}(k-1)+1$, 其逆或 LU 分解的计算代价很高, 因为矩阵的稀疏性会被破坏. 为克服这一困难, 我们对带状矩阵 $I_{N^{2}}-A_{k}^{(m+1)}$ 计算无填充的不完全 LU 分解 (ILU(0)). 设 $L$ 与 $U$ 是其 ILU(0) 分解的因子, 则二者都只含 $2k-1$ 条非零对角线. 因此 $P_{k}^{(m+1)}=LU$ 也是稀疏矩阵. 由于 $P_{k}^{(m+1)}$ 可以是 $I_{N^{2}}-A_{k}^{(m+1)}$ 的良好近似, 它也可作为系数矩阵 $\underline{\mathbf{A}}^{(m+1)}$ 的良好近似. 对矩阵 $P_{k}^{(m+1)}$ 及其逆 $(P_{k}^{(m+1)})^{-1}$ 作用于向量 $\mathbf{x}\in\R^{N^{2}}$, 有
\[
  P_{k}^{(m+1)}\mathbf{x}=LU\mathbf{x},\qquad
  \bigl(P_{k}^{(m+1)}\bigr)^{-1}\mathbf{x}=U^{-1}L^{-1}\mathbf{x}.
\]
于是, 预条件子 $P_{k}^{(m+1)}$ 是通过稀疏带状矩阵 $I_{N^{2}}-A_{k}^{(m+1)}$ 的不完全 LU 分解隐式构造的.

\subsection{预处理 GMRES($\rho$) 方法}
\label{sec:10-3-2}

现已为 $\underline{\mathbf{A}}^{(m+1)}$ 构造了预条件子 $P_{k}^{(m+1)}$. 我们不去直接求解线性方程组 \eqref{eq:10-4}, 而是求解如下预处理后的线性方程组
\begin{equation}
  \bigl(P_{k}^{(m+1)}\bigr)^{-1}\underline{\mathbf{A}}^{(m+1)}\mathbf{u}^{(m+1)}
  =\bigl(P_{k}^{(m+1)}\bigr)^{-1}\mathbf{b}^{(m+1)}.
\label{eq:10-5}
\end{equation}
我们把第 6 章介绍的 GMRES 方法与重启策略结合起来求解上述预处理后的线性方程组. 这里用 $\mathrm{GMRES}(\rho)$ 表示在 $\rho$ 步迭代后重启的重启 GMRES 方法. 关于重启 GMRES 方法的介绍, 可参阅 \cite{src:70}. 为完整起见, 我们对预处理后的线性方程组 \eqref{eq:10-5} 陈述在每次迭代时检验残差的预处理 $\mathrm{GMRES}(\rho)$ 方法 \cite{src:59}. 这里使用如下记号: 符号 $\mathbf{r}_{t}$ 与 $\mathbf{r}_{p}$ 分别表示原线性方程组与预处理后线性方程组的当前残差. 用 \textit{IterMax} 表示最大迭代次数, $\epsilon>0$ 表示给定的残差相对精度.

在每个时间层 $t^{(m+1)}$, 选取 $u_{0}$ 作为 $\mathbf{u}^{(m+1)}$ 的初始猜测. 计算 $I_{N^{2}}-A_{k}^{(m+1)}$ 的 ILU(0) 分解: $P^{(m+1)}=LU$. 然后计算残差 $\mathbf{r}=\mathbf{b}^{(m+1)}-\underline{\mathbf{A}}^{(m+1)}\mathbf{u}^{(m+1)}$, 并令 $\mathbf{r}_{t}=\mathbf{r}$. 初始时 $\mu=0$. 在迭代步中, 有

% 算法 10.1 是不能断行的大块伪代码，若作普通段落会整体跳到次页，
% 以致上一页底部只剩标题（孤行）且页面被 \flushbottom 强行拉伸。
% 故包进 figure 浮动体（无 \caption，不影响算法编号），让后文回填该页。
\begin{figure}[!t]
\begin{algorithm}{预处理 GMRES($\rho$) 方法}
  \label{alg:10-1}
  \[
    \left\{
    \begin{array}{l}
      \textbf{While}\ \mu\leqslant \textit{IterMax}\ \text{and}\ \norm{\mathbf{r}_{t}}_{2}/\norm{\mathbf{b}^{(m+1)}}_{2}>\epsilon\\
      \quad \mu=\mu+1\\
      \quad \mathbf{r}_{p}=U^{-1}L^{-1}\mathbf{r};\quad \beta=\norm{\mathbf{r}_{p}};\quad \mathbf{v}_{1}=\mathbf{r}_{p}/\beta\\
      \quad j=0\\
      \quad \textbf{While}\ j\leqslant\rho\ \text{and}\ \norm{\mathbf{r}_{t}}_{2}/\norm{\mathbf{b}^{(m+1)}}_{2}>\epsilon\\
      \qquad j=j+1\\
      \qquad \mathbf{w}_{j}=U^{-1}L^{-1}\underline{\mathbf{A}}^{(m+1)}\mathbf{v}_{j}\\
      \qquad \textbf{for}\ i=1:j\\
      \qquad\quad h_{ij}=\mathbf{w}_{j}^{T}\mathbf{v}_{i}\\
      \qquad\quad \mathbf{w}_{j}=\mathbf{w}_{j}-h_{ij}\mathbf{v}_{i}\\
      \qquad \textbf{end}\\
      \qquad h_{j+1,j}=\norm{\mathbf{w}_{j}}_{2}\\
      \qquad \mathbf{v}_{j+1}=\mathbf{w}_{j}/h_{j+1,j}\\
      \qquad \textbf{compute}\ \mathbf{y}_{j}\ \text{the minimizer of}\ \norm{\beta\mathbf{e}_{1}-\bar{H}_{j}\mathbf{y}}_{2}\\
      \qquad \mathbf{r}_{t}=\mathbf{r}-LUV_{j+1}\bar{H}_{j}\mathbf{y}_{j}\\
      \quad \textbf{end}\\
      \quad \mathbf{r}=\mathbf{r}_{t}\\
      \quad \mathbf{u}^{(m+1)}=\mathbf{u}^{(m+1)}+V_{j}\mathbf{y}_{j}\\
      \textbf{end}
    \end{array}
    \right.
  \]
\end{algorithm}
\end{figure}

其中 $\bar{H}_{j}=(h_{ik})_{1\leqslant i\leqslant j+1,\,1\leqslant k\leqslant j}$, $\mathbf{e}_{1}=[1,0,\ldots,0]^{T}\in\R^{j+1}$. 算法中的矩阵 $V_{j}=[\mathbf{v}_{1}\mid\mathbf{v}_{2}\mid\cdots\mid\mathbf{v}_{j}]\in\R^{N^{2}\times j}$ ($j\leqslant N^{2}$) 是列标准正交的矩阵.

在预处理 $\mathrm{GMRES}(\rho)$ 方法中, 每次迭代的主要计算工作如下:
\begin{enumerate}
  \item[(1)] 计算 $\mathbf{w}=U^{-1}L^{-1}\underline{\mathbf{A}}^{(m+1)}\mathbf{v}_{j}$;
  \item[(2)] 计算 $\mathbf{r}_{t}=\mathbf{r}-LUV_{j+1}\bar{H}_{j}\mathbf{y}_{j}$.
\end{enumerate}

由于 $\underline{\mathbf{A}}^{(m+1)}$ 具有特殊的类 Toeplitz 结构, 对 $\mathbf{x}\in\R^{N^{2}}$ 的矩阵-向量乘法 $\underline{\mathbf{A}}^{(m+1)}\mathbf{x}$ 可以非常高效地计算 \cite{src:16,src:43}. 注意
\[
\begin{array}{rl}
  \underline{\mathbf{A}}^{(m+1)}\mathbf{x}
  = & \mathbf{x}-D_{+}^{(m+1)}\bigl(I_{N_{2}}\otimes A_{L}^{\alpha,N_{1}}\bigr)\mathbf{x}
    -D_{-}^{(m+1)}\bigl(I_{N_{2}}\otimes A_{R}^{\alpha,N_{1}}\bigr)\mathbf{x}\\[1.5mm]
  & -E_{+}^{(m+1)}\bigl(A_{L}^{\beta,N_{2}}\otimes I_{N_{1}}\bigr)\mathbf{x}
    -E_{-}^{(m+1)}\bigl(A_{R}^{\beta,N_{2}}\otimes I_{N_{1}}\bigr)\mathbf{x},
\end{array}
\]
以及
\[
  I_{N_{2}}\otimes A_{R}^{\alpha,N_{1}}=\bigl(I_{N_{2}}\otimes A_{L}^{\alpha,N_{1}}\bigr)^{T},\qquad
  A_{R}^{\beta,N_{2}}\otimes I_{N_{1}}=\bigl(A_{L}^{\beta,N_{2}}\otimes I_{N_{1}}\bigr)^{T}.
\]
计算 $\underline{\mathbf{A}}^{(m+1)}\mathbf{x}$ 的主要开销在于
\[
  \bigl(I_{N_{2}}\otimes A_{L}^{\alpha,N_{1}}\bigr)\mathbf{x},\quad
  \bigl(I_{N_{2}}\otimes A_{R}^{\alpha,N_{1}}\bigr)\mathbf{x},\quad
  \bigl(A_{L}^{\beta,N_{2}}\otimes I_{N_{1}}\bigr)\mathbf{x},\quad
  \bigl(A_{R}^{\beta,N_{2}}\otimes I_{N_{1}}\bigr)\mathbf{x}.
\]
乘法 $(I_{N_{2}}\otimes A_{L}^{\alpha,N_{1}})\mathbf{x}$ 与 $(I_{N_{2}}\otimes A_{R}^{\alpha,N_{1}})\mathbf{x}$ 可通过计算 $(I_{N_{2}}\otimes A_{L}^{\alpha,N_{1}})(\mathbf{x}+\sqrt{-1}\,\breve{\mathbf{x}})$ 同时得到, 其中 $\breve{\mathbf{x}}=[x_{N},x_{N-1},\ldots,x_{1}]^{T}$ \cite{src:66}. 类似地, 也可同时计算 $(A_{L}^{\beta,N_{2}}\otimes I_{N_{1}})\mathbf{x}$ 与 $(A_{R}^{\beta,N_{2}}\otimes I_{N_{1}})\mathbf{x}$. 利用 FFT, 矩阵-向量乘法 $\underline{\mathbf{A}}^{(m+1)}\mathbf{x}$ 只需 $\mathcal{O}(N^{2}\log N)$ 次运算.

此外, 由于 $L$ 与 $U$ 都是只含 $2k-1$ 条非零对角线的稀疏三角矩阵, 计算 $L\mathbf{x}$, $U\mathbf{x}$, $L^{-1}\mathbf{x}$ 与 $U^{-1}\mathbf{x}$ 只需 $\mathcal{O}(N^{2}k)$ 次运算. 因此, 计算 $\mathbf{w}=U^{-1}L^{-1}\underline{\mathbf{A}}^{(m+1)}\mathbf{v}_{j}$ 需要 $\mathcal{O}(N^{2}\log N)$ 次运算. 至于 $\mathbf{r}_{t}$ 的计算, 由于 $V_{j+1}$ 与 $\bar{H}_{j}$ 分别是 $N^{2}\times(j+1)$ 与 $(j+1)\times j$ 矩阵, 可在 $\mathcal{O}(N^{2}(j+k))$ 次运算内得到. 于是, 对大的 $N$, 预处理 GMRES 方法每个时间层的计算量为 $\mathcal{O}(N^{2}\log N)$.

\subsection{数值结果}
\label{sec:10-3-3}

本小节给出一个数值例子, 以说明预处理 $\mathrm{GMRES}(\rho)$ 方法相对于无预条件子的 $\mathrm{GMRES}(\rho)$ 方法的表现. 对 CN-WSGD 格式 \eqref{eq:10-4}, 选取初始猜测为
\[
  \mathbf{u}_{0}=
  \left\{
  \begin{array}{ll}
    \mathbf{u}^{(0)}, & m=0,\\[1mm]
    \tilde{\mathbf{u}}^{(m+1)}\equiv 2\mathbf{u}^{(m)}-\mathbf{u}^{(m-1)}, & m>0.
  \end{array}
  \right.
\]
由于 $\tilde{\mathbf{u}}^{(m+1)}$ 是 $\mathbf{u}^{(m+1)}$ 的良好近似 \cite{src:90}, 我们取其为初始猜测. 停机准则为
\[
  \frac{\norm{\mathbf{r}_{j}}_{2}}{\norm{\mathbf{b}^{(m+1)}}_{2}}<10^{-7},
\]
其中 $\mathbf{r}_{j}$ 是第 $j$ 步迭代后的残差向量. 在数值实验中, 取参数 $k=4$. 此外, 对 $\mathrm{GMRES}(\rho)$ 方法取 $\rho=20$.

下表中, 用 $N_{1}$ 与 $N_{2}$ 分别表示 $x$ 方向与 $y$ 方向的空间剖分数, 用 $M$ 表示时间剖分数. 最后一个时间层上真解与数值解之间的最大误差记为 Error. 此外, 预处理 $\mathrm{GMRES}(\rho)$ 方法记为 $\mathrm{PGMRES}(\rho)$.

\begin{example}
\label{ex:10-1}
考虑方程 \eqref{eq:10-1}, 其中 $\alpha=1.8$, $\beta=1.9$. 左, 右扩散系数由
\[
\begin{array}{ll}
  d_{+}(x,y,t)=4(1+t)x^{1.8}(1+y), & d_{-}(x,y,t)=4(1+t)(1-x)^{1.8}(1+y),\\[1mm]
  e_{+}(x,y,t)=4(1+t)(1+x)y^{1.9}, & e_{-}(x,y,t)=4(1+t)(1+x)(1-y)^{1.9}
\end{array}
\]
给出. 空间区间为 $\Omega=(0,1)\times(0,1)$, 时间区间为 $[0,T]=[0,1]$. 初始条件为 $u_{0}(x,y)=x^{3}y^{3}(1-x)^{3}(1-y)^{3}$. 源项由
\[
\begin{array}{rcl}
  f(x,y,t) & = & -e^{-t}\Bigl[\bigl(x^{3}(1-x)^{3}y^{3}(1-y)^{3}\bigr)\\[1mm]
  & & +4(1+t)\Bigl(\dfrac{\Gamma(4)}{\Gamma(2.2)}\bigl(x^{3}+(1-x)^{3}\bigr)
        -\dfrac{3\Gamma(5)}{\Gamma(3.2)}\bigl(x^{4}+(1-x)^{4}\bigr)\\[1mm]
  & & +\dfrac{3\Gamma(6)}{\Gamma(4.2)}\bigl(x^{5}+(1-x)^{5}\bigr)
        -\dfrac{\Gamma(7)}{\Gamma(5.2)}\bigl(x^{6}+(1-x)^{6}\bigr)\Bigr)y^{3}(1-y)^{3}(1+y)\\[1mm]
  & & +4(1+t)\Bigl(\dfrac{\Gamma(4)}{\Gamma(2.1)}\bigl(y^{3}+(1-y)^{3}\bigr)
        -\dfrac{3\Gamma(5)}{\Gamma(3.1)}\bigl(y^{4}+(1-y)^{4}\bigr)\\[1mm]
  & & +\dfrac{3\Gamma(6)}{\Gamma(4.1)}\bigl(y^{5}+(1-y)^{5}\bigr)
        -\dfrac{\Gamma(7)}{\Gamma(5.1)}\bigl(y^{6}+(1-y)^{6}\bigr)\Bigr)x^{3}(1-x)^{3}(1+x)\Bigr].
\end{array}
\]
给出. 该分数阶扩散方程的解为
\[
  u(x,y,t)=e^{-t}x^{3}(1-x)^{3}y^{3}(1-y)^{3}.
\]
\end{example}

% 本章末页只余「一段收尾 + 表 10.1」，若沿用全书 \flushbottom 会被强行撑开
% （Underfull vbox badness 10000）。此处在分组内局部改为 \raggedbottom，
% 并在组内 \clearpage 让末页以自然高度出页，组结束后自动恢复全书的 \flushbottom。
\begingroup
\raggedbottom
\needspace{14\baselineskip}
例 \ref{ex:10-1} 的数值结果见表 \ref{tab:10-1}. 由表 \ref{tab:10-1} 可见, $\mathrm{GMRES}(20)$ 方法所需的迭代次数远多于 $\mathrm{PGMRES}(20)$ 方法. 随着 $N_{1}$ 与 $N_{2}$ 增大, $\mathrm{PGMRES}(20)$ 方法所需的迭代次数增长缓慢.

\begin{table}[htbp]
  \centering
  \caption{例 \ref{ex:10-1} 的数值结果}
  \label{tab:10-1}
  \begin{tabular}{cccc}
    \toprule
    $N_{1}=N_{2}=M$ & GMRES(20) & PGMRES(20) & Error \\
    \midrule
    $2^{4}$ & 48.750  & 11.000 & $1.1386\times10^{-6}$ \\
    $2^{5}$ & 81.594  & 12.406 & $3.0206\times10^{-7}$ \\
    $2^{6}$ & 157.750 & 14.250 & $7.6918\times10^{-8}$ \\
    $2^{7}$ & 273.914 & 17.055 & $1.9381\times10^{-8}$ \\
    $2^{8}$ & 416.625 & 20.477 & $4.8489\times10^{-9}$ \\
    \bottomrule
  \end{tabular}
\end{table}
\clearpage
\endgroup


% ---------- 附录（如有） ----------
% \appendix
% \include{chapters/chA}

% ---------- 参考文献 ----------
\makebibliography

% ---------- 符号表（后置材料）----------
% 与「术语索引」同为无编号章，进目录；内容见 back/notation.tex。
% back/notation.tex —— 符号表（后置材料）
% 依第 1 章 §1.1「基本记号」所列条目，并补入各章陆续引入的主要记号。
% 每条均在正文中有对应出现，不求穷尽，只收全书反复使用者；章次注于含义之后。
% 排版体例：与「术语索引」一致（longtable + booktabs，表头 \headfont）。
\cleardoublepage
\phantomsection
\chapter*{符号表}
\addcontentsline{toc}{chapter}{符号表}
\markboth{符号表}{符号表}

本书的记号在第 1 章开篇集中列出, 其余随各章内容陆续引入. 为便于查阅, 现将全书主要记号分类汇总于下; 含义之后括注该记号主要出现的章次, 未注者通贯全书.

\begingroup
\renewcommand{\arraystretch}{1.18}%
\begin{longtable}{@{}>{\raggedright\arraybackslash}p{0.27\textwidth}>{\raggedright\arraybackslash}p{0.63\textwidth}@{}}
  \toprule
  {\headfont 记号} & {\headfont 含义} \\
  \midrule
  \endfirsthead
  \toprule
  {\headfont 记号} & {\headfont 含义} \\
  \midrule
  \endhead
  \bottomrule
  \endlastfoot

  \multicolumn{2}{@{}l}{\headfont 一、数集与空间} \\*[3pt]
  $\R,\ \C$ & 实数域与复数域 \\
  $\R^n,\ \C^n$ & 实, 复 $n$ 维列向量空间 \\
  $\R^{m\times n},\ \C^{m\times n}$ & $m\times n$ 实, 复矩阵空间 \\
  $\ii$ & 虚数单位 $\sqrt{-1}$ \\
  $\mathcal{F}$ & 计算机浮点数集 (第 3 章) \\

  \addlinespace[6pt]
  \multicolumn{2}{@{}l}{\headfont 二、向量, 矩阵与基本运算} \\*[3pt]
  $A=[a_{ij}]$ & 第 $(i,j)$ 个元素为 $a_{ij}$ 的矩阵 \\
  $\mathbf{x},\mathbf{y},\mathbf{z}$ & 列向量, 以小写粗正体表示 \\
  $A^T,\ A^*$ & 转置, 共轭转置 \\
  $A^{-1},\ A^{\dagger}$ & 逆矩阵, Moore--Penrose 广义逆 (第 4 章) \\
  $I_n$ & $n$ 阶单位矩阵 \\
  $\mathbf{e}_i$ & 第 $i$ 个单位向量, 即 $I_n$ 的第 $i$ 列 \\
  $\mathbf{0}$ & 零矩阵或零向量, 尺寸视上下文而定 \\
  $|A|$ & 元素取绝对值的矩阵, $(|A|)_{ij}=|a_{ij}|$ \\
  $\det(A),\ \rank(A),\ \tr(A)$ & 行列式, 秩, 迹 \\
  $\dim(S)$ & 向量空间 $S$ 的维数 \\
  $\diag(\cdot)$ & 以给定元素为对角元的对角矩阵 \\
  $\spn\{\cdot\}$ & 由给定向量张成的子空间 \\
  $\operatorname{nullity}(A)$ & $A$ 的零度, 即 $\mathcal{N}(A)$ 的维数 (第 4 章) \\
  $A(i,j),\ A(i,:),\ A(:,j)$ & 仿 MATLAB 的元素, 行, 列记号 \\
  $A(k:l,p:q)$ & 由第 $k$ 至 $l$ 行, 第 $p$ 至 $q$ 列构成的子矩阵 \\

  \addlinespace[6pt]
  \multicolumn{2}{@{}l}{\headfont 三、范数与条件数} \\*[3pt]
  $\norm{\cdot}$ & 矩阵或向量的范数, 所指依上下文判定 \\
  $\norm{\cdot}_p$ & $p$ 范数 \\
  $\norm{\cdot}_1,\ \norm{\cdot}_2,\ \norm{\cdot}_\infty$ & 1 范数, 2 范数与 $\infty$ 范数 \\
  $\norm{A}_F$ & Frobenius 范数 \\
  $\triplenorm{\cdot}$ & 习题中表示矩阵范数的三竖线记号 \\
  $\kappa(A)$ & $A$ 的条件数, $\kappa(A)=\norm{A}\norm{A^{-1}}$ (第 3 章) \\
  $c(x)$ & 函数求值的条件数 (第 1 章) \\

  \addlinespace[6pt]
  \multicolumn{2}{@{}l}{\headfont 四、特征值, 奇异值与谱} \\*[3pt]
  $\lambda_i$ & 矩阵的第 $i$ 个特征值 \\
  $\Lambda$ & 以特征值为对角元的矩阵, $\Lambda=\diag(\lambda_1,\ldots,\lambda_n)$ \\
  $\rho(A)$ & $A$ 的谱半径, $\rho(A)=\max\{|\lambda|:\lambda\ \text{属于}\ A\ \text{的谱}\}$ (第 3 章) \\
  $\sigma_i$ & 第 $i$ 个奇异值 (第 8 章) \\
  $\Sigma,\ \Sigma_r$ & 奇异值构成的对角矩阵, $r=\rank(A)$ (第 8 章) \\
  $\mathbf{u}_i,\ \mathbf{v}_i$ & 第 $i$ 个左, 右奇异向量 (第 8 章) \\

  \addlinespace[6pt]
  \multicolumn{2}{@{}l}{\headfont 五、矩阵分解与标准形} \\*[3pt]
  $A=LU$ & LU 分解, $L$ 为下三角矩阵, $U$ 为上三角矩阵 (第 2 章) \\
  $A=QR$ & QR 分解, $Q$ 正交或酉, $R$ 为上三角矩阵 (第 4 章) \\
  $A=U\Sigma V^*$ & 奇异值分解 (第 4, 8 章) \\
  $Q^TAQ$ & 实对称矩阵的正交对角化, 即 $Q^TAQ=\diag(\lambda_1,\ldots,\lambda_n)$ (第 8 章) \\
  $A=X\Lambda X^{-1}$ & 可对角化矩阵的特征分解 (第 7 章) \\

  \addlinespace[6pt]
  \multicolumn{2}{@{}l}{\headfont 六、子空间与集合} \\*[3pt]
  $\mathcal{R}(A)$ & $A$ 的值域, 即列空间 (第 4 章) \\
  $\mathcal{N}(A)$ & $A$ 的零空间 (第 4 章) \\
  $\mathcal{S}^{\perp}$ & 子空间 $\mathcal{S}$ 的正交补 (第 4 章) \\
  $\mathcal{X}$ & 最小二乘问题的解集 (第 4 章) \\
  $\mathcal{P}_k$ & 次数不超过 $k$ 且 $P_k(0)=1$ 的多项式集合 (第 6 章) \\
  $\mathcal{M}_U$ & 可被酉矩阵 $U$ 同时对角化的矩阵集合 (第 6 章) \\
  $\mathcal{O}(\cdot)$ & 运算量或误差的阶 (大 $\mathcal{O}$ 记号) \\

  \addlinespace[6pt]
  \multicolumn{2}{@{}l}{\headfont 七、迭代法与 Krylov 子空间} \\*[3pt]
  $\mathbf{x}_k$ & 第 $k$ 步迭代向量 (第 5, 6 章) \\
  $\mathbf{r}_k$ & 第 $k$ 步残差, $\mathbf{r}_k=\mathbf{b}-A\mathbf{x}_k$ (第 6 章) \\
  $\mathcal{K}(A,\mathbf{r}_0,k)$ & 由 $A$ 与初始残差 $\mathbf{r}_0$ 生成的 $k$ 维 Krylov 子空间 (第 6 章) \\
  $\mathbf{r}_t,\ \mathbf{r}_p$ & 原方程组与预处理后方程组的当前残差 (第 10 章) \\
  $\rho$ & 重启 GMRES 方法的重启参数 (第 10 章) \\

  \addlinespace[6pt]
  \multicolumn{2}{@{}l}{\headfont 八、误差与浮点运算} \\*[3pt]
  $fl(x)$ & 实数 $x$ 的浮点数 (第 3 章) \\
  $\mu_\varepsilon$ & 机器精度 (第 3 章) \\
  $\epsilon$ & 小参数: 扰动幅度或残差相对精度 \\
  $\hat{\mathbf{x}}$ & 计算解, 即精确解的近似 (第 3 章) \\

  \addlinespace[6pt]
  \multicolumn{2}{@{}l}{\headfont 九、常微分方程与分数阶扩散方程} \\*[3pt]
  BVM,\ IVM & 边值法 (boundary value method) 与初值法 (initial value method) (第 9 章) \\
  DDE,\ LMF & 时滞微分方程 (delay differential equation) 与线性多步公式 (linear multistep formula) (第 9 章) \\
  $\mathcal{D}_{\nu_1,\nu_2}$ & BVM 的 $A_{\nu_1,\nu_2}$-稳定区域 (第 9 章) \\
  $\Gamma$ & 复平面中的一条曲线 (第 9 章) \\
  $\Gamma(z)$ & Gamma 函数, $\Gamma(z)=\int_0^\infty e^{-t}t^{z-1}\,\mathrm{d}t$ (第 10 章) \\
  ${}_{a}D_x^{\alpha},\ {}_{x}D_b^{\alpha}$ & $x$ 方向左, 右 Riemann--Liouville 分数阶导数 (第 10 章) \\
  ${}_{c}D_y^{\beta},\ {}_{y}D_d^{\beta}$ & $y$ 方向左, 右 Riemann--Liouville 分数阶导数 (第 10 章) \\
  $\alpha,\ \beta$ & 分数阶导数的阶, $1<\alpha,\beta<2$ (第 10 章) \\
  $\Omega,\ \partial\Omega$ & 求解区域及其边界 (第 10 章) \\
  $\tau$ & 时间步长, $\tau=T/M$ (第 10 章) \\
  $h_x,\ h_y$ & $x$, $y$ 方向的网格尺寸 (第 10 章) \\
  $N_1,\ N_2,\ M$ & 空间剖分数与时间剖分数 (第 10 章) \\
  $\underline{\mathbf{A}}^{(m+1)},\ \mathbf{u}^{(m+1)}$ & 第 $m+1$ 个时间层的系数矩阵与解向量 (第 10 章) \\

\end{longtable}
\endgroup


% ---------- 术语索引（必须最后） ----------
\printterminology

\end{document}
